Fix the tower, measures and matching data of \S\ref{newtheta:global}.
Theorem~\ref{tcr:global} constructs a supported test with nonzero
primed period and proves \(I(\Phi)=cI'(\Phi')\) for every matched
pair, at every allowed real signature. The scalar \(c\ne0\) belongs
to the chosen data. In its two rank recursions the comparison scalar
\(C\) is uniform in the admissible finite data once the adapted
upper and lower geometries, real matching maps and measures are
fixed. The proof includes the inert finite and real local comparisons.
The evaluation lattice used in Section~\ref{main:assembly} requires
more than proportionality of these two linear functionals. On
marked components the doubled kernel gives spanning over the
coefficient field (Propositions~\ref{dg:identity} and~\ref{munf:span}).
Theorem~\ref{bfam:oncharts} controls the finite denominators and
Petersson norms of spanning sections under its chart and cusp
hypotheses. The final change of weight also involves the integral
determinants of two conjugate Hodge eigenspaces. On the indefinite
selected family their common local orders and the CM evaluation
in Theorem~\ref{cmeig:splitmodel} give this change at every weight.

\subsection{Local matching in split matrix coordinates}\label{newtheta:local}

We compute the two local functionals used at a selected split place. This
also gives test vectors on which the functions agree at \emph{every} group
element. Let \(k\) be a nonarchimedean local field, \(\mathcal O\) its
valuation ring, \(\pi\) a uniformizer, and \(|\pi|=|\mathcal O/\pi|^{-1}\).
Put \(d=a+a'\), with \(a,a'>0\), \(G_j=\mathrm{GL}_j(k)\),
\(K_j=\mathrm{GL}_j(\mathcal O)\), and give \(G_j\) Haar measure
such that \(K_j\) has volume one. Use the row matrices
\(X_0=[I_a\;0]\) and \(Y_0=[0\;I_{a'}]\).
The characters \(\alpha,\beta:k^\times\to\mathbb C^\times\) are smooth
and unitary, and \(\nu=\alpha/\beta\). In the two polarizations the
orbital functions, on locally constant functions compactly supported in
the respective full-row-rank sets, are
\begin{align}
 B_\Phi(h)&=\beta(\det h)|\det h|^{-a/2}
 \int_{G_a}\Phi(gX_0h^{-t})\,\nu(\det g)|\det g|^{d/2}\,dg ,
 \label{newtheta:big}\\
 S_\Psi(h)&=\alpha(\det h)|\det h|^{a'/2}
 \int_{G_{a'}}\Psi(gY_0h)\,\nu(\det g)|\det g|^{d/2}\,dg .
 \label{newtheta:small}
\end{align}
Here \(h^{-t}=(h^{-1})^t\). For any fixed full-rank \(X\) the map
\(g\mapsto gX\) is a homeomorphism onto the full-rank matrices with
the same row space as \(X\). The image is closed in the full-rank set,
since having rows in a fixed subspace is closed. The inverse is right
multiplication by any right inverse of \(X\). Thus the integrals here
have compact support in \(g\). The large-side definition in the original
model is a sum of integrals over Cartan double cosets. A compact subset of
\(G_a\) meets only finitely many of its disjoint open double cosets, so
the two definitions agree in this case.
\src{src/Development/AbelianVarietyTheory/Definitions/Divisors/RankEqCardIffExistsFrameMatrixFirstColumnBlock.lean}{703}{The two orbital functions and their determinant factors}.

\begin{proposition}[The common image of full-rank split tests]\label{newtheta:split}
Let \(P\subset G_d\) be the block upper triangular parabolic with
diagonal blocks of sizes \(a,a'\). The image of each of
\eqref{newtheta:big} and \eqref{newtheta:small} on these test spaces is
the full space of smooth functions \(f:G_d\to\mathbb C\) satisfying
\[
 f\left(\begin{pmatrix}A&*\\0&D'\end{pmatrix}h\right)
 =\alpha(\det A)\beta(\det D')|\det A|^{a'/2}
             |\det D'|^{-a/2}f(h).
\]
In particular there are matched tests with value one at the identity.
\end{proposition}
\begin{proof}
For \(p\in P\), \(X_0p^{-t}=A^{-t}X_0\) and \(Y_0p=D'Y_0\).
In the first integral substitute \(gA^{-t}\) for \(g\), and in the
second \(gD'\) for \(g\). Right translations preserve Haar measure on
\(G_j\), as is also seen from its density
\(|\det g|^{-j}\prod_{i,l}dg_{il}\) in additive matrix coordinates.
Using \(\alpha=\nu\beta\) gives the displayed character in both cases.
Right translation of the output is obtained by the corresponding linear
substitution in the test times its determinant prefactor. Compact support
and local constancy give invariance under some open subgroup on the right.

Let \(K_j(e)\) be the kernel of reduction modulo \(\pi^e\). Take
\(e\ge1\) so large that both characters are trivial on
\(1+\pi^e\mathcal O\). For \(\Phi_e,\Psi_e\) take the indicators of
\(X_0+\pi^eM_{a,d}(\mathcal O)\), \(Y_0+\pi^eM_{a',d}(\mathcal O)\)
divided by the volumes of \(K_a(e)\), \(K_{a'}(e)\), respectively.
Both supports are compact and have full row rank. For \(h\in K_d\),
the integrality and full rank modulo \(\pi\) of \(gX_0h^{-t}\) on
the first support force \(g\in K_a\), by multiplying on the right by
\(h^t\). The same argument with \(h^{-1}\) gives \(g\in K_{a'}\)
on the second support. Modulo \(\pi^e\), the equations become
\([g\;0]=X_0h^t\) and \([0\;g]h=Y_0\), respectively.
Each has solutions precisely when the lower left block of \(h\)
is zero modulo \(\pi^e\). In this case the diagonal blocks \(A,D'\)
are invertible modulo \(\pi^e\). The solutions are one congruence
coset, given by \(g\equiv A^t\) or \(g\equiv (D')^{-1}\).
The two normalized values are consequently
\(\beta(\det h)\nu(\det A)\) and
\(\alpha(\det h)\nu(\det D')^{-1}\). They agree because
\(\det h\equiv\det A\det D'\pmod{\pi^e}\).
In particular the common function is one at \(h=1\).

Write \(K=K_d\) and \(U=(P\cap K)K_d(e)\). This is the inverse image
of the block upper triangular subgroup modulo \(\pi^e\): one lifts
the two invertible diagonal blocks and the upper right block to obtain
an element of \(P\cap K\). Our common function on \(K\) is supported
on \(U\) and is right \(K_d(e)\)-invariant. Also \(G_d=PK\). Indeed
the intersection of any \(a'\)-dimensional row space with \(\mathcal O^d\)
is a saturated free summand of rank \(a'\) by the elementary divisor
theorem. Extending an integral basis of that summand gives a matrix of
\(K\) with the same last-row space as the given matrix of \(G_d\).
Their quotient belongs to \(P\).

For any smooth \(f\) with the stated covariance choose \(e\) still
larger so that \(f|_K\) is right \(K_d(e)\)-invariant, possible by
compactness. On the finitely many cosets \(U k_j\) in \(K\), the
restrictions of \(f\) are determined by \(f(k_j)\). In fact \(K_d(e)\)
is normal in \(K\), so \(f(uk_j)\) transforms by the prescribed
character on \(P\cap K\) and trivially on the congruence kernel.
The right translates \(h\mapsto B_{\Phi_e}(hk_j^{-1})
=S_{\Psi_e}(hk_j^{-1})\) have exactly these transformation laws,
disjoint supports on \(K\), and value one at \(k_j\). Their linear
combination with coefficients \(f(k_j)\) agrees with \(f\) first on
\(K\) and then on \(PK\). Each translate comes from a test of the
same allowed kind. This proves surjectivity on both sides.
\end{proof}

The ramified split condition asks for a large-side value even on
tests meeting the singular locus. There is a direct justification for
the finite shell prescription in this case.
\begin{lemma}[Cancellation near the singular strata]\label{newtheta:shell}
Suppose \(\nu\) is nontrivial on \(\mathcal O^\times\). For every
locally constant compactly supported \(\Phi\) on \(M_{a,d}(k)\) and
every fixed full-rank \(X\), only finitely many double cosets
\(K_a\operatorname{diag}(\pi^{m_i})K_a\) have a nonzero integral of
\(\Phi(gX)\nu(\det g)|\det g|^{d/2}\).
\end{lemma}
\begin{proof}
Choose a right inverse \(X^\dagger\). The identity
\(g=(gX)X^\dagger\) and compact support give a common lower bound for
the valuations of the entries of all contributing \(g\). Multiplication
by \(K_a\) preserves the supremum norm of matrix entries in either
direction, so the bound applies to each \(m_i\). Choose \(u\in
\mathcal O^\times\) with \(\nu(u)\ne1\). Compact support and local
constancy give a common additive period lattice
\(\pi^tM_{a,d}(\mathcal O)\) for \(\Phi\), by covering its compact
support by finitely many balls of constancy. If \(m_i\) is large,
replacing \(k_1\) by \(k_1\operatorname{diag}(1,\ldots,u,\ldots,1)\)
in \(g=k_1\operatorname{diag}(\pi^{m_j})k_2\), with \(u\) at
position \(i\), changes \(gX\) by an element of this period lattice
uniformly in \(k_1,k_2\in K_a\). The determinant character is
multiplied by \(\nu(u)\), while its absolute value is unchanged.
The normalized measure on a double coset is the pushforward of product
probability Haar measure on \(K_a^2\): both measures are invariant
under its transitive left and right \(K_a\) actions. Averaging over
\(k_1\) thus makes this shell integral zero. With all \(m_i\) bounded
above and below, only finitely many integer tuples remain.
\end{proof}

\begin{proposition}[The whole kernel on the smaller split side]\label{split:kernel}
Put \(n=a'<a\) and \(s=d/2>n\), and set
\(\lambda(g)=\nu(\det g)|\det g|^s\) on \(G=\mathrm{GL}_n(k)\).
On all locally constant compactly supported functions on \(M_{n,d}(k)\),
the integral \eqref{newtheta:small} converges absolutely and its kernel
is precisely the span of
\[
        \Phi(g\,\cdot)-\lambda(g)^{-1}\Phi,\qquad g\in G .
\]
For unramified \(\nu\), the value at the identity of the standard integral
matrix indicator is \(\prod_{j=0}^{n-1}(1-\nu(\pi)q^{j-s})^{-1}\ne0\),
where \(q=|\mathcal O/\pi|\).
\end{proposition}
\begin{proof}
For any full-rank \(X\), boundedness of \(gX\) bounds \(g\) by a right
inverse of \(X\). On that bounded set the absolute integration density
is a constant times \(|\det g|^{s-n}\) in additive matrix coordinates.
It is bounded, so the integral converges. Haar substitution shows that
the displayed differences have zero integral.

On the full-rank open set the map taking the row span is a principal
\(G\)-bundle over a Grassmannian, trivial on each pivot chart by
row reduction. Subdivide the projection of a compact support into
finitely many disjoint compact open subsets of such charts. In chart
coordinates each test is a finite sum of products of a base test and
a test on \(G\). This follows by covering its compact support by
compact open rectangles of constancy and taking a common refinement.
On the fibre \(G\) the quotient by the displayed differences is
one-dimensional, detected by integration against \(\lambda(g)\,dg\).
In fact, for any finite list of tests take a compact open subgroup
\(J\subseteq\ker\lambda\) fixing all of them on the right. Modulo
differences \(1_{tJ}\) equals \(\lambda(t)1_J\); its integral is
\(\lambda(t)\operatorname{vol}(J)\). This proves the assertion on one
fibre. A base refinement where the base factors are constant gives the
same kernel description throughout the full-rank open set. All its
row spaces are tested by \(Y_0h\) as \(h\) varies.

We show that rank \(r<n\) adds no classes to this quotient. Over
a pivot chart of the row-span Grassmannian the rank \(r\) stratum has
fibre the space of injective \(n\)-by-\(r\) matrices. Here a basis of
the row space identifies the original matrix with its matrix of row
coefficients. The source group acts transitively on this fibre by left
multiplication. Suppose a linear functional \(D\) there satisfies
\(D(f(g\,\cdot))=\lambda(g)^{-1}D(f)\).
Around \(X_r=(I_r,0)^t\), take
\(V=X_r+\pi^t M_{n,r}(\mathcal O)\) for \(t\ge1\) sufficiently large
that \(\nu\) is trivial on \(1+\pi^t\mathcal O\). The congruence
subgroup \(K_n(t)\) acts transitively on \(V\), by completing its
first \(r\) columns, and preserves additive volume on \(V\).
Its invariant functionals on tests supported in \(V\) are multiples of
that volume: averaging any such test over the compact group gives the
constant average on \(V\), and is a finite average by smoothness.
Take \(t_0=\operatorname{diag}(I_r,\pi I_{n-r})\), which fixes \(X_r\)
and sends \(V\) into \(V\). On the one hand
\(D(1_{t_0V})=|\det t_0|^r D(1_V)\), by additive volume columnwise.
On the other hand equivariance gives \(\lambda(t_0)D(1_V)\).
Since \(|\lambda(t_0)|=|\det t_0|^s\) and \(r<s\), we have \(D(1_V)=0\).
The translates of \(V\) cover the fibre, so \(D=0\). For \(r=0\)
the fibre is a point and the same character test gives zero directly.
Testing products with each base indicator proves vanishing of every
equivariant functional on the stratum. Its quotient by differences
is zero since linear functionals separate points of any vector space.

Filter the closed rank-deficient locus by the finitely many closed
sets \(\{\operatorname{rank}\le r\}\). Restriction of tests to a closed
subset is onto: a compact open subset of it can be extended to a finite
union of compact open neighborhoods in the ambient space, chosen to
avoid its closed complement within the subset. The kernel consists of
tests compactly supported on the open complement. Taking quotients by
differences is right exact. Induction therefore gives zero quotient for
the whole rank-deficient locus. Consequently each test on the matrix
space is congruent modulo differences to one supported on the full-rank
open set. Their orbital integrals agree, so the full-rank calculation
proves the asserted equality of kernels.

Finally the integral at the identity for the spherical indicator is over
the integral nonsingular matrices. Its right \(K_n\)-cosets correspond
to the finite-index sublattices of \(\mathcal O^n\) and have volume one.
For unramified \(\nu\), the weight of an index \(q^b\) lattice is
\((\nu(\pi)q^{-s})^b\). In upper triangular column Hermite form with
diagonal \(\pi^{b_i}\), the entries in row \(i\) above the diagonal
are taken modulo \(\pi^{b_i}\). There are
\(\prod_i q^{(n-i)b_i}\) choices and \(b=\sum b_i\).
Summing the \(n\) convergent geometric series gives the displayed product.
\end{proof}
\begin{corollary}[Both orbital maps on ramified split tests]\label{split:ram}
If \(\nu|_{\mathcal O^\times}\ne1\), the larger-side shell prescription
has kernel the corresponding span of source-action differences on
\emph{all} tests. Both full test spaces have the common image of
Proposition~\ref{newtheta:split}.
\end{corollary}
\begin{proof}
In the larger integral first replace \(d/2\) in the integration weight
by \(z\) with \(\Re z>a\). Then it is absolutely convergent by the
bounded-density argument above, and change of variables under every
source translation is valid. For fixed \(\Phi,X\), the additive-period
argument of Lemma~\ref{newtheta:shell} kills all but finitely many
Cartan shells, independently of \(z\), since the determinant norm is
constant on each shell. Their sum is a Laurent polynomial in \(q^{-z}\).
The identity for a given translate, including its factor
\(\nu(\det g)^{-1}|\det g|^{-z}\), consequently holds at \(z=d/2\)
as well. The same change of variables gives parabolic covariance at
that parameter. A compact open subgroup fixes the test on the right
in the target action, which gives smoothness of its orbital function.

The rank argument in Proposition~\ref{split:kernel} applies with
\(n=a\); whenever \(r\ne d/2\), the absolute values in the test
still differ. For \(r=d/2<a\) fix \(u\in\mathcal O^\times\) with
\(\nu(u)\ne1\) and use \(\operatorname{diag}(I_r,u,1,\ldots,1)\)
in the stabilizer. This preserves the compact open ball \(V\) and
changes the functional on its indicator by the factor \(\nu(u)\).
It must therefore vanish there and on all translates. Thus every
singular stratum has zero quotient by differences. The full-rank
kernel calculation and the shell transform on the entire space finish
the proof exactly as on the smaller side.
\end{proof}
In particular every \(\lambda^{-1}\)-equivariant functional for the written
translations on the smaller split model factors through its orbital integral; the same statement applies on
the larger side at the indicated ramified places. Full-rank tests span
both indicated images. For the larger unramified side when \(\nu=1\)
and \(r=d/2\), the stabilizer test gives equality for every determinant.
The next lemma shows that the infinite type excludes this resonance
at the good places of the global construction.

\begin{lemma}[The character at a totally split rational prime]\label{split:ray}
Let \(\mathfrak w\) lie above a rational prime completely split in the
Galois CM field \(E\). If its restriction to \(F\) is outside the
ramification set \(S\) of the triple, then
\(t=\nu_{\mathfrak w}(\pi)\) has infinite order, for either choice of
the place above \(F\) and every uniformizer \(\pi\).
\end{lemma}
\begin{proof}
Fix \(E\subset\mathbb C\). The chosen embeddings above the real places
form a CM type \(\Phi\), and we can regard them as automorphisms of
\(E\) by normality. Put \(m_\tau=m_\alpha(\tau)-m_\beta(\tau)>d\).
The infinite component of \(\nu\) at \(x\in E^\times\) is
\(\prod_{\tau\in\Phi}(\tau x/|\tau x|)^{m_\tau}\).
Take a finite ray modulus for \(\nu\), prime to \(\mathfrak w\).
By finiteness of the ray class group there are \(h>0\) and \(x\in
E^\times\) such that \((x)=\mathfrak w^h\) and \(x\equiv1\) modulo
this modulus. All ramified finite components are one at \(x\);
at every other finite place except \(\mathfrak w\), \(x\) is a unit.
The continuous Hecke character is trivial on diagonal \(E^\times\).
Its value on \(x\), squared, therefore gives
\[
 t^{-2h}=\prod_{\tau\in\Phi}
            \bigl(\tau x/\overline{\tau x}\bigr)^{m_\tau}.
\]
The principal fractional ideal on the right has exponents \(hm_\tau\)
and \(-hm_\tau\) at \(\tau\mathfrak w\) and
\(\bar\tau\mathfrak w\), respectively. These primes are all distinct
because the decomposition group is trivial. The fractional ideal is
nontrivial, whereas every root of unity generates the unit ideal.
This proves the assertion. Unramifiedness gives independence of \(\pi\).
\end{proof}
\src{src/Development/AbelianVarietyTheory/Definitions/Foundations/NotIsSplitOfIsInertAt.lean}{266}{Continuity and triviality on the diagonal};
\src{src/Development/AbelianVarietyTheory/BlockTheory/DivisorBase/WeilTensorModels.lean}{114}{Infinite character on ideles with trivial finite part};
\src{src/Development/AbelianVarietyTheory/BlockTheory/DivisorBase/WeilTensorModels.lean}{374}{Local ramification and infinite exponents}.

\begin{proposition}[The larger orbital transform at a good split place]
\label{split:good}
Suppose \(\nu\) is unramified and \(t=\nu(\pi)\) has infinite order.
Set \(n=a\), \(s=d/2\). On every matrix test \(\Phi\), the larger
transform in \eqref{newtheta:big}, initially integrated with exponent
\(z\) for \(\Re z>n\), has a regular continuation at \(z=s\).
This transform and the smaller transform \eqref{newtheta:small}
have the same image on their whole test spaces, namely the induced
representation in Proposition~\ref{newtheta:split}. The kernel on
each side is the span of its source-action differences. For
unramified \(\alpha,\beta\) their integral-matrix indicators can be
normalized to have the same nonzero image.
\end{proposition}
\begin{proof}
Fix a full-rank \(X\). The function \(H(g)=\Phi(gX)\) on \(M_n(k)\)
is locally constant and compactly supported, since multiplying by a
right inverse of \(X\) bounds its support. In additive matrix
coordinates the integral is a constant times
\[
 Z_n(H,z)=\int_{M_n(k)} H(g)\nu(\det g)|\det g|^{z-n}\,dg .
\]
Assign the integrand value zero on singular matrices, a null set
by successive one-variable polynomial testing and Fubini.
For \(0\le j\le n\) use the same exponent \(z-n\) in \(Z_j\) on
\(M_j(k)\), with \(Z_0\) evaluation at its single point.
We show by induction that \(Z_j\) is rational in \(q^{-z}\) and
its only possible poles for finite \(z\) satisfy
\[
           1-t^i q^{-i(z-n+i)}=0,\qquad 1\le i\le j .
\]
Away from a neighborhood of the zero matrix, a compact support has
a finite disjoint compact open subdivision on each part of which a
chosen nonzero pivot \(b\) has constant valuation. Permute rows and
columns to make it the top left entry. For the other three blocks use
coordinates \((B,C,D-Cb^{-1}B)\) instead of \((B,C,D)\). This
preserves additive measure by translation in \(D\) and factors the
determinant as \(b\det(D-Cb^{-1}B)\), up to a unit sign.
In these coordinates a compactly supported locally constant test is
a finite sum of products on compact open rectangles, by a finite
refinement of a cover by rectangles of constancy. The pivot weight
is a Laurent monomial times a constant. The remaining integrals thus
have the asserted form by induction.

Near zero subtract the constant value of the test times the
indicator of a sufficiently small origin ball. Scalar dilation
reduces that ball to \(M_j(\mathcal O)\) up to a Laurent monomial.
In this unit ball the complement of \(\pi M_j(\mathcal O)\) is
covered by the preceding pivot argument. The integral on the
inner ball is \(t^j q^{-j^2-j(z-n)}\) times the unit-ball integral.
Solving this identity in the half-plane of convergence introduces
exactly the indicated possible denominator with \(i=j\). This
proves rational continuation. For real \(z=s\), a zero of a
denominator would force first \(s-n+i=0\) by \(|t|=1\), and then
\(t^i=1\). Hence every test is regular there.

Source equivariance holds first in the convergent half-plane by
Haar substitution and hence at \(s\) by continuation. The same
reasoning gives the displayed parabolic covariance of
Proposition~\ref{newtheta:split} at \(s\); smoothness follows from
the compact open stabilizer of the test. On full-rank tests the
integrals are entire and the image and kernel were computed above.
On the rank-\(r<n\) stratum the stabilizer test in
Proposition~\ref{split:kernel} would require
\[
          |\det t_0|^r=\nu(\det t_0)|\det t_0|^s,\qquad
          t_0=\operatorname{diag}(I_r,\pi I_{n-r}).
\]
The moduli differ for \(r\ne s\), and for \(r=s\) equality would
give \(t^{n-r}=1\). Thus every singular stratum has zero quotient by
source differences. The rank filtration and the full-rank kernel
calculation prove the asserted kernel on all tests. The smaller
side is Proposition~\ref{split:kernel}.

In the convergent region the Hermite count in
Proposition~\ref{split:kernel} gives the larger spherical value
\(\prod_{j=0}^{n-1}(1-tq^{j-z})^{-1}\). Its continuation at \(s\)
is finite and nonzero; the smaller value has the same formula
with \(n=a'\). If \(\alpha,\beta\) are unramified both output
functions are right \(K_d\)-invariant. Such an induced function is
determined by its value at one since \(G_d=PK_d\), so dividing each
map by its indicated value identifies the two spherical outputs.
\end{proof}

The covariance is \(\delta_P^{1/2}
((\alpha\circ\det)\boxtimes(\beta\circ\det))\):
\(\delta_P(A,D')=|\det A|^{a'}|\det D'|^{-a}\), from conjugation on
the upper right block. The two determinant characters give an
irreducible unitary representation of the Levi, since their external
product is one-dimensional and unitary. Bernstein's irreducibility
theorem for normalized induction of irreducible unitary representations
of a Levi in \(\mathrm{GL}_d(k)\) therefore proves that this common
image is irreducible, both here and in Corollary~\ref{split:ram}.
Bernstein,
\emph{P-invariant distributions on \(\mathrm{GL}(N)\) and the classification
of unitary representations of \(\mathrm{GL}(N)\) (non-Archimedean case)},
\S0.2, corollary, proved in \S8.2~\cite{bernsteinInvariant}.

Lemma~\ref{split:ray} supplies the character hypothesis for every
good totally split place of the triple, at either splitting.
The construction here evaluates a rational continuation on unramified
tests. For the integral-matrix indicator infinitely many Cartan shells
have nonzero integral, so the finite-support sum defining the prescribed
shell map cannot replace this continuation. The matching condition
that specifies that shell map is imposed at \(S\), where the earlier
ramified calculation applies. The results here concern the displayed
matrix models. Passing to the given adelic Weil actions requires the
identifications with these polarizations and their lattice vacua.
\src{src/Development/AbelianVarietyTheory/Definitions/Divisors/RankEqCardIffExistsFrameMatrixFirstColumnBlock.lean}{707}{Finite-support Cartan sum in the prescribed map};
\src{src/Development/AbelianVarietyTheory/BlockTheory/DivisorBase/WeilTensorModels.lean}{470}{Places at which equality with that map is required}.



\begin{lemma}[The finite epsilon factors in a descent]\label{newtheta:gauss}
Let \(q=|\mathcal O/\pi|\) and let the additive character \(\psi\)
have conductor \(\mathcal O\). Suppose the unitary character \(\nu\)
has conductor exponent \(c\ge1\). Put
\[
 G=\sum_{y\in(\mathcal O/\pi^c)^\times}
          \nu(y)^{-1}\psi(y/\pi^c).
\]
Then \(|G|^2=q^c\). In particular
\(q^{-cs}G\,\nu(\pi)^c\ne0\) for every integer \(s\).
\end{lemma}
\begin{proof}
The two characters on units in the sum do not depend on the chosen
lifts. Write \(R=\mathcal O/\pi^c\) and \(\Psi(x)=\psi(x/\pi^c)\).
The additive pairing \((t,x)\mapsto\Psi(tx)\) on \(R\) is perfect
by the conductor assumption. Extend \(\nu^{-1}\) by zero to the
nonunits. Its unnormalized Fourier transform at \(t\in R^\times\)
is \(\nu(t)G\) by changing \(y\) to \(ty\). For \(t\in\pi R\)
choose a unit \(v\equiv1\bmod\pi^{c-1}\) with \(\nu(v)\ne1\),
using primitivity (for \(c=1\) choose any such unit). Then \(tv=t\)
in \(R\), and multiplication of the summation variable by \(v\)
forces that Fourier coefficient to vanish. Summing the squared
moduli of all Fourier coefficients gives \(q^c|R^\times|\) by
expanding and applying additive orthogonality. The same sum is
\(|R^\times|\,|G|^2\), proving the claim.
\end{proof}

\subsection{The global comparison required of the local maps}\label{newtheta:global}

In the global triple, \(E/F\) is CM quadratic, \(E/\mathbb Q\) is
Galois, and \(E/F\) is unramified at finite primes. The additive
character on \(\mathbb A_F/F\) has infinite component
\(\exp(-2\pi i\operatorname{Tr}(x))\). It has conductor \(\mathcal O_v\)
at the designated split places and at each inert place, where the residue
characteristic is odd and the residue cardinality is greater than three.
The three spaces have dimensions \(a,a',d=a+a'\) and signatures
\((a-p,p),(q,a'-q),(d-p-q,p+q)\) at each real place, using \(i\)
times the skew-Hermitian form on the third space. The integral structures
are self-dual for the specified pairings (the skew lattice uses the
trace dual). Each comes with a Cartan involution and the selected local
bases. The rank and weight restrictions are
\[
\begin{gathered}
a'<a,\qquad d\equiv0\pmod2,\qquad 2(p+q)+1<\min\{(a-a')/2,a'\},\\
m_\alpha(w)-m_\beta(w)>d\qquad(w\mid\infty).
\end{gathered}
\]
The unitary Hecke characters \(\alpha,\beta\) of \(E\) have these
infinite exponents in the chosen embeddings, and restrict to
\(\eta_{E/F}^{\,a'},\eta_{E/F}^{\,a}\) on \(\mathbb A_F^\times\).
They are unramified outside the primes above \(S=S_0\sqcup Q\), a
subset of the designated finite places with a chosen splitting
\(E\hookrightarrow F_v\). Their quotient \(\nu\) is ramified at
every prime above \(S\). A further place \(v_1\notin S\) has
underlying rational prime totally split in \(E\), with a choice of
splitting at \(v_1\). The Weil actions use respectively the pairs
\((\beta,\beta/\alpha)\) and \((\alpha,\alpha/\beta)\) in their
prescribed polarizations. The local vacua belong to the prescribed
lattice eigenlines and are standard characteristic functions almost
everywhere.
\src{src/Development/AbelianVarietyTheory/BlockTheory/DivisorBase/WeilTensorModels.lean}{388}{Base, signatures, characters and Weil actions}.

In these models the missing global assertion has three parts.
The first is the simultaneous construction of nonzero irreducible
quotients at infinity and each finite place, as in
Proposition~\ref{geo:theta}. The maps kill the infinitesimal and compact
source actions at infinity (derivatives map to zero and compact actions
preserve the image of each vector) and kill the span of source-action
differences at finite places. Target Lie and compact actions, or target
group actions at finite places, preserve matching. The two vacua have
the same nonzero image at every \(v\notin S\). At \(v\in S\) equality
of images must be equivalent, for \emph{all} Schwartz--Bruhat inputs,
to equality of the two orbital functions, in the coordinates taking the
vacuum to the integral-matrix indicator. At good places (outside \(S\)
and above a totally split rational prime), the kernel on Schwartz tests
must be exactly the span of the source-action differences. Full-rank
primed tests in the conjugate polarization at either splitting there
must still surject. Proposition~\ref{newtheta:split} proves the
full-rank image calculation in split coordinates. Corollary~\ref{split:ram}
and Proposition~\ref{split:good} determine the whole split kernels at
the specified ramified and good places. Section~\ref{threal:main}
establishes convergence for the two adelic kernels and constructs the
remaining local quotients, including all real signatures.

The second assertion is global convergence and proportionality. On the
unprimed side every adelic Schwartz input must give an absolutely
integrable theta kernel in \(g\) for each \(h\). On the primed side
use the span of factorizable inputs whose factor at \(v_1\) vanishes
on rank-deficient matrices in the indicated polarization. This space
must be stable under the Weil action. For each such input and each
integer \(N\ge0\), there must be \(C,k\) such that
\[
 |\Theta_{\omega'}(g,h;\Phi')|\,H_{[U(V')]}(g)^N
                 \le C H_{U(W)}(h)^k .
\]
Here the quotient height is the infimum of the matrix height over
rational left translates. The primed integral must also converge and
have moderate growth. For \emph{every} choice of matching maps, one
nonzero scalar must compare the two integrals for all matched pairs,
including the limits specified in Proposition~\ref{geo:theta}, and
at least one matched pair must have a nonzero value.

The third assertion controls how these scalars vary when passing to
smaller spaces. Write \(k=(a-a')/2\). In the primed descent \(q>0\)
the dimensions drop from \((a,a',d)\) to \((a,a'-2,d-2)\).
In the unprimed descent \(q=0<p\) they drop to \((a-2,a',d-2)\).
The upper and lower models use adapted integral hyperbolic planes
in both changing spaces. If their fractional ideals along the first
isotropic lines are \(\mathfrak b,\mathfrak c\), set
\[
 F_X(j)=\nu(\mathfrak b)^{e_1}\nu(\mathfrak c)^{e_2}
    \prod_{v\in S}\nu_v(s)\,
       q_v^{-c_v j}G(\nu_v,\psi_v)\nu_v(\pi_v)^{c_v}.
\]
The ideal evaluations use the fixed local uniformizer values. Each
branch has one triple of signs \(e_1,e_2,s\in\{1,-1\}\) valid across
the fields and all admissible descents in that branch. For fixed
base, spaces, planes and real matching maps a further \(C\ne0\) must
work for \emph{all} allowed characters, Weil actions and finite
matching maps. For upper and lower constants \(c,c_1\), the relations are
\[
 cL^{S_0}(\nu,k+1)=\frac{c_1}{CF_X(-k)}
       \quad(q>0),\qquad
 c=CF_X(1-k)L^{S_0}(\nu,k)c_1\quad(q=0<p).
\]
Lemma~\ref{newtheta:gauss} proves nonvanishing of the finite factors.
For \(u=k,k+1\) one has \(u>1\), and the corresponding partial
Euler product over \(E\) is nonzero. Indeed each logarithmic factor
is bounded in absolute value by
\(\sum_{j\ge1}N\mathfrak p^{-ju}/j\).
There are at most \([E:\mathbb Q]\) primes above each rational
prime \(l\), all of norm at least \(l\), and
\(\sum_l l^{-u}<\infty\). This proves absolute convergence of the
sum of logarithms. These arguments justify division by the displayed
factors. Theorem~\ref{tcr:global} proves the common-constant assertion and the
identities comparing the automorphic functions through the descent.
In particular equality of a constant term detects a difference only up
to the kernel of that constant-term operator.
\src{src/Development/AbelianVarietyTheory/BlockTheory/DivisorBase/WeilTensorModels.lean}{552}{Quantifiers and the two scalar recursions}.

\paragraph{A published comparison at the anisotropic base.}
Let \(A\) have dimension \(m\ge1\) and be definite at a real place.
Any \(F\)-rational isotropic vector would stay isotropic under that
embedding, so \(A\) has Witt index zero over \(F\). For the doubled
space \(Y\oplus(-Y)\) with \(\dim_EY=n\ge1\), the diagonal copy of
\(Y\) is isotropic of dimension \(n\), half the total dimension.
In the standard Schr\"odinger model let \(\Phi\) be a Schwartz test,
\(\ell\) evaluation at zero, \(\omega\) the split Weil action, and
\(f_{\Phi,s}\) the standard flat Siegel section through
\(x\mapsto\ell(\omega(x)\Phi)\) at \(s_0=(m-n)/2\).
The restriction of its splitting character to \(\mathbb A_F^\times\)
is \(\eta_{E/F}^{\,m}\). For this section the Eisenstein series is
holomorphic at \(s_0\), and its value is
\[
 E(x,s_0,f_{\Phi,s_0})=
 \begin{cases}1&m>n,\\2&m\le n\end{cases}
 \int_{[U(A)]}\Theta(g,x;\Phi)\,d\bar g ,
 \qquad \operatorname{vol}_{d\bar g}([U(A)])=1.
\]
This is the unitary specialization of Gan--Qiu--Takeda, Theorem 7.1.
In their notation \(r=0\), \(d(n)=n\) and
\(\tau(U(A))/[E:F]=1\). Their integral in \S3.1 has probability
measure, which agrees with the average over the fundamental domain in
the doubled formula for any choice of Haar measure. These facts fix
both the sign of the critical parameter and the value at \(m=n\).
Gan, Qiu and Takeda, \S\S3.1 and 7.1~\cite{ganQiuTakeda}.

To apply this identity to an arbitrary doubled presentation one must
intertwine its Darboux coordinates, actions and functional with this
standard model, preserving the rational lattice sum. The defining
identities require that the functional be invariant under the
diagonal Lagrangian translations and the source group, and that its
translates along the opposite rational coordinates sum to the theta
distribution. The separate meromorphic continuation input for a
continuous \(K\)-finite section of the doubled group uses the Siegel
law \(\chi(\det t)|\det t|_{\mathbb A_E}^{\,s+n/2}\).
At each \(s_0\) it requires one pole order and a common neighborhood
on which the pole-cleared family is weakly holomorphic for every
functional bounded on each polynomial-growth unit ball. An independent
application of continuation must identify this topology as well as the
section. The next two subsections construct a supported nonzero
undoubled primed period when \(q=0\), compare both periods when
\(p=q=0\), and calculate the constant terms needed to pass between
ranks. The normalization of the smaller matching maps enters at
that last step.
\src{src/Development/WeilBound/ThetaLifts/LatticeBasics/Hypotheses2.lean}{63}{Doubled actions and summation identity};
\src{src/Development/WeilBound/ThetaLifts/LatticeBasics/HypothesesE.lean}{46}{Growth-functional continuation};
\src{src/Development/WeilBound/ThetaLifts/LatticeBasics/HypothesesS.lean}{11}{Doubled normalized value}.

\subsection{Absolute convergence, support at one split place, and the inert quotient}
\label{threal:main}

Write \(d=a+a'\), \(f=[F:\mathbb Q]\), and use the height \(H\) of
\S\ref{newtheta:global}. Put \([G]=G(F)\backslash G(\mathbb A_F)\).
There is a difference between the estimates needed to define the two
periods and the identity comparing their values. We prove the estimates
for the undoubled groups.

\begin{proposition}\label{threal:decay}
For the signatures and ranks of \S\ref{newtheta:global}, the unprimed
theta kernel is absolutely integrable over \([U(V)]\) at every \(h\)
for every adelic Schwartz function. For a primed test \(\Phi'\) in the
specified full-rank space at \(v_1\) and every \(N\geq0\), there exist
\(C_{\Phi',N}\) and an integer \(k\geq0\) such that
\[
 |\Theta_{\Phi'}(g,h)|\,H_{[U(V')]}(g)^N
                      \leq C_{\Phi',N} H(h)^k .
\]
The primed kernel is absolutely integrable in \(g\), its integral has
moderate growth in \(h\), and the full-rank test space is invariant
under the two unitary actions.
\end{proposition}
\begin{proof}
Use \(Y\) for either source, \(m=\dim_EY\), and let \(r\) be its
\(F\)-rational Witt index. Fix rational hyperbolic pairs
\((e_i,e_i^*)_{1\le i\le r}\) and their orthogonal complement \(Y_0\).
In \(Y\otimes_EW\), take one whole half of each hyperbolic pair
tensored with \(W\) as position coordinates and complete by a rational
Lagrangian in the underlying symplectic space of \(Y_0\otimes_E W\).
The theta sum now has positions
\[
        (x_1,\ldots,x_r,y)\in W(E)^r\times F^b,\qquad
             b=(m-2r)d ,
\]
where \(W(E)\) denotes the underlying \(E\)-vector space.
Choose the half for which the torus in the positive chamber acts on
functions by
\[
 \phi(x,y)\longmapsto
       \left(\prod_i T_i^{d/2}\right)
       \phi(t_1x_1,\ldots,t_rx_r,y),\qquad
          T_i=t_i^{2f},\quad t_i>0 ,
 \tag{*}
\]
up to a scalar of absolute value one. Indeed each of the \(d\)
\(E\)-coordinates contributes \(2f\) real dimensions and the
prefactor is the square root of the real change of volume. The torus
has value \(t_i\) at all real places and unit finite components.
With the parabolic stabilizing the expanding isotropic flag one uses
the opposite position half: the action on arguments is inverse to
the action on translations.

The change from the given Darboux coordinates is rational. Indeed the
images of a rational basis are rational by assumption, and their
determinant is nonzero. Rational linear changes, rational quadratic
phases and partial Fourier transforms preserve the absolute value of
the rational summation functional, by the product formula, triviality
of \(\psi\) on \(F\), and Poisson summation with self-dual measures.
These operations generate the rational symplectic group by the Siegel
Bruhat decomposition. They also show that the absolute value of the
kernel is invariant under rational left translates in either group.
These arguments apply to each allowed unitary implementation: two
unitary operators implementing the same action on the Heisenberg
operators differ by a scalar of modulus one. In Schr\"odinger
coordinates their quotient commutes with every modulation and hence
is multiplication on \(L^2\), and commuting with the translations
makes that multiplier constant. This argument also permits the
use of the usual tensor-product implementations for estimates at a
fixed adelic group element.

Apply adelic reduction theory to the connected reductive group \(U(Y)\).
Its center is the norm-one torus and has no \(F\)-split part.
A minimal \(F\)-parabolic is the stabilizer of the chosen flag.
The quotient is covered by a Siegel set with representatives
\(n m_0 t k_0\), with \(n,m_0,k_0\) in fixed compact sets of the
unipotent radical, norm-one part of the centralizing Levi, and
maximal compact subgroup, respectively. Here \(t_i/t_{i+1}\) and
\(t_r\) are bounded below by a positive constant. When \(r=0\) the
set is compact. These are the compactness and fundamental-set theorems
in adelic reduction theory, with inversion to pass from the right
quotient to the left quotient \cite[\S\S2.1--2.2]{springerReduction}.
The elements \(t^{-1}nt\) stay in a compact set, since the inverse
of every positive root is bounded in this chamber. Thus the
representatives have the form \(tc\), with \(c\) in a fixed compact
set. In the integration coordinates above the Haar density in \(t\)
is a constant times
\[
       \prod_{i=1}^{r} T_i^{-(m-2i+1)}\,\frac{dT_i}{T_i}.
\]
Indeed for each \(j>i\) the two relative roots with weights
\(t_i/t_j,t_it_j\) have parameter in \(E\); the root \(t_i\) has
parameters in the \(m-2r\) dimensional space \(Y_0\) over \(E\);
the root \(t_i^2\) has one parameter over \(F\). Their determinant
on the unipotent radical is
\(\prod_i T_i^{2(r-i)+(m-2r)+1}\).
Integration of a nonnegative rationally invariant function over a
covering Siegel set bounds its quotient integral from above.

Here are the bounds on the test in (*) as \(c,h\) vary. Fix coordinates
over \(\mathbb Q\) using bases of the number fields and put \(D=fmd\).
For a pure archimedean and finite tensor, the finite part of
\(\omega(c,h)\Phi\) is supported on \(M^{-1}\widehat{\mathbb Z}^{D}\)
and has supremum at most \(C M^{D/2}\), for an integer
\[
             1\le M\le C_\Phi H(h)^A .
\]
To see this, start with an integer \(M_0\) such that translations and
modulations by \(M_0\widehat{\mathbb Z}^{2D}\) fix the finite test,
using the standard \(\mathbb Q\)-adelic character and its dual
coordinates. Increasing \(M_0\) also fixes coordinate denominators.
After conjugation it suffices to have
\((c,h)^{-1}M\widehat{\mathbb Z}^{2D}\subset
 M_0\widehat{\mathbb Z}^{2D}\) in the symplectic representation.
At each place \(v\), put \(n_v\) equal to the largest denominator
order of the entries of \(h,h^{-1}\), or zero if all are integral,
using all factors of \(E\otimes F_v\). The extension \(E/F\) is
unramified at finite places. In fixed bases the entries on tensors
involve only matrix multiplication and fixed changes of coordinates.
All the \(c\)'s carry bounded denominators, supported at finitely many
primes. Thus \(M=C_0\prod_{p}p^{\max_{v\mid p}n_v}\) suffices for
a fixed positive integer \(C_0\). Since \(p\le q_v\) and the finite
height factor is at least \(q_v^{n_v}\), this proves the bound.
Invariance by modulations forces the stated support; invariance by
translations makes the function constant on cosets of
\(M\widehat{\mathbb Z}^{D}\), each of volume \(M^{-D}\) in product
measure of integral volume one. Unitarity gives the supremum bound.

At infinity, for every \(B\ge0\), the absolute value of the transformed
test at \(z\) is at most \(C_B H(h)^{A_B}(1+\|z\|)^{-B}\).
Indeed commute multiplication and differentiation through the unitary
operators using Heisenberg covariance. Each is a linear combination of
the original position and momentum operators, with coefficients
bounded by a fixed power of the maximum of the archimedean matrix
norms of \(h,h^{-1}\) and a constant for \(c\). Apply this identity to
words of length at most \(B+D+2\) (round \(B\) up to an integer).
Euclidean Sobolev embedding applied to each monomial times the
transformed function then gives the assertion. Finite sums of pure
tensors satisfy the same majorant. In these estimates \(A_B\) may
depend on the fixed spaces, \(B\), and the test. The specified height
uses the local spectral absolute values over \(F_v\), extending
\(|\pi_v|=q_v^{-1}\), and takes the maximum of \(1\) and the norms of
both matrices.
\src{src/Development/AbelianVarietyTheory/Definitions/Foundations/EventuallyNormFinLeOne.lean}{405}{Matrix height across the places}.

For the unprimed integral fix \(h\); then \(M\) is bounded and
\(t_i\gg1\). Summing the Schwartz majorant with \(B>D+1\) on
\(M^{-1}\mathbb Z^D\) gives \(O_{\Phi,h}(\prod_i T_i^{d/2})\).
For \(r\le p\) every exponent \(d/2-(a-2i+1)\) is strictly negative,
since \(a-d/2=(a-a')/2>2(p+q)+1\). Integration even over the larger
rectangle \(T_i\ge c_1>0\) converges.

For the primed integral we verify the support after the change of
polarization. At \(v_1\) take the specified conjugate split coordinates,
so the original position variable is an \(m\)-by-\(d\) matrix over
\(F_{v_1}\). The two unitary actions there are invertible row and
column changes times a scalar, since they preserve each of the two
split Lagrangians. Therefore they preserve full row rank. This also
proves invariance of the adelic span by factorization, since all other
factors stay Schwartz and a restricted adelic element fixes the
standard lattice factors up to scalar almost everywhere. For clarity,
write the conjugate split matrix first in the rational Witt row basis.
In a hyperbolic pair the two position rows are indexed by \(e_i,e_i^*\).
In the new model the \(E\otimes F_{v_1}\)-valued position block \(x_i\)
contains \emph{one unchanged row} and the Fourier dual of its paired
row. Indeed the tensor symplectic pairing pairs the \(e_i\)-row on
one split factor with the \(e_i^*\)-row on the other and conversely.
One exchanges the position and momentum in one of these two row
pairs, using an invertible linear identification with \(W\) from
its form. The remaining change of polarization takes place in the
\(Y_0\)-block. If the unchanged row vanishes, the original matrix has
deficient rank for every value integrated over in these partial
Fourier transforms. The transformed function therefore vanishes when
\(x_i=0\). This holds separately for each \(i\) and survives finite
sums. The coordinates called for in the support definition are unique
up to scalar on Schwartz--Bruhat functions: an intertwiner with the
stated translation and modulation laws is unique by the local
Schr\"odinger commutant. Algebraically this follows by averaging
modulations over compact open subgroups to project onto support in
any ball, and then translating the indicators of balls.

For a surviving rational summand \(\|x_i\|\ge c_2/M\): the rational
points in the indicated finite support lie in \(M^{-1}\mathbb Z^D\),
and the coordinates were chosen block by block. Given \(L\ge0\),
spend \(rL\) powers in the archimedean majorant to obtain
\(\prod_i (1+t_i\|x_i\|)^{-L}\ll M^{rL}\prod_i t_i^{-L}\).
Another \(D+2\) powers suffice to sum the remaining majorant over
\(M^{-1}\mathbb Z^D\), with bound \(O(M^D)\) by unit-cube counting.
Thus, for every \(J\ge0\),
\[
             |\Theta_{\Phi'}(tc,h)|
                    \ll_{\Phi',J} H(h)^{A_J}\prod_i T_i^{-J}.
\]
Here \(L\) was chosen to absorb both \(2fJ\) and the prefactor in (*).
The quotient height at any translate is at most \(H(tc)\), which on
this set is bounded by \(C\prod_i T_i^{A_0}\), by the matrix formula
for \(H\). Taking \(J\ge NA_0\) gives the claimed inequality.
Integrating the bound for any sufficiently large \(J\) proves
absolute convergence and moderate growth, including the case \(r=0\)
with empty products. All bounds use absolute values, so they apply
to the prescribed Haar fundamental domains.
\end{proof}

\begin{corollary}\label{threal:limits}
An identity with one fixed scalar for all polynomially matched pairs
whose primed input satisfies the support condition implies the same
identity for every matched pair in \S\ref{newtheta:global}.
\end{corollary}
\begin{proof}
Fix \(h\) and a finite Schwartz factor. The unprimed proof bounded the
absolute integral by finitely many archimedean Schwartz seminorms.
The same is true on the primed side for a fixed admissible finite
factor: the integer \(M\), the size of the finite factor and the
compact set used in reduction are fixed, and the Sobolev argument
uses finitely many seminorms of the varying archimedean factor.
For the primed approximations the individual finite factors in the
definition of matching need not themselves satisfy the support
condition. Let \(B\) be their finite-dimensional span and \(B^\circ\)
the subspace of the full finite Schwartz space spanned by admissible
finite tensors. For \(S=\mathcal S(\mathbb R^D)\) one has
\[
      (S\otimes B)\cap(S\otimes B^\circ)=S\otimes(B\cap B^\circ)
\]
inside the algebraic tensor product with all finite Schwartz tests,
as follows by extending a basis of \(B\cap B^\circ\) in each factor
space \(B,B^\circ\). Products of functions realize this tensor product
faithfully: evaluations span the dual of every finite-dimensional
space of functions. Express the approximants and their limit in a
basis of \(B\) adapted to \(B\cap B^\circ\). The coefficients outside
the intersection vanish and the others converge in \(S\), being fixed
finite linear combinations of the original coefficients. Continuity
on each resulting finite slice permits passage to the limit at \(h\).
\end{proof}

\begin{proposition}[Even source ranks at an inert place]\label{threal:inert}
Suppose \(v\) is inert and \(a=2u,\ a'=2w\). The two local
Schwartz--Bruhat spaces have a common irreducible \(U(W)(F_v)\)
quotient, invariant under their respective source groups. The maps
can be chosen to send the prescribed vacua to the same nonzero
spherical vector. Its unramified parameter, in the coordinates
\(q^{2\sigma_i}\) up to permutation and separate inversions, is
\[
      (q,q^3,\ldots,q^{2u-1})\ \sqcup\
      (q,q^3,\ldots,q^{2w-1}),\qquad q=|\mathcal O_F/v|.
\]
\end{proposition}
\begin{proof}
Put \(k=F_v,\ l=E_{\mathfrak w}\), and \(R=u+w\). This is an
unramified quadratic extension of odd residual characteristic, \(q>3\),
and \(\psi_v\) has conductor \(\mathcal O_k\). Since all ramification
of \(\alpha,\beta\) is split, their local components are unramified.
Both are trivial on \(k^\times\) by evenness of the ranks. A
uniformizer of \(k\) is also a uniformizer of \(l\), so both local
characters, and both splitting pairs, are trivial.

The two source lattices are perfect Hermitian lattices. The absolute
trace-dual condition on the target also gives a perfect skew-Hermitian
lattice here. Indeed the standard global additive character with the
specified infinite component has local annihilator of \(\mathcal O_k\)
equal to the absolute codifferent; this follows from the trace
pairing and the uniqueness of the character on \(\mathbb A_F/F\)
with that infinite component (the dual is \(F\)). The conductor
hypothesis makes this codifferent the unit ideal, and \(l/k\) adds
no different. A trace-zero unit converts a perfect skew-Hermitian
form into a Hermitian one. Such a unit is obtained by subtracting
the conjugate of a lift of a residue element outside the base field.
Every even-rank perfect Hermitian lattice here is integrally
hyperbolic. Over the residue fields diagonalize and use surjectivity
of the norm to obtain a nonzero isotropic vector in dimension at
least two. The derivative of the norm form there is the surjective
residue trace pairing with this vector. Hensel lifting gives a
primitive isotropic \(x\). Choose \(y\) pairing to one with \(x\)
and replace \(y\) by \(y-h(y,y)x/2\). The resulting hyperbolic plane
splits integrally, giving induction. In particular the target admits
the integral hyperbolic polarization and both source groups have the
quasi-split sign \(+\).

Use the ordinary Schr\"odinger model on \(Y^R\), for \(Y\) one of the
two Hermitian sources. The lattice vector is, up to scalar,
\(\mathbf1_{\Lambda_Y^R}\). To verify the identification with the
prescribed vacuum, apply its Heisenberg eigenline condition on each
of the two tensor lattice halves of the hyperbolic polarization of
\(W\). The quadratic phase is one on each half, which is isotropic
over \(l\) and stable under the trace-one element. The two conditions
are invariance under the translations by \(\Lambda_Y^R\) and the
modulations by its dual. They force support in \(\Lambda_Y^R\) and
constancy there. The change to these coordinates is a unitary
Heisenberg intertwiner and the original vacuum is nonzero.

We recall why the splitting, including its scalar character, agrees
with the one in the spherical calculation. The prescribed splitting
on each factor extends to its hyperbolic doubling. In the doubled
diagonal polarization evaluation at zero has on its Siegel stabilizer
the law \(\chi(\det t)|\det t|_l^{\dim_EY/2}\), with the oscillator
partner's dimension in the exponent if it is the source that is
doubled. This is precisely the standard doubled splitting with
character \(\chi\). To use the adelic law at a single place, take
evaluation in the diagonal polarization at all places. This is an
admissible functional: rational changes to that polarization are
compositions of the linear changes, phases and Fourier transforms
above, and evaluation is continuous on each fixed finite Schwartz
slice. The functional is nonzero. After unitary identification of the
doubled Heisenberg implementations, their quotient on elements
supported at \(v\) is a scalar character by the commutant argument.
The common evaluation law makes it trivial on the Siegel stabilizer.
Its conjugates generate the hyperbolic unitary group: the Siegel
Bruhat double cosets are represented by interchanging some of the
hyperbolic pairs. In each hyperbolic plane the interchange, up to
a diagonal element, is the product of three alternating upper and
lower root unipotents by the \(2\)-by-\(2\) identity in
\(\mathrm{SL}_2(k)\simeq SU(1,1)(k)\). Upper unipotents are in the
stabilizer and lower ones in its opposite conjugate. Thus the
character is trivial. Restrict on block elements \((g,1)\) of the
doubling and apply the defining product rule on pure tensors to
recover the given single-factor action. This argument works on both
factors. In particular it excludes a scalar character twist that
Heisenberg covariance alone would permit.
\src{src/Development/AbelianVarietyTheory/Definitions/Divisors/HilbertNinetyQuotientHomEqIff.lean}{435}{Doubled Siegel and block restriction laws}.

In Liu's notation take the larger skew-Hermitian group \(G_R\) and
the Hermitian space of dimension \(2s\), first \(s=u\), then \(s=w\).
His unramified additive character and trivial splitting character are
the ones just checked (a unit scaling in the trace pairing changes
only the choice of conductor-zero additive character or hyperbolic
basis). The trivial representation of \(U(2s)\) has parameters
\((q,q^3,\ldots,q^{2s-1})\), since the exponents of the normalized
Borel modulus on its \(l^\times\)-coordinates are
\(s-i+1/2\), for \(1\le i\le s\). In fact the root computation in
the preceding proof gives the modulus before taking its square root.
Liu's Theorem~6.2(2), Proposition~3.4 and the parameter maps of
\S2 say that, for \(R\ge s\) and sign \(+\), the small lift of the
spherical representation on \(G_R\) with the \(s\) parameters inverted
and \(q,q^3,\ldots,q^{2(R-s)-1}\) appended is the spherical
representation with those \(s\) parameters on \(U(2s)\).
The small-lift quotient sends \(\mathbf1_{\Lambda_Y^R}\) to a
nonzero vector, also explicitly in the proof of his Theorem~1.1.
The two applications give the displayed identical Weyl orbit.
Accordingly the oscillator has in each case a quotient onto
\(\pi\boxtimes{\bf1}\) with the same irreducible spherical \(\pi\)
on \(G_R\). The image of each vacuum belongs to the one-dimensional
hyperspecial fixed line; rescale one map to give equality. If maps
on the full space of all functions are requested, extend linearly
from Schwartz--Bruhat functions using any vector-space complement:
the conditions concern only Schwartz--Bruhat vectors.
For these choices of spherical vector, parameter and normalized
induction see \cite[\S\S2, 3 and 6]{liuThetaInert}.
\end{proof}

\begin{proposition}[All inert places, including the odd row ranks]
\label{threal:allinert}
At every inert place, in either parity, the two local Schwartz spaces
have a common irreducible target quotient annihilating the source
differences. The prescribed vacua have the same nonzero image after
one scalar normalization.
\end{proposition}
\begin{proof}
Write \(d=2R\), \(Y=V\) or \(V'\), \(m=\dim_EY\), and
\(\xi=\beta_v\) or \(\alpha_v\), respectively. These last two
characters of \(E_v^\times\) are the same: both are unramified
and their values at a uniformizer of \(F_v\), also a uniformizer
of \(E_v\), are \((-1)^m\). Thus \(\xi^2=1\). On the source the
splitting character is trivial in both cases. Use the integral
hyperbolic polarization of \(W\) in the preceding proof. In this
model \(U(Y)\) acts by linear substitution on \(Y^R\): in the
Siegel Levi implementing its action the determinant modulus is
one, with trivial splitting character. Evaluation at zero is
source invariant. Consequently
\[
        \phi\longmapsto
          \bigl(h\longmapsto(\omega(1,h)\phi)(0)\bigr)
\]
gives a nonzero equivariant map for the target into the normalized
Siegel induction \(I(s_m,\xi)\) of \(U(2R)(F_v)\) from
\((\xi\circ\det)|\det|_{E_v}^{s_m}\),
where \(s_m=(m-R)/2\). Indeed the Siegel unipotent acts by a
quadratic character taking value one at zero. The Levi acts at
zero by \(\xi(\det b)|\det b|_{E_v}^{m/2}\), and the half modulus
of this Siegel parabolic is \(|\det b|_{E_v}^{R/2}\).
The image is nonzero already on the lattice indicator at \(h=1\).
These substitutions are the standard Siegel splitting of the
polarized Weil representation. The doubled uniqueness argument
above identifies them with the actions specified here.

The image has finite length as a submodule of parabolic induction
of an admissible finite-length representation. Choose any
irreducible quotient of this nonzero image. The composite is a
nonzero map from the oscillator onto the trivial source times
this quotient. The unramified vacuum theorem says that such a map
takes the lattice vacuum to a nonzero vector in the hyperspecial
fixed line when the source representation is spherical. Its
hypotheses are precisely odd residue characteristic, an unramified
quadratic extension and additive character of conductor \(\mathcal
O_v\), self-dual lattices and unramified splitting characters. To
interchange the Hermitian and skew-Hermitian spaces one can use a
trace-zero unit and its inverse. The source vacuum in the given
coordinates is a nonzero multiple of this lattice vector by the
Heisenberg characterization above.
\cite[Appendix, Theorem~2.A.1]{harrisSpherical}

In particular this irreducible quotient has the spherical Hecke
character of \(I(s_m,\xi)\): hyperspecial invariants are exact by
compact averaging, and the induction has a single fixed line by
Iwasawa decomposition. For \(m'=2R-m\) one has \(s_{m'}=-s_m\).
The two inducing data are Weyl conjugate on their unramified
parameters: the long Siegel Weyl element takes \(b\) to its
inverse conjugate transpose, \(\xi\) is conjugation invariant
and self-inverse, and the \(R\) exponents of the trivial
\(\mathrm{GL}_R\) representation are symmetric about zero.
The Satake classification thus identifies the two spherical
irreducibles. Scaling the maps to their common fixed line gives
the desired normalization.
\end{proof}

\begin{proposition}[Ordinary relatively split places]
\label{threal:ordinarysplit}
At every place split in \(E/F\) outside \(S=S_0\sqcup Q\), there
is a common irreducible quotient with matching nonzero vacua,
including when the underlying rational prime does not split
completely in \(E\).
\end{proposition}
\begin{proof}
Use the two row and dual-row models of
Proposition~\ref{newtheta:split}. The relative splitting
\(E\otimes F_v=k\times k\) suffices for these models. In either
polarization the two idempotent summands of the tensor lattice give
translations by a full lattice and modulations by its character dual.
Indeed the source lattice is perfect for its \(E\)-valued form,
and the partner lattice is perfect for the absolute trace pairing.
Thus, in dual bases, the modulation coefficients range over the
inverse different of \(k/\mathbb Q_\ell\) times the ordinary dual
of the position lattice. This gives exactly the character dual for
the trace additive character. The latter is forced by its infinite
component and triviality on \(F\), since \(F\) is dense in its
finite adeles. On each idempotent summand separately the quadratic
phase in the vacuum condition is one by isotropy, also in residue
characteristic two. Translation invariance and modulation invariance
force any nonzero vacuum to be a constant times the indicator of
the position lattice. Choose position lattice bases and the
conjugate dual bases so its indicator is the integral-matrix test
on each side and the two target compact subgroups agree. Scaling
these local coordinate intertwiners identifies the given vacua
with the two indicators. The unitary linear action is the prescribed
one by the block and Siegel laws used above; conjugating to the dual
polarization sends the target action to inverse transpose as in
\eqref{newtheta:big}.

In the notation of Proposition~\ref{split:good} put \(n=a\),
\(t=\nu(\pi)\), \(s=d/2\) and
\(D_n(z)=\prod_{j=0}^{n-1}(1-tq^{j-z})\).
For any matrix test \(H\), the cleared determinant integral
\(D_n(z)Z_n(H,z)\) is a Laurent polynomial in \(q^{-z}\).
Here is the refinement needed when \(t\) has finite order. Keep the
integrand exponent \(z-n\) fixed as in that proposition. In size
\(j\) a ball about zero scales to \(M_j(\mathcal O_v)\). Its
denominator divides
\(\prod_{i=0}^{j-1}(1-tq^{n-j+i-z})\) by the same Hermite
count with parameter \(z-n+j\) and the conversion from additive
to group Haar measure. Its numerator is a constant times a
Laurent monomial. Subtract the value at zero times such a ball.
The remaining compact support is covered by pivot charts; the
Schur-complement argument reduces to size \(j-1\), whose denominator
is a subproduct of the one just displayed. Finite disjoint
subdivision and induction prove the claim starting with size zero.

Multiply the larger transform by \(D_n(z)\) and evaluate at \(s\).
The source and parabolic covariance pass from the right half-plane
to this value by the meromorphic identity principle, including if
the value is a residue. The character and modulus factors in that
covariance are holomorphic in \(z\). Smoothness in the target
follows from the compact open stabilizer of the test. The larger
vacuum has value one at the identity by the exact Hermite count,
with \(\operatorname{vol}(\mathrm{GL}_n(\mathcal O_v))=1\).
The smaller vacuum also has nonzero value there, since \(a'<s\),
and can be normalized to one. The two spherical vectors are now
equal. The common induced target of Proposition~\ref{newtheta:split}
is irreducible by Bernstein's theorem, so both nonzero maps are
onto. The stronger assertion on the entire source-difference
kernel at totally split good primes was proved separately in
Proposition~\ref{split:good}.
\end{proof}

\begin{proposition}[The real quotient for \(p=q=0\)]\label{threal:compact}
At each infinite place write \(A=m_\alpha,\ B=m_\beta\) and
\[
        j=(A-B-d)/2,\qquad l=(B+a)/2 .
\]
The two real polynomial modules have the same source-coinvariant
representation of the compact target \(U(d)\), with dominant highest
weight \(((l+j)^a,l^{a'})\). Superscripts indicate repetitions.
The tensor product of these representations over the real places
realizes the real matching maps for this branch.
\end{proposition}
\begin{proof}
The restrictions to \(\mathbb R^\times\) give \(A\equiv a'\),
\(B\equiv a\pmod2\). Thus \(j,l\) are integers, and \(j>0\) by the
weight condition. Let \(D_0=\mathbb C^d\), \(U_0=\mathbb C^a\),
\(U_1=\mathbb C^{a'}\) with the standard actions. In orthonormal
column frames the operators on the polynomial Fock modules are the
following two representations of source times target:
\[
\begin{split}
 &(\det_{U_0}^{-j}\boxtimes\det_{D_0}^{\,l})
               \otimes\operatorname{Sym}(U_0\otimes D_0),\\
 &(\det_{U_1}^{\,j}\boxtimes\det_{D_0}^{\,l+j})
               \otimes\operatorname{Sym}(U_1^*\otimes D_0^*).
\end{split}
\]
We verify the weights to fix the sign convention. In one unprimed
complex matrix entry the tensor form is
\(-\operatorname{Tr}_{\mathbb C/\mathbb R}(-i z\bar z')
=2\,\operatorname{Im}(\bar z z')\); it changes sign on the primed
side. The additive character is \(x\mapsto e^{-2\pi i x}\). In the
positive plane the Heisenberg derivatives for \(1,i\) can be taken
as \(\partial_x,-4\pi i x\). Acting on
\(\exp(-2\pi x^2)\), the second derivative operator is \(i\) times
the first. Thus a rotation acts on the creation operator
\(d\rho(1)-i\,d\rho(i)\) by its ordinary character. Reversing the
form gives the conjugate character. The infinitesimal metaplectic
rotation on the positive plane has eigenvalues \(i(n+1/2)\):
in the creation and annihilation operators \(b^\dagger,b\) with
\([b,b^\dagger]=1\) its generator is the symmetrized quadratic
\(i(b^\dagger b+bb^\dagger)/2\); the negative plane reverses this
sign. Therefore the metaplectic vacuum has half the sum of the signed
entry weights.

The doubled Siegel normalization contributes the additional half of
the splitting-character exponent on the determinant of each single
unitary factor. Indeed on the doubled Siegel Levi \(t\in
\mathrm{GL}_n(\mathbb C)\) the action of the metaplectic Lie algebra
on evaluation gives the positive modulus; its infinitesimal phase
is zero, as one also sees from the real change of coordinates.
The prescribed additional phase on this Levi is
\((\det t/|\det t|)^m\). The determinant of the doubled unitary
element on the same Levi is \(\det t/\overline{\det t}\); hence the
difference of the two infinitesimal implementations is
\((m/2)\operatorname{tr}\). They differ by a Lie algebra character,
and every such character on \(\mathfrak u(n,n)\) is a multiple of
trace since its derived algebra is \(\mathfrak{su}(n,n)\).
Restriction to the first block preserves the trace. Consequently the
source vacuum powers are \((d+B-A)/2=-j\) and
\((-d+A-B)/2=j\); the target powers are \((a+B)/2=l\) and
\((-a'+A)/2=l+j\). The unitary groups are connected, so the
infinitesimal action determines these integrated representations.
The Gaussian times the Hermite polynomials form the polynomial
module here: the annihilation equations for the decaying Gaussian
separate into the one-coordinate first-order equations above, with
the opposite choice of sign giving an increasing exponential, and
the Heisenberg derivatives generate all the creation monomials.

Use the Cauchy decomposition
\(\operatorname{Sym}(U\otimes D_0)=
\bigoplus_\lambda S_\lambda U\otimes S_\lambda D_0\);
this follows by applying Schur--Weyl duality to each homogeneous
degree and taking simultaneous permutation invariants.
The only partition for which \(S_\lambda U_0=\det_{U_0}^j\) is the
rectangle \((j^a)\), with multiplicity one. On the dual side the
rectangle is \((j^{a'})\). Thus the first invariant space is
\(\det^l S_{(j^a)}D_0\), and the second is
\(\det^{l+j} S_{(j^{a'})} D_0^*\). Their highest weights agree,
since duality reverses and negates the list of weights. Both spaces
are irreducible. Averaging over the source compact group projects
onto the invariants and kills precisely the span of the differences
on the polynomial module, by complete reducibility on each finite
degree. Composing with an identification of the two highest-weight
modules gives the claimed surjections. Their tensor product is
irreducible for the product of the compact targets. The maps may
be extended linearly off the polynomial spaces, since the
specified infinitesimal and compact transformation laws concern
polynomial inputs and these actions preserve that subspace.
For comparison with the source conventions see
\src{src/Development/WeilBound/ThetaLifts/LocalTheory/ToralCohomology/PiInfUnitNe.lean}{116}{The unprimed target highest weight}.
\end{proof}

\begin{proposition}[The real quotients for all signatures]\label{threal:allreal}
At each infinite place let \(A=m_\alpha,\ B=m_\beta\).
In Paul--Trapa's normalization the common real representation
is the cohomologically induced module on the ordered Levi
\[
 U(a-p,p)\times U(a'-q,q),\qquad
 A_{\mathfrak q}(\det^{A/2}\boxtimes\det^{B/2}).
\]
The half-determinant notation uses the covers in their normalization.
The output after twisting is an honest module for the target Lie
algebra and its maximal compact group.
\end{proposition}
\begin{proof}
Put \(r=d-p-q,\ t=p+q,\ k=A-B>d\). Work first with the
metaplectic oscillator, before the two splitting twists. With the
additive character and tensor form in Proposition~\ref{threal:compact},
a positive source row has vacuum determinant exponent \((r-t)/2\);
a negative source row has the opposite exponent. This follows by
summing the signed one-entry weights \(1/2\) there. It is the
convention of Paul--Trapa's joint harmonics, their formula (3.3).
In passing to the prescribed actions the source determinant exponents
are shifted by \((B-A)/2\) and \((A-B)/2\), respectively.
The doubled trace argument for that shift in
Proposition~\ref{threal:compact} applies to each signature, since
the doubled group and its Siegel Levi have the same trace character.
Thus the prescribed trivial source corresponds in the metaplectic
oscillators to \(\det^{k/2}\) and \(\det^{-k/2}\).

Since \(|r-t|<k\), subtracting the vacuum weights from
\((k/2,\ldots,k/2)\) gives positive entries on both compact
blocks of \(U(a-p,p)\). In the notation of Paul--Trapa's Lemma~3.2
the two joint-harmonic bounds are \(a-p\le r,\ p\le t\).
For the primed determinant all the entries after subtraction
are negative. The bounds for \(U(q,a'-q)\) become
\(a'-q\le r,\ q\le t\). They hold since \(r=(a-p)+(a'-q)\)
and \(t=p+q\). Thus each character occurs in the joint harmonics.
Part (1) of their Theorem~4.1 gives on the unprimed side
\(A_{\mathfrak q}(\det^{k/2}\boxtimes 1)\) with first block
\(U(a-p,p)\). Part (2) gives
\(A_{\mathfrak q}(1\boxtimes\det^{-k/2})\) with second block
\(U(a'-q,q)\) on the primed side. These are the same ordered
two-block parabolic before assigning its two inducing characters.
Finally the target splitting adds \(B/2\) to both block weights
on the first side and \(A/2\) on the second. Tensoring by this
determinant character commutes with cohomological induction, by
pulling the character through the induction and derived compact
invariant functors. Both modules therefore have the displayed
description. The shifts inherent in the authors' normalization
of \(A_{\mathfrak q}\) are identical on the two sides.
\cite[\S\S2.1, 2.3, 2.8, Lemma~3.2, Theorem~4.1 and Corollary~4.3]{paulTrapaHowe}

These are small theta quotients of the polynomial Fock modules,
hence irreducible nonzero targets of equivariant surjections.
Taking their algebraic tensor products over the real places gives
the quotient for the product Lie algebra and maximal compact.
Indeed the source polynomial module is spanned by finite products
of creation monomials from the separate places. A submodule of the
external product of simple Harish-Chandra modules is either zero
or the whole product, by the density theorem applied to finite
tuples in each factor (their commuting endomorphisms are scalars).
The quotient maps kill source derivatives and compact differences;
target derivatives and compact actions preserve matching. Extending
the maps linearly from polynomial vectors to their ambient
Schwartz spaces satisfies the stated real-map conditions, all of
which test polynomial inputs. For simultaneous archimedean
one-parameter groups one differentiates the finite product of
the individual actions, giving the sum of the placewise derivatives.
\end{proof}

\begin{theorem}[A nonzero supported undoubled period]\label{threal:nonzero}
For \(q=0\) and every allowed \(p\) the primed period is nonzero for some polynomial
Schwartz test satisfying the imposed full-rank condition at \(v_1\).
Both parities of \(a'\) are allowed.
\end{theorem}
\begin{proof}
Put \(Y=V'\), \(m=a'\), \(H=U(Y)\), \(J=U(W)\). The source quotient
is compact by the rank-zero case of reduction theory. Take probability
quotient measures for now; other Haar normalizations multiply the
periods by positive constants. We first prove the identity
\begin{equation}\label{threal:norm}
 \int_{[J]}\left|\int_{[H]}\Theta_\phi(g,h)\,dg\right|^2 dh
  = c_H\int_{H(\mathbb A_F)}
                \langle\omega(g,1)\phi,\phi\rangle\,dg,
                 \qquad 0<c_H<\infty .
\end{equation}
Here the scalar simply accommodates the choice of group Haar measure
on the right; it is one for quotient volume one with rational counting
measure. We use the usual self-dual Hilbert pairing, linear in the
first variable. For the proof it suffices to take finite pure tensors
with polynomial Fock factors at infinity.

On the doubled space \(Y\oplus(-Y)\), change from the product
polarization for \(\phi\otimes\bar\phi\) to the diagonal polarization.
The value at zero in the latter coordinates is
\(\int\phi(x)\overline{\phi(x)}\,dx\): the change is a difference and
average of the two position coordinates followed by Fourier transform
in the average variable. The Jacobian scalars can be fixed using
rational coordinates and globally self-dual measures, by the product
formula. At \((g_1,g_2)\) the analogous functional gives
\(\langle\omega(g_1,1)\phi,\omega(g_2,1)\phi\rangle\) times the
second-block determinant character of the doubled splitting. In the
product with the conjugate kernel that character is divided out.
Its exact conventional form is \(\chi_W(\det g_2)\) for the splitting
character on the \(H\)-action in the notation of the doubling
construction. Its modulus is one and it is trivial on rational points.
This block identity is also the change-of-model identity in
Gan--Qiu--Takeda, \S11.3, equations preceding (11.2)~\cite{ganQiuTakeda}.
The rational summation on the two product factors changes to the
doubled rational summation by Poisson summation.

For the given presentations, transport the standard implementations
using their rational Darboux coordinates. On each doubled factor the
possible discrepancy with the prescribed implementation on elements
supported at a single place is a scalar character trivial on the
Siegel stabilizer, by the argument of
Proposition~\ref{threal:inert}. This argument applies over the
real fields as well, using hyperbolic planes. At split places the
doubled group is a general linear group, generated by the block
parabolic and its conjugates by elementary elimination. Hence every
single-place discrepancy is trivial on each single factor by the
block rule. The global discrepancy is a continuous character:
compare the two implementations with a nonzero matrix coefficient
in a neighborhood of the identity, using unitarity and weak
continuity. It is trivial on the dense subgroup of elements supported
on finitely many places, hence everywhere. As a result the standard
automorphic theta and these periods agree after the rational change
of coordinates, up to the single scalar inherent in choosing the
unitary coordinate transform. Taking its modulus to be one suffices
in \eqref{threal:norm}.

To interchange the three integrations, use the chamber estimate of
Proposition~\ref{threal:decay} with \(W\) as integrating space of
dimension \(d\) and \(Y\) as partner of dimension \(m\). Its proof
applies to a skew-Hermitian integrating space as well. If a
Hermitian integrating space is used in the statement of the
Siegel--Weil formula, multiply its given skew-Hermitian form by
any \(0\ne\delta\in E\) with \(\bar\delta=-\delta\) and multiply
the other form by \(\delta^{-1}\). The product of the forms, the
underlying isometry groups, the flags and their moduli are preserved
in these conventions.
For \(g\) in compact source representatives the single kernel is
\(O(\prod T_i^{m/2})\), uniformly in \(g\). The squared bound is
integrable since the \(F\)-Witt index of \(W\) is at most \(p\)
and \(d-2p+1>m\). Now in the doubled Siegel--Weil formula integrate first over \([J]\),
with integrating space of dimension \(d\) and doubled group \(U(2m)\).
The parameters are \(s_0=(d-m)/2>m/2\), as \(a>a'\).
The coefficient of the probability theta integral is one by Weil's
convergent-range formula, applied with \(d>2m\); in the anisotropic
case this is also Theorem~7.1 of Gan--Qiu--Takeda. The Siegel Eisenstein
series with normalized induction is absolutely and locally uniformly
convergent in \(\Re s>m/2\); absolute convergence holds uniformly
on compact representatives of \([H]^2\) by majorization by the
positive spherical section at the real parameter. Thus integration
over \([H]^2\) permits absolute unfolding.
We use the probability normalization for the theta integral
\cite[introduction, Theorem~2 and \S3]{yamanaUnitary}.

Every maximal rational isotropic subspace of \(Y\oplus(-Y)\) projects
injectively onto each copy of \(Y\) since \(Y\) is anisotropic. Its
dimension is \(m\), so it is the graph of a rational isometry.
The set indexing the Siegel series is therefore one
\((H\times H)(F)\)-orbit with stabilizer \(\operatorname{diag} H(F)\).
After unfolding we integrate the matrix coefficient above, with its
second-block character cancelled, over
\(\operatorname{diag}H(F)\backslash H(\mathbb A_F)^2\).
The change \(g=g_2^{-1}g_1\) gives \eqref{threal:norm} and proves
absolute convergence of its right-hand integral. The same calculation
with any finite replacements of the test factors is valid.

We choose the factors so that the last integral is nonzero. At infinity
the primed source is negative definite. In the entry calculation of
Proposition~\ref{threal:compact} there are \(d-p\) negative planes
and \(p\) positive planes in each source row. Its vacuum determinant
power as a source representation is
\((A-B-d+2p)/2=j+p>0\). The \(d-p\) columns from the negative
planes are copies of the dual standard source representation.
Choose \(m\) of those columns, possible since \(d-p\ge m\), and
take their determinant to the power \(j+p\) times the vacuum.
This is a nonzero source-invariant polynomial vector. Its local matrix coefficient on
the compact source is the positive constant \(\|\phi_v\|^2\).
At a designated finite split place the conjugate polarization has
positions \(X\in M_{m,d}(F_v)\) and the source acts by left row
operations times a scalar character and the square root of the
modulus. Take the indicator of a sufficiently small compact open
ball about \((I_m\mid0)\). If \(g\) moves this ball to intersect
itself, the invertible leading block implies
\(g\in 1+\pi^b M_m(\mathcal O_v)\) for a prescribed \(b\);
this follows by writing \(g=X_1X_2^{-1}\) on the first blocks, or
its inverse if the action is inverse. Choose \(b\) so the character
is trivial there. The coefficient is nonnegative, vanishes off a
compact set and is positive on a neighborhood of \(1\).
This test satisfies the full-rank condition at \(v_1\).

At any specified inert place \(W\) has an integral hyperbolic
polarization by the lattice argument above, irrespective of the
parity of \(m\). The source acts on \(Y^{d/2}\) by linear changes
on all vectors times a scalar character. Now \(d/2>m\). Pick a small
ball about a tuple whose first \(m\) vectors are a basis of \(Y\).
The same leading-block argument inside \(\mathrm{GL}_m(E_{\mathfrak w})\)
forces overlap only in any prescribed neighborhood of \(1\) in
\(H(F_v)\). The scalar on this compact subgroup is trivial after
shrinking: a continuous character of a profinite group has open
kernel. The corresponding local integral again is strictly positive.

It remains to ensure that the unrestricted product of local
integrals is nonzero. Start with the factors just chosen at infinity
and \(v_1\), and the usual unit-norm lattice indicators at almost
all other places. There the hyperspecial group \(K_v\) fixes the
indicator: in the split linear model the unramified character is
trivial on units; in the inert hyperbolic polarization the source
acts by pullback with the unramified splitting. More explicitly the
character on the source for this polarization is trivial when the
partner is hyperbolic and its splitting character is trivial, as
is the case for the source splitting \(\alpha/\beta\) at inert places:
the two unramified characters have equal value on a uniformizer
by their common parity. The doubled Schr\"odinger model makes this
pullback implementation explicit, and the doubled uniqueness already
proved identifies it with the prescribed one. Take the local Haar
measures to give \(K_v\) mass one at almost all places. Write
\(b_v(g)=\langle\omega_v(g,1)\phi_v,\phi_v\rangle\) and
\(B_v=\int|b_v|\). Outside a finite set \(S\) one has \(b_v|_{K_v}=1\)
and \(B_v\ge1\). The integral of \(\prod_v|b_v|\) is finite by
the absolute unfolding, and is positive on some cylinder about
the identity. Exhaustion by finite restricted products and Tonelli
give \(\prod_{v\notin S} B_v<\infty\), hence
\(\sum_{v\notin S}(B_v-1)<\infty\). Thus
\[
        \left|\int_{H(F_v)} b_v(g)\,dg-1\right|
              \le B_v-1
\]
at these places. Only finitely many local integrals can vanish;
replace those and any exceptional finite factors by the compact
open tests just constructed. Absolute factorization now gives a
nonzero product. Formula~\eqref{threal:norm} yields a nonzero
undoubled primed period. The finitely modified tensor is a valid
supported test since the prescribed vacuum is standard almost
everywhere.
\end{proof}

\begin{theorem}[The scalar at the definite base]\label{threal:base}
Suppose \(p=q=0\). Both undoubled periods are nonzero, the primed
one on the required supported subspace. Common matching quotient
maps exist. For every choice of such maps there is a unique
\(c\in\mathbb C^\times\) such that \(I_\Phi=cI'_{\Phi'}\) for
all matched tests. Both parities are allowed. The inert and ordinary
split maps are furnished by Propositions~\ref{threal:allinert} and
\ref{threal:ordinarysplit}; the specified orbital maps at \(S\)
and the stronger kernel and support conditions at the totally split
good primes are those of Corollary~\ref{split:ram} and
Proposition~\ref{split:good} in the row polarizations. The
identification of those polarizations and their vacua is the one
in Proposition~\ref{threal:ordinarysplit}; at \(S\) use the
specified lattice bases. Proposition~\ref{threal:compact} gives
the archimedean maps.
\end{theorem}
\begin{proof}
First establish unprimed nonvanishing. Write \(n=a\),
\(H=U(V)\), \(J=U(W)\), \(\chi=\beta/\alpha=\nu^{-1}\),
\(T=S=S_0\sqcup Q\) the exact set above which \(\nu\) is ramified,
and \(s_0=(d-n)/2=a'/2\). Both adelic quotients are compact.
Form the Siegel section of \(U(V\oplus -V)\) through the
doubled Schwartz vector \(\phi\otimes\bar\phi\) at \(s_0\),
and extend it flat in normalized induction using the compact
subgroups specified below. Pair the resulting Eisenstein series
on the two copies of \([H]\) against the constants, cancelling the
second-block determinant character as above. Write this scalar
as \(Z(s)\). The graph orbit and the change of variables in
Theorem~\ref{threal:nonzero} express \(Z(s)\), for \(\Re s\)
large, as a positive measure constant times the product of the
local integrals of \(f_{v,s}(g,1)\) over \(H(F_v)\). Constants
are fixed by giving the hyperspecials volume one outside \(T\)
and choosing product Haar measure up to one global positive
factor. The meromorphically continued Eisenstein series can be
paired over the compact quotient: its meromorphic continuation
as smooth automorphic forms is locally uniform on compact sets
after multiplication by a local scalar denominator. At \(s_0\),
the anisotropic Siegel--Weil formula with integrating space \(W\)
and \(d>n\) gives \(Z(s_0)=C\|I_\phi\|^2\), \(C>0\), by Fubini.
No initial absolute convergence at \(s_0\) is used.
\cite[Theorem~7.1 and \S11.3]{ganQiuTakeda}

At each infinite place take a nonzero source-invariant polynomial
from Proposition~\ref{threal:compact}. Choose the maximal compact
in the doubled group containing \(H(F_v)^2\); hence
\(f_{v,s}(g,1)=\|\phi_v\|^2\) there. The doubled section is
compact-finite, since the doubled maximal compact at this place
is the product of the two compact groups. At \(v\in T\), which is
split because \(T\subset P\), take the small-ball test about a
full row frame in \(M_{n,d}(F_v)\) constructed above. On
\((g,1)\) the flat section equals the compactly supported matrix
coefficient at \(s_0\) times a positive Iwasawa height raised to
\(s-s_0\). Its integral is entire and strictly positive at \(s_0\).

At every other finite place the doubled section can be taken to
be the basic spherical section of value one at the identity.
At split places use the characteristic function of integral
matrices in the linear model. On the doubled row space the two
blocks together give the integral matrices with \(2n\) rows.
The determinant character is unramified and the integral general
linear group fixes their indicator. At inert places the extension
and characters are unramified, the additive character has conductor
\(\mathcal O_v\), and the residue characteristic is odd. The
lattices are perfect, \(W\) has even dimension and is integrally
hyperbolic, as proved above. In its polarization the doubled row
group acts linearly on \((V\oplus -V)^{d/2}\) with unramified
splitting, and the product of the two lattice indicators is
fixed by the stabilizer of the doubled self-dual lattice. This
works also for \(n\) odd. Indeed an integral hyperbolic plane
splits off repeatedly from a perfect Hermitian lattice of rank at
least two over an unramified quadratic extension of odd residue
characteristic by reduction and Hensel lifting, leaving a unit
line in odd rank. Thus the original row group too is unramified.
Normalize the indicators to have \(L^2\)-norm one. Evaluation on
the doubled diagonal then has value one. Transforming back to the
given presentation uses only finitely many departures from its
prescribed vacuum, exactly as above.

Let \(\eta=\eta_{E/F}\), and use a superscript \(T\) to omit finite
Euler factors over \(T\). The unramified doubling computation
with trivial representation of \(U(n)\) gives for the product of
the factors just described outside \(T\) and infinity
\[
 \frac{\displaystyle\prod_{i=0}^{n-1}
                 L^{T}_{E,\mathrm{fin}}(s+n/2-i,\chi)}
      {\displaystyle\prod_{i=0}^{n-1}
                 L^{T}_{F,\mathrm{fin}}(2s+n-i,\eta^{\,i})}.
\]
Here the base change of the spherical trivial representation has
the \(n\) exponents \((n-1)/2-i\): it is the spherical
constituent with trivial representation's Satake parameter,
the half sum of positive roots, restricted to the split group
over \(E\). In the inert odd case this list includes zero. The
doubling \(L\)-function uses \(s+1/2\), and \(\chi|_{\mathbb A_F^\times}=1\)
because the two ranks have equal parity. These give exactly the
shifts and the denominator above. See the unramified formula
\cite[\S4.2.1, (53)]{ehlsDoubling}. The product equals the
meromorphic continuation from its region of absolute convergence.
In particular no place with an unramified numerator has been
removed arbitrarily: at every prime of \(E\) over \(T\) the
character is ramified by the hypotheses, so its local factor is one.

At \(s_0\) the denominator arguments are at least \(a'+1>1\).
The numerator arguments are integers \(k=d/2-i\).
At each infinite place the angular exponent \(m\) of \(\chi\)
has \(|m|>d\); the gamma factors
\(\Gamma_{\mathbb C}(k+|m|/2)\) are finite and nonzero throughout
this interval. The Euler product gives nonvanishing when \(k>1\).
At \(1\), one elementary proof uses that both \(\chi\) and
\(\chi^2\) have nonzero infinity type. They have no Hecke pole.
For \(u>1\),
\(\zeta_E(u)^3|L_{E,\mathrm{fin}}(u,\chi)|^4
 |L_{E,\mathrm{fin}}(u,\chi^2)|\ge1\) off the ramified factors,
by \(3+4\cos t+\cos(2t)\ge0\) in the logarithmic Euler products.
The finite omitted factors of the three functions are bounded
above and below near \(1\). A zero would have order at least four
in this product, while \(\zeta_E^3\) has pole order three. For
\(k\le0\), apply the Hecke functional equation to the completed
function and use nonvanishing at \(1-k\) for \(\chi^{-1}\);
the conductor power and root number are nonzero
\cite{tateThesis}. This proves
\(Z(s_0)\ne0\). The primed supported period is already nonzero
by Theorem~\ref{threal:nonzero}.

We next compare the actual automorphic copies. On the definite
source the primed integral exists also for tests without the
support restriction. Use that extension in the following argument.
All three quotients are compact, with compact groups at infinity.
For any specified finite open compact level and infinite compact
type the space of automorphic vectors is finite-dimensional:
there are finitely many finite-adelic double cosets at that
level, and matrix coefficients of the prescribed compact type
span a finite-dimensional space on each. Consequently the unitary
spectrum is discrete with finite multiplicities. More explicitly
the orthogonal projection to any such level and type has finite
rank; these projections exhaust a dense subspace, and a continuous
direct integral of inequivalent unitary representations would have
infinite rank on any projection on which it is nonzero.

Let \(P_v,P'_v\) be the given local quotient maps and \(\pi\)
the restricted tensor product of the common irreducibles, using
their normalized vacuum images. At infinity take the entire
product of the compact factors; the polynomial calculation of
Proposition~\ref{threal:compact} gives its source-invariant
quotient with multiplicity one. For each of the two periods,
project to any irreducible automorphic summand of the target.
A nonzero projection, with all but one test factor fixed and
the other target factors paired with linear functionals, gives
an intertwiner from the local oscillator to the source trivial
representation times the local target constituent. At finite
places local Howe duality gives a unique such irreducible quotient
and a one-dimensional space of maps to it. This applies to
unitary dual pairs in characteristic zero and to the split
general linear dual pair; twisting an oscillator splitting by
characters preserves the assertion
\cite{ganTakedaHowe,minguezHowe}. At infinity the identical
conclusion follows from the compact calculation. Hence every
nonzero projection takes values in a copy of \(\pi\) and
annihilates each local kernel. Pairing is legitimate on algebraic
tensors: a smooth finite-level, compact-finite vector of an
irreducible automorphic summand belongs to its algebraic restricted
tensor product, and the pairings on the remaining factors separate
these vectors. At each finite tensor stage the kernel of the tensor
quotient is the sum of the subspaces having a kernel in one
factor, by exactness of tensor products over \(\mathbb C\).
This proves factorization of both periods. Smoothness then gives
pointwise equality whenever two resulting automorphic vectors are
equal in \(L^2\).

In fact \(\pi\) has multiplicity one here. Let \(i_0\) be the
unitary embedding defined by the nonzero primed image, using on
\(\pi\) the norm transported from that image. If another copy
exists, choose a unitary embedding \(i_1\) with image orthogonal
to \(i_0\). For any polynomial primed test \(u\) and compact-finite
smooth \(x\in\pi\), form the reverse lift
\[
  T_u(i_jx)(g)=\int_{[J]}(i_jx)(h)\,
                \overline{\Theta'_u(g,h)}\,dh
         \qquad (g\in U(V')(\mathbb A_F)).
\]
This function is constant in \(g\). In fact pairing against any
irreducible summand of the compact-quotient source spectrum and
varying \(x,u\) gives the same local intertwiners as above,
now with target constituent \(\pi_v\) fixed. Injectivity of
local Howe correspondence forces the source constituent to
be trivial at every place. At infinity the equivariant Hom in
the reverse direction also has this property: the Cauchy
decomposition in the compact calculation pairs distinct partitions
on the two compact groups. The globally trivial representation
occurs precisely on constant functions. Therefore \(T_u(i_1x)=0\):
its average on \([U(V')]\) is a positive normalization times
\(\langle i_1x,I'_u\rangle=0\).

On the other hand
\[
                 \|T_u(i_0x)\|_2^2=\|T_u(i_1x)\|_2^2.
\]
To see this form the flat doubled section \(f_{u,s}\) on
\(U(W\oplus -W)\) from \(u\otimes\bar u\), with parameter
\(s_*=(a'-d)/2=-a/2\). Pair its Eisenstein series, after the
second-block determinant correction, over \([J]^2\) against the
two copies of the vector \(i_jx\) and its conjugate. For \(\Re s\)
large there is just the graph orbit, since \(W\) is anisotropic.
Unfolding gives the integral over \(J(\mathbb A_F)\) of the
restricted section against the matrix coefficient of the right
regular action on \(i_jx\). In the convention with the section at
\((1,g)\) its scalar weight is \(f_{u,s}(1,g)\chi_*(\det g)^{-1}\),
where \(\chi_*=\alpha\) is the primed target splitting character.
The matrix coefficient is \(\langle R(g)i_jx,i_jx\rangle\), or
its conjugate upon conjugating the whole pairing. Either is
independent of \(j\) by isometry and equivariance. The two
meromorphic integrals thus agree identically. At \(s_*\) the
anisotropic Siegel--Weil formula is holomorphic, and its coefficient
is two for the probability theta integral because \(a'<d\).
The doubled product identity and Fubini over the compact quotients
give the two asserted squared norms times the same positive constant
\cite[Theorem~7.1 and \S11.3]{ganQiuTakeda}. Choose \(u\) with
\(I'_u\ne0\) and \(i_0x=I'_u\); then the average of
\(T_u(i_0x)\) is a positive multiple of \(\|I'_u\|^2\).
This contradicts \(T_u(i_1x)=0\). The two nonzero maps from the
common quotient now lie in the same one-dimensional intertwining
space by Schur's lemma. Restriction back to supported primed tests
and Corollary~\ref{threal:limits} give the required nonzero scalar
and its uniqueness.
\end{proof}

\begin{proposition}[The two periods against the cuspidal spectrum]\label{threal:orthA}
The unprimed period pairs to zero with every irreducible cuspidal
representation of \(J=U(W)\) whenever \(p>0\), for arbitrary \(q\).
If \(q=0<p\), the primed period also has this property, and on
removing a target hyperbolic plane the zero-vector orbit of the
unprimed line constant term vanishes identically.
\end{proposition}
\begin{proof}
Choose a rational prime completely split in \(E\) outside the finitely
many exceptional primes and the prime under \(v_1\), using the
splitting-prime theorem; let \(v\)
be a place above it. Write \(k=F_v\). By Lemma~\ref{split:ray}, the
unramified characters \(\alpha,\beta\) there have distinct values
on a uniformizer. The source-invariant quotient of the unprimed local
oscillator is the entire irreducible normalized induction
\[
       \pi_v=(\alpha\circ\det_a)\times(\beta\circ\det_{a'})
\]
by Proposition~\ref{split:good}. Thus every cuspidal constituent
pairing nontrivially with the unprimed period must have this local
factor. In detail the pairing, with other Schwartz coordinates and
other cuspidal coordinates fixed, gives an equivariant map into
the smooth dual of the conjugate of the local cusp representation.
Smoothness follows from invariance of the test under an open
compact subgroup, and the smooth dual in question is the original
representation by unitarity. Source invariance follows by changing
variables in the source integral.

For \(p>0\) the source has a rational hyperbolic plane: its form is
isotropic at every real place, it is split at the split finite places,
and at the inert places its perfect lattice has Witt index \(\lfloor
a/2\rfloor\). For the underlying \(F\)-quadratic norm form local
isotropy therefore gives global isotropy by Hasse--Minkowski, and
one completes the isotropic vector to a hyperbolic pair.
Write \(\Pi\) for such a cuspidal representation. Its adjoint lift
back to \(U(V)\), with conjugated oscillator, has a vector of
nonzero average, by Fubini. Indeed the absolute source theta sum in
Proposition~\ref{threal:decay}, integrated in \(g\), is bounded by
a power of \(H(h)\): in its unprimed estimate retain the factors
\(O(M^D H(h)^{A_B})\) before integrating the strictly negative
torus exponents. A cuspidal automorphic vector of finite compact
type decays faster than any power on a target Siegel set.
These estimates justify both the pairing and the interchange.

At any earlier source in the Witt tower, of dimension \(a-2b\)
with \(b\ge1\), no irreducible local representation can lift to
\(\pi_v\) in the forward oscillator. Indeed M\'inguez's split
correspondence adds the characters
\[
   \beta|\cdot|^{(1-l)/2},\
   \beta|\cdot|^{(3-l)/2},\ldots,\beta|\cdot|^{(l-1)/2},
               \qquad l=d-(a-2b)=a'+2b ,
\]
to the cuspidal support of the transformed source parameter.
The bare matrix oscillator adds the segment of the trivial
\(\mathrm{GL}_l\) representation. The scalar in the target row
action in \eqref{newtheta:big} is \(\beta(\det h)\) after removing
the unitary change-of-variables multiplier. The inverse transpose
on matrices interchanges the segment with its reversed dual, which
has the same multiset of real powers. This gives exactly the
displayed twist. The actual support of \(\pi_v\) has only \(a'\)
entries with unitary part \(\beta\). The other \(a\) entries have
unitary part \(\alpha\). This excludes the displayed segment.
\cite[Corollary~6.3(2)]{minguezHowe}
The first nonzero global adjoint lift of a cusp representation in a
Witt tower is cuspidal, by the tower constant-term formula
\cite[Proposition~3.1]{wuThetaTower}. If a lower lift were nonzero,
the first such lift would therefore have an irreducible cuspidal
constituent. Pairing against it and conjugating gives the excluded
forward local intertwiner. Thus the lift on \(U(V)\) is itself
cuspidal. The integral of a cusp representation against the constant
function on this positive-rank group is zero: a nonzero invariant
functional on a constituent would make it the trivial representation,
whose constant term along the given rational parabolic is nonzero.
This contradicts the preceding nonzero average.

For the primed period suppose \(q=0<p\) and let \(Y=V'\). A target
plane exists by the same norm-form argument, after multiplying its
skew form by a nonzero trace-zero element of \(E\).
The group \(U(Y)\) is anisotropic.
If the pairing with a cusp representation \(\Pi\) is nonzero, its
adjoint lift on \(U(Y)\) has nonzero projection to constants.
The entire adjoint lift is cuspidal by anisotropy, and is irreducible
by \cite[Theorem~5.3]{wuThetaTower}. Consequently it is the constant
representation. Its occurrence as functions consists of actual
constants, since right-invariant functions on \([U(Y)]\) are
constant. Wu's involution theorem now says that the forward lift of
these constants to \(W_1\) vanishes, where \(W_1\) is obtained from
\(W\) by deleting the rational hyperbolic plane
\cite[Theorem~5.1(3), Remark~5.2]{wuThetaTower}.
To check the twists in this use of the theorem, let \(d_X,d_Y\)
be its two splitting characters restricted to norm-one determinants
on the two groups \(X=W,Y\), respectively. Complex conjugation
of the oscillator replaces \(\psi,\chi_X,\chi_Y\) by their
inverses. Relative to Wu's oscillator for \(\psi^{-1}\) and the
original splitting characters it thus twists by
\(d_X^{-1}\boxtimes d_Y^{-1}\): on the doubled Siegel Levi this
follows from \(\chi(z/\bar z)=\chi(z)^2\); the doubled uniqueness
argument identifies the restrictions to the two factors. Wu's
reverse input is therefore \(d_X^{-1}\Pi\) and its output is
\(d_Y\) times the constants. The factor \(d_Y^{-1}\) in his
Theorem~5.1 cancels it in the forward lift.

This lower forward lift is nonzero by the calculation in
Theorem~\ref{threal:nonzero}. Here \(Y\) is unchanged, and the
new target has dimension \(d-2>2a'\), signature
\((d-p-1,p-1)\) with \(d-p-1\ge a'\), and is still hyperbolic
at the inert finite places, by Witt cancellation. The convergence
inequality there is \((d-2)-2(p-1)+1=d-2p+1>a'\).
These give respectively the absolute doubling identity, its real
source-invariant polynomial, and the required finite tests in that
proof. The splitting characters have unchanged parity on deleting
a plane. This contradiction proves the assertion for the primed
period as well.

Finally the zero-vector term on the unprimed side uses the
oscillator for the same \(V\) and the smaller \(W_1\). In the
polarization with position variable \(x\in V\) for the removed
target plane, setting \(x=0\) leaves the theta sum for
\(V\otimes W_1\); the extra source determinant character is
trivial for this added hyperbolic plane, whose source action on
the \(x\)-block is ordinary pullback. The source integral converges
by the proof of Proposition~\ref{threal:decay} with \(d-2\).
At a real place, the positive-character joint-harmonic condition in
Proposition~\ref{threal:allreal} now requires \(p\le p-1\).
Hence this local oscillator has zero trivial-source coinvariant
module: if its big theta module were nonzero, its unique irreducible
quotient would occur in the local correspondence. The source
integral at each \(h\) kills all the Lie algebra and compact-group
differences and so vanishes on the polynomial module. Continuity
and Hermite approximation give the same conclusion for all Schwartz
inputs. For the primed source the only rational isotropic vector is
zero, since \(Y\) is definite. Thus it has no nonzero-vector
orbit in this constant term.
\end{proof}

\begin{proposition}[Unfolding the larger nonzero orbit]\label{threal:coneA}
Let \(p>0\). Choose rational hyperbolic pairs \(x,y\) in \(V\)
and \(x',y'\) in \(W\), with orthogonal complements \(V_1,W_1\).
Write \(m(t,h_1)\) for the target line Levi and \(k=(a-a')/2>1\).
Use rational mixed Schr\"odinger coordinates and probability
measure on the target unipotent quotient. Let \(e_A\) be evaluation
in the coordinates dual to the modulation space
\[
        S_A= V\otimes_E Ex'\ \oplus\ Ex\otimes_E W_1.
\]
With \(K\) the source compact stabilizer and \(\rho\) the Heisenberg
action, put
\[
 T_A\Phi=\int_K\int_{\mathbb A_E^\times}
  e_A\rho(l\,x\otimes y')\omega(k_0,1)\Phi\
             \nu(l)|l|_E^k\,d^\times l\,dk_0 .
\]
The entire contribution of the rational nonzero isotropic vectors
to the unnormalized line constant term, at \(m(t,h_1)h\), is
\[
   c_A\,\alpha(t)|t|_E^{a'/2}\,
      I_{V_1,W_1}\big(T_A(\omega(1,h)\Phi)\big)(h_1).
\]
Here \(c_A\ne0\) depends on the rational models, the evaluations
and Haar measures; it is independent of the splitting characters.
At \(q=0<p\) this is the entire unprimed line constant term.
The entire primed constant term there is
\(\alpha(t)|t|_E^{a'/2}\) times its smaller zero-evaluation
period, with coefficient one for coherent mixed models and the same
source Haar measure.
\end{proposition}
\begin{proof}
First use \(V\otimes Ex'\) as modulation space. In its dual
coordinate the sum runs over \(z\in V(E)\). For fixed source and
target translates the sums are absolutely uniform on the compact
unipotent quotient. The center acts by an additive character of a
rational nonzero multiple of \(\langle z,z\rangle_V\); averaging
it leaves exactly the null vectors. The remaining target unipotent
coordinates have quotient \(W_1(\mathbb A_F)/W_1(F)\), writing
\(W_1\) here for its underlying \(F\)-vector group. They act on
the residual oscillator by the Heisenberg translations or
modulations \(z\otimes w\). This follows by applying the unipotent
matrix to the target coordinates \(x',W_1,y'\): the \(W_1\) shift
of \(y'\) is \(w\), and the quadratic correction in the \(x'\)
coordinate multiplies \(\langle z,z\rangle_V\), hence gives no
phase on this cone.

All nonzero rational null vectors form one \(U(V)(F)\)-orbit by
Witt extension. At the representative \(z=x\), use \(Ex\otimes W_1\)
as a further modulation space in the residual oscillator.
The last compact additive average kills every nonzero dual rational
coordinate by character orthogonality. Its surviving value is
therefore
\[
    \sum_{u\ {\rm rational}}
      (e_A\rho(x\otimes y')\,\omega(g,h)\Phi)(u),
\]
where \(u\) belongs to a position space for \(V_1\otimes W_1\).
The changes of rational polarization here preserve the theta
summation by Poisson summation with self-dual measures. Dual hyperbolic
bases may introduce a nonzero rational \(E\)-multiple of \(x\)
in this expression. Replacing the representative by that multiple
does not change the orbit, and its idele norm and Hecke character
are one.

Unfold the orbit to \(H_x(F)\backslash H(\mathbb A_F)\),
\(H=U(V)\), where the stabilizer of the vector is
\(H_x=N_x\rtimes U(V_1)\); the line stabilizer has the additional
Levi factor \(E^\times\). Write \(g=n\,m(b,g_1)k_0\) in source
Iwasawa coordinates. The adelic compact in this decomposition can
be taken to be the prescribed \(K\). At split finite places this
is integral row reduction. At an inert finite place one scales the
line generator to a primitive lattice vector and pairs it to one
by self-duality. The second vector can be made isotropic by the
surjectivity of the integral trace in the unramified extension.
Their unimodular hyperbolic plane splits off. Any two unitary
lattice complements here are isometric: a nonsingular finite-field
Hermitian space has a vector of nonzero norm, so integral
orthogonalization gives a unit diagonal, and every base-field unit
is an integral norm. The skew self-dual case is reduced to the
Hermitian case by the local different generator. These facts give
transitivity on isotropic lines. At infinity the two compact unitary
factors carry any isotropic line to the chosen line, by moving its
two components of equal length. The resulting local decompositions
are restricted products.

On the evaluated residual oscillator \(N_x\) acts trivially:
it fixes \(x\), acts with determinant one on \(S_A\), and acts
trivially on \(S_A^\perp/S_A=V_1\otimes W_1\). In mixed coordinates
the off-diagonal entries produce Heisenberg operators with parameters
linear in the dual \(S_A\)-coordinate, and quadratic phases in that
coordinate; both disappear at zero. It also fixes the displacement
\(x\otimes y'\). The source line Levi contributes
\(\nu(b)^{-1}|b|_E^{(d-2)/2}\) and the residual action of \(g_1\):
split off the whole target plane first, on which the source acts
by pullback on \(V\); on the remaining oscillator the source
Siegel character is \(\beta/\alpha\) and its modulation block
\(Ex\otimes W_1\) has dimension \(d-2\). Commuting the Levi left
past the displacement replaces \(x\otimes y'\) by
\(b^{-1}x\otimes y'\). The source Haar density on the line Levi
contains \(|b|_E^{1-a}\), because \(N_x\) has \(a-2\) ordinary
\(E\)-coordinates and one central \(F\)-coordinate of weight
\(b\bar b\). Integration over \([N_x]\) contributes its volume.
After the substitution \(l=b^{-1}\) the total character and modulus
are \(\nu(l)|l|_E^k\), since \(a-1-(d-2)/2=k\).
Integrating over \([U(V_1)]\) gives the smaller period.
The quotient integration constant and the unipotent volume are
positive for positive compatible Haar measures.

For justification of this unfolding one can integrate the
\emph{absolute value of the single orbit summand} over
\(H_x(F)\backslash H(\mathbb A_F)\) first. In the coordinates just
used its integral is bounded by the \(K,\mathbb A_E^\times\)
integral with weight \(|l|_E^k\) of the absolute smaller theta
sum integrated over \([U(V_1)]\). Choose the rational
source-adapted polarization of Proposition~\ref{threal:decay} for
this last sum. Its absolute integral is bounded by finitely many
Schwartz seminorms for each fixed finite test. The gap in dimension
is now \(k-1\) and the rational source Witt index is at most \(p-1\),
so the same negative-exponent estimates apply. In mixed coordinates
the new test is obtained by restriction to the linear coordinates
\((l,u)\); the partial Fourier transforms for changing the smaller
polarization act on \(u\) alone. At infinity these seminorms are
bounded by arbitrarily rapidly decreasing functions of \(l\),
uniformly on the compact \(K\)-orbit. At finitely many modified
finite places the restrictions are finite sums of products of
compactly supported tests in \(l,u\). At every other place they
are the indicator of the integral \(l\) times the smaller lattice
test. The absolute product is dominated by the Dedekind zeta
Euler product at \(k>1\); the finitely many local integrals at zero
also converge since \(k>0\). Tonelli applied to the translates of
the absolute orbit summand now justifies its global unfolding.
The same bounds prove that \(T_A\Phi\) is Schwartz and justify
interchanging its evaluation with the smaller period.
With multiplicative unit volume one at the basic places, the actual
scalar there is the absolutely convergent product of
\((1-\nu_w(\pi_w)(Nw)^{-k})^{-1}\), by the integral of the
lattice indicator along the displacement line.

On the target line Levi the evaluation scalar is
\(\beta(t)|t|_E^{a/2}\) because the whole \(V\otimes Ex'\)
block was removed first. Moving the Levi to the left past the
displacement changes \(l\) to \(l\bar t\); hence substitution in
the Mellin integral multiplies this scalar by
\(\nu(\bar t)^{-1}|t|_E^{-k}=\nu(t)|t|_E^{-k}\).
The target character is thus \(\alpha(t)|t|_E^{a'/2}\).
An evaluation with the same specified Heisenberg laws changes the
formula by a fixed scalar. Locally, invariance under the modulations
factors evaluation through restriction to zero: at a finite place
average on balls; at infinity differentiate the modulation law
to annihilate the coordinate multiples, which generate Schwartz
functions vanishing on that subspace (use Taylor's formula near it
and a cutoff away from it). The remaining intertwiner commutes with
all smaller translations and modulations, so is scalar. The
prescribed vacuum images fix the local normalizations at almost
every place. Thus changing the evaluations or the product measures
leaves the asserted independence from the splitting characters.
Finally Proposition~\ref{threal:orthA} removes the other unprimed
orbit when \(q=0\). On the primed side definiteness leaves only
zero, and the same first mixed evaluation gives the stated formula
with exponent \(a'/2\) and character \(\alpha\).
\end{proof}

\begin{proposition}[Detection by the entire line constant term]\label{threal:detectline}
Let \(f\) be a linear combination of the two periods of polynomial
tests, with the required support on the primed tests. In all
signatures, vanishing of the constant term of \(f\) along one
rational isotropic line, as a function on \(J(\mathbb A_F)\),
implies that \(f\) is cuspidal. For \(p+q>0\) this implies \(f=0\).
\end{proposition}
\begin{proof}
Witt extension conjugates any two rational isotropic subspaces of the
same dimension. If \(f\) is not cuspidal, choose the smallest \(r\)
for which its constant term along the stabilizer \(P_r\) of a rational
isotropic \(r\)-plane is nonzero. These are all the conjugacy classes
of maximal proper rational parabolics. By assumption \(r\ge2\).
On the Levi of \(P_r\), fix the other factor and vary the first
factor \(\mathrm{GL}_r(\mathbb A_E)\). Some resulting function,
allowing a compact right translate of \(f\) by Iwasawa decomposition,
is nonzero. It is cuspidal on this general linear group: the further
constant term along a \(j\)-plane in this block is the constant term
along the \((j,r)\)-flag, computed by unipotent Fubini, and this
also factors through the vanishing \(j\)-plane term for \(0<j<r\).

Here is the spectrum to which this cuspidal function belongs.
The original polynomial periods are compact-finite at infinity
and finite under the center of the target enveloping algebra:
source invariance factors the polynomial oscillator at each place
through its big theta module of the trivial representation, which
is of finite length by real Howe duality \cite{howeTranscending}.
Moderate growth follows from Proposition~\ref{threal:decay}.
The constant term restricted to its Levi is again center-finite.
Indeed on left unipotent-invariant functions restriction of a
central right differential operator gives its relative
Harish-Chandra projection (up to the modulus twist for the
unnormalized constant term). The Levi center is finite over this
image, since the invariants of the Levi Weyl group in the Cartan
polynomials are finite over the full Weyl invariants. This also
gives center-finiteness on the first direct factor of the Levi.

One can thus take a nonzero eigenfunction for the positive real
split center in the finite-dimensional span of its translates
(equivalently, extract an exponential coefficient on its Lie
algebra). These operations commute with finite-place translations.
Project also to a character on the compact norm-one idele class
center, and twist by a common real power of the idele determinant
to make the resulting central character unitary. Some such projection
is nonzero by compact Fourier decomposition. The resulting smooth,
compact-finite, center-finite cusp form is square-integrable modulo
the center by rapid decay on Siegel sets. It has a nonzero projection
to an irreducible unitary cuspidal representation of
\(\mathrm{GL}_r(\mathbb A_E)\). Its local factors are generic by
the cuspidal Whittaker expansion \cite{shalikaGeneric}. Thus at
each place above the \(v\) used in Proposition~\ref{threal:orthA}
the constant-term space must admit a generic irreducible subquotient
unitarizable after a determinant twist. Equivalently one can carry
out the center projections on the finite-place submodule generated
by the original slice: its split-center translates span a finite
number of copies of that module, and restriction to a central
eigenspace preserves its possible finite local subquotients.

At \(v\) each period factors through \(\pi_v\) by the full split
kernel computations, whether or not its coefficients at the other
places factor through their small theta quotients. The local
\(P_r\) has block sizes \((r,d-2r,r)\). Evaluation on the Levi after
taking the constant term factors through its unipotent coinvariants,
since a right unipotent translation becomes a translation of the
compact integration quotient. Normalizing the Jacquet functor only
twists each general linear block by a common real determinant power.
The geometric lemma computes the first \(\mathrm{GL}_r(k)\) action
in the resulting filtration, with possible subquotients of
\[
 \left(\alpha|\cdot|^{-(a-i)/2}\right)\circ\det_i
       \ \times\
 \left(\beta|\cdot|^{-(a'-j)/2}\right)\circ\det_j,\qquad i+j=r .
\]
Indeed a shuffle selects an initial block of length \(i\) of the
first determinant representation and of length \(j\) of the
second. The inverse square root of the modulus of their internal
parabolics gives the displayed exponents. The other end of the
isotropic plane accounts for the last, separate block of size \(r\).
\cite[geometric lemma, \S2.12]{bzgLemma}
The first block contributes a nondegenerate Whittaker functional
only if \(i,j\le1\): the highest derivative of the induced
representation is the tensor product of the highest derivatives
of its factors, and a determinant character of dimension \(>1\)
vanishes under that functor. The Whittaker coinvariant functor is
exact, so the same exclusion holds for every subquotient
\cite[\S4]{bzgLemma}.

It remains to examine \(r=2,\ i=j=1\). This normalized principal
series of \(\mathrm{GL}_2(k)\) is irreducible since its character
ratio has unitary part \(\alpha/\beta\) of infinite order and
therefore is neither \(|\cdot|\) nor \(|\cdot|^{-1}\).
The two Satake eigenvalues have distinct arguments and their
absolute values differ by the factor
\(|\pi_k|^{-(a-a')/2}\ne1\), where \(\pi_k\) is a uniformizer.
For unitarity the multiset must be invariant under
\(z\mapsto1/\bar z\), by the isomorphism to the conjugate dual.
That involution preserves the argument of each eigenvalue.
Consequently their distinct arguments force both absolute values
to be one. No common change of modulus can achieve this, giving the
required contradiction. Thus \(f\) is cuspidal. For \(p+q>0\)
Propositions~\ref{threal:orthA}, \ref{tpd:zero} and \ref{tpd:orth}
cover the two periods in every signature. Thus \(f\) is orthogonal
to the cuspidal spectrum and hence to itself; continuity gives
pointwise zero.
\end{proof}

\begin{corollary}[The detector in the \(a\)-descent]\label{threal:endA}
For \(q=0<p\), suppose that for one fixed \(c\) the
\(\alpha|\cdot|_E^{a'/2}\) components of the two periods agree up to
\(c\) on \(m(1,h_1)\) for every \(h_1\) and every polynomially
matched pair with supported primed test. Then those periods agree
up to \(c\) on all of \(J(\mathbb A_F)\) for every such pair.
\end{corollary}
\begin{proof}
By Proposition~\ref{threal:coneA} each entire line constant term
is just its indicated component. The matching condition survives
right translation by the target compact group: its real compact
action preserves the polynomial module, and its finite action
preserves full rank at \(v_1\) and intertwines the matching maps.
The hypothesis therefore applies also to \(m(1,h_1)k_0\) for each
target compact \(k_0\). The common Levi character extends this to
\(m(t,h_1)k_0\); left unipotent invariance and Iwasawa decomposition
cover the whole adelic group. Proposition~\ref{threal:detectline}
now applies to the difference.
\end{proof}

\begin{proposition}[The detector in the \(a'\)-descent]\label{threal:endB}
Let \(q>0\). Suppose a difference \(f\) of the two polynomial
periods, with supported primed test, has its full line constant
term transforming by \(\alpha(t)|t|_E^{a'/2}\) alone. Then \(f=0\).
In particular equality of the \(\beta(t)|t|_E^{a/2}\) components
of two matched periods, with one scalar for all the tests on
\(m(1,h_1)\), implies equality with that scalar on \(J(\mathbb A_F)\).
\end{proposition}
\begin{proof}
Induct on \(p\), keeping \(q>0\). For \(p=0\) the constant terms
of both periods have only the \(\beta\)-character by
Propositions~\ref{tpd:zero} and~\ref{tpd:unfold}. The two line
characters have distinct moduli since \(a>a'\), so the constant
term vanishes and Proposition~\ref{threal:detectline} applies.

For \(p>0\), split off planes in \(V,W\). At the real places one
can take compact groups adapted to the chosen planes. In detail,
choose on each hyperbolic plane a Cartan involution exchanging its
null lines and choose one on its orthogonal complement. The
resulting positive structure is conjugate to the original by a
real isometry. Translating the test by this isometry transports its
polynomial Fock module to the new one. Source translation leaves
the period fixed and simultaneous target translation replaces it
by a right translate, preserving purity and the zero assertion
to be proved. The lattices for the unused smaller finite models
can be chosen split along the rational decomposition at almost all
places, and compatible with the specified split polarizations at
\(v_1\). The induction uses only these rational theta models, their
finite tests, the prescribed signatures and the gap inequalities.
These hold in the smaller ranks \(a-2,a',d-2\), with indices
\(p-1,q\).

For each fixed target compact translate, restriction of this
pure constant term to the smaller unitary group is a difference
\(f_1\) of its two polynomial periods, with supported primed test.
Indeed the \(\alpha\)-component of the unprimed form is precisely
the cone in Proposition~\ref{threal:coneA}; the zero evaluation
has character \(\beta(t)|t|^{a/2}_E\). The smaller unprimed test
is polynomial at infinity by the orthogonal Fock factorization
of Lemma~\ref{tpd:mellin}, now with \(S_A\).
On the primed side the supported zero evaluation of
Lemma~\ref{tpd:so} gives precisely the \(\alpha\)-component.
The adjustment by source differences in that lemma may be made
anew for each fixed translate. Each resulting \(f_1\) in turn
has a two-character line constant term in its own ranks, by the
same calculations (and their definite-source versions). For an
arbitrary upper target translate, source Iwasawa is not needed
again: upper target Iwasawa and the pure character express its
slice as a scalar times a smaller target translate of a
compact-translate slice. Thus the same two line laws hold for all
slices used below.

We prove that the \(\beta\)-component of every such smaller
slice is zero. Choose any rational isotropic line in \(W_1\);
together with the first line it spans a rational isotropic
two-plane \(D\) in \(W\). For any fixed upper right translate
take the \(D\)-parabolic constant term of \(f\), and restrict its
first Levi block to \(\mathrm{SL}_2(\mathbb A_E)\), writing \(F\)
for the resulting automorphic function. Its ordinary upper
triangular constant term on this group is the average for the
isotropic \((1,2)\)-flag, by Fubini over the compact unipotent
quotients. Compute the same flag average in the other order by
first taking the outer line constant term and then the smaller
line term. These are averages along the two filtrations of the
\emph{same flag radical}. On \(d(t)=\operatorname{diag}(t,t^{-1})\)
the outcome, for each fixed right translate \(h\), is
\[
              F_N(d(t)h)=A(h)+\nu(t)|t|_E^{1-k} B(h),
                 \qquad k=(a-a')/2.
\]
The first outer factor is always \(\alpha(t)|t|_E^{a'/2}\).
The inner factors are \(\alpha(t^{-1})|t|_E^{-a'/2}\) and
\(\beta(t^{-1})|t|_E^{-(a-2)/2}\), respectively. These give
the two terms displayed. Their coefficients are continuous and
bounded on compact sets: they can be recovered by evaluating at
two fixed \(t\)'s with different idele norms.

In particular \(F\) belongs to
\(L^2(\mathrm{SL}_2(E)\backslash\mathrm{SL}_2(\mathbb A_E))\).
To check this, the displayed constant term is bounded on a Siegel
set since \(1-k<0\). Its difference from \(F\) decays faster than
any power of the split chamber parameter there. Indeed the
height exponent in Proposition~\ref{threal:decay} can be held fixed
for all right derivatives of the original period: choose once
the seminorm orders needed for integration; differentiating first
changes the test in these seminorms, not the powers of height in
the conjugation bounds. Integrating over the compact \(D\)-radical
therefore gives uniform moderate growth of derivatives of \(F\).
At fixed finite level the upper additive unipotent quotient is a
compact real torus. Subtract the mean of \(n\mapsto F(n d(t)c)\)
and use Fourier series there to bound its value by a Sobolev norm
of its positive-order derivatives of any sufficiently high order.
Along the archimedean positive chamber each unipotent derivative
gains the inverse positive root of \(d(t)\), as a right derivative
on \(F\). Compact factors \(c\) only change the fixed bounds. This
is the estimate in Corollary~\ref{tpd:detect} on the rank-one
group, with the bounded mean retained. Boundedness and finite
quotient volume give the assertion. This applies also to every
fixed finite local translate of the vector giving \(F\).

Use the good split place \(v\) of Proposition~\ref{threal:orthA}.
At one of its two \(E\)-places \(w\) the right
\(\mathrm{SL}_2(E_w)\) action on \(F\) factors through restriction
of the first block in the \(r=2\) Jacquet filtration computed in
Proposition~\ref{threal:detectline}. Its possible factors before
restriction are determinant characters and the one principal series
with \(i=j=1\). The latter has no unitarizable restriction
constituent. Here is a verification. Write \(G=\mathrm{GL}_2(E_w)\)
and \(G^+=Z(G)\mathrm{SL}_2(E_w)\). The quotient \(G/G^+\) is
the finite group of square classes. The irreducible principal
series \(\tau\) is finitely generated over \(G^+\) by coset
representatives applied to a generating vector. Choose a maximal
proper \(G^+\)-submodule; its finitely many \(G\)-conjugates
intersect trivially by \(G\)-irreducibility. The restriction
embeds in a finite sum of simple quotients, so is semisimple, and
its simple factors are \(G\)-conjugate. The center acts by a scalar
character. After twisting \(\tau\) by a real determinant power this
central character is unitary. If any factor on \(\mathrm{SL}_2(E_w)\)
had a positive invariant Hermitian form, so would every conjugate,
and their direct sum would give a \(G^+\)-invariant form.
Averaging it over \(G/G^+\) would unitarize the twist of \(\tau\),
contrary to the Satake calculation in Proposition~\ref{threal:detectline}.

The local subspace under consideration is a quotient of a
finite-length admissible \(\mathrm{SL}_2(E_w)\)-module: in the
Jacquet module take invariants of compact open subgroups fixing
the vector in all the other Levi blocks. Their invariant spaces
in the finite filtration are finite-dimensional on those blocks.
The first block restricts as just described. A finite-length
admissible submodule of a unitary Hilbert representation has
unitarizable simple constituents. In fact the closure of any of
its submodules has exactly that submodule's smooth vectors within
it, by compact-open averaging and finite-dimensionality at each
level; orthogonal projection then splits the submodule. Hence
only the trivial \(\mathrm{SL}_2(E_w)\) factors remain in the
space of translates of \(F\), and its local action is trivial.
In particular \(F_N(d(t))=F_N(1)\) if \(t\) is a uniformizer
at \(w\) and one elsewhere. Its modulus makes
\(\nu(t)|t|_E^{1-k}\ne1\). The flag formula now implies \(B(1)=0\).
The upper right translate was arbitrary, so the entire smaller
\(\beta\)-component vanishes.

The smaller slice \(f_1\) is thus pure for the \(\alpha\)-character.
Induction gives \(f_1=0\) for each upper compact translate.
The character law, upper unipotent invariance and Iwasawa
decomposition give vanishing of the full upper line constant term.
Proposition~\ref{threal:detectline} gives \(f=0\).
Finally equality on \(m(1,h_1)\) for all matched pairs kills
their entire \(\beta\)-component by compact stability of matching
and Iwasawa, as in Corollary~\ref{threal:endA}. The preceding
assertion applies to the remaining pure difference.
\end{proof}

\begin{proposition}[Matching all tests at an integral good split place]
\label{threal:localcone}
Let \(v\) lie above a completely split rational prime outside the
ramification set. Use unramified characters and a pair of integral
hyperbolic splittings in which the source, target and smaller basic
lattices decompose as direct sums and the displacement line is
primitive. Normalize the unit groups and \(K_v\) to volume one.
Write \(P_{1,A},P_{1,B}\) for the two smaller quotient maps,
normalized on their common vacuum. For every pair of locally
matched tests \(\phi,\phi'\), in the \(a\)-step one has
\[
   P_{1,A}(T_{A,v}\phi)=
    \prod_{w\mid v}L_w(\nu,k)\, P_{1,B}(e'_{0,v}\phi').
\]
In the \(a'\)-step the analogous equality is
\[
   P_{1,B}(T_{B,v}\phi')=
    \prod_{w\mid v}L_w(\nu,k+1)\, P_{1,A}(e_{0,v}\phi).
\]
Here the evaluations and basic vectors in the mixed coordinates
are normalized compatibly: evaluation of a basic vector is the
smaller basic vector, and its displacement evaluates to zero
outside the integral line. \(T_{B,v}\) is the local operator in
Lemma~\ref{tpd:mellin}. Both products range over the two places
of \(E\) above \(v\).
At a designated ramified split place the same pairs of comparison
maps span a space of dimension at most one in each direction.
\end{proposition}
\begin{proof}
Set \(K_0=F_v\). By Propositions~\ref{split:kernel} and
\ref{split:good}, or by Corollary~\ref{split:ram} at a designated
ramified place, the two source coinvariants are the same
irreducible \(\mathrm{GL}_d(K_0)\) representation \(\pi_v\).
In particular these are the precise kernels of any of the
surjective local maps in the matching data: a map killing the
source differences induces a nonzero map out of this irreducible
module. We show that the two maps to the smaller output after
applying a cone or zero evaluation factor through these
coinvariants.

For the zero evaluation, equivariance for the unchanged source
and its action on the smaller quotient give this immediately.
For \(A\), project
\(e_A\rho(x\otimes y')\omega(g,1)\phi\) to the smaller quotient.
It is left invariant by the full \(K_0\)-point vector stabilizer
\(H_x\), since its unipotent radical acts trivially as in
Proposition~\ref{threal:coneA}, and its remaining source factor
acts trivially after projection. This stabilizer is unimodular:
the action of the residual general linear block on the row and
column root groups has opposite determinant moduli, and is trivial
on the center. Integration over \(H_x\backslash H(K_0)\) therefore
gives a right \(H(K_0)\)-invariant map. Its Iwasawa coordinates
give precisely \(P_{1,A}T_{A,v}\), up to a nonzero Haar factor,
by the source Levi calculation above. The integral is absolute
since \(k>0\); the locally constant compactly supported restrictions
in the displacement and smaller position variables span a finite
number of smaller tests. For \(B\) use the vector stabilizer of
\(x\) in the primed source. It fixes the \(S_B\)-reduced quotient
and the class of \(y\otimes x'\) modulo \(S_B\), with no central
Heisenberg phase: the difference lies in
\((V'_1+Ex)\otimes Ex'\), isotropic together with
\(y\otimes x'\). The residual \(U(V'_1)\) again disappears after
projection. The Iwasawa calculation in Lemma~\ref{tpd:mellin}
gives \(P_{1,B}T_{B,v}\) by this quotient integral, now with
positive exponent \(k+1\).

Both resulting maps from \(\pi_v\) are invariant under the
unipotent radical of the target line parabolic. For \(A\), the
displacement changes there by an element of
\(Ex\otimes(W_1+Ex')\) in \(S_A\), with no Heisenberg phase
because \(x\) is isotropic. For \(B\), the displacement is fixed;
the target unipotent preserves \(S_B\), has determinant one there,
and acts trivially on the smaller tensor quotient. For the zero
evaluation the assertion holds already before the smaller quotient.
The actions on the reduced spaces in these checks are the mixed
ones at zero, as in Proposition~\ref{threal:coneA}. Moreover the
two maps in either comparison are equivariant for the smaller
target group and have the same line character, by the two Levi
computations.

At this split place the line parabolic has block sizes
\((1,d-2,1)\). Write its outer entries as \(u,z\); the element of
\(E_v^\times\) acting on the isotropic \(E_v\)-line is
\((u,z^{-1})\). The unnormalized line character in the \(a\)-step
is therefore \(\alpha(u)\alpha(z)|u/z|^{a'/2}\). Dividing by
\(\delta_P^{1/2}=|u/z|^{(d-1)/2}\) prescribes the first and last
characters of the normalized Jacquet module to be
\[
       \alpha|\cdot|^{-(a-1)/2},\qquad
       \alpha|\cdot|^{(a-1)/2}.
\]
The geometric lemma applied to the two determinant blocks of
\(\pi_v\) has exactly one such subquotient: the first and last
rows must both come from the \(\alpha\)-block, since
\(\alpha/\beta\) has infinite order on a uniformizer in the good
case, and is nontrivial on units in the ramified case. They leave
the normalized induction of the determinant characters of lengths
\(a-2,a'\) on the middle block. Its \(\alpha\)-exponent has no
shift because one row was removed at either end. It is irreducible
by the unitary induction theorem used in Proposition~\ref{split:good}.
In the \(a'\)-step the line character is \(\beta(t)|t|_E^{a/2}\);
the prescribed ends are \(\beta|\cdot|^{\mp(a'-1)/2}\). Both
come from the other block, leaving lengths \(a,a'-2\). Thus the
equivariant Hom into the stipulated smaller representation and
the two given characters has dimension at most one in each case,
by the Jacquet filtration and Schur's lemma. \cite[\S2.12]{bzgLemma}

At the stated good places evaluate the two maps on the upper
common spherical vector.
The zero evaluation gives the nonzero smaller vacuum. In the
cone integral the compact group fixes the test, and the shifted
evaluation gives exactly the smaller vacuum for integral displacement.
At each of the two places \(w\) the resulting geometric series
is \(L_w(\nu,k)\) or \(L_w(\nu,k+1)\), respectively. The product
is nonzero, fixing the scalar on every matched test.
\end{proof}

\begin{corollary}[The ramified unprimed comparison on all tests]\label{threal:ramA}
At a designated ramified split place, use the local character
\(\chi=\nu_w\), conductor \(\pi^c\), residue cardinality \(Q\)
and measures of Proposition~\ref{tpd:ram}. For all upper matched
tests the scalar in the \(a\)-step, multiplying the smaller primed
zero evaluation, is
\[
       Q^{c(k-1)}\chi(\pi)^c
       \sum_{r\in(\mathcal O/\pi^c)^\times}
                         \chi(r)^{-1}\psi_v(r/\pi^c).
\]
\end{corollary}
\begin{proof}
Use the pivot-determinant tests of Proposition~\ref{tpd:ram},
now removing the first and last of the \(a\) target pivot columns
on the unprimed side. The upper and lower pairs are matched by the
same integral row computation there, and both orbital functions
have value one at the identity before applying the reductions.
The primed zero evaluation integrates one nonpivot column with
\(a'\) entries in \(\pi^c\mathcal O\), giving
\(Q^{-a'c}\phi'_1\). On the unprimed side write
\(l=(s,t)\) in the two split components. The shift \(l x\otimes y'\)
has position entry \(s\) at the first row and last pivot column
and momentum entry \(-t\) at the last row and first pivot column:
in the chosen idempotent the coordinates of \(x,y'\) are the
first and last basis vectors; pairing their conjugates into the
dual basis gives the last and first vectors, with the minus sign
from the skew target pair. The two entries have zero mutual
Heisenberg half-pairing.
\src{src/Development/WeilBound/ThetaLifts/LocalTheory/StandardValuations/AdelicInclusionMapPlaceSection.lean}{218}{Paired-coordinate identities for the unprimed displacement}.

In the \(S_A\)-evaluation one integrates the first pivot column
and the first source row except for the last pivot column. This
follows by using the translations in \(V\otimes Ex'\) first
and then those in \(Ex\otimes W_1\); the dual modulation
coordinates are set equal to zero. Thus the pivot matrix has
the form
\[
      \begin{pmatrix} *&*&s\\ *&D_1&0\\ u&0&0\end{pmatrix}
\]
with Fourier factor \(\psi_v(-tu)\). Its determinant is
\(-su\det D_1\). Apart from \(u\), the integrated pivot entries
have unit volume, and the \(a'\) integrated nonpivot entries of
the first row give \(Q^{-a'c}\). Replacing \(u\) by \(-u\)
removes the two minus signs. With
\(D_\chi=\mathbf1_{\mathcal O^\times}\chi^{-1}\), this yields
\(Q^{-a'c}D_\chi(s)\widehat D_\chi(t)\phi_1\).
The compact source average fixes the test. The Mellin weight is
\(\chi(s)\chi(t)^{-1}|st|^k\). The \(s\)-integral is one; the
\(t\)-integral is supported on \(t=\pi^{-c}\mathcal O^\times\),
where \(\widehat D_\chi(\pi^{-c}u)=
 Q^{-c}\chi(u)G(\chi,\psi_v)\) by the unit residue sum.
Consequently this integral is \(Q^{c(k-1)}\chi(\pi)^c
G(\chi,\psi_v)\). Cancel the common factor \(Q^{-a'c}\) and use
Proposition~\ref{threal:localcone} to pass from this nonzero
smaller vector to all tests.
\end{proof}

For an adapted integral plane let \(\mathfrak b x\) and
\(\mathfrak c x'\) be its two first-line lattices. The duality
condition on the skew lattice, using the absolute trace, gives
the second target line lattice
\(\mathfrak D_{E/\mathbb Q}^{-1}\bar{\mathfrak c}^{-1}y'\).
Indeed a coefficient must pair with every element of
\(\mathfrak c x'\) into the codifferent, and the other integral
summands are orthogonal. Thus the finite displacement
\(l x\otimes y'\) is integral precisely for \(l\) in
\(\mathfrak b\mathfrak D_{E/\mathbb Q}^{-1}
  \bar{\mathfrak c}^{-1}\widehat{\mathcal O}_E\).
These ideals are units at the specified ramified places by the
prescribed paired integral bases and total rational splitting.
The relative different for \(E/F\) is trivial and
\(\nu|_{\mathbb A_F^\times}=1\), by the common parity of the
two splitting characters. Hence the character of the codifferent
is one, and that of the displacement ideal is
\(\nu(\mathfrak b)\nu(\mathfrak c)\); the norm factor with
exponent \(-k\) depends only on the chosen lattices. This gives
the two positive ideal exponents for the \(a\)-step. In particular
the splitting-prime calculation of the two comparison factors also
allows nonprimitive line generators by the corresponding change of
variable in the local Mellin test of the lattice vacuum.
\src{src/Development/WeilBound/ThetaLifts/LocalTheory/ConeShifts/IdealId0.lean}{49}{The two displacement ideals and their character values}.

\paragraph{The scalar comparison for \(p+q>0\).}
The common local maps give an algebraic test for the
global discrepancy. Fix polynomial real and finite tests and take the
restricted tensor products with the specified vacuum lines. Let
\(P,P'\) be the tensor maps to the common quotient, using the
full-rank restriction on the primed domain. Suppose the quotient
and full-rank surjectivities in \S\ref{newtheta:global} have been
established at the places in question. Then \(P,P'\) are onto:
each pure tensor has finitely many non-vacuum factors; choose lifts
factor by factor and use the prescribed vacua elsewhere. Work with
the rational summation functional averaged over the source, denoted
by \(I_h,I'_h\). For fixed \(h\) its action on source differences
vanishes by change of variables wherever the summation descends to
the source quotient. The first additional global requirement is
\[
            I_h(\ker P)=0,\qquad I'_h(\ker P')=0 .
\]
In particular a local quotient that kills the source differences
does not identify the annihilator of the global integral unless its
extra kernel is killed too. Once these equalities hold, the two
periods give well-defined linear maps \(A,A'\) from the common tensor
module to functions of \(h\). Proportionality is equivalent, when
\(A'\ne0\), to the following directly testable equalities for all
module vectors \(x,y\) and adelic \(h,h_0\):
\[
                   A(x)(h)\,A'(y)(h_0)
                    =A'(x)(h)\,A(y)(h_0).
\]
Choose \(y,h_0\) with \(A'(y)(h_0)\ne0\) to prove this equivalence;
the scalar is nonzero exactly when \(A\ne0\). Corollary~\ref{threal:limits}
then gives the assertion on the entire matching space. Conversely a
nonzero primed period on some supported Schwartz test already gives
one on a polynomial test: keep its finite factors fixed, approximate
each archimedean Schwartz coefficient by its finite Hermite
expansions, and use the continuity proved above. Hermite expansions
converge in all Schwartz seminorms, by iterating the harmonic
oscillator on the coefficients and using the creation and annihilation
operators to control the seminorms of its eigenfunctions.

Theorem~\ref{threal:base} establishes the scalar comparison at the
definite base. Theorem~\ref{threal:nonzero} proves \(A'\ne0\) also
for \(q=0<p\) whenever the factorization exists, and gives a nonzero
supported period there before factorization. The graph-orbit proof
of the base comparison changes once \(W\) has isotropic lines:
maximal isotropic subspaces of \(W\oplus -W\) can project with
nonzero kernels, and their terms survive in the unfolding unless
one proves the necessary vanishing. Proposition~\ref{threal:allreal}
gives the real quotient in these signatures, and
Proposition~\ref{threal:detectline} reduces the global comparison to
one entire line constant term for all target translates.

The rank restriction on other available multiplicity formulas
matters here. The determinant characters on the two inducing blocks
give the two Satake strings of lengths \(a,a'>1\) at the completely
split good primes. The natural Arthur parameter with these strings
has two character blocks with nontrivial Arthur \(\mathrm{SL}_2\);
the generic multiplicity-free formula does not apply to it.
Chen--Zou's Theorem~7.7 treats elliptic parameters at \(F\)-rank
one and also those for which each nontrivial Arthur block has
cuspidal degree greater than the rational Witt index. Here those
degrees are one. Moreover this theorem describes a discrete
spectrum in terms of packets, defined as multisets; using it on a
specific period would require proving discrete realization and
identifying its packet member. Even a multiplicity of one does not
give the ratio of functionals in different ranks.
\cite[\S5.3, Assumption~7.2 and Theorem~7.7]{chenZouMultiplicity}

In the descent in \(a\), the actual comparison of tests must send
the larger cone integral to the smaller zero evaluation times
\(C_A\,\mathrm{stepFac}_A\,L^{S_0}(\nu,(a-a')/2)\).
Proposition~\ref{threal:coneA} identifies the two global constant
terms with those operations and Corollary~\ref{threal:endA} proves
the required detection. Identifying their images in the smaller
matching quotients with the uniform scalar is a separate assertion.
In the descent in \(a'\), the smaller cone integral must match the
larger zero evaluation times
\(C_B\,\mathrm{stepFac}_B\,L^{S_0}(\nu,(a-a')/2+1)\), and must
land in the smaller supported space. Lemma~\ref{tpd:mellin}
proves that support assertion, Proposition~\ref{tpd:unfold}
identifies the cone component, and Proposition~\ref{threal:endB}
proves detection by it in every signature. The signs in the two
finite products and the conductors are as in \S\ref{newtheta:global}.
The constants \(C_A,C_B\) precede the quantifiers over the characters,
both Weil systems, and both sets of finite quotient maps once the
evaluations, real maps and measures are fixed.
Proposition~\ref{threal:localcone} determines both ratios on all
tests at the integral good split places in its hypotheses.
Proposition~\ref{tpd:ordinary} extends these identities to every
unramified place split in \(E/F\), including resonant characters
and the factors from the adapted line lattices.
Corollary~\ref{threal:ramA} and Proposition~\ref{tpd:ram} determine
the ratios at the designated ramified split places.
At inert places Proposition~\ref{threal:allinert} gives the
irreducible quotient and its common vacuum. Propositions~\ref{tcr:descend} and \ref{tcr:unique} prove that the
cone and zero operations factor through the upper quotient on every
local test and compute their ratio after projecting to the smaller quotient. At the real places
Proposition~\ref{threal:allreal} identifies the two modules in
each fixed rank. Proposition~\ref{tcr:unique} compares the compact-averaged
Mellin operators and evaluations between the two ranks on the
polynomial vectors for the fixed real quotient maps. Theorem~\ref{tcr:global}
assembles the comparisons for the prescribed polarizations, lattice vacua
and product measures, giving the uniform constants and nonvanishing in
positive signature.
\src{src/Development/WeilBound/ThetaLifts/ConeHypotheses/PolDense.lean}{46}{The two normalized cone comparisons};
\src{src/Development/WeilBound/ThetaLifts/ConeHypotheses/PolDense.lean}{32}{Detection by the selected constant-term component}.

\subsection{The positive primed Witt index}
\label{tpd:main}

Put \(k=(a-a')/2\), and let \(H\) denote a hyperbolic plane of the
appropriate sign. In the step \(q\mapsto q-1\), the unchanged source
is \(V\), the other source changes from \(V'_1\oplus H\) to \(V'_1\),
and the target from \(W=W_1\oplus H\) to \(W_1\). We use the
splitting characters of \S\ref{newtheta:global} throughout this tower.
Write \(P=MN\) for the target parabolic stabilizing the first
isotropic line of \(H\), and normalize \([N]\) to have volume one.
Write \(m(t,h_1)=\operatorname{diag}(t,h_1,\bar t^{-1})\).

\begin{proposition}[The unprimed coefficient when \(p=0\)]
\label{tpd:zero}
Suppose \(p=0\) and \(q>0\). In rational mixed coordinates, let
\(e_0\) be evaluation at zero in the new \(V(\mathbb A_F)\)
variable. For the coherent mixed realizations and the same Haar
measure on the unchanged source, the full, unnormalized constant term is
\[
 (I_\Phi)_P(m(t,h_1)h)=
 \beta(t)|t|_{\mathbb A_E}^{a/2}\,
 I^{W_1}_{e_0(\omega(1,h)\Phi)}(h_1).
\]
Some polynomial archimedean test with an allowable finite factor has
nonzero coefficient. Every unprimed lift with polynomial archimedean
test is orthogonal to the cuspidal spectrum of \(U(W)\).
\end{proposition}
\begin{proof}
In the mixed model the rational summation variables are \((x,y)\),
where \(x\in V(E)\) and \(y\) belongs to a rational position space
for \(V\otimes_E W_1\). The center of \(N\) acts through
\(\psi(bc\langle x,x\rangle_V)\), with \(b\in\mathbb A_F\) and
a fixed nonzero rational factor \(c\in F\) determined by the paired
basis of \(H\). Its integral over \(F\backslash\mathbb A_F\) selects
\(\langle x,x\rangle_V=0\). Definiteness gives \(x=0\).
The other coordinates of \(N\) act on the smaller oscillator through
Heisenberg operators whose translation and modulation parameters are
linear in \(x\), with quadratic central phase zero at \(x=0\).
They therefore act trivially on this value. Both unipotent and source
quotients are compact, and the rational Schwartz sum is absolutely
and uniformly convergent on compact sets of representatives; this
justifies both integrations. The line Levi acts at zero by
\(\beta(t)|t|_{\mathbb A_E}^{a/2}\), the square root of the modulus
of scalar substitution on the \(a\) coordinates times the prescribed
Siegel character. The smaller target uses the same splitting; on the
source the added plane acts on this variable by pullback. This is
also the mixed-model tower formula for a cuspidal source, applied to
the trivial representation of anisotropic \(U(V)\).
\cite[Proposition~3.1, (3.2)--(3.4)]{wuThetaTower}

There are \(q\) rational hyperbolic planes in \(W\). At split finite
places it is split, at inert places it is hyperbolic by the perfect
lattice calculation of Proposition~\ref{threal:allinert}, and at
every real place its Witt index is \(q\). To split off one plane,
multiply the skew form by a nonzero trace-zero element of \(E\) and
apply Hasse--Minkowski to the resulting quadratic norm form on the
underlying \(F\)-space. Complete an isotropic vector to a hyperbolic
pair. Witt cancellation preserves the remaining local indices, so
the argument can be iterated. The bottom \(X\) is definite of
dimension \(d_0=d-2q\).

The unprimed nonvanishing calculation in Theorem~\ref{threal:base}
applies to \((V,X)\): its argument up to the evaluation of
\(Z((d_0-a)/2)\) uses only these two spaces. Indeed
\(a'-2q>1\), \(d_0\) is even and the infinite angular exponent still
exceeds \(d_0\). At inert places \(X\) is hyperbolic and admits a
self-dual lattice, and at split places one may take dual bases.
These lattices agree almost everywhere with the lattice obtained by
a rational orthogonal decomposition. The denominator arguments
are at least \(a'-2q+1>1\) and the numerator integers are
\(d_0/2-i\) for \(0\leq i<a\); the same functional equation and
nonvanishing at one apply. At infinity the unprimed invariant
determinant used there has exponent \((m_\alpha-m_\beta-d_0)/2>0\)
and there are \(d_0\geq a\) columns. This gives a nonzero bottom
period.

Evaluation in each new variable is onto the smaller Schwartz
space by tensoring with a test of value one at zero. These extensions
are restricted tensors, since away from finitely many places the
chosen rational hyperbolic basis and its inverse are integral. The
constant-term formula gives nonvanishing successively in the tower.
With the finite test fixed, the resulting nonzero coefficient at
a fixed Levi element is a continuous functional of the infinite
Schwartz test, by compactness of the integrating quotients.
Polynomial times Gaussian tests in the original coordinates are
dense in the Schwartz topology, giving the asserted polynomial test.

For orthogonality we keep the unconjugated convention of Wu, in which
the two directions use inverse additive characters. Set
\(Y=V\), and let \(\chi_1,\chi_2\) be the splitting characters on
the target \(U(X)\) and the source \(U(Y)\), respectively, in
Wu's labeling (the twists below mean composition with determinant).
The bottom lift \(\sigma=\theta^X_{Y,\psi}(\mathbf1)\) is nonzero
and cuspidal since \(X\) is anisotropic. The input is itself an
irreducible cuspidal representation. Wu's Theorem~5.5 gives
irreducibility of \(\sigma\) and Theorem~5.4(1), with the two
additive characters in their indicated order, gives
\[
 \theta^Y_{X,\psi^{-1}}(\chi_1^{-1}\sigma)=\chi_2\mathbf1.
\]
The input and output in this last equality are again nonzero,
irreducible and cuspidal. Applying Theorem~5.1(2) with
\(\pi=\chi_1^{-1}\sigma\) therefore shows that
\(\theta^{X\oplus H^q}_{Y,\psi}
 (\chi_2^{-1}(\chi_2\mathbf1))\) pairs to zero with every cusp
form for \(q>0\). This is the asserted unprimed lift.
\cite[Theorems~5.1, 5.4 and 5.5]{wuThetaTower}
\end{proof}

\begin{proposition}[Cuspidal orthogonality of the supported primed lift]
\label{tpd:orth}
For every \(q>0\), including \(p>0\), every supported primed
theta integral pairs to zero with all cuspidal forms on \(U(W)\).
\end{proposition}
\begin{proof}
First take a compact-finite automorphic vector \(u\) in a single
irreducible cuspidal constituent on \(J=U(W)\). Such vectors span
a dense subspace of the cuspidal spectrum. Set
\[
 B(g)=\int_{[J]}\Theta_{\Phi'}(g,h)\overline{u(h)}\,dh,
        \qquad g\in U(V')(\mathbb A_F).
\]
Use Siegel representatives for \(h\) in this integral. Rapid decay of
\(u\) there dominates every fixed positive power of \(H(h)\), by
the rapid decrease theorem for cusp forms. Proposition~\ref{threal:decay}
then permits Fubini, and \(B\) is a rapidly decreasing, square-integrable
reverse theta lift. More explicitly its proof gives a bound
\(C_{L,u}\prod_i T_i^{-L}\) on each source Siegel set, with the
same bound when the right compact set of representatives is enlarged.
The finite supported condition is preserved by translations in either
group, so the estimate also applies to every fixed right translate.

Consider any rational maximal parabolic of \(U(V')\), stabilizing a
totally isotropic \(j\)-plane. In the reverse direction the input
\(\bar u\) is cuspidal. Proposition~3.1 of Wu applies to this
\(j\)-plane, with the roles of the Hermitian and skew-Hermitian
spaces exchanged (equivalently rescale the two forms by reciprocal
trace-zero scalars). Its unnormalized constant term at
\(m(t,h_0)g_0\), where \(t\in\mathrm{GL}_j(\mathbb A_E)\), is a
function of \(h_0,g_0\) times
\(\xi(\det t)|\det t|_{\mathbb A_E}^{d/2}\) for a unitary
splitting character \(\xi\). Take \(t\) to be the scalar \(R I_j\)
at each real place and the identity at finite places and let
\(R\to\infty\). On this positive chamber ray the integral over the
unipotent quotient tends to zero: its representatives can be taken
in a compact set which stays bounded after conjugation by
\(m(t,1)^{-1}\) on the left and \(m(t,1)\) on the right, and the
preceding bound applies. The modulus of the displayed character
grows as \(R^{[F:\mathbb Q]jd}\), forcing its coefficient to
vanish for every \(h_0,g_0\). Thus \(B\) is cuspidal.
\cite[Proposition~3.1]{wuThetaTower}

The group \(U(V')\) has a proper rational parabolic. Indeed the
given real Witt index is \(q>0\); at inert places the perfect
Hermitian lattice has Witt index \(\lfloor a'/2\rfloor\), and at
the remaining finite places it is split. The same norm-form
argument as above supplies a rational isotropic line.
The orthogonal projection in \(L^2([U(V')])\) of the constant
function to the closed cuspidal subspace is right invariant, hence
constant. Its constant term at the line parabolic must vanish, so
this projection is zero. It follows that \(\int B(g)\,dg=0\).
Fubini identifies this with the asserted pairing.
\end{proof}

\begin{lemma}[Removal of the primed zero orbit]\label{tpd:primezero}
Suppose \(p=0<q\). In taking the constant term of the supported
primed integral along \(P\), one can choose a test giving the same
automorphic function for which the \(x=0\) summand vanishes
identically, also after any target translation.
\end{lemma}
\begin{proof}
At one real place average the upper test over \(K_+=U(q)\), the
positive compact block of \(U(V')\). The unchanged source integral
follows by right invariance of Haar measure, compact averaging and
Proposition~\ref{threal:decay}. The full-rank condition at \(v_1\)
is preserved. Mixed evaluation at \(x=0\) is \(U(V')\)-equivariant
to the smaller oscillator for \((V',W_1)\): on the extra variable
the group acts by pullback, using the coherent splitting on decreasing
the partner by a hyperbolic plane. Thus its value on the averaged
test belongs to the \(K_+\)-invariants of that oscillator.

Write \((r_1,t_1)=(d-q-1,q-1)\) for the signature of \(W_1\).
As in the entry calculation of Proposition~\ref{threal:allreal},
the \(K_+\)-vacuum character in the primed oscillator is
\(\det^{b}\), where \(b=(m_\alpha-m_\beta+r_1-t_1)/2>0\).
There are \(r_1\) standard and \(t_1\) dual column variables
for this block. By the Cauchy decomposition the compact types can
only come from
\(\det^b\otimes S_\lambda(\mathbb C^q)
                    \otimes S_\mu(\mathbb C^q)^*\),
with \(\ell(\mu)\leq t_1=q-1\); variables from other rows are
trivial on \(K_+\). An invariant would equate the highest weights
\(\mu=\lambda+(b,\ldots,b)\), impossible at the last entry.
Hence the compact average is zero on the polynomial module, and
by density and continuity on the Schwartz space. Target translations
commute with compact source averaging. The smaller evaluation still
vanishes after applying any source element by equivariance.

On a compact truncation of the integrating source, average the
center of \(N\) first. Additive orthogonality leaves the sum over
the rational null cone, and the zero summand just treated vanishes
pointwise. Taking an exhaustion of the source quotient gives the
constant term as the limit of these nonzero-orbit expressions;
the original kernels are dominated by Proposition~\ref{threal:decay}
uniformly over compact representatives for \([N]\). This formulation
requires no separate integrability assertion for an unrestricted
smaller zero-orbit period.
\end{proof}

\begin{corollary}[Detection for a definite unprimed source]\label{tpd:detect}
If \(p=0,q>0\), two matched polynomial tests give proportional
automorphic functions with scalar \(c\) as soon as their full line
constant terms are proportional with that scalar at every target
translate. For the unprimed form this constant term is the smaller
period of Proposition~\ref{tpd:zero}; for the primed form it consists
only of the nonzero orbit of Lemma~\ref{tpd:primezero}.
\end{corollary}
\begin{proof}
For any \(q\) this follows from Proposition~\ref{threal:detectline}.
We give a direct argument in rational Witt index one.
Suppose \(q=1\).
The rational Witt index of \(W\) is one by its signatures and the
isotropy proved above. Every proper rational parabolic stabilizes a
line and these lines are rationally conjugate by Witt's theorem. The
difference \(D\) of the forms thus has every proper constant term zero.
Here one can see square integrability directly. The seminorm bounds
in Proposition~\ref{threal:decay} give a single moderate growth
exponent for all right archimedean derivatives of \(D\). Indeed
choose the source decay order in that proof once, large enough for
source integration. Applying a right derivative replaces the infinite
test by its derived oscillator action. The orders of the seminorms
used to bound the translated test and the powers of the target
height in those bounds depend on this fixed decay order and the
dimensions. Only the seminorm values of the initial test change.

On a rank-one Siegel set write the target representative as \(tc\)
with positive chamber coordinate \(R\ge1\) and bounded \(c\).
At fixed finite level the compact unipotent quotient in its constant
term reduces to a compact real nilmanifold. Indeed strong
approximation on the unipotent group follows successively for its
additive root spaces, and the finite open stabilizer can be taken
common to the bounded finite components. On this nilmanifold the
function \(n\mapsto D(ntc)\) has mean zero. Let \(\Delta\) be
the sum of negative squares of a left-invariant real basis; Haar
volume is invariant and its kernel consists of constants. Elliptic
Sobolev embedding and the spectral gap on the mean-zero subspace
bound a point value by \(C_j\|\Delta^jD(\cdot\,tc)\|_{L^2}\)
for every sufficiently large \(j\). Each basis derivative becomes
a right derivative on \(D\) with Lie algebra argument
\(\operatorname{Ad}((tc)^{-1})X=O(R^{-1})\), on choosing the
coordinate with smallest positive root \(R\). The uniform growth
bound on the compact set of unipotent representatives now gives
\(D(tc)=O_j(R^{B-2j})\) for a fixed \(B\). Thus \(D\) is rapidly
decreasing and belongs to the cuspidal subspace of \(L^2\).
Propositions~\ref{tpd:zero} and~\ref{tpd:orth} make it orthogonal
to the cuspidal spectrum, hence to itself.
\end{proof}

\begin{lemma}[The Mellin exponent and the good Euler factors]
\label{tpd:mellin}
Choose paired bases \(x,y\) and \(x',y'\) of the source and target
hyperbolic planes, respectively. In the rational mixed polarization
for simultaneous removal of these two planes, take the following
isotropic space to act by modulations:
\[
 S=Ex\otimes_E W\ \oplus\ V'_1\otimes_E Ex'.
\]
Write \(e'\) for evaluation at zero in its dual position coordinates,
and \(\rho\) for the Heisenberg action. With \(K\) the product of the
specified source compact stabilizers, the integral
\[
 T\Phi'(z_1)=\int_K\int_{\mathbb A_E^\times}
   \big[e'\rho(\bar l\, y\otimes x')\omega'(k_0,1)\Phi'\big](z_1)
           \,\nu(l)|l|_E^{k+1}\,d^\times l\,dk_0
\]
is absolutely convergent and defines a smaller Schwartz function.
It preserves the required full-rank support at \(v_1\).
Its target line character is \(\beta(t)|t|_E^{a/2}\).
With local unit groups of multiplicative mass one almost everywhere,
the scalar from the unmodified unramified factors outside any
sufficiently large finite set is the corresponding partial Euler
product for \(L_E(\nu,k+1)\).
\end{lemma}
\begin{proof}
The chosen subspace has dimension \(d+a'-2\) over \(E\);
its symplectic perpendicular modulo itself is the smaller tensor
space \(V'_1\otimes_EW_1\), by the two orthogonal decompositions.
We can choose its dual coordinate space to contain the translation
line \(E(y\otimes x')\), paired nondegenerately with
\(E(x\otimes y')\). Consequently the bracket in the integral is the
restriction of a mixed-model Schwartz function to a linear subspace
with coordinates \((\bar l,z_1)\), allowing a fixed sign or unit
change in the first coordinate. A rational choice of these
coordinates is integral with integral inverse almost everywhere,
where the vacuum is a nonzero multiple of the product lattice
indicator and the compact source fixes it. Choosing each original
basic vector and its image as the smaller basic vector removes the
same multiple on both sides. For a fixed finite test the compact
finite orbit is finite, by smoothness; at infinity the compact
orbit is bounded in every Schwartz seminorm. Finite-place restrictions
are finite sums of products of locally constant compactly supported
functions of \(l\) and \(z_1\) off the basic places. At infinity
all smaller Schwartz seminorms and derivatives in \(z_1\) are
dominated by arbitrarily rapidly decreasing majorants in \(l\).
Since \(k+1>1\), the integrals at zero converge. At a basic place
\(w\) of \(E\) the integral is the geometric series
\(\sum_{j\ge0}\nu_w(\pi_w)^j(Nw)^{-j(k+1)}\) times the smaller
basic vector. The absolute products converge since the Dedekind
zeta series converges at \(k+1\). This proves the analytic and Euler
factor assertions, including differentiation under the integrals.

At \(v_1\), choose the two conjugate idempotent factors of the
split tensor space as position and momentum in the specified primed
matrix model. Order the rows by \((x,V'_1,y)\) and the columns
by \((x',W_1,y')\) on the position side. The complementary momentum
basis pairs \(x\) with \(y\) and \(x'\) with \(y'\). In passing to
the above mixed evaluation, the position vectors belonging to \(S\)
are integrated out by partial Fourier transformation. They are the
first row and the first column of the middle rows. The momentum
vectors belonging to \(S\) require zero in the last row and the last
column of the middle rows. Translation by \(\bar l\,y\otimes x'\)
puts a nonzero entry in the first position of the last row; its other
idempotent component contributes only a Fourier phase. Thus, for
middle position matrix \(Z\), the original test is evaluated solely
on matrices of the form
\[
   \begin{pmatrix} *&*&*\\ *&Z&0\\ b&0&0\end{pmatrix},\qquad
                                      b\in F_{v_1}^{\times}.
\]
Changing the chosen idempotent uses the corresponding conjugate
bases and the same pattern. Clearing the middle entries of the
first column by the last row bounds the full rank by
\(2+\operatorname{rank}Z\). The upper supported test, also after
the compact source action by invertible row operations, vanishes
on all these matrices if \(\operatorname{rank}Z<a'-2\).
Every partial Fourier integral then vanishes at \(Z\), proving
the smaller support assertion.

At infinity an adapted Cartan involution preserves the orthogonal
decomposition into the tensor blocks formed from the two planes
and their complements. In particular the smaller symplectic tensor
is a complex-orthogonal summand for the compatible positive
structure. The Fock vacuum factors over this summand and its
complement; the polynomial module is spanned by finite products of
their creation monomials. Compact source translates of a polynomial
test remain in the finite span of fixed degrees. The evaluation and
translation in the integral take place only in the complement.
The smaller real output is therefore polynomial for polynomial
input, in the coordinates transported by the smaller Cartan
structure. The moments in the complement converge by the same
Gaussian bounds.

Furthermore fixing a nonzero common pair of real matching maps fixes
the infinite component of \(\nu\). Given two allowed systems for
these same maps, their primed source actions on a compact unitary
block at one real place differ by the determinant character
\(\det^{(m_\nu-m_{\widetilde\nu})/2}\), by the doubled splitting
calculation in Proposition~\ref{threal:allreal}. Both actions leave
the same quotient map unchanged. Apply that map to a polynomial
with nonzero image and translate by this compact block. The
determinant character is identically one since the block determinant
ranges over the circle, so \(m_\nu=m_{\widetilde\nu}\) there.
The restriction of each character to \(\mathbb A_F^\times\) fixes
its radial part at infinity. Varying the place proves the claim.
In particular the real Mellin operator itself is fixed across
the allowed characters and actions for fixed real maps, after the
same coordinates and Haar measures are chosen.

The source line Levi \(m(l,1)\) has determinant modulus
\(|l|_E^d\) on \(S\), since it scales the first summand by \(l\).
The corresponding evaluation character is \(\nu(l)|l|_E^{d/2}\).
Indeed first split the source plane, with modulation space
\(Ex\otimes W\) and its prescribed Siegel character \(\nu\),
then split the target plane in \(V'_1\otimes W\). On the source
unipotent radical the line Levi acts on \(a'-2\) coordinates over
\(E\) with determinant modulus \(|l|_E^{a'-2}\), and on the central
\(F\)-coordinate by \(l\bar l\), adding the factor \(|l|_E\).
Its inverse modular character in Iwasawa integration is therefore
\(|l|_E^{1-a'}\); together these give \(d/2+1-a'=k+1\).
The target line has determinant modulus \(|t|_E^{a'-2}\) on \(S\):
its action on the full \(W\) in the first summand has modulus one.
At evaluation its scalar is \(\alpha(t)|t|_E^{(a'-2)/2}\), from
the residual target Siegel action (keep the target splitting
character on reducing the source by a plane). Moving this target
element past the translation replaces \(l\) by \(l/\bar t\).
In the Mellin integral substitute \(l=u\bar t\). Since
\(\nu(\bar t)=\nu(t)^{-1}\), the resulting phase is
\(\alpha(t)/\nu(t)=\beta(t)\) and the exponent is
\((a'-2)/2+k+1=a/2\), as claimed. This computation explains both
the choice of the translation line and the shift by one in this
branch.
\src{src/Development/AbelianVarietyTheory/Definitions/Divisors/CTSmallCone.lean}{57}{The two-plane evaluation, cone and translation};
\src{src/Development/WeilBound/ThetaLifts/SmallCones/ConeLawAB.lean}{115}{The primed two-Levi character}.
\end{proof}

The two fractional-ideal exponents in this branch can also be fixed.
By adaptedness the line \(Ex\) has lattice \(\mathfrak b x\) and
the line \(Ex'\) has lattice \(\mathfrak c x'\). The source lattice
is Hermitian self-dual and the two lines and the orthogonal smaller
lattice split integrally. Pairing with \(\mathfrak b x\) therefore
shows that the lattice of \(Ey\) is \(\bar{\mathfrak b}^{-1}y\).
It follows by taking direct summands in the tensor lattice that the
finite translations \(\bar l\,y\otimes x'\) belonging to it have
\[
                l\in\mathfrak a\widehat{\mathcal O}_E,qquad
                \mathfrak a=\mathfrak b^{-1}\bar{\mathfrak c}.
\]
Both ideals \(\mathfrak b,\mathfrak c\) are prime to the conductor
of \(\nu\). At the chosen half of a ramified split place, this
follows because both paired vectors occur in the specified integral
basis. At the conjugate half, apply duality to the second vector in
that basis; the target's trace codifferent is a unit at the
designated totally split rational primes. At unramified places
the value of \(\nu\) on a uniformizer is independent of its unit
multiple, and \(\nu(\bar l)=\nu(l)^{-1}\). Consequently
\(\nu(\mathfrak a)=\nu(\mathfrak b)^{-1}
\nu(\mathfrak c)^{-1}\). Integrating the line-lattice indicator
outside the conductor contributes this character value, its absolute
norm to the power \(-(k+1)\), and the elementary Euler factors.
The norm power depends only on the fixed lattice and ranks.
Thus the lattice calculation gives \(e_1=e_2=-1\).
\src{src/Development/WeilBound/ThetaLifts/LocalTheory/ConeShifts/IdealId0.lean}{49}{The character on the translated lattice line}.

\begin{proposition}[Unfolding the supported nonzero orbit]
\label{tpd:unfold}
Suppose \(q>0\). In the rational mixed models the
\(\beta|\cdot|_E^{a/2}\) component of the line constant term of
any supported primed test equals
\[
 C_1\,\beta(t)|t|_E^{a/2}
 I'_1\bigl(T(\omega'(1,h)\Phi')\bigr)(h_1).
\]
Here the line Levi element is \(m(t,h_1)\). The difference from
the entire constant term transforms by \(\alpha|\cdot|_E^{a'/2}\).
It vanishes identically if \(p=0\).
For self-dual Fourier changes \(C_1>0\) depends only on the quotient
and Iwasawa measures and the fixed spaces. It is independent of the
characters and the matching maps.
\end{proposition}
\begin{proof}
We first compute at the target identity for a supported test
whose zero evaluation in the target-first model vanishes identically.
Write this condition as \(e_{W,0}\Phi'=0\). When \(p=0\) the real
average of Lemma~\ref{tpd:primezero} supplies such a test and leaves
both the upper integral and \(T\Phi'\) unchanged, the latter by
right invariance of Haar measure on \(K\).
Put \(Y=V'\), \(Y_1=V'_1\). At a fixed
source element \(g\), integration over the target unipotent center
in the target-first model leaves only null vectors in \(Y(E)\).
Their nonzero terms group by their rational lines. Fix one such
line \(Ex\); its nonzero vectors are \(u x\), \(u\in E^\times\).
Having chosen its companion \(y\), polarize the remaining tensor
\(Y\otimes_EW_1\) using \(Ex\otimes_E W_1\) as modulation space
and a fixed position space for \(Y_1\otimes_EW_1\).
For this vector \(ux\), the remaining target radical acts through
the Heisenberg modulation paired with \(ux\otimes w\),
\(w\in W_1(\mathbb A_F)\). Integrating it over the additive quotient
sets the corresponding rational \(W_1\) position to zero, since
\(u\ne0\). This is the character orthogonality in
\cite[(3.15)--(3.20)]{wuThetaTower}. The line contributes
\(\sum_{u\in E^\times}\sum_{z_1} A_g(u,z_1)\), where \(z_1\)
runs over the rational smaller positions and \(A_g\) is the
restriction of the resulting Schwartz test to the indicated positions.
The source element acts on the test before these operations.

There are now two natural polarizations of the removed coordinates.
The modulation space of this last evaluation is
\[
 S_A=Y\otimes Ex'\ \oplus\ Ex\otimes W_1 .
\]
Its common subspace with \(S\) of Lemma~\ref{tpd:mellin} has
codimension one over \(E\) in both. The two remaining dual directions
are \(y\otimes x'\) and \(x\otimes y'\). Thus replacing \(S_A\)
by \(S\) is a partial Fourier transform in a single \(E\)-coordinate
and leaves \(z_1\) unchanged. Explicitly, on restricting all common
evaluation coordinates to zero, its Fourier transform is
\[
 B_g(l,z_1)=[e'\rho(\bar l\,y\otimes x')\omega'(g,1)\Phi'](z_1).
\]
The trace pairing on the two lines is a nondegenerate rational
pairing; a fixed rational change of variable or sign permutes
\(E^\times\) and has adelic modulus one. With self-dual measures
Poisson summation therefore identifies the sums over \(E\) of
\(A_g\) and \(B_g\), separately at each \(z_1\).
Both missing zero terms vanish. For \(A_g(0,z_1)\) this follows
from \(e_{W,0}\Phi'=0\), since this evaluation intertwines the
source action before the second polarization and evaluation.
For \(B_g(0,z_1)\) use the
matrix in Lemma~\ref{tpd:mellin} with \(b=0\): its last row is
identically zero at \(v_1\). This remains so for the test translated
by \(g\), since source translations preserve full-rank support.
Consequently the sums over \(E^\times\) agree. At fixed \(g\)
the original grouping over lines is legitimate by absolute
Schwartz summation and compact integration in the target radical.

To unfold these line sums set \(D_x=\operatorname{Stab}_{U(Y)}(x)\).
It is the semidirect product of its line unipotent radical \(N_Y\)
and \(U(Y_1)\). The source group is transitive on its rational
nonzero null vectors by Witt extension. In the new expression put
\(f(g)=\sum_{z_1} B_g(1,z_1)\). This function is invariant under
\(N_Y(\mathbb A_F)\): an element of this group changes \(y\)
by a vector in \(Ex+Y_1\), so the change in the translation lies
in \(S\). Its Heisenberg commutator is trivial because \(Ex'\)
is isotropic; evaluation itself is invariant under the unipotent.
For the last assertion split the source plane first: the unipotent
acts on the smaller oscillator through Heisenberg operators with
parameters linear in its extra position, and its center by a
quadratic phase, all trivial at zero. The following evaluation on
the smaller target plane preserves this invariance.
The \(U(Y_1)(F)\) action preserves the rational summation.
For the source line Levi \(m_Y(l,g_1)\), moving it past
\(\rho(y\otimes x')\) replaces \(y\) by \(\bar l\,y\).
The simultaneous evaluation therefore gives
\[
 f(m_Y(l,g_1)k_0)=\nu(l)|l|_E^{d/2}\,
 \Theta_1'\bigl(g_1,1;
    e'\rho(\bar l\,y\otimes x')\omega'(k_0,1)\Phi'\bigr).
\]
For rational \(l\) the scalar is one. Hence summing over rational
lines and their nonzero scalar parameters gives exactly the
periodization of \(f\) over \(D_x(F)\backslash U(Y)(F)\).

We check absolute integrability of this periodization at the level
of the signed smaller theta sum. Integrate first on
\(D_x(F)\backslash U(Y)(\mathbb A_F)\). Iwasawa decomposition for
the source line and the chosen maximal compact gives, up to fixed
positive normalization, integration over \([N_Y]\),
\([U(Y_1)]\), \(\mathbb A_E^\times\) and \(K\), with density
\(|l|_E^{1-a'}\). The finite compact is in good position: it is
hyperspecial for the self-dual source lattice at the inert places,
and is an integral general linear compact at the split places.
The quotient \([N_Y]\) is compact with fixed volume. The absolute
integral remaining is bounded by the integral of the absolute
smaller theta kernel against \(|l|_E^{k+1}\). At the distinguished
finite place each smaller test in this integral has full-rank
support by Lemma~\ref{tpd:mellin}. More explicitly the compact
locally constant test in \((l,z_1)\) has only finitely many slices
in the first variable at each exceptional finite place. At \(v_1\)
all slices occurring for invertible \(l\) have the required support.
At infinity all seminorms needed in Proposition~\ref{threal:decay}
for the smaller source are rapidly decreasing in \(l\) and uniformly
bounded near zero and in \(k_0\). The smaller ranks satisfy its
hypotheses, since \(q\) has decreased by one and \(a'\) by two.
The absolute Euler products outside the exceptional primes are
bounded by the corresponding factors of \(\zeta_E(k+1)\), as
in Lemma~\ref{tpd:mellin}. These bounds prove the required
integrability, justify unfolding the periodization, and allow the
smaller period to be interchanged with the Mellin and compact
integrals. This gives \(C_1 I'_1(T\Phi')(1)\).

Write \(A(t)=\alpha(t)|t|_E^{a'/2}\),
\(B(t)=\beta(t)|t|_E^{a/2}\), and let \(J(\Phi')\) be the entire
constant term at the identity. For arbitrary \(p\) and arbitrary
supported \(\Psi\), apply the computation just proved to
\[
   \Phi'_t=\omega'(1,m(t,1))\Psi-A(t)\Psi .
\]
This test satisfies \(e_{W,0}\Phi'_t=0\): in the target-first
model scalar substitution on the \(a'\) coordinates at zero gives
exactly \(A(t)\). Both summands retain the required split support.
Writing \(J_B(\Psi)=C_1 I'_1(T\Psi)(1)\), the result is
\[
 J(\omega'(1,m(t,1))\Psi)-A(t)J(\Psi)
                   =(B(t)-A(t))J_B(\Psi),
\]
using the Mellin covariance in Lemma~\ref{tpd:mellin}. Thus
\(J-J_B\) transforms by \(A\), whereas \(J_B\) transforms by \(B\).
These components are unique since \(a>a'\) and the splitting
characters are unitary. If \(p=0\), averaging any test first shows
directly that \(J=J_B\). Taking \(\Psi=\omega'(1,h)\Phi'\) and
using intertwining for the smaller target factor gives the asserted
formula at \(m(t,h_1)h\). All measure constants occurred in the fixed
Haar decompositions or unipotent quotients.
\end{proof}

\begin{lemma}[A supported representative for the other coefficient]
\label{tpd:so}
Let \(q>0\). For each supported primed test \(\Psi\), one may
replace \(\Psi\) by a supported \(\widetilde\Psi\) having the
same upper period, so that its one-plane zero evaluation
\(e_{W,0}\widetilde\Psi\) is full-rank supported for the unchanged
source \(V'\) and the target \(W_1\). In the coherent one-plane
model the \(\alpha|\cdot|_E^{a'/2}\) component at \(m(t,h_1)\)
is
\[
   \alpha(t)|t|_E^{a'/2}
       I_{V',W_1}(e_{W,0}\widetilde\Psi)(h_1).
\]
For a fixed further right target translate apply the construction
to the already translated upper test.
\end{lemma}
\begin{proof}
Put \(m=a'\) and work at \(v_1\) in the primed row coordinates.
Choose dual target bases with paired first and last columns. In
this model \(e_{W,0}\phi(Z)=\int\phi(u,Z,0)\,du\), with the
corresponding conjugate choice if the position idempotent is
interchanged. In fact one of the removed columns acts by translations
and is integrated, while the paired modulation column must vanish.
An upper test of full row rank therefore maps to a compactly
supported function on the open region \(D=\{\operatorname{rank}
Z\ge m-1\}\). The map is onto that space. At a given \(Z_0\in D\)
choose \(u_0\) completing its columns to rank \(m\), and small
compact open balls in \(u,Z\) on which rank stays full. A unit-mass
indicator in the first ball and a small last-column indicator of
value one at zero lift every test supported on the second ball.
A finite disjoint refinement covers the support of any test in \(D\).

Modulo compactly supported functions on the open full-rank part
of \(D\), these lower tests restrict to all compactly supported
functions on its relatively closed rank \(m-1\) stratum. Extension
follows by compact open local charts and disjoint refinement.
The stratum's source coinvariants are zero by the rank stabilizer
calculation of Proposition~\ref{split:kernel}, applied with
\((d-2)/2=m+k-1>m\). This calculation concerns the stratum itself
and applies with the same unitary phase and the lowered modulus.
Thus, up to a full-rank lower test, \(e_{W,0}\phi\) is a finite sum
of lower source-action differences of tests in \(D\). Lift each
such test through the surjection above, and subtract the
corresponding upper differences from \(\phi\). Equivariance of
evaluation for the unchanged source proves the required support.
Use individual supported local factors and extend by linearity;
at the other places all factors stay fixed. These upper differences
integrate to zero by right translation on the source, with absolute
convergence from Proposition~\ref{threal:decay}.

In the expansion of the upper constant term the zero-vector
contribution of this representative now converges by the
full-rank estimate applied to \((V',W_1)\). Its torus character is
\(\alpha(t)|t|_E^{m/2}\); at \(t=1\) evaluation gives precisely
the smaller displayed period, using the same source measure.
For each fixed source translate the \(A_g(0,z_1)\) of
Proposition~\ref{tpd:unfold} vanishes even though
\(e_{W,0}\widetilde\Psi\) need not be zero. In fact apply the
further source-plane evaluation \(Ex\otimes W_1\) at the split
place to this lower full-rank test. All entries in one row must be
zero and the opposite row is integrated, so every value in the
integrand vanishes. This remains true for any rational null line,
by changing row basis, and for translates by the unchanged source
and smaller target, which preserve matrix rank. The other Poisson
zero term still vanishes on the upper full-rank test. The entire
nonzero orbit therefore has the \(\beta\) law by the same absolutely
convergent unfolding. Uniqueness of the two characters proves the
assertion, including independence of the resulting component from
the choice of the lift of the source differences.
\end{proof}

At a designated split place the upper to lower comparison space is
at most one-dimensional by Proposition~\ref{threal:localcone}.
Its unprimed zero-evaluation map is nonzero, since evaluation is
onto and the lower quotient is nonzero. We next compute the scalar
on one pair of tests with nonzero smaller image. This determines
the comparison on every test at this place.

\begin{proposition}[The ramified primed comparison on all tests]
\label{tpd:ram}
At a designated ramified split place let \(K_0=F_v\) have residue
size \(Q\), let \(\chi=\nu_w\) on its chosen factor of \(E\),
and let \(c>0\) be the conductor exponent of \(\chi\).
Use the prescribed bases and row polarizations, \(\operatorname{vol}
(\mathcal O)=1\) for the additive character of conductor \(\mathcal
O\), and give \(\mathcal O^\times\) and the maximal source compact
masses one. Use the integral-vacuum-normalized evaluations. On
every upper pair matched at this place, the ratio of the lower
primed quotient on \(T_v\phi'\) to the lower unprimed quotient on
\(e_{0,v}\phi\), as a scalar identity of maps, is
\[
   \epsilon_v(-k)=
      Q^{ck}\chi(\pi)^c
       \sum_{r\in(\mathcal O/\pi^c)^\times}
                         \chi(r)^{-1}\psi_v(r/\pi^c).
\]
\end{proposition}
\begin{proof}
Put \(m=a'\) and \(D_\chi(u)=\mathbf1_{\mathcal O^\times}(u)
\chi(u)^{-1}\), extended by zero at \(u=0\).
For an \(n\)-by-\(d\) matrix choose a block of \(n\) pivot columns.
Let its test be \(D_\chi\) of the pivot determinant times the
indicators that the pivot entries belong to \(\mathcal O\) and
all other entries to \(\pi^c\mathcal O\). Take the first \(a\)
columns for \(\phi\), the last \(m\) for \(\phi'\), and the same
rule for \(\phi_1,\phi'_1\) after deleting the first and last of
the \(m\) target columns. The source compact fixes each test:
integral invertible row operations preserve its entry conditions,
and the source factor cancels the inverse determinant character.

The upper tests are matched, as are the two lower tests. To check
this directly, restrict the orbital functions
\eqref{newtheta:big}--\eqref{newtheta:small} to
\(h=\left(\begin{smallmatrix}A&B\\ C&D\end{smallmatrix}\right)
\in\mathrm{GL}_d(\mathcal O)\).
A contributing row transformation must belong to the integral
general linear group: the original row frame is primitive and
the pivot of the transformed test is invertible modulo \(\pi\).
On both sides support requires \(C=0\bmod\pi^c\), in which case
the diagonal blocks are invertible. The determinant character of
the row transformation cancels throughout that integral compact.
The two values are \(\beta_w(\det h)\chi(\det A)\) and
\(\alpha_w(\det h)\chi(\det D)^{-1}\), equal by
\(\det h\equiv\det A\det D\bmod\pi^c\). Iwasawa decomposition
extends equality using their common parabolic law. The common value
at the identity is one. The same calculation in the smaller ranks
gives a nonzero common lower quotient vector.

Unprimed evaluation integrates the \(a\) entries of one nonpivot
column, setting the other removed nonpivot column to zero. Thus
\(e_{0,v}\phi=Q^{-ac}\phi_1\). For the primed calculation write
\(l=(s,t)\in K_0^\times\times K_0^\times\), the first factor
being the chosen place \(w\). Order the pivot rows and columns as
first, middle, last in the prescribed bases. In the conjugate row
polarization the position and modulation parts of
\(\bar l\,y\otimes x'\) are respectively the first-row last-pivot
entry \(s\) and the last-row first-pivot entry \(-t\).
Indeed switching from the ordinary to the conjugate row polarization
interchanges momentum and position with one minus sign; the source
pair has Hermitian off-diagonal entries \(1,1\), the target pair
has skew entries \(1,-1\), and conjugating \(l\) swaps its two
components. These coordinates occupy distinct matrix entries, so
the half-pairing in the Heisenberg action is zero.
\src{src/Development/WeilBound/ThetaLifts/LocalTheory/StandardValuations/AdelicInclusionMapPlaceSection.lean}{228}{The two entries for the specified Fourier signs}.

In this coordinate ordering the primed mixed evaluation integrates
the last row and the last pivot column in the middle rows. The
other removed entries vanish except for the position entry \(s\).
This is the matrix calculation of Lemma~\ref{tpd:mellin}, expressed
in the prescribed conjugate bases. The resulting pivot matrix is
\[
     \begin{pmatrix}0&0&s\\0&D_1&*\\ x&*&*\end{pmatrix},
       \qquad \det=-sx\det D_1 ,
\]
and the Fourier factor is \(\psi_v(-tx)\). The \(a\) integrated
nonpivot entries of the last row have joint volume \(Q^{-ac}\).
The other integrated entries, apart from \(x\), range independently
over \(\mathcal O\) of mass one. Since \(s,x\) and the entries
of \(D_1\) must be integral, the determinant is a unit precisely
when each of \(s,x,\det D_1\) is a unit. The minus sign in the
determinant is removed together with the Fourier minus sign by
\(x\mapsto-x\). The evaluated test at \(l\) is therefore
\[
 Q^{-ac}D_\chi(s)\widehat D_\chi(t)\phi'_1,\qquad
 \widehat D_\chi(t)=
       \int_{\mathcal O^\times}\chi(x)^{-1}\psi_v(tx)\,dx .
\]
The complementary place has character \(\chi^{-1}\), so the Mellin
weight is \(\chi(s)\chi(t)^{-1}|st|^{k+1}\).
Integration in \(s\) gives one. The Fourier transform vanishes
outside \(v(t)=-c\): below this shell additive integration
cancels on cosets modulo \(\pi^c\), and above it multiplicative
translation by a unit trivial on the additive factor but
nontrivial under \(\chi\) forces zero. At \(t=\pi^{-c}u\) its
value is \(\chi(u)Q^{-c}G(\chi,\psi_v)\) by finite residue
summation. Integration on that multiplicative shell therefore
gives \(\chi(\pi)^c Q^{ck}G(\chi,\psi_v)\).
The source compact averaging has no further effect on the chosen
test. The factors \(Q^{-ac}\) on the two sides cancel, and
Proposition~\ref{threal:localcone} extends the resulting scalar
identity to all matched upper tests. The nonvanishing of the
Gauss sum is Lemma~\ref{newtheta:gauss}.
\end{proof}

In the other descent Corollary~\ref{threal:ramA} gives the
corresponding factor \(\epsilon_v(1-k)\) with the same unit sign.

In this unprimed descent the integral line of \(l x\otimes y'\)
has coefficient ideal
\(\mathfrak b\,\mathfrak D_E^{-1}\bar{\mathfrak c}^{-1}\),
where \(\mathfrak D_E^{-1}\) is the absolute codifferent. Indeed
the trace-perfect target lattice has this ideal as its duality
range for the skew-Hermitian form; pairing \(y'\) against
\(\mathfrak c x'\) gives the indicated factor. Since \(E/F\)
is unramified at finite places the absolute codifferent extends
from \(F\), and its character value is one by
\(\nu|_{\mathbb A_F^\times}=1\). The two ideal exponents in the
unprimed branch are therefore \(+1,+1\), and the preceding
Fourier computation gives \(s=+1\) also here.

\begin{proposition}[The comparison at all relatively split good places]
\label{tpd:ordinary}
Let \(v\notin S\) split in \(E/F\), and choose either factor as
the first copy of \(K_0=F_v\). Use the integral-vacuum-normalized
local evaluations, compact mass one, and unit mass one on each
multiplicative group. Write \(z=k+1\) in the primed descent and
\(z=k\) in the unprimed descent. The lower quotient of the local
cone operator is the following scalar times the lower quotient of
the opposite zero evaluation, on every upper matched pair:
\[
   \nu_v(\mathfrak a_v)\,|\mathfrak a_v|_v^z L_{E,v}(\nu,z).
\]
Here \(\mathfrak a_v\) is the corresponding translation-line ideal
computed above, its modulus is that of an idele generating its two
local factors, and \(\nu_v\) is evaluated on such a generator.
The assertion includes \(\nu_w(\pi)=1\).
\end{proposition}
\begin{proof}
The source compact fixes the vacuum at this unramified split place:
in integral row bases its action is linear substitution with
unramified determinant character. In the two-plane evaluation the
translated vacuum maps to the smaller vacuum for
\(l\in\mathfrak a_v\) and to zero otherwise. To see the normalization,
use the two idempotent position and momentum groups, as in
Proposition~\ref{threal:ordinarysplit}. The vacuum is proportional
to the position lattice indicator, and its modulation lattice is
the character dual. Adaptedness decomposes both lattices into the
paired hyperbolic directions and the smaller tensor. Each exchanged
position entry is integrated by partial Fourier transformation at
zero. Translation in the nonintegrated entry leaves its indicator
nonzero precisely on that entry's lattice. Modulation on the
integrated entry gives its original integral on the character-dual
lattice and zero elsewhere. The two entries are distinct, so the
half-pairing vanishes also in residue characteristic two. Together
the two conditions say exactly that the displacement belongs to
the tensor lattice, equivalently \(l\in\mathfrak a_v\). All other
entries give the untranslated evaluation, normalized to the smaller
vacuum. This uses the character dual for the actual additive character,
including its codifferent. Integrating the two multiplicative shells
gives the displayed scalar. The opposite one-plane evaluation sends
vacuum to vacuum.

Let \(u=\nu_w(\pi)\) vary on the unit circle, keeping \(\beta_w\)
fixed and setting \(\alpha_w=\beta_w\nu_w\) with \(\nu_w\) unramified.
On the other split factor take the inverse characters. This family
is local. For \(u\) of infinite order the full source coinvariant
calculations of Propositions~\ref{split:kernel} and~\ref{split:good}
apply over \(K_0\) itself. Their proofs use the local orbital
functions and the infinite-order condition. The quotient-integral
and Jacquet argument of Proposition~\ref{threal:localcone} also uses
only those local facts and the split line decomposition. The smaller
zero evaluation is onto. Thus the two maps from the upper common
quotient to the smaller quotient are proportional for each such
\(u\). The vacuum calculation determines the scalar.

To extend the identity to \(u_0\) on the circle, realize the upper
and lower normalized induced targets by restriction of their functions
to the integral compact. The compact pictures are independent of
\(u\), since all varying inducing characters are unramified. Use
the orbital maps normalized to send both vacua to the spherical
vector of value one. For each fixed test they are continuous in
\(u\) on the circle. For the larger row map this follows by clearing
the determinant denominator as in Proposition~\ref{threal:ordinarysplit}.
The pivot-chart proof there works with the unramified value as a
variable: integration on a ball at zero gives a rational function
with denominator the indicated subproduct of
\(\prod_{j=0}^{n-1}(1-u Q^{j-d/2})\); a pivot of fixed valuation
contributes a Laurent monomial, and the other charts give the
lower-dimensional subproducts. Here \(Q\) is the residue size
and \(n\) the upper row rank. Clearing this denominator and normalizing
the lattice value gives Laurent polynomials in \(u\) up to a nonzero
fixed constant. A basis change with nonunit determinant introduces
a further Laurent monomial. The smaller row map is absolutely
convergent at its parameter and its spherical value has no zero or
pole on the circle. Use the corresponding dimensions and ranks in
the two lower maps. Compact-open invariance reduces their restrictions
to finitely many compact-picture evaluations at fixed level. Thus
continuity includes the normalized residue at a pole.

The zero evaluation is independent of \(u\). The local cone operator
also varies continuously on each test in a fixed finite-dimensional
span of smaller tests. The compact source action is independent of an
unramified change of character, and the restrictions to the two
displacement coordinates are compactly supported locally constant
functions with finitely many slices. On each compact open slice
away from zero the shell integrals are finite sums; near zero the
geometric series converge uniformly since \(z>0\). This applies
independently to the two local factors of \(E\).

For tests \(x,y\) with equal normalized upper images at \(u_0\)
denote the two upper maps by \(P_u,P'_u\). Choose a compact open
subgroup \(J\subset K_d\) fixing their images for every \(u\);
the upper compact actions on the tests are independent of \(u\).
The \(J\)-fixed space of the compact picture is finite-dimensional
and independent of \(u\). By surjectivity of the smaller-row map
\(P'_{u_0}\) from Proposition~\ref{threal:ordinarysplit}, choose
test preimages of a basis and average them over \(J\).
Surjectivity at the critical parameter uses the common irreducible
normalized unitary induction and the nonzero spherical image.
The matrix of \(P'_u\) on the span of these averaged preimages
stays invertible near \(u_0\); denote its inverse by \(H_u\).
The tests
\[
             y(u)=y+H_u(P_u x-P'_u y)
\]
are matched to \(x\) there, and converge to \(y\) in a fixed
finite-dimensional test space. For \(v=v_1\) the argument can
first be made on the full spaces and then restricted to supported
pairs. Apply the equality at infinite-order parameters tending
to \(u_0\) and pass to the limit on the two lower compact pictures.
The computed vacuum scalar is continuous as well. This only uses
unitary parameters.

Finally any allowed irreducible common quotient at this place is
isomorphic to the canonical one, and each map to it for trivial
source is unique up to scalar by the split local Howe theorem.
\cite[Theorem~1 and Corollary~6.3]{minguezHowe}
Equality of the two nonzero vacuum images forces the two scalar
changes to agree, in both the upper and lower models. Thus the
identity holds for every allowed choice of quotients.
\end{proof}

For comparison with the normalized recursion, denote the smaller
periods by \(I_1,I'_1\) and the primed, nonzero-cone test by
\(T\Phi'\). The comparison needed at the identity of the line
torus has three constants. With fixed models and measures they are
specified by the following identities, in which \(b,b'\) denote
the upper components with character \(\beta|\cdot|_E^{a/2}\):
\[
 b=C_0 I_1(e_0\Phi),\qquad b'=C_1 I'_1(T\Phi').
\]
The third requirement is that the following pair be matched on the
smaller spaces, with its second entry supported:
\[
 \big(C_2 F_X(-k)L^{S_0}(\nu,k+1)e_0\Phi,\ T\Phi'\big).
\]
If these identities hold and the smaller scalar is \(c_1\), then
\[
      C_0 c_1 b'=C_2 C_1 F_X(-k)L^{S_0}(\nu,k+1)b .
\]
Proposition~\ref{threal:endB} gives detection of the functions.
The upper scalar under the third requirement is therefore
\[
 c=\frac{c_1}{C F_X(-k)L^{S_0}(\nu,k+1)},\qquad
                  C=\frac{C_2C_1}{C_0}.
\]
Proposition~\ref{tpd:unfold} supplies the second identity for every
\(q>0\). For \(p>0\) on the unprimed side,
Proposition~\ref{threal:coneA} places the entire nonzero orbit in
the other character; the zero orbit gives the first identity by
evaluation, with a convergent smaller unprimed period by the
estimates of Proposition~\ref{threal:decay} with \(d\) decreased
by two. Thus \(C_0=1\) throughout for coherent mixed models. The
remaining test comparison is the third identity on the two smaller
matching quotients with one common normalization.
Propositions~\ref{tpd:ram} and~\ref{tpd:ordinary} give it at the
ramified and all relatively split unramified places, including the
unit sign \(s=+1\) and the factors of the two adapted line ideals.
The vacuum calculation and the partial Fourier transform supply the input for Propositions~\ref{tcr:descend} and \ref{tcr:unique}, which extend the inert identity to all smooth tests and compare the real operators on the polynomial module. Theorem~\ref{tcr:global} assembles their normalized tensor product.
\src{src/Development/WeilBound/ThetaLifts/LocalTheory/StandardValuations/ExistsVol.lean}{162}{Evaluation on a ramified matching pair};
\src{src/Development/WeilBound/ThetaLifts/LocalTheory/LocBOf.lean}{17}{Passage from these tests to local matching}.
The required \(C\) is common to the admissible characters and
the choices of finite matching maps.
\src{src/Development/AbelianVarietyTheory/Definitions/Divisors/CTSmallCone.lean}{239}{The two constant-term evaluations};
\src{src/Development/WeilBound/ThetaLifts/ConeHypotheses/KeyAOf.lean}{37}{Assembly of the three constants};
\src{src/Development/WeilBound/ThetaLifts/ConeHypotheses/PolDense.lean}{46}{The supported smaller matching identity}.

\paragraph{Removing a hyperbolic plane on the given lattice.}
The integral part of the change of model admits the following direct
construction. It will allow the local descent to use a rational plane
while retaining the lattice at every finite place.

\begin{proposition}\label{famtheta:plane}
Let \(E/F\) be the CM extension in the theta tower, unramified at
every finite place. Let \(L\) be a lattice in an \(E\)-space of
dimension \(n\ge3\), with nondegenerate Hermitian or skew-Hermitian
form \(b\), linear in its first argument. Suppose \(L\) is self-dual
for an invertible conjugation-stable fractional ideal \(\mathfrak d\)
and the form has both signs at each infinite place (after multiplication
by an imaginary scalar in the skew case). There are rational isotropic
vectors \(x,y\) with \(b(x,y)=1\) for which
\[
 L=(L\cap Ex)\ \oplus\ (L\cap Ey)\ \perp\ L_0 .
\]
Here the first two lines are paired perfectly with values in
\(\mathfrak d\), and \(L_0\) is self-dual for \(\mathfrak d\).
At any prescribed finite set of split places where
\(\mathfrak d_w=\mathcal O_{E,w}\) on both sides, the two vectors
may simultaneously be taken as local generators of their lines.
\end{proposition}
\begin{proof}
In the skew case multiply \(b\) temporarily by a nonzero element
negated by conjugation. The associated \(F\)-quadratic form
\(v\mapsto b(v,v)\) for the resulting Hermitian space is nondegenerate
of dimension \(2n\). It is indefinite at every real place and isotropic
at every finite place, since quadratic forms over a finite extension
of \(\mathbf Q_p\) of dimension at least five are isotropic (also for
\(p=2\)). Hasse--Minkowski supplies a rational isotropic \(E\)-line
\(U\) for the original form. The intersection \(I=L\cap U\) is
saturated and of rank one. Both \(I\) and \(L/I\) are torsion-free
finite modules over the Dedekind ring \(\mathcal O_E\), hence projective.
In particular \(I\) is a direct summand as a module.

Put \(J=\operatorname{Hom}_{\mathcal O_E}(\bar I,\mathfrak d)\),
where the bar denotes the conjugate module. Self-duality identifies
\(L\) with \(\operatorname{Hom}(\bar L,\mathfrak d)\) by the pairing.
Restriction to \(\bar I\) is onto by the direct summand property.
Choose an integral section \(s:J\to L\) of this restriction,
using projectivity of \(J\). The pairing of \(sJ\) and \(I\)
is perfect in both directions, by biduality of the invertible modules
and conjugation stability of \(\mathfrak d\).
Choose \(\theta\in\mathcal O_E\) of trace one. Indeed the trace ideal
is the unit ideal locally at every finite prime: at a split prime
this is addition, and in an unramified quadratic extension it follows
by reduction from the surjective residue-field trace. This argument
also covers residue characteristic two.
There is a unique linear map \(C:J\to I\) with
\[
 b(Cj,sj')=\theta\, b(sj,sj')\qquad (j,j'\in J)
\]
by the perfect pairing. Hermitian symmetry with its given sign implies
\(b(sj,Cj')=\bar\theta b(sj,sj')\). Since \(b(I,I)=0\), the map
\(s'=s-C\) is another section and \(b(s'J,s'J)=0\).
The Gram pairing on \(I\oplus s'J\) has zero diagonal blocks and
perfect off-diagonal blocks. It is self-dual for \(\mathfrak d\).
For any \(z\in L\), its pairing functional on this submodule is thus
realized by a unique vector inside it. Subtracting that vector gives
the orthogonal projection to \(L_0\). This proves the integral
decomposition and self-duality of the complement. The two summands
of the plane are precisely the full lattice intersections with their
rational spans.

Choose a rational generator \(x\) of \(I\otimes E\); by weak
approximation scale it to generate \(I\) at all the indicated primes.
There is a unique \(y\in(s'J)\otimes E\) with \(b(x,y)=1\). Perfection
and \(\mathfrak d_w=\mathcal O_{E,w}\) make it a generator at the
same places, including the conjugate primes. Complete these two vectors
by bases of \(L_0\) at the split places and put them in the designated
columns. At each real place the plane has signature \((1,1)\);
choosing a Cartan involution on it and on its orthogonal complement
gives the adapted marking. The lower signatures lose one of each sign.
\end{proof}
\src{src/Development/WeilBound/ThetaLifts/TowerModels/PlaneMain.lean}{47}{rational isotropy and integral trace one},
\src{src/Development/WeilBound/ThetaLifts/TowerModels/PlaneSkew.lean}{18}{skew form and codifferent}, and
\src{src/Development/WeilBound/ThetaLifts/TowerModels/TowerCodimensionTwoOf.lean}{32}{marked planes for the two towers}.

\subsection{Local descent on an adapted pair of theta towers}
\label{tcr:main}

Write \(d=a+a'\), \(k=(a-a')/2\), and \(s=p+q\). In the
\(A\)-step take \(q=0<p\), remove a plane from \(V,W\), and put
\(z=k,\ \xi=\alpha|\cdot|_E^{a'/2}\). In the \(B\)-step take
\(q>0\), remove a plane from \(V',W\), and put
\(z=k+1,\ \xi=\beta|\cdot|_E^{a/2}\).
The one-plane evaluation \(e_0\) belongs to the unchanged source
and the two-plane cone operator \(T\) to the changed source. Subscripts
on these maps denote their local factors. The adaptedness assumptions
are those of \S\ref{newtheta:global}: the two integral lattices
split into their two lines and their orthogonal complements, the
finite designated bases have the prescribed columns, and the Cartan
involutions pair the two lines and restrict to the lower involutions.

\begin{lemma}[Square integrability and kernels]\label{tcr:kernels}
Each upper polynomial period, with the stipulated primed support,
belongs to \(L^2([U(W)])\). For any of the matching maps, the period
on each side factors through the restricted tensor product of its
small local theta quotients. A nonzero such period map is injective
on that tensor product.
\end{lemma}
\begin{proof}
The periods are smooth, of uniform moderate growth, finite under the
enveloping center and compact-finite at infinity, as in the proof
of Proposition~\ref{threal:detectline}. Apply the cuspidal-exponent
criterion for square integrability, with the constant term divided
by the square root of the modulus of its parabolic
\cite[I.4.11]{moeglinWaldspurger}. The center of the unitary group
has compact adelic quotient.

Consider a cuspidal exponent on a Levi of a flag of isotropic
subspaces. If the first subspace has dimension \(r>1\), its first
Levi block is \(\mathrm{GL}_r(\mathbb A_E)\). Any cuspidal projection
there is generic and unitary after a central real twist. Here a
cuspidal projection is computed by pairing the flag constant term
on \(M(F)\backslash M(\mathbb A_F)^1\), where all rational
character moduli are one, with a rapidly decreasing cuspidal test
on the Levi \(M\). Uniform moderate growth suffices for convergence;
the central height variables remain free for extracting exponents.
Fix the other cuspidal factors and vary the first cusp vector. A
nonzero pairing gives an equivariant functional with values in
the smooth contragredient of that cuspidal factor, using finite
compact invariance of the local test. An action off the norm-one
subgroup is accounted for by a compensating split-central
translation, hence by the extracted central exponent.
At a completely
split unramified prime with \(\nu_w(\pi)\) of infinite order it has
to be a subquotient of the first block in the Jacquet module of the
upper small quotient. Indeed the full flag constant term can be taken
first along the \(r\)-plane and then in its remaining unitary factor.
Integration there and cuspidal projection commute with the first
block action. The local calculation in Proposition~\ref{threal:detectline}
rules out precisely these subquotients for \(r>1\). This argument
uses the full split source-coinvariant calculation at that one
prime and does not require the period to be in \(L^2\).

If \(r=1\), the two line components are smaller periods by
Propositions~\ref{threal:coneA}, \ref{tpd:unfold} and
Lemma~\ref{tpd:so}, including their zero-orbit formulas. A missing
source plane gives no corresponding nonzero orbit. The same
calculation applies to the remaining flag on each contributing
smaller period. For this calculation rational hyperbolic coordinates
suffice: at the places where the resulting bases are not integral
they change finitely many local tests. At infinity the estimates of
Proposition~\ref{threal:decay} apply to Schwartz tests and their
derivatives. Center-finiteness of the successive coefficient as an
automorphic function follows also by Harish-Chandra projection from
the preceding coefficient. Thus every general linear block of a
contributing cuspidal exponent is a line. If there are \(j\) such
blocks, write \(T_i=|t_i|_E\) for their parameters in flag order.
The real exponents \(b_i\) of the unnormalized term satisfy
\(b_i\le a/2\), since at every stage they are half the rank of one
of the two current sources. The square root of the modulus has
exponents \((d-2i+1)/2\): the line radical for a space of dimension
\(d-2i+2\) has determinant modulus \(T_i^{d-2i+1}\) (its middle
module has dimension \(d-2i\) over \(E\), and its central
\(F\)-line has modulus \(T_i\)). Hence
\[
        b_i-(d-2i+1)/2\ \le\ -(a'-2i+1)/2<0
        \qquad(1\le i\le j\le s).
\]
Every nonzero generating ray of the closed positive chamber
\(\log T_1\ge\cdots\ge\log T_j\ge0\) has strictly negative
normalized exponent. The criterion gives the assertion.

Fix the other local tests. Source invariance factors the resulting
map into \(L^2\) through the big theta module of the trivial
representation, which has finite length and a unique simple quotient
with multiplicity one in its cosocle
\cite{howeTranscending,ganTakedaHowe,minguezHowe}. An admissible
finite-length image in a unitary representation is semisimple.
For a submodule take its Hilbert closure: the projection to a
compact-open fixed space at a finite place or a compact isotypic
space at a real place shows that the closure intersects the
algebraic image in the submodule. These spaces are finite-dimensional,
so orthogonal complements taken on them give an algebraic splitting;
at infinity it is stable under derivatives by skew adjointness on
smooth vectors. Thus each local map kills the kernel of the prescribed
small quotient. At the support place \(v_1\) on the primed side one
uses the supported coinvariants, which map isomorphically to this
same small quotient by Proposition~\ref{split:kernel}. Applying
the universal property of a tensor product to finite sums and then
enlarging the finite set of places gives the factorization.

The kernel of the map on the resulting restricted product is stable
under independent finite-place translations and the real Lie and
compact actions. A submodule containing a nonzero vector contains
the full finite tensor product supporting that vector: the individual
simple factors have scalar endomorphisms, and the algebraic density
theorem gives simplicity of their finite external product. Enlarging
the finite set, using the nonzero spherical vectors at the new
places, gives the whole restricted product. This proves injectivity
whenever the period map is nonzero.
\end{proof}

\begin{lemma}[Nonzero smaller images at infinity]\label{tcr:bump}
Suppose the smaller period on the changed-source side is nonzero
on polynomial tests. Then \(P_{1,\infty}T_\infty\) is nonzero on
upper polynomial tests, where \(P_{1,\infty}\) is its real matching
map. For the other side \(P'_{1,\infty}e_{0,\infty}\) is nonzero
whenever its smaller real quotient is nonzero.
\end{lemma}
\begin{proof}
Choose smaller finite tests and a target argument at which the
smaller period is a nonzero functional \(D\) of the smaller real
test. It is continuous for the Schwartz topology by
Corollary~\ref{threal:limits} and source-invariant. Its restriction
to polynomial tests factors through the smaller real quotient by
Lemma~\ref{tcr:kernels}. For \(e_0\), multiply any smaller Fock
polynomial by the vacuum on the complementary summand, whose
evaluation is nonzero. The positive complex structure preserves
this decomposition by adaptedness. Thus the one-plane evaluation
is onto on polynomials.

For \(T\), use the mixed coordinates with modulation space
\(Y\otimes Ex'\oplus Ex\otimes W_1\), where \(Y\) is the changed
source. Before the second evaluation the first position is
\(v\in Y_\infty\). For a real compact source element \(h\) denote
the evaluation at \(v=u x\), followed by the second evaluation
and \(D\), by \(A_h(u)\), where \(u\in\prod_{w\mid\infty}\mathbb C\).
In these coordinates multiplication of the initial test by a
scalar test \(b(v)\) multiplies \(A_h(u)\) by \(b(h^{-1}ux)\).
Indeed the source acts on the first position by pullback and on
the residual oscillator by its Weil action. With \(h=1,u=1\),
the evaluation onto the smaller Schwartz space is surjective
by extension in the omitted linear coordinates.

In the \(A\)-step the scalar under the integral for \(DT_\infty\)
is \(A_h(u)\) times the everywhere nonzero Mellin density on
\((\mathbb C^\times)^{[F:\mathbb Q]}\). In the \(B\)-step one
instead pairs the one-coordinate Fourier transform of \(A_h\)
at each place with the Mellin density of exponent \(k+1\), by
the partial Fourier identity in Proposition~\ref{tpd:unfold}.
We describe the transformed density on \(u\ne0\). Relative to
additive area on \(\mathbb C\), the original density is a nonzero
constant times
\[
          (l/|l|)^{m_w}|l|^{2k},\qquad |m_w|=m_{\alpha,w}-m_{\beta,w}>d .
\]
Complex conjugation of the displacement or reversal of the Fourier
pairing can reverse the angular weight. The Fourier transform of
this tempered distribution has a smooth nowhere-zero density on
\(\mathbb C^\times\). Indeed it is equivariant under rotations and
positive dilations. On their free transitive orbit such a distribution
is a scalar times the corresponding character times Haar density:
divide by the character and use translation invariance in angle and
logarithmic radius. If its scalar vanished it would be supported
at the origin, hence be a finite sum of derivatives of the delta
distribution. Its inverse Fourier transform would then be a
polynomial of homogeneous degree \(2k\) and angular weight of
absolute value \(>2k\), which is impossible by expanding in
\(l,\bar l\). This argument applies separately at every real place.

Here is why integration over the compact directions cannot cancel
this density. In the first mixed coordinates
\((h,u)\mapsto h^{-1}ux\) maps the compact group times the
nonzero scalars onto the product of punctured null cones. At each
place the two compact unitary groups are transitive on the two
unit spheres, and the lengths of the two components of a null
vector are equal. The fibers are the compact line stabilizers,
acting by \((h,u)\mapsto(jh,\eta u)\) for \(jx=\eta x\).
The factor acting on the smaller oscillator is the smaller compact
source and disappears after \(D\). The scalar on this stabilizer
also cancels: for the unchanged radial measure the source determinant
factor on the line is \(\nu^{-1}(\eta)\) in the \(A\)-step and
\(\nu(\eta)\) in the \(B\)-step, as in the source Levi calculations
of Proposition~\ref{threal:coneA} and Lemma~\ref{tpd:mellin}.
In the first chart the weight on \(u\) is \(\nu(u)\) in the
\(A\)-step. In the other step the Fourier transform carries a
rotation of \(u\) to its dual rotation, so the transformed angular
weight is \(\nu^{-1}(u)\). Explicitly the \(B\)-chart functions
after applying \(D\) satisfy
\[
        B_{jh}(l)=\nu(\eta)B_h(\eta l),\qquad |\eta|=1.
\]
The argument \(\eta l\) follows because the displacement is
\(\bar l\,y\otimes x'\): the Fourier kernel pairs \(u\) with \(l\)
by a real trace of a fixed nonzero multiple of \(ul\).
Substitution in the defining Mellin integral with weight \(\nu(l)\)
cancels the factor \(\nu(\eta)\). This proves invariance of the integrated density
along the compact fibers in either chart.

Choose first a test with \(A_1(1)\ne0\), then multiply it by a
nonnegative bump \(b(v)\) supported sufficiently near \(x\) on
each punctured cone and positive at \(x\). The pullback of that
support has all \(|u_w|\) bounded away from zero and infinity.
The Fourier distribution may therefore be evaluated there using
its smooth density. After compact-fiber integration the remaining
integrand, apart from \(b\), is continuous and nonzero at \(x\);
in a small neighborhood a fixed phase makes its real part positive.
This proves \(DT_\infty\ne0\) on Schwartz tests. Compact averaging,
Fourier transformation, restriction and the absolutely convergent
Mellin moment are continuous in the Schwartz seminorms, as in
Lemma~\ref{tpd:mellin}. Hermite expansions approximate any
Schwartz test by polynomials times the vacuum in these seminorms.
The output for an upper polynomial is again lower polynomial:
decompose it over the smaller positive summand and its orthogonal
complement; the compact orbit spans a finite-dimensional polynomial
space and all evaluation and displacement act in the complement.
Consequently some upper polynomial has \(DT_\infty\ne0\), as claimed.
\end{proof}

\begin{proposition}[Descending the test maps]\label{tcr:descend}
Suppose both smaller period maps are nonzero. After application of
the smaller matching map, each of \(T_v,e_{0,v}\) factors through
the upper matching map, at every inert place and at the product of
the real places (on polynomial vectors there). Both resulting maps
intertwine the target line parabolic with
\(\xi\boxtimes\pi_1\) inflated trivially on its radical.
At infinity this means its Lie algebra and the compact part of its
Levi. The two real maps so obtained are nonzero.
\end{proposition}
\begin{proof}
First produce auxiliary upper tests at every other place for each
of the two operations, such that their smaller images are nonzero.
At infinity use Lemma~\ref{tcr:bump}. At a ramified place use the
standard matched pair of Proposition~\ref{tpd:ram} or
Corollary~\ref{threal:ramA}. At an unramified split place use the
vacuum in Proposition~\ref{tpd:ordinary}. At \(v_1\), if a
supported primed test is required, use the surjection of supported
tests onto the whole upper quotient. All the indicated local maps
at that split place already factor through the upper quotient
there by those propositions (one may compare with the zero test on
the opposite side). In the \(A\)-step the smaller unchanged primed
source is compact at infinity and needs no support for integration.
When the smaller supported space is required for matching, one
can also use the representatives of Lemma~\ref{tpd:so}: the same
matrix proof of its lifting assertion works for \(q=0\).
In the \(B\)-step the changed-source operator itself preserves the
full-rank condition by Lemma~\ref{tpd:mellin}.

At an inert place the following vacuum calculation requires no
factorization. The tensor lattice at the place is self-dual for
the additive-character symplectic pairing. The prime is nondyadic.
In the two-plane symplectic decomposition its summands in the
removed tensor blocks have the following description. Each isotropic
line tensor a residual space pairs perfectly with the opposite
isotropic line tensor that space. The four products of the two pairs
of lines split into two such dual pairs. Adaptedness and self-duality
make these pairs of lattices mutually dual, including any fractional
line ideals. Take the removed modulation space and these opposite
summands as its dual position space. In those positions the basic
vector is the lattice indicator times the smaller basic vector,
up to scalar: invariance under the modulation lattice imposes
support on its position dual and invariance under the full dual
translation lattice forces constancy there. The tensor factor on
the residual space is itself the smaller basic line by the same
Heisenberg characterization. This argument also applies to the
one-plane evaluation. Normalize each evaluation on the vacuum.
The source compact fixes the basic vector, by the unramified
oscillator vacuum calculation in Proposition~\ref{threal:allinert}.
The displacement is a single position coordinate in the respective
chart. Thus shifted evaluation gives precisely the smaller basic
vector times \(1_{\mathfrak a_v}(l)\), where the integral translation
ideal \(\mathfrak a_v\) is the localization of
\[
 \mathfrak b\,\mathfrak D_E^{-1}\bar{\mathfrak c}^{-1}
       \quad(A),\qquad \mathfrak b^{-1}\bar{\mathfrak c}\quad(B).
\]
With multiplicative unit volume one its Mellin integral is
\[
  u_v=\nu_v(b_v)|b_v|_{E_v}^{z}
           (1-\nu_v(\pi_v)|\pi_v|_{E_v}^{z})^{-1},
        \qquad b_v\mathcal O_{E_v}=\mathfrak a_v .
\]
This is nonzero, including when \(\nu_v(\pi_v)=1\). All but
finitely many \(b_v\) are units, and the product of these basic
factors converges absolutely to a nonzero number since \(z>1\).

Now put a vector in the kernel of an upper matching map at the
place under consideration and the chosen auxiliary tests everywhere
else. Lemma~\ref{tcr:kernels} kills its upper period identically.
In particular its \(\xi\) coefficient is zero. This coefficient
is a fixed nonzero Haar multiple of the smaller period on the
tensor of the transformed tests, by the constant-term formulas.
For the primed cone formula for a general upper test, first take
the Levi difference of Proposition~\ref{tpd:unfold}; on the
\(\beta\)-coefficient it acts by the difference of the two
characters at the chosen idele, a nonzero scalar, and on \(T\)
it acts by that same scalar. For a zero evaluation with supported
representative, make the correction at \(v_1\) before reading the
other local factor. By smaller injectivity
(Lemma~\ref{tcr:kernels}) the tensor of lower images vanishes.
All its other factors are nonzero, so the one at the place under
consideration vanishes. This proves the claimed factorizations.

The target equivariance can be read from the same formulas. For
a local element of the line Levi or radical take its right translate
in the upper constant term. At the identity the radical acts
trivially by left invariance, the line acts by \(\xi\), and the
smaller target group translates the smaller period. Subtract the
corresponding two smaller expressions and again cancel the fixed
nonzero auxiliary factors by smaller injectivity. At infinity
one takes compact Levi translates and differentiates the identity
for one-parameter subgroups at zero; these operations preserve the
polynomial domain in the required differentiable sense. Continuity
for passing the derivative through the periods follows from
Proposition~\ref{threal:decay}. For a one-parameter element outside
the compact the identity before differentiation holds for Schwartz
tests by the same absolutely convergent integrals. The nonzero
real images follow from Lemma~\ref{tcr:bump}.
\end{proof}

\begin{proposition}[Inert and real uniqueness]\label{tcr:unique}
Under the assumptions of Proposition~\ref{tcr:descend}, the two
descended maps at an inert place satisfy \(D_{T,v}=u_vD_{0,v}\).
At infinity \(D_{T,\infty}=c_\infty D_{0,\infty}\) for a unique
nonzero scalar. For fixed real models, upper and lower real
matching maps and measures, this number is independent of every
finite character, matching quotient and choice of admissible
Weil actions in the two towers.
\end{proposition}
\begin{proof}
At an inert place the upper spherical vector generates its
irreducible representation. The line parabolic and the lattice
stabilizer satisfy \(G=PK\) by the integral Witt argument in
Proposition~\ref{threal:coneA}. Hence the spherical vector generates
even under \(P\). The comparison of its two images in the preceding
proof gives the identity on the whole representation.

Here is the compact weight calculation at one real place.
Write \(R=d-s,S=s,P=a-p,Q=p,A=m_\alpha,B=m_\beta\), and
\(\Delta=A-B>d\). Put
\[
 c_\pm=\tfrac12(B\pm(P-Q)),\qquad
 h_+=c_++\tfrac12(\Delta-R+S),\quad
 h_-=c_-+\tfrac12(\Delta+R-S).
\]
The bottom \(U(R)\times U(S)\)-type of
Proposition~\ref{threal:allreal} has weights
\[
 ((h_+)^P,(c_+)^{R-P};(h_-)^Q,(c_-)^{S-Q}),
\]
where superscripts denote repetitions. To see the normalization,
start from the two target compact vacuum weights \(\pm(P-Q)/2\);
the source determinant in the joint harmonics adds
\((\Delta-R+S)/2\) in the first \(P\) positions and
\((\Delta+R-S)/2\) in the first \(Q\) positions, and the target
splitting adds \(B/2\) throughout.
Alternatively the bottom-layer formula for the ordered two-block
parabolic is its inducing differential plus
\(\rho(\mathfrak u\cap\mathfrak p)-\rho(\mathfrak u\cap\mathfrak k)\)
in Paul--Trapa's cover normalization: in the first positive block
this correction is \(((S-Q)-(R-P))/2\), in the first negative
block its negative, and in the remaining positive and negative
blocks it is respectively \((P-Q)/2\) and its negative.
The inducing parameter is in the good range by \(\Delta>d\);
the bottom type occurs once by the bottom-layer formula for
cohomological induction (the inducing Levi type is a character).
In ordinary, uncovered normalization one subtracts \(\rho(\mathfrak u)\)
from the inducing differential. The two character exponents then
become \((A-a')/2,(B+a)/2\); their difference is
\((\Delta-d)/2>0\), so the bottom-type theorem applies.
\cite[\S2 and Theorem~4.1]{paulTrapaHowe}
\cite[Theorem~5.3 and Proposition~6.1]{voganZuckerman}

The compact Levi on an isotropic line acts by the same circle
\(\eta\) in one positive and one negative coordinate, and by
\(U(R-1)\times U(S-1)\) on the other coordinates. Indeed a null
vector is a sum of two vectors of equal norm, and an element
carrying it to \(\eta\) times itself acts by \(\eta\) on each.
For a \(U(n)\) weight \(\lambda_1\ge\cdots\ge\lambda_n\) the
\(U(n-1)\)-weights interlace it with multiplicity one, and their
omitted circle weight is \(\sum\lambda_i-\sum\mu_i\). Its maximum
\(\lambda_1\) forces \(\mu_i=\lambda_{i+1}\) for all \(i\);
its minimum \(\lambda_n\) forces \(\mu_i=\lambda_i\).
In the \(A\)-step both high blocks occur and \(h_++h_-=A\) is
the maximum diagonal-circle weight. In the \(B\)-step both low
blocks occur and \(c_++c_-=B\) is the minimum. These are exactly
the respective angular characters of \(\xi\). The extreme block
in the upper bottom type thus consists of a single irreducible
type of the smaller compact, with one entry deleted in each sign.
It is the lower bottom type by the same weights formula:
\(P-Q\) and \(R-S\) do not change. It occurs once in \(\pi_1\).
The compact restrictions of the two parabolic maps on the upper
bottom type are therefore proportional. The Iwasawa Lie algebra
identity \(\mathfrak g=\mathfrak p+\mathfrak k\) and the
Poincar\'e--Birkhoff--Witt theorem show that \(U(\mathfrak p_{\mathbb C})\)
acting on this compact type generates the entire irreducible module.
For the product of real places take the exterior tensor product
of the bottom types and apply the same argument componentwise.
This proves dimension at most one for the required space of maps.
Both are nonzero, so the scalar is unique and nonzero.

For fixed base data the infinite characters are fixed by their
specified exponents \(m_\alpha,m_\beta\). The doubled Siegel law
fixes the resulting real oscillator actions: the ratio of two
Heisenberg implementations with the same doubled Levi law is a
character trivial on that Levi and on the root unipotents, hence
trivial on the doubled group, as in Proposition~\ref{threal:inert}.
The compact integral defining \(T_\infty\) and its Mellin weight
therefore agree for all the allowed systems. The fixed realizations
of the upper and lower real maps give the same proportionality
equation on every polynomially matched real pair in all these
systems. Since some real test has nonzero zero-evaluation image
and its upper image can be matched on the other side, this equation
determines the same number \(c_\infty\) for all systems. If
there are no such systems the uniform comparison is vacuous.
\end{proof}

\begin{theorem}[The scalar for the two undoubled periods]\label{tcr:global}
For every triple of \S\ref{newtheta:global} and every allowed set
of real and finite matching maps, the primed period is nonzero
on some supported test, and there is a unique \(c\ne0\) such that
\(I(\Phi)=cI'(\Phi')\) on all matched pairs. For given adapted upper
and lower geometries and given upper and lower real matching maps,
there is \(C\ne0\) independent of all admissible \(X\), Weil actions
and finite matching maps such that
\[
 \begin{array}{ll}
 c=C\,F_X(1-k)L^{S_0}(\nu,k)c_1,
       &(q=0<p;\ e_1=e_2=+1,\ s_{\rm unit}=+1),\\[2pt]
 c\,L^{S_0}(\nu,k+1)=c_1/(C F_X(-k)),
       &(q>0;\ e_1=e_2=-1,\ s_{\rm unit}=+1).
 \end{array}
\]
Here \(F_X\) is the finite factor of \S\ref{newtheta:global},
and \(c_1\) is the scalar for the specified smaller matching maps.
\end{theorem}
\begin{proof}
First work on an adapted pair of geometries and suppose the
assertions of nonvanishing and proportionality hold in the smaller
rank. In particular both smaller period maps are nonzero: supported
primed vectors are matched by surjectivity of the opposite quotient
and the distinguished lower full-rank quotient. Use polynomial
inputs initially; nonzero polynomial inputs are available by
Corollary~\ref{threal:limits} and density. For \(T\) use the
integral with multiplicative units of volume one and local source
compacts of volume one. Normalize finite mixed evaluations on
their vacua. Rational coherent evaluations differ from this system
by fixed scalars with almost every scalar one
(Proposition~\ref{threal:coneA}). At infinity keep the real
evaluations fixed.

On upper matched tests, the restricted tensor product of lower
images of \(T\) is a scalar times that of \(e_0\). At the real
places use Proposition~\ref{tcr:unique}; at the finite places use
Propositions~\ref{tpd:ordinary}, \ref{tpd:ram},
Corollary~\ref{threal:ramA} and Proposition~\ref{tcr:unique}.
For a finite tensor sum this follows by applying the tensor product
of the descended linear maps to the equality of upper tensors in
the definition of polynomial matching. Outside the finite test set
one must retain the product of the vacuum multipliers: these
converge absolutely and have nonzero product. The translation
ideal contributes \(N(\mathfrak a)^{-z}\nu(\mathfrak a)\).
It is a unit at the ramified factors. As computed after
Lemma~\ref{tpd:mellin} and Proposition~\ref{tpd:ram}, its
character value is \(\nu(\mathfrak b)\nu(\mathfrak c)\) in \(A\)
and its inverse in \(B\). The partial Euler factors at the primes
of \(Q\) are one since \(\nu\) is ramified there. The ramified
factors are the stated Gauss factors, with arguments \(1-k,-k\)
respectively and with unit sign \(+1\) in both cases.
All these factors are nonzero. The remaining factors consist of
\(c_\infty\), the absolute ideal norm power, the fixed changes of
evaluation, and the quotient of the two Haar normalizations in
the orbit formulas. They give a single \(C\ne0\) for the stated
fixed data by Proposition~\ref{tcr:unique}. The Haar measures
are part of the models; no equality of Haar masses in different
dimensions is needed.

Use the smaller identity on these two tensors. In \(A\), the
unprimed \(\xi\)-coefficient is the cone and the primed one is the
zero evaluation, giving the first formula. In \(B\), the primed
coefficient is the cone and the unprimed one is the zero evaluation,
giving the reciprocal in the second formula.
Propositions~\ref{threal:endA} and~\ref{threal:endB} show that
the resulting scalar compares the upper periods at every target
argument. The zero-coefficient side is nonzero for some upper
test: choose its real polynomial and finite factors as in
Proposition~\ref{tcr:descend}, then apply smaller injectivity to
their nonzero tensor. Choose a matching vector on the other side,
using the supported surjection at \(v_1\) when the other side is
primed. The comparison just proved gives nonzero upper periods
on both tests. This also proves uniqueness of \(c\). Passage to
every matched limit is Corollary~\ref{threal:limits}.

To apply this argument by induction on \(p+q\), start at
Theorem~\ref{threal:base}. For the induction step the integral
plane construction of Proposition~\ref{famtheta:plane} supplies
the hyperbolic summands with their ideal lattices in the changing
source and target. On a given rational hyperbolic plane choose
Cartan involutions interchanging the two null lines at infinity
and an arbitrary compatible Cartan on its complement. At the
designated split places take the new adapted integral bases in the
prescribed column order. These choices do not change the rational
spaces, the lattices, the Haar measures or their theta integrals.
The smaller tensor has a rational symplectic basis by symplectic
elimination, giving its usual rational Schr\"odinger model.
Its basic vector is the lattice Heisenberg eigenline. It can also
be obtained by partial evaluation of the upper basic vector in the
adapted coordinates above: the removed mutually dual lattices
give the position indicator in those coordinates, leaving the
smaller eigenline. At a split dyadic place one uses the two isotropic
idempotent halves as in Proposition~\ref{tpd:ordinary}; their
individual quadratic phases vanish. In rational coordinates the
basic lattice gives the ordinary integral indicator almost everywhere.
Use any fixed lower source Haar measures. For the two lower
oscillators use the same pairs of splitting characters
\((\beta,\beta/\alpha)\), \((\alpha,\alpha/\beta)\).
Subtracting two from either rank preserves the quadratic character
restriction required for each unitary splitting. Use the local
splitting and its doubling law
\cite[(1.16) and Proposition~2.2]{harrisKudlaSweet}, and their
adelic product for the two Hecke characters
\cite[\S\S2--3]{wuThetaTower}. On rational elements the character
factors, absolute determinant factors and Weil indices have product
one, giving the rational theta action through Poisson summation.
Equivalently the doubling in a rational diagonal polarization
uses the prescribed determinant character on its Siegel Levi
and evaluation at zero; restricting to the two isometric summands
recovers the chosen single Weil action. All smaller gap inequalities
hold since \(p+q\) decreases by one, \(a'\) decreases by at most
two and \(k\) decreases by at most one. The local matching results
in \S\ref{threal:main} apply at both levels. The induction
hypothesis and the preceding adapted argument now apply.

We finally transport the existence and uniqueness assertion to the
original markings. At infinity choose separately in each source
and commonly in the target real isometries conjugating the original
Cartan involutions to the adapted ones. Their oscillator operators
carry the Gaussian line and its creation derivatives onto the new
ones by Heisenberg covariance. On the period the source translation
disappears by Haar invariance and the common target translation
right translates both functions by the same element. Under this
transport the original and new pairs of real simple quotient maps
differ, after identifying their targets, by one nonzero scalar
for each side, by real Howe uniqueness. The same comparison at
the finitely many remarked split places follows by finite Howe
uniqueness. At all other places one can keep the original maps.
The support condition uses matrix rank in the polarization
specified by the chosen split embedding and is invariant under
invertible row and column changes. Therefore a pair matched with
the original maps is transported to an adapted matched pair after
rescaling one test by the product of these finitely many scalar
ratios. The two upper period functions are proportional also for
the original data, and a supported nonzero matched pair transports
back. This transport is needed only for the existence assertion
for an arbitrary marking. For the uniform recursion with fixed,
already adapted geometries use the direct calculation in the
first part of the proof. The upper period scalar may depend on
the chosen matching maps and on \(X\).
\end{proof}
\src{src/Development/WeilBound/ThetaLifts/LocalTheory/LocAOf.lean}{17}{The unprimed test comparison},
\src{src/Development/WeilBound/ThetaLifts/LocalTheory/LocBOf.lean}{17}{the primed test comparison}, and
\src{src/Development/WeilBound/ThetaLifts/LocalTheory/StepFacEq.lean}{85}{the quantifiers for a single descent constant}.


\subsection{The geometric degree share and its arithmetic cost}\label{newtheta:HS}

Fix one component \(Z\) of a good deep cover and write \(D\) for the
dimension of its flag space, \(T>D\) for the final weight index,
\(\delta\) for the positive geometric degree of its birational image,
\(h\) for its normalized Chow height, \(d_0\) for the rank of the
terminal cusp lattice, and \(g\) for its arithmetic degree before the
associated-bundle transfer. The coefficient field is totally complex.
The fraction \(\binom{T-1}{D}\delta/d_0\) has a geometric meaning.

\begin{lemma}[Independent columns on a transverse section]\label{newtheta:filtration}
Over an infinite splitting field \(L\), let \(y_1,\ldots,y_\delta\)
be nonzero homogeneous-coordinate vectors in \(L^N\) for the points
of a transverse section by \(D\ge1\) hyperplanes \(s_1,\ldots,s_D\).
Let the terminal section space have dimension \(d_0\), and let
\(U_0\supset\cdots\supset U_{T-D}\) be a filtration by subspaces of
it. Suppose that degree-zero sections surject by restriction onto
\(L^\delta\), and that for \(0\le j<T-D\) the symbol on \(U_j\)
has kernel \(U_{j+1}\) and entries indexed by pairs \((\gamma,p)\)
with \(\gamma\in\mathbb N^D,\ |\gamma|=j,\ 1\le p\le\delta\).
For a degree-zero section \(m\), a linear form \(l\), and each such
\(\gamma\), suppose \(s^\gamma l^{T-j}m\in U_j\) and its symbol at
\((\gamma',p)\) is
\(\mathbf1_{\gamma=\gamma'}l(y_p)^{T-j}m(y_p)\). Then
\[
 d_0\ \ge\ B:=\delta\binom{T-1}{D}.
\]
There are \(B\) independent terminal columns of this kind. On the
associated graded their determinant in the point and monomial
coordinates is, up to ordering of the bases,
\[
 \left(\prod_{p=1}^{\delta}l(y_p)\right)^{
       \sum_{j=0}^{T-D-1}\binom{j+D-1}{D-1}(T-j)} .
\]
\end{lemma}
\begin{proof}
Choose \(l\) avoiding the finitely many proper hyperplanes
\(l(y_p)=0\), possible over an infinite field. Lift the standard
coordinate vectors of \(L^\delta\) to \(m_1,\ldots,m_\delta\).
For fixed \(j,\gamma,p\), the symbol of
\(s^\gamma l^{T-j}m_p\) has just one nonzero entry,
\(l(y_p)^{T-j}\). These symbols form a basis of the indicated
coordinate space. A relation between all the chosen columns, inspected
at the least \(j\) with a nonzero coefficient, would therefore give
a nontrivial relation between these independent symbols. The number
of \(\gamma\) at level \(j\) is \(\binom{j+D-1}{D-1}\), by placing
\(D-1\) separators in a row of \(j\) equal marks. Summing from \(0\)
to \(T-D-1\) gives \(\binom{T-1}{D}\) by Pascal's identity.
Multiplication of the displayed diagonal entries gives the determinant.
\end{proof}

This computation takes place over a field. In the integral problem one
needs the dimensions, determinants \emph{and} the norms of the columns.
In the original argument a nonzero polynomial \(Q\) in the \(DN\)
hyperplane coefficients selects transverse sections. On \(Q(s)\ne0\)
the specialized Chow form is a nonzero scalar times
\(\prod_p l(y_p)\) after passing to a splitting field. Degree-zero
restriction, the kernel of the symbol at each level, its action on
monomials, and generation of \(U_{T-D}\) by the images \(s^\gamma V_D\)
for \(|\gamma|=T-D\) all enter. Here \(V_j\) is the space of
weight-\(j\) sections in the chosen path of weights. The geometric
proof uses nefness and bigness of both line bundles and coherent
restriction and vanishing on the modified flag space. A rational
choice of lifts for the symbol bases alone gives no upper bound for
their norms or denominators.
\src{src/Development/AbelianVarietyTheory/Definitions/Divisors/HermSysHypOneLeLambdaZero.lean}{199}{Filtration, residue and multiplication hypotheses};
\src{src/Development/AbelianVarietyTheory/Proofs/DivisorProofs/HilbertSamuelFiltrationAtGoodSections.lean}{51}{Geometric construction of the filtration}.

\paragraph{Finite-degree estimate still needed.}
For the same cover, level, weight path and component, the numerical
statement at the point of use is the following. Put \(P_0\ge1\) for
the family size parameter and \(\kappa\ge0\) for its weight parameter.
There must be \(n,\Lambda,\beta,\Lambda_D,\Lambda_T\ge1\) and \(J\in
\mathbb R\) such that
\begin{align}
 \max(\log n,\log\beta,\log\delta)&\le P_0,&
 \log\Lambda&\le P_0(1+\kappa/T),\nonumber\\
 \max(\log\Lambda_D,\log\Lambda_T)&\le P_0\kappa,&
 J&\le 4DP_0(1+\kappa/T)+4TP_0+2DP_0\kappa/T .
 \label{newtheta:budgets}
\end{align}
Define
\[
 u=1-B/d_0,\qquad
 L_h=\log\Lambda+26\log\!\big(nD\delta(2+|h|+\log\Lambda)\big),
 \qquad G=\tfrac{D+1}{2}(\log n+1)+1 .
\]
The required estimate is \(0\le u\le1\) and
\begin{align}
 \frac{g}{d_0}
 &\ge(1-u)\big[Th-3(D+1)|h|-J-26D^2(L_h+\log T)
          \nonumber\\[-2pt]
 &\hspace{39mm}-(T+D+1)\log\beta-(T+D^2)(G+2)\big]
 \nonumber\\
 &\hspace{5mm}-u\big[T(\log2+L_h+\tfrac12\log\beta)
                    +\log\Lambda_D+\log\Lambda_T\big].
 \label{newtheta:finiteHS}
\end{align}
Here \(d_0>0\). The choice of auxiliary parameters may depend on \(Z\).
For each fixed admissible family and accuracy, the threshold in its
scale \(x\) must work simultaneously for all its late admissible
scales, indices, points, stages, good covers under consideration and
components. In particular the exponent budgets and constants in
\eqref{newtheta:budgets} are common. The subsequent comparisons also
require \((1-z)d_0\le B\) and a transfer cost \(c\ge0\) such that
the degree \(g_M\) at the weight of the actual bundle is at least
\(g-d_0c\).
\src{src/Formalism/LocalToGlobal/IndefiniteForms/IndefiniteTree/RealCubicMeans/LogMeanShares/LogMeanStatements/HSClaim.lean}{16}{Exact finite-degree inequality};
\src{src/Development/AutomorphicArithmetic/ThetaCorrespondence/ThetaPairLifts/ThirdPlaceSeries/ExistsDegreeParametersBound.lean}{34}{Family quantifiers and budgets}.

In the construction, \(n\) counts columns of an integral cusp frame
at the first weight, \(\Lambda\) bounds every column norm at every
embedding, and \(\beta\) bounds the Bergman kernel on every eligible
frame of the component. Both auxiliary weights must define
birational linear systems there. Multiplication by an integral frame
column must preserve integral lattices in all consecutive degrees.
At each embedding, multiplying a complex section by coefficients
\(v\) must have norm at most \(\sqrt\beta\|v\|_{\rm frame}\) times
the section norm. Integral sections spanning the spaces over the
coefficient field must be available with norms bounded at every
embedding by \(\exp(DP_0(1+\kappa/T))\) in degree zero and by
\(\exp(P_0\kappa)\) along the path through \(T\). The latter choice
works for both \(\Lambda_D,\Lambda_T\). The former gives
\(J=DP_0(1+\kappa/T)\). One also needs the normalized height of the
birational image in exactly the frame used for the multiplication
norms. The argument behind \eqref{newtheta:finiteHS} then estimates
the height of a determinant polynomial off the exceptional locus
\(Q=0\) and transports its slope bound back to the component
lattice. We give that argument below. Its geometric hypotheses must
still be realized uniformly on the actual covers; a fixed-variety
Hilbert polynomial specifies no bounds for the multiplication norms.
\src{src/Development/AutomorphicArithmetic/ThetaCorrespondence/ThetaPairLifts/HodgeBundleWeights/ExistsDegreeParametersBound.lean}{33}{Simultaneous frame, birational and short-section hypotheses};
\src{src/Development/AbelianVarietyTheory/Proofs/DivisorProofs/ExplicitArithmeticHilbertSamuelBound.lean}{26}{Determinant and Chow-height steps}.

\paragraph{Construction of the common determinant.}
We can in fact derive \eqref{newtheta:finiteHS} from the geometric
restriction, multiplication and metric conditions just identified.
Here are the exact hypotheses for this derivation. Over a number
field \(K\), let \(V_j\) be finite-dimensional section spaces with
full \(\mathcal O_K\)-lattices \(M_j\) and Hermitian norms at every
embedding. Let \(K^N\) act by commuting linear multiplication operators
of degree one. Assume a spanning frame \(w_1,\ldots,w_n\) of \(K^N\)
has norms at most \(\Lambda\) at each embedding and its multiplications
send \(M_j\) to \(M_{j+1}\) for \(0\le j<T\).
In complex coordinates let the input norm be
\(\|v\|_\sigma=\|A_\sigma^{-1}v\|_2\), \(A_\sigma\in\mathrm{GL}_N(\mathbb C)\).
The operator for \(v\) must have norm at most
\(\sqrt\beta\|v\|_\sigma\) in each of these degrees. For \(j=0,D,T\)
assume \(V_j\) has a \(K\)-basis contained in \(M_j\) with all
embedding norms bounded respectively by \(\Lambda_0,\Lambda_D,\Lambda_T\),
each at least one.

Let \(f(s,l)\) be the irreducible Chow polynomial in \(D+1\)
blocks of \(N\) variables, of degree \(\delta\ge1\) in each block.
For every rational \(s=(s_1,\ldots,s_D)\) outside a fixed hypersurface
\(Q(s)=0\), \(Q\ne0\), require the filtration and symbol rules of
Lemma~\ref{newtheta:filtration} after scalar extension, with
\(f(s,l)\) a nonzero constant times \(\prod_p l(y_p)\). The scalar
extension must preserve independence of rational sections in \(V_T\);
the deepest space is generated by \(s^\gamma V_D\) for
\(|\gamma|=T-D\). These conditions involve every good rational
specialization. The norm and lattice conditions involve the original
field and do not require norm control on the splitting fields.

\begin{proposition}[The filtration on the adjoint model]\label{newtheta:koszul}
Let \(X/K\) be a smooth geometrically integral projective variety of
dimension \(D\) over a number field, \(L\) a globally generated line
bundle with generating coordinates \(b_1,\ldots,b_N\), and \(A,B_2\)
nef and big line bundles. Choose a section
\(\xi\in H^0(X,A\otimes(B_2\otimes L^D)^{-1})\) which is invertible
on a nonempty open \(X^\circ\). Write
\[
 W_j=H^0(X,\omega_X\otimes A\otimes L^j),\qquad
 W_j\subset V_j\subset H^0(X^\circ,\omega_X\otimes A\otimes L^j).
\]
Assume \(V_0=W_0\), the \(V_j\) are stable under multiplication by the
coordinates through degree \(T\), and these inclusions persist after
the field extensions used to split the sections. Let \(D\) linear
combinations \(s_i\) of the \(b_j\) cut out a finite reduced transverse
complete intersection \(Y\subset X^\circ\) of degree \(\delta\).
Then the filtration and residue assertions of
Lemma~\ref{newtheta:filtration}, including deepest-space generation as
used in Lemma~\ref{newtheta:power}, hold over any splitting field of \(Y\).
\end{proposition}
\begin{proof}
For \(M=A\) or \(B_2\), and \(k\ge0\), \(M\otimes L^k\) is nef and
big. Multiplication by a fixed nonzero section of \(L^k\), raised to
the \(n\)th power, embeds \(H^0(X,M^n)\) into
\(H^0(X,(M\otimes L^k)^n)\) for every \(n\), proving bigness.
Kawamata--Viehweg vanishing gives
\(H^i(X,\omega_X\otimes M\otimes L^k)=0\) for \(i>0\). Its smooth
projective characteristic-zero hypotheses hold after embedding \(K\)
into \(\mathbb C\), and coherent cohomology commutes with field extension.
This proves the same vanishing over \(K\) and its finite extensions.
Kawamata, \emph{A generalization of Kodaira--Ramanujam's vanishing theorem}, 1982~\cite{kawamataVanishing}.

Let \(\mathcal I\) be the ideal sheaf generated by the \(s_i\) in
local trivializations of \(L\). The Koszul resolution of
\(\mathcal O_Y\) has terms \((L^{-i})^{\binom Di}\). Twisting it by
\(\omega_X\otimes B_2\otimes L^D\) gives terms with exponents
\(D-i\ge0\), all acyclic for global sections by the preceding
vanishing. The global sections of this complex thus compute the
cohomology on \(Y\), in particular the restriction to \(Y\) is
surjective. Multiplication by \(\xi\) identifies the two bundles
over \(Y\), so the restriction of \(W_0\) is surjective there too.
Similarly, twisting by \(\omega_X\otimes A\otimes L^k\) for \(k\ge D\)
gives exactness at the penultimate term:
\[
 \ker\big(W_k\longrightarrow H^0(Y,(\omega_X A L^k)|_Y)\big)
           =\sum_{i=1}^{D}s_i W_{k-1}.
\]
Here juxtaposition of line bundles denotes tensor product.

Work over a splitting field. At a point of \(Y\), the local equations
of the \(s_i\) form a regular system of parameters by transversality.
Thus the associated graded of the local ring is the polynomial ring
in their initial forms over the residue field. Equivalently, monomials
of total degree \(j\) in the \(s_i\) give the coordinates of the
degree-\(j\) Taylor term. This follows also by passing to the completion,
a formal power series ring in these parameters at a smooth rational
point. Trivialize the line bundles compatibly at each point. The kernel
of this Taylor term on sections vanishing to order \(j\) consists of
those vanishing to order \(j+1\), and the symbol of a product is the
product of its initial forms. This gives the diagonal symbol formula
of Lemma~\ref{newtheta:filtration} with \(y_p\) the vector
of the \(b_i(p)\) in the chosen trivialization.

For \(k-j\ge D\), any section in \(W_k\) vanishing to order \(j\)
on \(Y\) belongs to \(\sum_{|\gamma|=j}s^\gamma W_{k-j}\).
Induct on \(j\), the claim for \(j=0\) being immediate. At the next
step represent the section by
\(\sum_{|\gamma|=j}s^\gamma m_\gamma\) using induction. Vanishing
of its degree-\(j\) Taylor term implies \(m_\gamma|_Y=0\) for
every \(\gamma\), by the independence of the monomials at every
point. The kernel formula expresses \(m_\gamma\) as
\(\sum_i s_i m_{\gamma,i}\). Collecting equal monomials gives the
assertion in degree \(j+1\). The degree used in the kernel formula
is \(k-j\ge D\) throughout.

Set \(k_0=T-D\) and
\(V_*=\sum_{|\gamma|=k_0}s^\gamma V_D\subset V_T\).
Filter \(W_T+V_*\) by order of vanishing along \(Y\), with
\(U_j\) the subspace of sections of order at least \(j\).
Since \(V_*\) vanishes to order \(k_0\), every
\(w+v\in U_{k_0}\) with \(w\in W_T,\ v\in V_*\) has \(w\) vanishing
to that order. The preceding generation formula puts \(w\) in
\(\sum s^\gamma W_D\subset V_*\). It follows that \(U_{k_0}=V_*\).
The degree-zero vectors and their products occurring in the symbol
formula lie in \(W_T\), since \(V_0=W_0\). This verifies the
required symbol and deepest-space relations.
\end{proof}
The two adjoint bundles in the family have exactly the powers \(j\)
and \(j+D\) used here. The section \(\xi\) is required to stay
invertible on the chosen open.
\src{src/Development/AbelianVarietyTheory/Definitions/Divisors/HasBoundedPolesGlobalData.lean}{207}{Adjoint models and canonical multiplier};
\src{src/Development/AbelianVarietyTheory/Proofs/DivisorProofs/VanishesToOneOfStalkMap.lean}{133}{Surjectivity for the second bundle};
\src{src/Development/AbelianVarietyTheory/Proofs/DivisorProofs/RestrictedSectionChartCoords.lean}{180}{First-order kernel for the first bundle}.

\begin{proposition}[Boundary vanishing and the adjoint inclusion]\label{newtheta:modified}
Let \(\bar X\) be a smooth integral projective variety over a number
field, \(B\) a reduced effective Cartier divisor, \(U=\bar X-B\),
and \(D=\dim\bar X\). Let \(M_1,M_2,C\) be line bundles on \(\bar X\).
Assume the canonical comparison on \(U\) extends to an isomorphism
\(C\simeq\omega_{\bar X}(B)\). For \(0\le j\le T\), set
\(N_j=M_2 M_1^{D+j}\), and assume that the given space of cusp forms
in weight \(j\), as a space of sections on \(U\), is exactly the
restriction of \(H^0(\bar X,CN_j(-B))\).
Let \(W_i'\subset H^0(\bar X,M_i)\) be finite-dimensional systems
with no base points on \(U\). Let \(\pi:X\to\bar X\) be a proper
birational morphism from a smooth variety, an isomorphism over \(U\),
on which both base ideals become invertible ideals
\(\mathcal O_X(-E_i)\). Put \(L=\pi^*M_1(-E_1)\),
\(B_2=\pi^*M_2(-E_2)\), \(A=\pi^*(M_2M_1^D)\).
Then \(L,B_2\) are globally generated. There is a section
\(\xi\) of \(A(B_2L^D)^{-1}\) invertible on \(U\).
In the adjoint coordinates on \(U\), the cusp spaces satisfy
\[
 V_j=H^0(X,\omega_X A L^j(jE_1))\big|_U,\qquad
 W_j=H^0(X,\omega_X A L^j)\big|_U\ \subset V_j,\qquad W_0=V_0 .
\]
\end{proposition}
\begin{proof}
The local coefficients of \(W_i'\) generate the base ideal.
After pullback, division by a local generator of that invertible
ideal produces functions generating the unit ideal. These functions
are the local coefficients of sections generating the moving part.
The \(E_i\) are supported over \(B\). The canonical section of
\(\mathcal O_X(E_2+DE_1)=A(B_2L^D)^{-1}\) therefore gives \(\xi\).

Under \(C\simeq\omega_{\bar X}(B)\) the cusp space is
\(H^0(\bar X,\omega_{\bar X}N_j)\) on \(U\). For any line bundle \(N\)
on \(\bar X\), rational pullback identifies
\[
 H^0(\bar X,\omega_{\bar X}N)=H^0(X,\omega_X\pi^*N).
\]
Indeed, pullback of a regular top differential is regular by the
map on differentials. Conversely a top form regular on \(X\) defines
a rational top form on \(\bar X\), since the function fields agree.
At each codimension-one point \(y\) the inverse rational map to \(X\)
extends to \(\operatorname{Spec}\mathcal O_{\bar X,y}\) by properness.
Pulling the form back along this morphism puts it in
\((\omega_{\bar X}N)_y\), as differentials commute with localization.
A rational section of a line bundle on a normal variety regular at
all codimension-one points is regular everywhere, by the intersection
description of a normal noetherian domain inside its fraction field.
This proves the equality. Now \(\pi^*N_j=A L^j(jE_1)\).
Multiplication by the canonical section of \(\mathcal O_X(jE_1)\)
is an inclusion and is invertible over \(U\). It is the identity for
\(j=0\). Using this section to identify the two open bundles gives
the displayed comparisons.
\end{proof}

On the toroidal flag component the first hypothesis above requires
the determinant of the Kodaira--Spencer comparison to have local
coefficient \(u/t\), where \(t\) cuts out the \emph{reduced}
boundary and \(u\) is a unit. In smooth toroidal coordinates where
the boundary is \(x_1\cdots x_b=0\), a basis of logarithmic
differentials is \(dx_i/x_i\) for \(i\le b\) and \(dx_i\) for
\(i>b\). Wedging proves the pole formula if the automorphic
comparison matrix in this basis has unit determinant. A comparison
known only on \(U\) does not give that last unit condition. In one
variable, multiplication by \(x\) is invertible on \(x\ne0\) and
changes the boundary order by one. The second hypothesis identifies
\emph{every} form vanishing on the boundary with an automorphic cusp
form of the specified weight and conversely. These two hypotheses
give the inclusion for every \(j\) on the same modification without
any estimate involving \(j\). Generation of the two systems on \(U\),
bigness of their moving parts where used, and nefness and bigness of
\(A\) in Proposition~\ref{newtheta:koszul} require geometric input.
\src{src/Development/AutomorphicArithmetic/ThetaCorrespondence/ThetaPairLifts/HodgeBundleWeights/ExistsWeightHodgeBundleGenerated.lean}{30}{Both base systems and the common modification};
\src{src/Development/AutomorphicArithmetic/ThetaCorrespondence/ThetaPairLifts/HodgeBundleWeights/ExistsSimpleLogPoleAlongBoundary.lean}{38}{Logarithmic Kodaira--Spencer determinant at the boundary};
\src{src/Development/AbelianVarietyTheory/Definitions/Divisors/ExistsResolProjSpace.lean}{203}{Fixed divisors and the moving adjoint coordinates}.

There is a primary comparison theorem that supplies the unit condition
on the standard PEL model. On the characteristic-zero fiber of a
neat unitary PEL moduli space with a compatible admissible smooth
projective cone decomposition, write \(S\) for its smooth toroidal
compactification and \(B_S\) for the reduced boundary. For the
universal semiabelian family the polarized, endomorphism-balanced
Kodaira--Spencer sheaf is isomorphic by its extended Kodaira--Spencer
morphism to \(\Omega^1_S(\log B_S)\). The reduced normal crossings
boundary and the isomorphism are parts (3) and (4) of the same
theorem. Lan,
\emph{Toroidal compactifications of PEL-type Kuga families},
Definition 1.40 and Theorem 1.41~\cite{lanToroidal}.

Here is the determinant on the flag bundle for a CM field action
with conjugation as Rosati involution. Work over a splitting field
for this action and use the polarization isogeny to identify the two
Hodge bundles over the characteristic-zero base. Write \(H_\tau\)
for the dual of the \(\tau\)-part of the Lie bundle, of rank \(m_\tau\).
In Lan's balanced tensor quotient, the relation for an element
\(a\) of the CM field kills the tensor of the \(\tau\) and
\(\nu\) parts unless \(\nu=\bar\tau\), since some \(a\) has
\(\bar\tau(a)\ne\nu(a)\) otherwise. The polarization symmetry
identifies each surviving ordered pair with its transposed pair.
As \(\tau\ne\bar\tau\), the quotient has precisely one summand
\(H_\tau\otimes H_{\bar\tau}\) per unordered pair. This also follows
by taking the projection onto that summand on one ordering and its
transpose on the other, giving an explicit inverse to the inclusion.

Let \(p:\bar X\to S\) be the bundle of complete increasing flags
in the Lie eigenbundles. Let \(Q_{\tau,i}\) denote the
\emph{dual} of the \(i\)-th graded Lie line, numbered
\(0\le i<m_\tau\). The smooth projection gives an exact sequence
for logarithmic cotangents with terms \(p^*\Omega^1_S(\log B_S)\)
and \(\Omega^1_{\bar X/S}\), as is seen on a trivializing flag
chart by adding the flag coordinates to local coordinates on \(S\).
In the determinant the first term contributes \(m_{\bar\tau}\) to
the exponent of every \(Q_{\tau,i}\), by the tensor determinant
formula above. For the relative tangent the graded factors are
\(\operatorname{Hom}(Q_{\tau,i}^{\vee},Q_{\tau,j}^{\vee})\),
\(i<j\): an infinitesimal deformation of a flag is an
endomorphism modulo those preserving its filtration.
The determinant of the relative \emph{cotangent} thus contributes
\(i-(m_\tau-1-i)=2i+1-m_\tau\). Consequently
\[
 \omega_{\bar X}(p^{-1}B_S)
 \simeq\bigotimes_{\tau,i}
 Q_{\tau,i}^{\,m_{\bar\tau}+2i+1-m_\tau}.
\]
This is the character called \(\mathrm{detTanEx}\) in the source;
the central exponent is zero. Thus when the flag and line bundles
in Proposition~\ref{newtheta:modified} are the indicated PEL bundles
and their canonical comparison on the interior is the determinant
of this map up to a nonzero constant, the required boundary order
follows. Checking this comparison uses the identification of the
original framed moduli problem, flags and character rule with this
model. Multiplication by a varying invertible function on the open
would require a separate calculation of its boundary orders.
No metric constant enters this calculation.
\src{src/Development/AbelianVarietyTheory/Definitions/Divisors/ExistsUniqueIsTanFrameOfIsEmpty.lean}{276}{Canonical weight and determinant exponents}.

On the same standard model, let \(\mathcal L\) be the line bundle of
a flag character and set \(\mathcal W=p_*\mathcal L\). Over a
trivialization of the Hodge bundles this direct image is the constant
bundle with fiber the sections of the character line on the product
of flag varieties: formation of \(H^0\) commutes with the flat base
change from the coefficient field. Its gluing is exactly the change
of Hodge frame in that representation. By the projection formula
\[
 H^0(\bar X,\mathcal L(-p^{-1}B_S))
       = H^0(S,\mathcal W(-B_S)).
\]
Thus the cusp comparison reduces to identifying the coefficient
space with the rational form of the usual holomorphic cusp space in
this representation. Over \(\mathbb C\), the latter has the standard
description by sections of the subcanonical bundle
\(\mathcal W(-B_S)\). Over \(K\), equality of two \(K\)-subspaces
in a common space follows from equality after \(K\hookrightarrow\mathbb C\)
by faithful flatness. This last assertion requires both embeddings
of the spaces in the same ambient space of sections on the open.
It gives no comparison of integral lattices or of Petersson norms.
Lan and Stroh,
\emph{Relative cohomology of cuspidal forms on PEL-type Shimura varieties},
introduction preceding Theorem 1.1~\cite{lanStroh}.

\begin{lemma}[One open set of pencils]\label{newtheta:pencil}
In Proposition~\ref{newtheta:koszul}, suppose further that the \(b_i\)
define a birational morphism \(X\to Z\subset\mathbb P^{N-1}\), and
let \(f\) be the Chow form of the geometrically integral image \(Z\)
of degree \(\delta\). There is a nonzero polynomial \(Q(s)\) over
\(K\) such that at every \(s\) with \(Q(s)\ne0\) the required transverse
section lies in \(X^\circ\) and its Chow specialization is a scalar
times \(\prod_p l(y_p)\) over a splitting field.
\end{lemma}
\begin{proof}
There is a dense smooth open of \(Z\) on which the morphism is an
isomorphism. Shrink it to avoid the image of \(X\setminus X^\circ\);
this image is closed by properness and has dimension at most \(D-1\).
In the affine parameter space of \(D\) hyperplane coefficient vectors
the condition that all hyperplanes pass through a fixed point imposes
\(D\) independent equations. Thus the incidence over the excluded
locus has dimension at most \(DN-1\). Over the retained smooth locus,
the hyperplanes through a fixed point induce arbitrary covectors on
its tangent space, since it is embedded in projective space. Linear
dependence of their \(D\) covectors imposes one further nontrivial
determinant equation. The incidence of nontransversality likewise has
dimension at most \(DN-1\). The closures of their parameter images
are therefore proper closed sets. On the remaining open intersect also
the open of proper linear sections of degree \(\delta\) from the
definition of projective degree. It contains a principal open \(Q\ne0\)
defined over \(K\), since all the conditions descend to \(K\).
Each such section has \(\delta\) distinct geometric points and the
same inverse image on \(X\). For fixed \(s\), the defining incidence
property of the Chow form says exactly that \(f(s,l)=0\) when the
last hyperplane passes through one of these points. The resulting
linear forms in \(l\) are pairwise nonproportional. Each divides the
nonzero polynomial \(f(s,l)\), whose degree in \(l\) is \(\delta\).
Their product therefore equals it up to a nonzero scalar. The points
and scalar lie in a finite extension.
\end{proof}

\begin{lemma}[One determinant divisible by a specified Chow power]\label{newtheta:power}
Set \(R_0=\binom{T-1}{D}\), \(B=R_0\delta\), and
\[
 b_*=\sum_{j=0}^{T-D-1}\binom{j+D-1}{D-1}(T-j)
     =R_0\frac{T+D(D+1)}{D+1}.
\]
There are \(0\le y\le d_0-B\) and a single nonzero determinant
polynomial \(\Delta=f^{b_*}H\) over \(K\). Its columns comprise
\(B\) degree-zero short vectors multiplied \(T\) times,
\(y\) degree-\(D\) short vectors multiplied \(T-D\) times,
and \(d_0-B-y\) degree-\(T\) short vectors. The cofactor \(H\)
is multihomogeneous of total degree
\[
 e=BT+y(T-D)-b_*(D+1)\delta\le T(d_0-B).
\]
\end{lemma}
\begin{proof}
The first equality for \(b_*\) counts the exponent in
Lemma~\ref{newtheta:filtration}. Write \(T-j=D+(T-D-j)\).
Summing the latter factor by counting one copy of each term for
every \(j\le i<T-D\) gives \(\sum_{i=0}^{T-D-1}\binom{i+D}{D}
=\binom{T}{D+1}\). Thus \(b_*=D R_0+\binom{T}{D+1}\),
which gives the stated expression.

Choose a rational \(s^0\) outside \(Q=0\) maximizing the rank \(y\)
of the deepest space. Such a point exists since \(K\) is infinite
and the ranks form a nonempty finite set. At \(s^0\), choose
\(\delta\) vectors \(m_r\) from the short degree-zero basis whose
residues form a basis, and \(y\) monomials \((s^0)^\eta m\),
\(|\eta|=T-D\), with \(m\) in the short degree-\(D\) basis,
that span the deepest space. This is possible since the proposed
generators span after base change. Use all the columns
\(s^\gamma l^{T-j}m_r\), \(|\gamma|=j<T-D\), together with
the selected \(s^\eta m\). At \(s=s^0\) choose a rational \(l\)
with all \(l(y_p)\ne0\); finitely many nonzero linear forms over
an extension cannot vanish everywhere on \(K^N\). The jet symbols
and the deepest columns give \(B+y\) independent vectors by
Lemma~\ref{newtheta:filtration}. Complete them using the short
degree-\(T\) basis. With the choices of short vectors and exponents
now fixed, take the determinant in an arbitrary \(K\)-basis of
\(V_T\). Its specialization just constructed is nonzero.

At any other good \(s\) the \(y\) selected deepest columns are
dependent or span the deepest space, by maximality of \(y\). In
the former case the specialized determinant is identically zero.
In the latter take coordinates adapted to the filtration and complete
them by a basis modulo \(U_0\). The matrix is block triangular in
the successive symbol bases, with the deepest block and the block
modulo \(U_0\) independent of \(l\). For each \(j,\gamma\) the
diagonal symbol block is
\(\operatorname{diag}(l(y_p)^{T-j})_{p}\,(m_r(y_p))_{p,r}\).
The symbol rule and surjectivity of restriction show that each graded
piece has exactly the indicated size. Taking determinants gives
\(\Delta(s,l)=c_s f(s,l)^{b_*}\), allowing \(c_s=0\) and working
over the splitting field.

Choose \(l_0\in K^N\) such that \(f(s,l_0)\) is a nonzero
polynomial in \(s\). Clearing the two specialized identities gives
\(\Delta(s,l) f(s,l_0)^{b_*}=f(s,l)^{b_*}\Delta(s,l_0)\)
on \(Q(s)\ne0\), hence identically by polynomial density over the
infinite field. Since \(f\) is irreducible and has positive
\(l\)-degree, it is coprime to \(f(s,l_0)\). Unique factorization
gives \(f^{b_*}\mid\Delta\). Each column has fixed multidegree, so
the determinant has their sum as multidegree, and the quotient by
a nonzero multihomogeneous factor is multihomogeneous. Total
degrees give the stated \(e\): in particular \(e=y(T-D)-BD(D+1)\)
and \(y\le d_0-B\).
\end{proof}

\paragraph{Height of the determinant.}
For a multiform \(P\ne0\) put \(r=[K:\mathbb Q]\) and define
\[
 \mathcal H(P)=\frac{1}{r}\left[
 \sum_{\mathfrak p}\log\|P(w\cdot)\|_{\mathfrak p,\mathrm{Gauss}}+
 \sum_{\sigma:K\hookrightarrow\mathbb C}
 \int\log|P^\sigma(A_\sigma z)|\,dm(z)\right].
\]
The substitution \(w\) sends an \(n\)-tuple to \(\sum_i w_i t_i\)
in each input block. The finite norm is the maximum coefficient
absolute value with \(|x|_{\mathfrak p}=N\mathfrak p^{-v_{\mathfrak p}(x)}\).
The integral uses independent uniform unit vectors in \(\mathbb C^N\)
in all the blocks. At a finite prime the Gauss norm is multiplicative:
divide each polynomial by a coefficient of largest absolute value and
reduce modulo the prime, where products of nonzero polynomials are
nonzero. At infinity logarithms add, so \(\mathcal H\) is additive
on products. The normalized component height used below satisfies
\[
 h=\frac{\mathcal H(f)}{(D+1)\delta}+\tfrac12 H_{N-1},
 \qquad H_j=\sum_{i=1}^{j}\frac1i,\quad H_0=0.
\]
This is the Chow-height normalization for the \emph{same} integral
frame and Hermitian matrices. Identifying the given component height
with this expression requires its birational presentation.
\src{src/Development/AbelianVarietyTheory/Definitions/Divisors/PointHeight.lean}{30}{Frame-substituted Chow height}.

\begin{lemma}[Multiform and determinant height bounds]\label{newtheta:heights}
Write \(C=\log(n\Lambda)+(\log n+1)/2\). For every nonzero
multiform \(P\) of total degree \(v\),
\(\mathcal H(P)\ge-Cv\). For the polynomial of
Lemma~\ref{newtheta:power},
\[
 g\ge\mathcal H(\Delta)
 -B\log(\beta^{T/2}\Lambda_0)
 -y\log(\beta^{(T-D)/2}\Lambda_D)
 -(d_0-B-y)\log\Lambda_T .
\]
\end{lemma}
\begin{proof}
For a homogeneous polynomial \(p\) of degree \(v\) and a unit vector
\(x\), its sphere mean satisfies
\(\int\log|p(z)|\,dm(z)\ge\log|p(x)|-vH_{N-1}/2\) whenever
\(p(x)\ne0\). Rotate coordinates so \(x\) is the first coordinate vector.
Condition on \(z_1\) and average the common phase of the orthogonal
coordinates. Jensen's inequality for the logarithm of the modulus of
a holomorphic polynomial on a disk gives a lower bound
\(\log|z_1^v p(x)|\). For \(N>1\), \(|z_1|^2\) has density
\((N-1)(1-t)^{N-2}\) on \(0<t<1\): a normalized vector of independent
standard complex Gaussians is uniform on the sphere and its squared
coordinates before normalization are independent exponential variables.
Thus \(-\mathbb E\log |z_1|^2=\int_0^1[1-(1-t)^{N-1}]\,dt/t
=H_{N-1}\). For \(N=1\) the cost is zero. Applying the same phase
averages blockwise gives the assertion for multiforms with \(v\) the
sum of the block degrees. The logarithm is bounded above on the
compact product of spheres, and the lower bound starting at any
nonzero value makes its negative part integrable.

The frame spans, so \(P(w\cdot)\ne0\). Choose any nonzero coefficient
\(c\) of it. At finite places \(|c|_{\mathfrak p}\le
\|P(w\cdot)\|_{\mathfrak p,\mathrm{Gauss}}\). At infinity, Cauchy's
formula bounds \(|\sigma c|\) by the supremum of \(|P^\sigma|\)
on tuples of frame combinations with coefficients of modulus at most
one. Each combination has frame norm at most \(n\Lambda\). By the
sphere estimate and homogeneity its logarithm is at most the sphere
mean plus \(v(\log(n\Lambda)+H_{N-1}/2)\). The product formula for
\(c\) and \(H_{N-1}\le1+\log n\) (as \(N\le n\)) give the first
bound.

Fix a rational basis of \(V_T\), and let \(I\) be its determinant
fractional ideal for \(M_T\), generated by determinants of integral
columns. At every finite prime the coefficients of each constructed
column after frame substitution lie in \(M_T\), by multinomial expansion
and stability under the \(w_i\). All coefficients of \(\Delta(w\cdot)\)
therefore lie in \(I\). Its finite log-norm sum is at most
\(-\log N(I)\). At infinity let \(G_\sigma\) be the positive definite
Gram matrix in the same basis. Hadamard's inequality bounds the volume
\(|\Delta^\sigma|\sqrt{\det G_\sigma}\) by the product of the column
norms. Evaluating with each input on its unit sphere, iterating the
multiplication estimate gives the three norm bounds in the statement.
Integrate the logarithm and sum over embeddings. In this basis the
normalized arithmetic degree is
\(g=-r^{-1}(\log N(I)+\tfrac12\sum_\sigma\log\det G_\sigma)\):
the finite contribution is the local determinant ideal and the infinite
contribution is the logarithmic volume of this basis, with a minus
sign. The formula is basis-independent by the product formula. It
yields the second inequality.
\end{proof}

\begin{corollary}[The effective inequality from the normed filtration]\label{newtheta:effective}
Under the common-determinant hypotheses and height identification above,
\eqref{newtheta:finiteHS} holds with \(J=\log\Lambda_0\).
\end{corollary}
\begin{proof}
Put \(h'=\mathcal H(f)/((D+1)\delta)\). By
Lemma~\ref{newtheta:heights}, \(h'\ge-C\) and
\(\mathcal H(H)\ge-Ce\ge-CT(d_0-B)\).
Also \(b_*\mathcal H(f)=B(T+D(D+1))h'\).
Add these estimates in \(\mathcal H(\Delta)\) and insert the
three column costs in the same lemma. Since \(y\le d_0-B\), all
the costs are nonnegative, and \(u=1-B/d_0\), this gives
\begin{align*}
 \frac{g}{d_0}&\ge(1-u)
 [Th'-D(D+1)C-J-\tfrac T2\log\beta]\\
 &\hspace{8mm}-u[TC+\tfrac T2\log\beta
                        +\log\Lambda_D+\log\Lambda_T].
\end{align*}
Here \(0\le u\le1\) by Lemma~\ref{newtheta:filtration}.
The height normalization gives \(h'\ge h-(1+\log n)/2\).
Since \(n,D,\Lambda,\delta\ge1\),
\(L_h\ge\log\Lambda+26\log(2n)\ge C\).
Thus \(D(D+1)C\le26D^2(L_h+\log T)\), and
\(T(1+\log n)/2\le(T+D^2)(G+2)\).
Enlarging the nonnegative costs, including the term
\(3(D+1)|h|\) in the first bracket, now gives
\eqref{newtheta:finiteHS} term by term. The dependence on \(T,D,n\)
is explicit, so the proof applies simultaneously to all components
carrying the stated inputs.
\end{proof}

In particular, this derivation removes a separate appeal to an effective
Hilbert--Samuel estimate \emph{after} the normed filtration and Chow
height have been identified. No bound on \(s^0\) or on its splitting
field enters Lemma~\ref{newtheta:power}. The controlled norms belong
to the short coefficient vectors of a polynomial in variable \(s,l\).
The unproved geometric step here is producing, on each actual modified
flag component at the common family threshold, the restriction
surjectivity at degree zero, the successive symbols and deepest-space
generation on the complement of \(Q=0\), the stated short bases and
multiplication operator norms, and the same-frame Chow-height
identification. Propositions~\ref{newtheta:koszul} and
Lemma~\ref{newtheta:pencil} give the restriction, symbol and generation
statements once the smooth projective model, adjoint inclusions,
nef and big bundles, nonvanishing multiplier, and birational
pencil are realized there. The numerical estimates on \(P_0,\kappa\)
are needed at that threshold. For example, surjectivity of restriction after
tensoring with the coefficient field does not bound the norms of a
basis of the original integral lattice. The rational filtration and
the metric hypotheses address different necessary parts of this step.
\src{src/Development/AbelianVarietyTheory/Proofs/DivisorProofs/SlopeLowerBoundTopLatticeOffHypersurface.lean}{123}{Choice of a maximal-rank section and common determinant};
\src{src/Development/AbelianVarietyTheory/Proofs/DivisorProofs/MultiformHeightLowerBound.lean}{115}{Multiform cost};
\src{src/Development/AbelianVarietyTheory/Proofs/DivisorProofs/LatticeDegreeGeHeightPolynomialVectorsDet.lean}{91}{Determinant ideal and degree}.

The required uniformity of the first frame can be expressed one level
earlier. At an admissible stage write \(Z_0=f_0r_0L_0\) for the
product of the host parameter, rank parameter and scale, and
\(m_1=\sum_{\tau,k}|u_1(\tau,k)|\) for its first weight size.
The short-generator comparison and the uniform Bergman estimate
require fixed \(C_s,C_b\ge0\), independent of the component and the
late stage, for which
\[
 \log n\le C_s Z_0^3,\quad
 \log\Lambda\le C_s(Z_0^3m_1+Z_0^5),\quad
 \sup_{\text{eligible frames}}\operatorname{Berg}(u_1)
                         \le e^{C_bZ_0^3}.
\]
Here the short generators must span the integral frame lattice with
the recorded embedding norms; the logarithmic estimate for
\(\Lambda\) allows enlarging it to at least one.
At the same stage the size bounds give \(Z_0^{12}\le P_0\) and
\(m_1\le Z_0^2(1+\kappa/T)\).
Consequently, as soon as \(Z_0\ge2C_s+C_b+1\),
\(\log n,\log\beta\le P_0\) with \(\beta=e^{C_bZ_0^3}\),
and \(\log\Lambda\le2C_sZ_0^5(1+\kappa/T)\le
P_0(1+\kappa/T)\). This proves the claimed uniform absorption of
these three costs using a single threshold. The existence of the short
generators with that dependence and the componentwise supremum bound
are the upstream metric estimates. A pointwise bound for any fixed
level with an unspecified level-dependent constant cannot provide the
single threshold.
\src{src/Development/AutomorphicArithmetic/ThetaCorrespondence/ThetaPairLifts/ThirdPlaceSeries/HilbertSamuelBergmanFrameBound.lean}{137}{Short frame and uniform cubic and quintic costs}.

\begin{lemma}[From field-spanning sections to integral generators]\label{newtheta:short}
Let \(K\) be a number field of degree \(d\) and absolute
discriminant \(\Delta\). Let \(V\) be an \(R\)-dimensional \(K\)-space
and \(M\subset V\) a finitely generated \(\mathcal O_K\)-module.
At each embedding \(\sigma\), give \(V\) a seminorm \(N_\sigma\)
with \(N_\sigma(ax)=|\sigma(a)|N_\sigma(x)\).
Assume a subset of \(M\) spans \(V\) over \(K\) and all its elements
satisfy \(N_\sigma(x)\le H\) for every \(\sigma\), where \(H\ge0\).
Then \(M\) has at most \(2dR\) generators over \(\mathcal O_K\), all
with every seminorm at most \(\max(1,d^4R\Delta^{1/d}H)\).
\end{lemma}
\begin{proof}
Use the Euclidean structure
\(\|z\|^2=\sum_\sigma|z_\sigma|^2\) on the real subspace of
\(\mathbb C^d\) where \(z_{\bar\sigma}=\overline{z_\sigma}\).
The dual to the embedded lattice \(\mathcal O_K\) consists of the
coordinatewise conjugates of the images of the inverse different
\(\mathfrak D^{-1}\). Indeed its pairing with \(a\in\mathcal O_K\)
then equals \(\operatorname{Tr}_{K/\mathbb Q}(ay)\); the trace
pairing is nondegenerate, and the integral-valued linear functionals
on the full lattice are rational on its rational span.
For \(0\ne y\in\mathfrak D^{-1}\), the ideal \(y\mathfrak D\)
is integral, so \(|N_{K/\mathbb Q}(y)|\Delta\ge1\).
Arithmetic--geometric mean on the \(d\) squared absolute values gives
\(\|y\|\ge\sqrt d\,\Delta^{-1/d}\).
The Euclidean transference theorem
\(\lambda_d(\Gamma)\lambda_1(\Gamma^*)\le d\), applied to this
full real lattice with this dual, now gives \(d\) independent integers
\(\omega_i\) with \(|\sigma(\omega_i)|\le\sqrt d\,\Delta^{1/d}\)
at every embedding. Here \(\lambda_i\) denotes the least radius
of a closed Euclidean ball containing \(i\) independent lattice vectors.
Banaszczyk,
\emph{New bounds in some transference theorems in the geometry of numbers}, Theorem 2.1~\cite{banaszczyk}.

If \(R=0\) the empty set suffices. Otherwise choose a \(K\)-basis
\(b_j\) from the spanning sections. Then \(e_{ij}=\omega_i b_j\)
are in \(M\), form a rational basis of \(V\), and all their
seminorms are at most \(\sqrt d\,\Delta^{1/d}H\).
As an abelian group \(M\) is finitely generated and torsion free
of rank \(dR\). For each element \(v\) of an integral basis of \(M\),
write \(v=\sum q_{ij}e_{ij}\), \(q_{ij}\in\mathbb Q\), and take
\(v'=v-\sum\lfloor q_{ij}\rfloor e_{ij}\).
The triangle inequality gives
\(N_\sigma(v')\le dR\sqrt d\,\Delta^{1/d}H\).
The \(e_{ij}\) and all these \(v'\) generate \(M\) over \(\mathbb Z\),
hence over \(\mathcal O_K\). Their count and the asserted weaker
fourth-power bound follow.
\end{proof}

Let \(R>0\) be the component rank. Suppose
\[
 \log d,\ \log(\Delta^{1/d})\le391 Z_0^2,\qquad \log R\le C_b Z_0^3.
\]
Let \(H=\exp(C_2(Z_0^3m_1+Z_0^5))\) bound the integral
field-spanning sections, with \(C_2\ge0\).
Lemma~\ref{newtheta:short} gives the required count and height of
integral generators with \(C_s=2000+C_b+C_2\). Indeed
\(\log(2dR)\le\log 2+391 Z_0^2+C_bZ_0^3\), and
\(\log(d^4R\Delta^{1/d}H)\le1955 Z_0^2+C_bZ_0^3+\log H\),
where \(Z_0\ge1\). This removes the need to bound an integral change
of basis for the cusp lattice.
\src{src/Development/AbelianVarietyTheory/Proofs/DivisorProofs/ExistsShortGeneratorsOfShortSpanningFamily.lean}{13}{Totally bounded generators of the whole module};
\src{src/Development/AutomorphicArithmetic/ThetaCorrespondence/ThetaPairLifts/HodgeBundleWeights/ShortGeneratorNormComparison.lean}{25}{Uniform field and rank costs}.

The remaining field-spanning assertion concerns the actual deep
cover, coefficient field, integral lattice and Petersson norm on
each component. In the source the weight \(\lambda\) is decreasing
at every embedding, its gap between the last of the \(p\) coordinates
and the first of the \(q\) coordinates is at least \(8f_0r_0+10\),
\(p,q\ge1\), and \(\log\sum_{\tau,k}|\lambda_{\tau,k}|\le10L_0\).
The auxiliary even theta rank used for this conversion is
\(n_\theta=8(f_0r_0+1)\). The span at that even rank is obtained
from three assertions about a common family of theta series:
field-spanning at the given level, an archimedean bound at
\emph{all} embeddings and components, and a positive integer
clearing the finite denominators of the specified rescaled series.
The costs in the logarithms of the archimedean bound and of the
integer must each be at most a fixed constant times
\(Z_0^2\,\mathrm{sizeK}\), where \(\mathrm{sizeK}\) is the
theta weight size
\(\sum_{\tau}\sum_{0\le i<m_\tau}a_{\tau,i}
 +f_0r_0 n_\theta/2\) for the nonnegative polynomial weights
\(a_{\tau,i}\) of the series. The conversion to \(\lambda\) additionally
uses two twisting sections integral at all finite places, with
controlled archimedean norms, whose product induces the prescribed
change of weights. Qualitative generation over the coefficient field
does not bound either cost. For instance, the lattice
\(\mathbb Z\subset\mathbb Q\) is spanned over \(\mathbb Q\) by
\(1/m\), and its least integer denominator is \(m\), arbitrarily
large even with the usual norm at most one. For the actual theta
family Theorem~\ref{intjet:semi} supplies a common finite multiplier,
and Proposition~\ref{intjet:arch} gives a complex estimate in terms
of the shortest period and the norms of the generator, directions
and polynomial. Proposition~\ref{bfam:harm} provides short data for the polynomial span and
Corollary~\ref{bfam:metric} gives the normalized metric bound under its
stated analytic hypotheses. The final identification of their image
with cusp sections on the given charts follows from Propositions~\ref{bfam:unitread}, \ref{dg:identity} and~\ref{munf:span}. In the final weights one must in addition bound the parallel determinant twist of \eqref{ctwist:degree} uniformly on the actual integral components.
\src{src/Development/AutomorphicArithmetic/ThetaCorrespondence/ThetaPairLifts/HodgeBundleWeights/ExistsGoodCoverAtMinimalLevel.lean}{25}{Three theta-family inputs for a single even weight};
\src{src/Development/AutomorphicArithmetic/ThetaCorrespondence/ThetaPairLifts/HodgeBundleWeights/ShortSpanCompositionClosed.lean}{50}{Change of weight and controlled twisting sections}.

\paragraph{Cancellation of the lattice index in a rescaled series.}
Here the auxiliary integer \(N\) can be arbitrarily large for a
fixed level. Write \(O=\mathcal O_E\), \([E:\mathbb Q]=2f\),
\(r=p+q\), \(M\) for the positive integer level, and
\(\Lambda\subset E^n\) for the self-dual lattice for
\(\langle a,b\rangle=\sum\bar a_i b_i\).
The construction takes \(N\) odd, \(M\mid N\),
\(NO^n\subset\Lambda_0=\Lambda\cap O^n\). On a geometric
characteristic-zero PEL fiber let \(A\) be the abelian variety with
symmetric polarizing line \(L\) and theta section \(s_0\).
Its period lattice \(\Pi\) is projective of rank \(r\) over \(O\).
Put \(d_P=\deg P\) for a nonzero polynomial in the integral harmonic
space of the series. Write \(\mathcal Q\) for the \(n\)-fold level
quotient and \(C_M=|\mathcal Q|=M^{frn}\). In fact the datum's
projective lattice of rank \(r\) is reduced modulo the half-level
ideal \(\mathfrak m\), with \(\mathfrak m\bar{\mathfrak m}=(M)\).
Hence \(\mathrm N\mathfrak m=M^f\), and local freeness, prime by
prime, gives the asserted cardinality. The Fourier weight
\(\omega_{x_0}(u)=\zeta_M^{\langle x_0,u\rangle_M}\) is a unit at
every finite place. Here the exponent is the integral level pairing.

The torsion set \(Z_\Lambda\subset A[N^2]^n\) consists of the
tuples \(z\) satisfying \((N\mu)z=0\) for all rows
\(\mu\in\operatorname{Hom}_O(\Lambda,O)\subset E^{1\times n}\).
The entries of \(N\mu\) are integral. For \(u\in\mathcal Q\), fix the
specified \(N\)-division point \(t_u\in A[N^2]^n\) of its level
point. Form the pullback of \(s_0^{\boxtimes n}\) by multiplication
by \(N\) on \(A^n\), then translate by the normalized theta lift of
\(z\), then by that of \(t_u\). Denote by \(J_P(u,z)\) the jet
evaluation of this section in the given Hodge frame. Explicitly,
for the individual translated sections \(F_b\) in the \(n\) factors
it is
\[
 J_P(u,z)=\sum_{\mathbf k} P_{\mathbf k}
       \prod_{b=1}^{n}\nabla_{\mathbf k_b}F_b .
\]
Here \(\mathbf k_b\) is the column of the monomial's multi-index at
\(b\); \(\nabla_{\mathbf k_b}\) is the prescribed linear theta
jet in that frame. The expression is multilinear in the factors,
hence depends on their tensor. For the raw theta section \(y\)
and the chosen scalar \(s\), the actual finite sum is
\[
 (s y)(x_0)=\frac{s}{C_M\,m_\Lambda N^{d_P}}
     \sum_{u\in\mathcal Q}\omega_{x_0}(u)
                      \sum_{z\in Z_\Lambda}J_P(u,z),\qquad
 m_\Lambda=[\Lambda_0:NO^n]^r .
\]
The notation suppresses the geometric point and Hodge frame. The
auxiliary determinant scalar \(\beta_H\) in the short evaluation
of \(\beta_H^{-1}y\) cancels when one returns to \(y\).
\src{src/Formalism/LocalToGlobal/ProjectiveLattices/ThShort.lean}{244}{Finite torsion evaluation with all three denominators};
\src{src/Development/AbelianVarietyTheory/Definitions/PlaceAdaptation/AdaptStd.lean}{107}{Harmonic evaluation on translated sections}.

\begin{lemma}[Exact cancellation on every level class]\label{newtheta:index}
Put \(K_0=Z_\Lambda\cap A[N]^n\). Then \(|K_0|=m_\Lambda\) and
\(J_P(u,z+v)=J_P(u,z)\) for every \(v\in K_0\). In particular
\(m_\Lambda^{-1}\sum_{z\in Z_\Lambda}J_P(u,z)\) is precisely
the sum over \(Z_\Lambda/K_0\), for each \(u\) separately.
\end{lemma}
\begin{proof}
Work first over \(\mathbb C\). Write \(z=x/N\bmod\Pi^n\).
The defining conditions say that \(\mu x\in\Pi\) for all these
rows, that is, \(x\in\Lambda\otimes_O\Pi\) inside
\((E\otimes_O\Pi)^n\). This last equality can be checked after
localizing at every prime of \(O\), where both lattices are free
and evaluation against the dual basis recovers all coordinates.
Self-duality and \(NO^n\subset\Lambda\) give
\(\Lambda\subset N^{-1}O^n\), so \(N^2z=0\) already holds.
The additional condition \(Nz=0\) is \(x\in\Pi^n\).
Since \(\Pi\) is locally free of rank \(r\), this identifies \(K_0\)
with \((\Lambda_0/NO^n)\otimes_O\Pi\), whose order is \(m_\Lambda\).

Let \(E_L\) be the integral alternating Riemann form of the
polarization, extended rationally. For \(x/N,x'/N\) in this
description the commutator of the two translations on
\((L^{\boxtimes n})^{N^2}\) is
\(\exp(\pm2\pi i\sum_i E_L(x_i,x'_i))\), by the Weil pairing
formula for the Riemann form. For \(x_i=a_i h,\ x'_i=b_i h'\),
with \(a,b\in\Lambda,\ h,h'\in\Pi\), the exponent before
\(\pm2\pi i\) is \(E_L(h,\langle a,b\rangle h')\) by the
Rosati involution. This is integral by self-duality; bilinearity
gives the same result for sums of tensors. Hence the lifts commute
on \(Z_\Lambda\).

For clarity, on odd \(N^2\)-torsion the normalized lift has
\(N^2\)-th power one and inversion sends it to its inverse.
Two such lifts at the same point differ by a scalar \(c\) with
\(c^{N^2}=c^2=1\), hence coincide. On a commuting subgroup the
product satisfies both normalizations, so the normalized lift there
is multiplicative. On \(A[N]^n\), use the canonical identification
of the pullback by \([N]\) with \((L^{\boxtimes n})^{N^2}\).
Translation by \(v\in A[N]^n\) on the pullback lifts by the
identity on the fiber at \([N]x=[N](x+v)\). This lift is
normalized, as its powers and inversion are the corresponding
translations on a pullback. It fixes the pulled-back section.
Consequently translation by \(z+v\) has exactly the same effect
as translation by \(z\) on that section. The subsequent translation
by \(t_u\) and the linear jet evaluation preserve equality.
The identities descend to arbitrary characteristic-zero geometric
fibers by defining the abelian variety, the finitely many points,
the line, the sections and the required finite jets over a finitely
generated subfield of characteristic zero and embedding it into
\(\mathbb C\). Mumford,
\emph{On the equations defining abelian varieties I}, \S1~\cite{mumfordEquations}.
\end{proof}

The scalar \(s\) in the same formula also admits a simultaneous
choice at the finite and infinite places. Write
\(I=[\Lambda:\Lambda_0]\), let \(F\) be the totally imaginary
coefficient field of degree \(h\) containing the coefficient
embeddings \(\tau:E\to F\), and set
\(\mathfrak a=\det(\Lambda_0)\det(\Lambda)^{-1}\), using the
fractional determinant ideals in \(E\). Then
\(\mathfrak a\subset O\) and \(\mathrm N\mathfrak a=I\):
over each discrete valuation ring the inclusion has elementary
divisors and both statements follow by multiplying those divisors.
Put
\[
 \mathfrak b=\prod_\tau
       \big(\tau(\mathfrak a)\mathcal O_F\big)^{m_\tau}.
\]
The finite scalar condition is \(s\in\mathfrak b^{-1}\), since
the quotient of the two type norms of determinants is exactly
\(\mathfrak b\). Each extended ideal in the product has norm
\(I^{h/[E:\mathbb Q]}\), by the residue-degree and ramification
identity above each prime. Since \(\sum m_\tau=fr\), we have
\(\mathrm N\mathfrak b=I^{hr/2}\).
The image of \(\mathfrak b^{-1}\) in \(\mathbb C^{h/2}\) has real
covolume \(2^{-h/2}\sqrt{|\Delta_F|}/\mathrm N\mathfrak b\),
by the discriminant determinant for an integral basis and the
fractional ideal index. Take the product of closed disks of common
radius \(R_s=|\Delta_F|^{1/(2h)} I^{-r/2}\). Its volume
\(\pi^{h/2}R_s^h\) exceeds \(2^h\) times that covolume because
\(\pi>2\). Minkowski's convex body theorem supplies a single nonzero
\(s\in\mathfrak b^{-1}\) with \(|\sigma(s)|\le R_s\) at every
embedding, including the conjugates. At every level prime
\(\ell\mid M\), local standardness of \(\Lambda\) gives
\((\Lambda_0)_\ell=\Lambda_\ell\); hence \(\mathfrak b\) is a unit
at all primes over \(\ell\), and this same \(s\) is integral there.
This proves the determinant rescaling bound on the actual family,
including its index dependence. It applies before the jet estimates.
\src{src/Development/AbelianVarietyTheory/Proofs/DivisorProofs/ExistsSmallScalarReflexIdealSelfDualLattice.lean}{29}{Single rescaling scalar for all embeddings and the determinant ideals}.

In fact \(|Z_\Lambda/K_0|=I^r\) by the same locally free tensor
calculation. Cauchy--Schwarz on the \(C_M I^r\) remaining terms
and \(|\sigma(\omega_{x_0}(u))|=1\) thus give at that frame
\[
 |\sigma(sy)(x_0)|^2
 \le \frac{|\Delta_F|^{1/h}}{C_M N^{2d_P}}
       \sum_{u,\,\bar z}|J_{P^\sigma}(u,\bar z)|^2 .
\]
The right side uses the conjugate geometric data at \(\sigma\).
In particular no separate factor \(I^r\) is lost in estimating the
sum by the squared jet values.

There is an exact test for the finite denominator.
Write \(w=\sum_{\tau,i}a_{\tau,i}\),
\(\delta_E=|\operatorname{disc}E|\), \(S=w+frn/2\), and
\[
 D_M=C_M(M\delta_E)^w.
\]
This is the positive integer \(\mathrm{intZ}\) of the construction,
whose proposed common multiplier is \(D_M^2\).
For any valuation \(v\) of an integral framed test, extended to the
finite torsion field, put
\[
 T_{x_0}=s\sum_{u\in\mathcal Q}\omega_{x_0}(u)
                   \sum_{\bar z\in Z_\Lambda/K_0}J_P(u,\bar z).
\]
At that test the desired integrality is \emph{equivalent} to
\[
 v(T_{x_0})\ge d_P v(N)+v(C_M)-2v(D_M).
\]
Thus at a prime over \(\ell\mid N\) with
\(\ell\nmid M\delta_E\), the required divisibility is by
\(N^{d_P}\) locally, even though \(\ell\) need not divide the level.
At a level prime the allowable loss is exactly
\(2v(D_M)-v(C_M)\) relative to \(d_Pv(N)\).
Integrality of the roots of unity and the geometric coset
cancellation impose no such divisibility on the harmonic jets.
Checking all these tests is stronger than checking an expansion at
one cusp unless a corresponding integral expansion principle for
the chosen lattice is supplied.

\begin{proposition}[The auxiliary integer at good abelian reduction]
\label{newtheta:goodjet}
Let \(\ell>\max(2,d_P)\), \(\ell\nmid M\delta_E\), and
\(\Lambda\otimes\mathbb Z_\ell=O^n\otimes\mathbb Z_\ell\).
At a mixed-characteristic discrete valuation ring \(R\) of residue
characteristic \(\ell\), take a framed test extending to an abelian
scheme equipped integrally with the rigidified symmetric line,
theta section, integral Hodge directions, polarized biextension
and its rigidified identification with \(L\) on the graph of the
specified endomorphism \(w_0\). Then for every allowed \(N\),
every \(u,z\), and the actual translated sections,
\[
                 v(J_P(u,z))\ \ge\ d_P v(N).
\]
The valuation may be extended to define the torsion points.
This proves the preceding denominator test at these places.
\end{proposition}
\begin{proof}
Put \(N_0=\operatorname{lcm}(M,\exp(O^n/\Lambda_0))\), taking
the exponent of the trivial group to be one. Then \(N_0\mid N\),
\(\ell\nmid N_0\), and \(k=N/N_0\) sends \(t_u,z\) to \(kt_u,kz\).
The first maps under \(N_0\) to the same level point. The second
belongs to \(Z_{\Lambda,N_0}\) by the description \(z=x/N\)
in Lemma~\ref{newtheta:index}. Both images are killed by \(N_0^2\).
For either translation by \(t\), lift translation by \(kt\) on
\(L^{N_0^2}\) and pull its line isomorphism back along \([k]\),
using \([k]\circ T_t=T_{kt}\circ[k]\).
The resulting lift at \(t\) is symmetric. Its \(N^2\)-th power is
the identity, since \(N^2t=0\) and \(N_0^2\mid N^2\).
By uniqueness of the normalized odd lift this is precisely the
one used in the test. The rigidified identifications for successive
multiplications compose, so both translations, in their prescribed
order, carry the pullback of \([N_0]^*s_0\) to the pullback of
the corresponding translated section.

Here is the jet calculation, including the frame correction.
In formal coordinates write \(e(x,y)\) for a normalized frame of
the polarized biextension. Its logarithmic correction \(h(x,y)\)
vanishes on the axes. The frame \(e\exp(-h)\) multiplies exactly
in each argument, since the two multiplication factors of \(e\)
are the exponentials of the respective additive defects of \(h\).
Pulling this corrected frame to \((x,w_0x)\) yields the frame in
which the jet of a section of \(L^m\) with raw coefficient \(f\)
is obtained by differentiating \(f\exp(mh(x,w_0x))\).
In changing the raw biextension frame to \(e b\), one can take
\(h+\log b\) for the correction; the raw coefficient changes to
\(f b(x,w_0x)^{-m}\), leaving the corrected germ unchanged.
Any further choice differs by a biadditive Rosati-compatible term,
whose contribution to the harmonic jet vanishes as shown below.
Biadditivity identifies the corrected frame at \((kx,kw_0x)\)
with the \(k^2\)-th tensor power of that at \((x,w_0x)\).
The induced isomorphism of lines is the rigidified one: it is
constructed by iterating the two global biextension laws and is
the identity in the rigidification at zero. In formal logarithm
coordinates \(x\), multiplication by \(k\) is \(x\mapsto kx\).
Each homogeneous derivative of degree \(j\) therefore gains \(k^j\).
The prescribed row sums of every monomial of \(P\) give the same total
degree \(d_P=\sum_{\tau,i}a_{\tau,i}\).
Suppressing \(u\) and displaying both translation points gives the
exact equality on the generic fiber
\[
       J_{P,N}(t_u,z)=k^{d_P} J_{P,N_0}(kt_u,kz).
\]

It remains to check the smaller-level jet over \(R\). All translations
now have order prime to \(\ell\). After finite extension their points
extend by properness. Each translated line is isomorphic to the
original line over \(R\): their quotient represents zero in the
relative Picard scheme by the polarization and \(N_0^2\)-torsion,
and a line from the local base is free. Any integral lift has
\(N_0^2\)-th power a unit. The scalar that changes it to the
normalized lift likewise has valuation zero, by taking that power.
Thus each of the smaller-level sections has integral raw germ in
an integral frame.

Work through total degree \(d_P\). The formal logarithm and its
inverse are integral to this order: integrating the closed invariant
differentials term by term divides each homogeneous component only
by its positive total degree (contract the form with the Euler
vector field), and inversion of a series with identity linear term
requires no division. In these coordinates the group law is
addition. Restrict both arguments \(X,Y\) to the linear combinations
of the integral eigen-directions of the frame, with variables
indexed by \((\tau,i)\). Write \(h\) also for this restriction.
The three series
\[
 \begin{split}
 h(X,Y)+h(X',Y)-h(X+X',Y),\qquad&
 h(X,Y)+h(X,Y')-h(X,Y+Y'),\\
 &h(aX,Y)-h(X,\bar aY)\qquad (a\in O)
 \end{split}
\]
are integral through that order. Indeed they are the logarithms
of the integral multiplication and Rosati adjunction factors of the
biextension frame, each of constant coefficient one; the truncated
logarithm has unit denominators. For a coefficient of \(X^\alpha
Y^\gamma\) with \(|\alpha|\ge2\), split off one variable \(X'_j\)
with \(\alpha_j>0\) in the first defect. Its coefficient is
\(-\alpha_j\) times the one sought, with \(\alpha_j\) a unit.
The second defect handles \(|\gamma|\ge2\). The axes vanish.
In bidegree \((1,1)\), the last defect multiplies the coefficient
in types \((\tau,\nu)\) by \(\tau(a)-\nu(\bar a)\).
For \(\tau\ne\bar\nu\) some \(a\in O\) makes this a unit:
since \(\ell\nmid\delta_E\), the distinct embeddings have distinct
reductions in the common splitting residue field. Thus, to this
order, \(h\) is integral up to a bilinear form supported on
the conjugate pairs.

Write \(X_b\) for the variables in factor \(b\).
After evaluation at \((X_b,w_0X_b)\) and summation over \(b\),
the possible nonintegral part is a linear combination \(Q\) of
\(\sum_b X_{\tau i b}X_{\bar\tau j b}\).
Under the pairing \(P(\partial)F(0)\), multiplication by such a
quadratic is adjoint to \(\sum_b\partial_{\tau i b}
\partial_{\bar\tau j b}\) on \(P\), by the monomial factorial formula.
These operators annihilate \(P\) by its defining harmonic equations.
Consequently the factor \(\exp(N_0^2 Q)\) contributes only its
constant term. The remaining exponentials, raw germs and coefficients
of \(P\) are integral through degree \(d_P\). For the coefficients
of \(P\) this uses local standardness of \(\Lambda\) in the
integral linear generators of its polynomial ring.
We obtain \(v(J_{P,N_0})\ge0\), and \(v(k)=v(N)\).
The same local standardness makes \(\mathfrak b\) a unit at this
place, so \(s\) is integral as well.
\end{proof}
For the standard lattice this works at every such good test with
\(N_0=M\), even when \(\ell\) divides \(N\) to arbitrary order.
We next remove the restrictions on the residue characteristic and
the local lattice at abelian reduction. The archimedean norm
of the smaller-level sections is a separate estimate.
\src{src/Formalism/LocalToGlobal/ProjectiveLattices/RowDeg.lean}{43}{Harmonic trace equations with identity column matrix};
\src{src/ReciprocityTheory/ReductiveGroups/ArithmeticGeometry/BirationalGeometry/ToroidalCompactification/CocycleRightZeroThd.lean}{174}{The two biextension defects and the adjunction factor};
\src{src/Development/AbelianVarietyTheory/Proofs/JetHarmonics/HarmonicValueNormalisedPullbackChainRule.lean}{516}{Normalized isogeny on harmonic evaluations}.

\paragraph{Finite denominators at abelian reduction.}
\label{intjet:sec}
Write \(d=\deg P=\sum a_{\tau i}\), \(\Delta=|\operatorname{disc}E|\),
and \(v\) for the additive valuation of a mixed-characteristic
discrete valuation ring \(R\). Valuations and all the inclusions below
may be tested after finite extension. We retain the integral symmetric
line, its rigidification, the section, the Hodge directions and the
polarized biextension on an abelian scheme \(A/R\) as in
Proposition~\ref{newtheta:goodjet}. The graph of \(w_0\in O_E\)
in that biextension gives the line. In a semiabelian degeneration the
generic abelian scheme need not extend in this form.

\begin{lemma}[Integral Taylor derivatives of harmonic jets]\label{intjet:der}
In the preceding formal chart take an integral raw coefficient,
any nonnegative tensor power of the line, and an integral homogeneous
harmonic polynomial of degree \(d\) at the standard lattice.
Its evaluated jet belongs to \(\Delta^{-\lfloor d/2\rfloor}R\).
The assertion holds in every residue characteristic.
\end{lemma}
\begin{proof}
Work over the fraction field in logarithm coordinates. For a series
\(F\) define its Taylor values to be the mixed derivatives
\(\partial^\alpha F(0)\), with no division by \(\alpha!\).
If \(F\) comes from an integral function on the formal group,
on integral tangent vectors these values lie in \(R\). Indeed
each vector extends by translation to an invariant integral
derivation of the formal group ring. In logarithm coordinates it
is the constant vector field, so applying a sequence of these
derivations and evaluating at zero gives exactly the desired value.
This argument works on a product and on any list of integral
directions, whether or not they form an integral basis.

Series with integral Taylor values form a ring by Leibniz's formula.
An inverse with unit constant term has integral Taylor values by the
same formula and induction. For \(g(0)=1\) and integral Taylor
values, those of \(\log g\) are integral by
\(\partial_i\log g=g^{-1}\partial_i g\). For \(h(0)=0\) with
integral Taylor values, use
\(\partial_i\exp h=(\exp h)\partial_i h\) and induction.
These are identities over the fraction field and give integrality
of derivatives without an estimate for any inverse factorial.

Let \(h(X,Y)\) be the logarithmic correction of the biextension in
the chosen frame. Its two additive defects and its \(O_E\)-adjunction
defect are logarithms of integral units with constant term one, as
in Proposition~\ref{newtheta:goodjet}. For an index \(\alpha\) of
degree at least two write \(\alpha=\alpha'+e_j\), \(|\alpha'|>0\).
Apply \(\partial_X^{\alpha'}\partial_{X'_j}\partial_Y^\beta\) to
\(h(X,Y)+h(X',Y)-h(X+X',Y)\) at zero. The result is
\(-\partial_X^\alpha\partial_Y^\beta h(0)\). This proves its
integrality. Interchanging the two arguments treats \(|\beta|\ge2\).
The two axes vanish. In bidegree \((1,1)\) write \(c_{\tau i,\nu j}\)
for the derivative along eigenvectors of types \(\tau,\nu\).
The adjunction defect gives
\((\tau(a)-\nu(\bar a))c_{\tau i,\nu j}\in R\) for every \(a\in O_E\).
If \(\tau\ne\bar\nu\), then \(\Delta c_{\tau i,\nu j}\in R\).
Indeed choose an integral basis \(b_i\) and its trace-dual basis
\(b_i^\vee\). The embedding matrices are inverse transposes, so
\(\sum_i\tau(b_i^\vee)\sigma(b_i)\) is one for \(\sigma=\tau\)
and zero otherwise. All \(\Delta b_i^\vee\) are integral by the
adjugate of the trace matrix. Apply this equality to
\(\sigma=\tau,\bar\nu\) and subtract.

Thus the gauge on the eigen-directions is the sum of a series with
integral Taylor values, a bilinear form with coefficients in
\(\Delta^{-1}R\) on the nonconjugate pairs, and an unrestricted
bilinear form on the conjugate pairs. On the graph of \(w_0\) these
are quadratic terms, summed over the product factors. The last
quadratics are killed by harmonicity as in
Proposition~\ref{newtheta:goodjet}. In the exponential of the
nonconjugate quadratic each derivative of order \(j\) at zero is
a sum over partitions into pairs of the \(j\) labelled derivative
positions. This follows by differentiating \(\partial_i\exp
Q=(\exp Q)\partial_i Q\); a singleton has zero value and derivatives
of \(Q\) of order greater than two vanish. Each summand costs at
most \(\lfloor j/2\rfloor v(\Delta)\). Leibniz's formula for the
remaining integral factors gives the claimed bound.
\end{proof}

\begin{proposition}[All abelian tests at a locally standard lattice]
\label{intjet:standard}
Suppose \(\Lambda_\ell=(O_E^n)_\ell\) and the integral abelian data
above are given. For every allowed \(N\) and each pair of translation
points in the original order one has
\[
 v(J_{P,N}(u,z))\ \ge\
 d\bigl(v(N)-v(M)\bigr)-\lfloor d/2\rfloor v(\Delta).
\]
In particular \(D_M^2\) clears the rescaled sum at every such test,
including the level primes, \(2\), and primes ramified in \(E\).
\end{proposition}
\begin{proof}
Let \(N_0=\operatorname{lcm}(M,\exp(O_E^n/\Lambda_0))\) as before.
Now \(v(N_0)=v(M)\). The identity
\(J_{P,N}= (N/N_0)^d J_{P,N_0}\) on the generic fibre proved in
Proposition~\ref{newtheta:goodjet} works whether or not \(\ell\mid N_0\).
All torsion points of the smaller test extend after finite extension
by properness. For \(t\in A[N_0^2]\) the line \(T_t^*L^{N_0^2}\)
is integrally isomorphic to \(L^{N_0^2}\). The class of their
ratio in the dual abelian scheme is zero since \(N_0^2t=0\),
and an invertible module on the local base is free. Any integral
translation lift has its \(N_0^2\)-th power in \(R^\times\).
The generic scalar adjusting it to the normalized lift therefore
has valuation zero. This applies to both translations in succession,
also for torsion of residue-characteristic-power order.
The smaller-level section has integral raw germ. Apply
Lemma~\ref{intjet:der}. Local standardness makes \(P\) integral in
the standard columns and \(s\) integral at this place. Sum after the
index cancellation of Lemma~\ref{newtheta:index} and use
\(D_M=C_M(M\Delta)^d\).
\end{proof}

\begin{theorem}[The trace at an arbitrary self-dual lattice at abelian reduction]
\label{intjet:trace}
Suppose the abelian data above extend integrally. Take any
\(\Lambda,N,P,s\) in the rescaled series, with \(\Lambda\) standard
at the primes of \(M\), and the prescribed level points \(t_u\).
At a prime \(\ell\nmid M\),
\[
 v\left(s\sum_{\bar z\in Z_\Lambda/K_0}J_{P,N}(u,\bar z)\right)
 \ \ge\ d\,v(N)-\lfloor d/2\rfloor v(\Delta)
 \qquad(u\in\mathcal Q).
\]
Together with Proposition~\ref{intjet:standard}, this proves the
finite denominator test at every residue characteristic on these
abelian models.
\end{theorem}
\begin{proof}
Let \(O=O_E\), \(B=A^n\), \(B_0=\Lambda_0\otimes_O A\) and
\(B_1=\Lambda\otimes_O A\). A full lattice over the Dedekind
ring \(O\) is finitely generated projective. Its tensor with \(A\)
is obtained by splitting an idempotent on a product of copies of
\(A\). This gives abelian schemes over \(R\). Their Lie algebras
are obtained by taking the same tensor with \(\operatorname{Lie}A\).
Consider the maps induced by
\[
 O^n\xrightarrow{\,N\,}\Lambda_0
 \xrightarrow{\,i\,}\Lambda,\qquad
 \Lambda_0\xrightarrow{\,j\,}O^n
\]
where \(i,j\) are inclusions. Write \(a_0,i,j\) also for the
resulting isogenies, and \(a=i a_0\). Then \(j a_0=[N]\).
These maps are isogenies on every fibre since suitable integer
multiplications factor through them in both directions. They
are finite flat by properness and flatness of a quasi-finite map
between smooth schemes of the same relative dimension.

Give each tensor scheme its rigidified line from the graph of
\(w_0\) in the polarized biextension for the restricted Hermitian
form. Denote the lines on \(B,B_1\) by \(\mathcal L,\mathcal L_1\).
Then \(j^*\mathcal L=i^*\mathcal L_1\) and
\(a^*\mathcal L_1=[N]^*\mathcal L=\mathcal L^{N^2}\)
by rigidified identifications. To see the normalization, on free
modules a matrix \(U\) pulls the biextension for a Hermitian
matrix \(G\) back to that for \(U^\dagger G U\). Expand the
biextension in its two arguments and move each first-argument
entry across by Rosati adjunction. These identifications are
the identity on the axes. Restriction to direct summands proves
the assertion for projective modules as well. This construction
uses the line on the graph, so no choice among symmetric lines
inducing the same polarization is needed. See Amir-Khosravi,
\S\S1.2--1.4~\cite{serreTensor}.

On the geometric generic fibre, the period lattice calculation of
Lemma~\ref{newtheta:index} gives \(\ker a=Z_\Lambda\) and
\(\ker a_0=K_0\). The natural action of \(\ker a\) on the pullback
line is exactly its normalized lift. It has \(N^2\)-th power one,
and inversion on the symmetric pullback sends the lift to its
inverse. These conditions determine the lift since \(N\) is odd.
In particular they agree with the products of the individual
normalized lifts on the \(n\) factors by the same uniqueness.
Put \(t_0=j^*(s_0^{\boxtimes n})\), viewed as a section of
\(i^*\mathcal L_1\). Before the level translation the sum of
the \(|Z_\Lambda/K_0|\) translated sections on \(B\) is
\[
                 a^*\operatorname{Tr}_i(t_0).
\]
Indeed a finite étale isogeny on the generic fibre has trace
equal, after pullback, to the sum over its kernel. The input
on \(B\) was \(a_0^*t_0=[N]^*(s_0^{\boxtimes n})\).

The determinant of \(di\) generates the ideal
\(\mathfrak b R=\prod_\tau
 (\tau(\det\Lambda_0\,\det\Lambda^{-1})R)^{m_\tau}\).
Choose bases of the two lattices over the semilocal Dedekind
ring \(O\otimes\mathbf Z_{(\ell)}\), where they are free.
The determinant on each generic eigenspace is the indicated
embedding of the determinant of the inclusion, repeated \(m_\tau\)
times. The chosen bases give integral changes of basis on the
Lie lattices, so this calculates their determinant ideal even
when the eigenspaces do not split integrally.
Translation of invariant differential frames carries this same
different ideal to every point. For a finite flat map between smooth
schemes of equal relative dimension the different is generated by
the determinant of its differential. Locally in the finite
algebra the trace functional is this generator times a generator
of the dualizing module, by the complete intersection formula
for the different~\cite{tateDifferent}. Trivializing \(\mathcal L_1\)
shows that \(\operatorname{Tr}_i(t_0)\) is divisible by
\(\mathfrak b R\) as a section. Thus \(S=s\operatorname{Tr}_i(t_0)\)
is integral, including when \(v(s)<0\).

It remains to apply the actual final translation. Its point
\(t=t_u\) is killed by a power of the half-level ideal, hence by
a power of \(M\); this extra condition is part of the evaluation
rule. In particular its order \(h\) is prime to \(\ell\).
\src{src/Formalism/LocalToGlobal/ProjectiveLattices/ThShort.lean}{252}{The support of the prescribed division point}.
Set \(b=a(t)\). Rigidify the line
\(\mathcal D=T_b^*\mathcal L_1\otimes\mathcal L_1^{-1}\) at zero
by choosing an integral generator of its fibre. It is an element
of \(\operatorname{Pic}^0(B_1)[h]\), since \(b\mapsto
T_b^*\mathcal L_1\otimes\mathcal L_1^{-1}\) is a homomorphism
to the dual abelian scheme. Thus \(\mathcal D^h\) has an integral
rigidified trivialization. On the formal neighborhood of zero
there is a unique \(h\)-th root frame with prescribed value at
zero and this trivialization as its \(h\)-th power. Coefficient
recursion only divides by the unit \(h\).

The normalized lift at \(t\) identifies the translation of
\(a^*S\) with the pullback of \(T_b^*S\), followed by an
isomorphism \(a^*\mathcal D\to\mathcal O_B\). This isomorphism
is integral by the translation-lift argument of
Proposition~\ref{intjet:standard}. Its \(h\)-th power and the
pullback of the chosen trivialization differ by a global unit
of the abelian scheme, hence by an element of \(R^\times\).
It takes the pullback of the root frame to a constant unit:
over the fraction field a formal unit whose \(h\)-th power is
constant is constant. In particular, in the formal neighborhood
the translated section is, up to a constant unit, \(a^*\) of
an integral germ of \(\mathcal L_1\), obtained by dividing the
germ \(T_b^*S\) by the root frame. This explains why the final
translation does not change the trace divisibility.

To compute the jet of this pullback, choose a split surjection
\(U:O^t\to\Lambda\) with section \(V:\Lambda\to O^t\),
extended rationally, and put \(G=U^\dagger U\). Its entries
are integral by self-duality. View \(U\) as an \(n\)-by-\(t\)
matrix in \(E^n\). Then \(UV=I_n\), \(V^\dagger G V=I_n\),
and \(NV\) has integral entries. The map \(a\) factors as \(NV\)
followed by \(U\). Pull the integral germ back by \(U\); its raw
coefficient in an integral frame of the line for \(G\) is integral.
In the corrected frame the chain rule for \(NV\) gives \(N^d\)
times the differential operator \(P_V\) obtained from \(P\)
by sending a column-direction \(v\) of type \(\tau\) to
\(\tau(V)v\). Indeed in logarithm coordinates the map is
linear on each eigenspace. The corrected biextension frame
multiplies in both arguments and intertwines the adjoint
\(O\)-actions by its defining identities. Expanding the matrix
entries in these identities identifies its pullback with the
corrected frame for \((NV)^\dagger G(NV)=N^2 I_n\).
Thus the comparison differentiates the full section with the
same correction as in the original jet.

The coefficients of \(P_V\) are integral. The defining integral
polynomial ring is generated by the linear directions from
\(\Lambda\), and \(V\Lambda\subset O^t\).
\src{src/Formalism/LocalToGlobal/ProjectiveLattices/ThShort.lean}{235}{Integral polynomial directions of a nonstandard lattice}.
For the Hermitian matrix \(G\) the quadratic ambiguity in the
corrected log frame comes from summing
\(h(X_j,w_0G_{jk}X_k)\) over its columns. For a conjugate pair
of types its contraction on \(P_V\) uses the coefficients
\(\bar\tau(G_{jk})\). By the ordinary chain rule and
\(V^\dagger G V=I_n\), it is the image under \(V\) of the standard
contraction on \(P\); hence it vanishes. All other terms obey
Lemma~\ref{intjet:der} because the entries of \(G\) and \(w_0\)
are integral. This costs at most \(\lfloor d/2\rfloor v(\Delta)\).
The factor \(N^d\) proves the stated inequality.
\end{proof}

\paragraph{The norm of the generator at semiabelian reduction.}
Let \(\mathcal A/R\) be the semiabelian model in an integral test,
with abelian generic fibre \(A\), rigidified cubical line
\(\mathcal L\), integral section \(s_0\), and integral tangent
vectors. Its principal polarization is \(\lambda\). Write
\(\mathcal P\) for the rigidified Poincaré line of \(A\times A^t\).
In this test \(L=(1,\lambda w_0)^*\mathcal P\), with
\(w_0+\bar w_0=1\) and \(\bar w_0\) the Rosati adjoint of \(w_0\).
Indeed these are the graph-line and adjunction conditions of the
datum. The graph identification can be made rigidified by its value
at zero. A rigidified isomorphism of cubical lines on an abelian
variety respects the cube structure by uniqueness.
\src{src/ReciprocityTheory/ReductiveGroups/ArithmeticGeometry/BirationalGeometry/ToroidalCompactification/TwoWeightEqOneAddRoot.lean}{82}{Graph-line condition};
\src{src/Formalism/LocalToGlobal/ProjectiveLattices/VAtVNeBot.lean}{53}{The integral model, eigendirections and theta generator}.

\begin{lemma}[The maximum of the canonical theta norm]\label{intjet:maximum}
On the analytification over the completed algebraic closure, the
canonical norm of \(s_0\) is at most one everywhere. Its maximum
is attained at the point of the skeleton belonging to the identity
component of the special fibre.
\end{lemma}
\begin{proof}
Use the valuation \(v=-\log|\cdot|\) with normalized uniformizer.
One may make a valued extension for this calculation. For the split
Raynaud uniformization let \(Y\) be the period lattice,
\(\Phi:Y\longrightarrow X\) its identification with the character
lattice under the principal polarization, and \(b\) the positive
definite monodromy form on \(Y_{\mathbb R}\).
For any rigidified ample line \(H\) giving \(\lambda\) take its
descent data on the Raynaud extension. Denote the constant of descent
at \(u\in Y\) by \(c_H(u)\) and the period pairing by \(t\), using
the convention
\[
 c_H(u+v)=c_H(u)c_H(v)t(u,\Phi v),\qquad
 -\log\|t(u,\Phi v)\|=b(u,v).
\]
Products in this identity use the biextension identifications of
the fibres. The Poincaré line pulled back by \(1\times\lambda\)
is \(m^*H\otimes p_1^*H^{-1}\otimes p_2^*H^{-1}\), with its
rigidification. Pulling back by the graph of \(w_0\) and taking
the quotient of the three descent constants therefore gives
\[
 c_L(u)=\frac{c_H(u+w_0u)}{c_H(u)c_H(w_0u)}
       =t(u,\Phi w_0u),\qquad
 -\log\|c_L(u)\|=b(u,w_0u)=\tfrac12 b(u,u).
\]
The last equality uses \(b(w_0u,v)=b(u,\bar w_0v)\), by functoriality
of the dual period pairing, and symmetry of \(b\).

Here are the precise standard uniformization results used to pass
from this computation to a metric. The balanced translate in
de Jong--Shokrieh, Proposition~8.3, has descent valuation
\(b(u,u)/2\). Translation by a point of tropical coordinate \(z\)
adds \(b(u,z)\) to the descent valuation (Lemma~8.1).
Thus the balancing point for the present triple has \(z=0\).
It acts trivially on the skeleton. The symmetry of both lines makes
its image a two-torsion point, since \(T_y^*L\simeq T_{-y}^*L\)
implies \(\lambda(2y)=0\). Their canonical metrics correspond under
translation up to a constant from rigidification (Lemma~5.1).
For the balanced translate Proposition~9.1 and the tropical
Riemann theta formula give on \(Y_{\mathbb R}/Y\)
\[
        -\log\|s_0(\nu)\|=\mathrm{constant}
                      +\tfrac12\min_{u\in Y} b(u+\nu,u+\nu).
\]
Its minimum occurs at \(\nu=0\bmod Y\). Corollary~10.4 gives
the same maximum of the norm over the whole analytification.
Proposition~10.3 at the identity component identifies the negative
log norm there with the order of the original section along that
component in its cubical model. This order is nonnegative by
integrality. This last application uses the original rigidification,
so no unspecified metric constant enters the inequality.
See de Jong and Shokrieh, \S\S6--11~\cite{deJongShokrieh}, and the
Raynaud theta descent in Foster, Rabinoff, Shokrieh and Soto,
\S\S3--4~\cite{fosterTheta}.
\end{proof}

\begin{lemma}[Integral traces on identity discs]\label{intjet:disc}
Let \(i:C_0\to C_1\) be an isogeny of abelian generic fibres induced
by an inclusion of full projective \(O_E\)-lattices tensored with
\(\mathcal A\). Write \(\mathfrak d R\) for the determinant ideal
of its differential on the integral Lie algebras. Let \(L_1\)
be a symmetric cubical line on \(C_1\) with its canonical metric
and integral cubical extension to the identity model. Suppose
\(u\) is a section of \(i^*L_1\) bounded by one everywhere
in the pullback metric. Write \(T=\operatorname{Tr}_i u\).
Take \(b\in C_1[h]\), where \(h\) is odd and prime to the residue
characteristic, choose \(c\) with \(hc=b\) and \(h^2c=0\),
and suppose the normalized translation at \(c\) exists on
\([h]^*L_1\). On the formal disc at zero the raw coefficient of
\(\widetilde T_c[h]^*T\) in an integral frame of \([h]^*L_1\)
is divisible by \(\mathfrak d R\).
\end{lemma}
\begin{proof}
Extend the ground field finitely to contain the points of the
relevant finite fibres. The homomorphism of completed identity
discs of the two models is finite flat. For details, each smooth
disc has ring \(R[[X_1,\ldots,X_g]]\). On special fibres the
kernel of \(i\) is finite since an integer multiplication factors
through \(i\) and acts as an isogeny of semiabelian varieties.
Consequently the images of target parameters and a uniformizer
generate an ideal primary to the maximal ideal of the source.
Lifting a basis of the finite quotient by this ideal and iterating
in the ideal gives finite generation over the complete target ring.
These \(g+1\) parameters form a regular sequence in the regular
local source. The local criterion for flatness, successively along
the parameters, gives a finite free module. This finite, finitely
presented flat map is a local complete intersection at the closed
point: its closed fibre is cut out by the regular sequence of \(g\)
specialized parameters in the smooth special-fibre disc. The complete
intersection locus is open, hence contains the whole local spectrum.
The finite module of relative differentials is computed by the
formal Jacobian. Indeed derivations into finite modules are
continuous for the maximal-ideal topology by the product rule and
these modules are complete. In integral translation-invariant
differential bases the Jacobian determinant has constant ideal
\(\mathfrak d R\). Thus the trace of any power series in the source
lies in \(\mathfrak d R[[Y]]\) by the different formula in
Theorem~\ref{intjet:trace}.

The rigid generic fibre of either completed local ring is its open
unit polydisc. The disc map just constructed is surjective there:
the finite free map has positive rank, and points lying above a
point reducing to zero reduce to zero in this local map. More
explicitly the primary-ideal identity gives, for some \(e\),
\(X_j^e\in(\pi,i_1(X),\ldots,i_g(X))\); thus every point of the
finite formal fibre above such a point has \(|X_j|<1\). Finite
base change, or evaluation first on closed subdiscs, gives the
analytic covering. It is étale in characteristic zero, and each
fibre is a coset of the subgroup \(K^\circ\subset\ker i\) lying in
the identity disc. The full preimage over that disc in the abelian
generic fibre is the disjoint union of its translates indexed by
\(\ker i/K^\circ\). Indeed choose a disc preimage of a given
point. All other preimages differ by elements of \(\ker i\), and
two translated discs meet precisely when their translates differ by
an element of \(K^\circ\).

An integral frame of a symmetric cubical line on the identity disc
has canonical norm one. To see the metric comparison, on the generic
fibre of the formal completion along the entire special fibre take
the model metric. This domain is covered by finitely many compact
affinoids and is stable under doubling. The rigidified cubical
isomorphism for \([2]\) implies that the logarithm \(F\) of the
ratio of canonical and model metrics satisfies \(F(2x)=4F(x)\).
It is bounded on this domain, so iteration gives \(F=0\).
An integral formal frame differing from a model frame by a power
series with unit constant term has the same norm on the identity
disc. The canonical metric is functorial for rigidified isogeny
pullbacks by the same doubling argument on the compact generic
fibre.

A finite-order lift \(U\) of translation at any torsion point
preserves the canonical metric on its symmetric line. Choose
\(m>1\) congruent to one modulo the order of the point. Translation
commutes with \([m]\); comparison of the two line isomorphisms
on the proper connected generic fibre differs by a scalar.
The continuous logarithmic distortion therefore satisfies
\(F([m]x)=m^2F(x)+\mathrm{constant}\). Iteration on the compact
space makes \(F\) constant, and the finite order of \(U\) makes it
zero. This also applies to the normalized lift at \(c\).

In the identity frame on the target, view the translated pullback
as the trace above \(y\mapsto b+[h]y\), using that lift to identify
the line. Multiplication by \(h\) is an isomorphism on either
identity disc, by the integral formal inverse function theorem.
For each class in \(i^{-1}(b)/K^\circ\), choose a representative
\(x\). The preimages on its sheet are parametrized by
\(x+[h]z\), with \(i(z)=y\). On this disc the coefficient of
\(u(x+[h]z)\) in the pulled-back translated frame is analytic of
absolute value at most one, by the metric calculation. Every such
function on the open unit polydisc has an integral power series:
apply the Gauss norm bound on closed polydiscs of radii tending to
one, after a valued extension if needed. The coefficients already
belong to the field of definition. The sum over this sheet is the
trace of that power series under the finite free formal map \(i\).
On geometric generic fibres this is the usual sum over the sheet,
which identifies the two analytic traces. Each is divisible by
\(\mathfrak d R\), so their finite sum has the same divisibility.
\end{proof}

\begin{theorem}[Finite denominators at semiabelian tests]\label{intjet:semi}
At every integral test with the stated semiabelian, cubical and
Hodge data, \(D_M^2\) clears the rescaled finite theta sum. The
valuation inequalities of Proposition~\ref{intjet:standard} at
locally standard lattices and of Theorem~\ref{intjet:trace} away
from \(M\) hold also at semiabelian reduction.
\end{theorem}
\begin{proof}
At a locally standard lattice use \(N_0\) and the generic-fibre
identity of Proposition~\ref{intjet:standard}. The norm of the
generator is at most one everywhere by Lemma~\ref{intjet:maximum}.
Its normalized pullback by \([N_0]\) and its two finite-order
translates have the same bound, by functoriality and the isometry
proved in Lemma~\ref{intjet:disc}. In the integral formal frame
at zero their raw germs therefore have integral coefficients by
the polydisc argument. Lemma~\ref{intjet:der} applies on a smooth
semiabelian formal group verbatim: it uses the formal group law,
the integral tangent directions and the cubical adjunction.
This proves the individual inequality, also when \(\ell\mid M\).

For \(\ell\nmid M\) use the notations \(a_0,i,j,a\) in
Theorem~\ref{intjet:trace}. The tensor models are direct summands
of products of the smooth semiabelian scheme. An idempotent on a
smooth group gives a product decomposition by its image and the
image of its complement: these are the closed kernels of the
opposite projections, and addition and the pair of projections
are mutual inverses. Each factor is smooth by the lifting
criterion and has semiabelian fibres as a direct factor of such a
group. The graph-line construction on these models uses only the
integral cubical biextension and its \(O_E\)-adjunction. On the
generic fibre the identities of pullback lines and normalized
translations give, after summing the first translation,
\(a^*\operatorname{Tr}_i(j^*s_0^{\boxtimes n})\), as before.
The section inside this trace has norm at most one everywhere,
since the rigidified line identifications are isometries.

The final point \(t\) has odd order \(h\) prime to \(\ell\).
Take \(t'\) with \(ht'=t\), and set \(c=a(t')\), \(b=a(t)\).
The normalized lift at \(c\) on \([h]^*\mathcal L_1\) exists:
\(h^2c=0\) and \([h]^*\mathcal L_1=\mathcal L_1^{h^2}\).
Set \(T=\operatorname{Tr}_i(j^*s_0^{\boxtimes n})\).
The sections
\[
       [h]^*\widetilde T_t\,a^*T,\qquad
           a^*\widetilde T_c\,[h]^*T
\]
agree up to a root of unity. Both arise from the same translation
\(t'\) on \([h]^*a^*\mathcal L_1\), by pulling back the respective
lifts and using \([h]T_{t'}=T_t[h]\) and \(aT_{t'}=T_ca\).
These two lifts have finite order: a power divisible by the point
order and the order of the original lift is the identity. Their
ratio is a scalar on the proper generic fibre, hence is a root
of unity. Lemma~\ref{intjet:disc} gives divisibility by
\(\mathfrak b R\) for the section on the right before applying
\(a^*\), since the determinant on integral Lie algebras has
exactly the same ideal computed in Theorem~\ref{intjet:trace}.
Multiplication by \(s\in\mathfrak b^{-1}\) gives an integral raw
germ. Apply the integral polynomial and harmonic chain rule for
\(a=U(NV)\) in that theorem. The extra pullback \([h]\) on the
left multiplies a homogeneous corrected jet by \(h^d\), a unit.
The factor \(N^d\) and the loss of at most
\(\lfloor d/2\rfloor v(\Delta)\) follow as claimed. The index
cancellation and the level factor \(D_M=C_M(M\Delta)^d\) finish
the original finite-place test.
\end{proof}


The size of the multiplier itself is unconditional under the stage
bounds \(\log M\le30L_0,\ \log\delta_E\le7fL_0\),
\(f,r,L_0\ge1\). Indeed \(frn+w\le2S,\ w\le S\), so
\[
 \log D_M^2
 =2(frn+w)\log M+2w\log\delta_E
 \le134 f L_0 S\le134(frL_0)^2 S .
\]
This bound is independent of \(N,\Lambda,P,x_0\) for the fixed
weight and level. Theorem~\ref{intjet:semi} supplies the valuation
test for the jets with the scalar \(s\) allowed by the determinant
ideals. In the source the proof uses valuations of theta
derivatives in integral semiabelian charts and a local isogeny on
the primary torsion subgroup.
\src{src/Development/AbelianVarietyTheory/Proofs/DivisorProofs/LogShortThetaSeriesDenominatorBoundLe.lean}{102}{Size of the common clearing integer};
\src{src/Development/AbelianVarietyTheory/Proofs/DivisorProofs/SmulIntZSqNonnegOfSqDvd.lean}{71}{Normalized jet sum at a square divisor of the level};
\src{src/Development/AbelianVarietyTheory/Proofs/RingsOfIntegers/SumHvPrimaryKernel.lean}{221}{Local isogeny and integral formal theta jets}.

\paragraph{An archimedean estimate on each complex theta fibre.}
Fix a complex embedding \(\sigma\) of the coefficient field \(F\),
of degree \(h\) over \(\mathbf Q\). Use the cubical metrics
rigidified to have norm one at zero on the complex graph lines and
Haar probability on each abelian variety. Put \(g=fr=\dim A\),
\(m=gn\), \(d=\deg P\), and \(\Delta=|\operatorname{disc}E|\).
Let \(\lambda>0\) be the shortest nonzero period of \(A\) in the
positive Hermitian metric \(H\) of its principal polarization,
normalized by the Riemann form \(\Im H\). Take an orthogonal Hodge
eigenframe whose individual direction vectors have length at most
\(U\), with \(U>0\). Use the row copies of that frame in \(A^n\).
Let \(\|P^\sigma\|_1\) denote the sum of absolute values of its
coefficients in these directions. The integer \(N\) and the self-dual
Hermitian column lattice \(\Lambda\) can vary at this fixed frame.

\begin{proposition}[A complex fibre bound independent of the auxiliary level]
\label{intjet:arch}
With the rescaling \(s\) chosen above by Minkowski's theorem, put
\[
 C=2fn\Delta,\qquad
 R=\frac{\min(1,\lambda/C)}{4\sqrt m}.
\]
Then the cubically normalized value of the rescaled theta series satisfies
\[
 \frac{|\sigma(sy)(x_0)|}{\|s_0\|_2^n}
 \le |\Delta_F|^{1/(2h)}e^{\pi/2}
       \|P^\sigma\|_1\,d!\,(\sqrt2\,U)^d R^{-(d+m)}.
\]
This estimate is uniform in \(N,\Lambda\), the choice of division
points and the level character.
\end{proposition}
\begin{proof}
Use the isogeny \(a:B=A^n\to B_1=\Lambda\otimes_O A\) from
Theorem~\ref{intjet:trace}; in rational column coordinates its
differential is \(N\) times the identity. Both graph lines on
\(B\) and \(B_1\) induce principal polarizations. For the latter
this follows by tensoring the principal polarization of \(A\) with
the isomorphism from \(\Lambda\) to its conjugate \(O\)-dual
given by the unimodular Hermitian pairing. The equality
\(a^*\mathcal L_1=[N]^*\mathcal L\) is isometric: the two
curvature forms agree and the ratio of metrics on a compact complex
torus is then constant, fixed at zero. A finite-order normalized
translation lift is also an isometry, since the ratio is constant
and a positive real number of finite multiplicative order is one.
All the \(L^2\) pullbacks under isogenies preserve norms because
the pushforward of Haar probability is Haar probability.

Write \(v=[N]^*s_0^{\boxtimes n}\), and \(U_z\) for the normalized
translation. The vectors \(U_zv\), \(z\in Z_\Lambda/K_0\), are
pairwise orthogonal. Indeed \(v\) is fixed by the deck translations
\(b\in B[N]\), and \(U_zv\) is then an eigenvector for each
\(U_b\), with character given by the theta commutator with \(z\).
The annihilator of \(B[N]\) inside \(B[N^2]\) for this commutator
is \(B[N]\) itself. In a symplectic basis for the principal
polarization, the pairing is the standard perfect alternating
pairing modulo \(N^2\), and \(B[N]\) consists of all multiples of
\(N\). Distinct classes modulo \(K_0=Z_\Lambda\cap B[N]\) therefore
give distinct characters of commuting unitaries. The trace identity
of Theorem~\ref{intjet:trace} gives
\[
 a^*T=\sum_{z\in Z_\Lambda/K_0}U_zv,\qquad
 \|T\|_2= I^{r/2}\|s_0\|_2^n,\qquad I=[\Lambda:\Lambda_0].
\]
In particular \(\|\sigma(s)T\|_2\le
|\Delta_F|^{1/(2h)}\|s_0\|_2^n\). This equality of squared norms
uses the original order of the translations. The \(t_u\) act
after this calculation.

We bound the shortest period on \(B_1\). On \(\Lambda\) take the
real trace norm \(\|v\|_{\rm tr}^2=
\operatorname{Tr}_{E/\mathbf Q}\langle v,v\rangle\).
For \(v\ne0\) its squared norm is at least one, since the trace
is a positive rational integer. The trace-dual lattice equals
\(\mathfrak D_E^{-1}\Lambda\), by \(O\)-duality and the definition
of the codifferent. Multiplication by \(\Delta\) sends it into
\(\Lambda\), so its minimum is at least \(1/\Delta\).
Euclidean transference in dimension \(2fn\) gives \(2fn\) independent
vectors of \(\Lambda\) of trace norm at most \(C\)
\cite[Theorem 2.1]{banaszczyk}. Contractions with those vectors
separate every nonzero \(x\in\Lambda\otimes_O\Pi\), where
\(\Pi\) is the period lattice of \(A\): their rational span is the
whole column space over \(\mathbf Q\), and the pairing is perfect
over \(E\). Each contraction belongs to \(\Pi\). Its operator norm
on the complex tangent space with product Hermitian norm is at
most \(C\), by Cauchy--Schwarz on each eigenblock; on a given block
the sum of the squared column coefficients is at most the squared
trace norm. The positive form on \(B_1\) in rational column
coordinates is this same product form by the isometry of pullbacks
and \(da=N\,\mathrm{id}\). Thus every nonzero period on \(B_1\)
has norm at least \(\lambda/C\).

Use orthonormal eigen-coordinates on the universal cover of \(B_1\)
centered at any prescribed point. The polydisc of coordinate radius
\(R\) injects, by the period bound. In it a holomorphic line frame
of squared norm \(\exp(-\pi\|z\|^2)\) exists and is unique up to
unit phase. Indeed its prescribed logarithmic norm has the required
curvature, and a real pluriharmonic discrepancy on the simply
connected chart is the real part of a holomorphic function.
The Euclidean covolume of the periods in these coordinates is one:
the real volume form is \((\Im H_1)^m/m!\) up to orientation,
and its period integral for a principal polarization is one.
The same holds for any translate of the line and any choice of
center. If \(G=\sum c_\alpha z^\alpha\) is the coefficient of a
section of \(L^2\)-norm \(D\) in such a frame, then
\[
 |c_\alpha|^2\,
 \frac{\pi^mR^{2|\alpha|+2m}}{\prod_i(\alpha_i+1)}
 \le \int_{\{|z_i|<R\}}|G|^2\,dz\le e^\pi D^2 .
\]
The first inequality is angular orthogonality followed by radial
integration. The second uses \(mR^2\le1\) and injectivity. Since
\(\prod(\alpha_i+1)\le2^{|\alpha|}\), derivatives of degree \(d\)
in the row frame are bounded by
\(e^{\pi/2}D\,d!(\sqrt2 U)^d R^{-d-m}\).
Here each row direction is a multiple of an orthonormal coordinate
direction. Summing with \(P\) introduces \(\|P^\sigma\|_1\).

To compare this local derivative with the harmonic jet, on the
vector-space cover choose the holomorphic frame of the polarized
biextension with squared norm
\(\exp[-\pi(H(z,w)+H(w,z))]\), normalized at the origin.
Its curvature is that of the polarized biextension, by taking the
pullback under addition and dividing by the two projections.
Addition in either variable and Rosati adjunction preserve this
frame: their metric factors agree by sesquilinearity, and the ratio
is holomorphic of unit modulus and is one on the axes. On the graph
\(w=w_0z\) its squared norm is \(\exp(-\pi H(z,z))\) since
\(w_0+w_0^\dagger=1\). Make this frame comparison first on the
single factor \(A\times A\). Any permitted additive frame there
differs by the exponential of a Rosati-compatible complex bilinear
form, hence only couples conjugate types. On \(B=A^n\) the jet
uses the exterior product of the same single-factor frame in all
columns, followed by its \(N^2\)-th tensor power. Its graph correction
is therefore a sum of scalar multiples of
\(\sum_b z_{\tau i b}z_{\bar\tau j b}\). Each such sum disappears
under the harmonic equations, as in Lemma~\ref{intjet:der}.
For the final translation
by \(t_u\), apply the centered frame at \(a(t_u)\) to the translated
germ. Its pullback has squared frame norm
\(\exp(-\pi N^2\|z\|^2)\); the isometry \(U_{t_u}\) identifies it
with the frame at zero up to unit phase, by uniqueness. The derivative
then gains precisely \(N^d\), which cancels the normalizing
denominator in the series. Finally average over \(u\). Its Fourier
weights have modulus one and the divisor is their number \(C_M\).
This proves the inequality.
\end{proof}

Write \(S=d+m/2\) and \(Z_0=frL_0\). Suppose the frame satisfies
\[
 \max\{\log^+(\lambda^{-1}),\log^+ U,\log^+\|s_0\|_2\}\ll Z_0^2,
 \qquad \log^+\|P^\sigma\|_1\ll Z_0^2S .
\]
Then the log of the absolute bound is \(O(Z_0^2 S)\) under the
stage bounds above.
Here \(\log^+x=\max(0,\log x)\) for \(x>0\); a zero polynomial is
immediate. In fact \(m\le2S\), \(d\le S\),
\(\log^+(d!)\le d\log(1+d)\), \(\log(2S)\ll L_0\),
\(\log\Delta\le7fL_0\) and \(h^{-1}\log|\Delta_F|\ll Z_0^2\).
These bound every factor in the displayed inequality, including
\(\log(1/R)\ll Z_0^2\). To pass from this complex fibre estimate
to the simultaneous section norms requires the normalized algebraic
and analytic readings on every component, short polynomials spanning
the actual space, and the indicated frame control at its half
maxima. The normalized norm of the chosen generator matters when
that generator is identified with a frame of the determinant line.
\src{src/Development/AbelianVarietyTheory/Proofs/CuspFormEvaluations/NormalisedThetaValueBound.lean}{146}{Normalization on the tensor torus};
\src{src/Development/AbelianVarietyTheory/Proofs/DivisorProofs/NormRawThetaSeriesIndexBoundLeSqrt.lean}{18}{Square-root index in the archimedean estimate}.


\paragraph{The component Bergman bound and its volume normalization.}
Fix one good complex embedding of the coefficient field. Write
\(f=[E:\mathbb Q]/2\), \(\Delta=|\operatorname{disc}E|\),
\(R_E=\Delta^{1/(2f)}\), \(r\ge2\) for the Hermitian rank,
\(d\) for the positive host integer and \(M\) for the deep level.
Let \(w\ge1\) bound the sum, over the CM type, of the spreads
\(\max_i\lambda_{\tau i}-\min_i\lambda_{\tau i}\) of the flag weights.
The complex dimension of the product of full flag domains is
\(b=fr(r-1)/2\). The constants in the component estimate are
\begin{align*}
 \rho&=[200r^3(w+f)]^{-1},&
 \varepsilon&=\min\{1,\pi/[2M\sqrt d\,(fw+2)\log2]\},\\
 A&=(1+2(2fr)^2\Delta/\varepsilon)^{2fr^2},&
 V_*&=(15M)^{2fr^2}[3(2f)^2r^4R_E^2]^{6fr^2},\\
 C_{\rm loc}&=\frac{8\,b!}{(\pi\rho^2)^b}\,A,&
 \beta_*&=C_{\rm loc}V_* .
\end{align*}
Here \(V_*\) is the upper bound called \(\mathrm{massLev2}\) in
the Tamagawa comparison. The normalizing mass itself is the sum
of the volumes of the chosen arithmetic fundamental sets.
\src{src/Development/AbelianVarietyTheory/Definitions/Divisors/UnifOfSpaceEq.lean}{103}{Flag radius, return bound and local factor};
\src{src/Development/WeilBound/RepresentationTheoryGL/TubeMeasures/DisjointTubes/MassLev2.lean}{13}{Level-dependent total mass bound};
\src{src/Development/WeilBound/RepresentationTheoryGL/TubeMeasures/DisjointTubes/ConstantSupremum/SupremumConstantLogBound.lean}{14}{Product defining the normalized supremum constant}.

\begin{lemma}[Evaluation on a flag ball]\label{bergvol:local}
At an orthonormal frame \(x\), let \(B_x\) be the lower unipotent
flag chart of Euclidean radius \(\rho\). Write \(d\mu\) for the flag
measure specified below. For every holomorphic flag section \(P\),
\[
        |P(x)|_{\rm Pet}^2\le
        \frac{2b!}{(\pi\rho^2)^b}\int_{B_x}|P|_{\rm Pet}^2\,d\mu .
\]
\end{lemma}
\begin{proof}
In the frame \(x\) each factor has a diagonal Hermitian matrix
\(J\), with the positive signs first and the negative signs last.
Write the chart matrix as \(U=I+Z\) with \(Z\) strictly lower
triangular; the entries over all factors are the Euclidean coordinates.
In particular \(\|Z\|_{\rm op}\le\|Z\|_{\rm HS}\le\rho\).
Let \(D_{\tau i}\) be the absolute determinant of the leading
\(i\)-by-\(i\) Gram block of \(U^*JU\). The endpoint determinants
\(D_{\tau0},D_{\tau r}\) equal one. On the chart the squared norm
and invariant measure in question are
\[
 |P|_{\rm Pet}^2=|P(xU)|^2
 \prod_{\tau}\prod_{i=0}^{r-1} (D_{\tau,i+1}/D_{\tau i})^{-\lambda_{\tau i}},
 \qquad
 d\mu=\left(\prod_{\tau}\prod_{i=0}^{r-1}D_{\tau i}^{-2}\right)dV .
\]
The difference of each Gram block from the corresponding sign matrix
has operator norm at most \(2\rho+\rho^2\le3\rho\). Its singular
values therefore lie in \([1-3\rho,1+3\rho]\), and the block has
the required inertia by continuity. Thus the whole ball belongs to
the flag domain and \(|\log D_{\tau i}|\le6r\rho\).
Subtract the midpoint of the extreme weights from each row of
weights at \(\tau\). This leaves the logarithm of the norm factor
unchanged by telescoping the endpoint determinants. Its absolute value
is at most \(6r^2\rho w\le1/20\); the absolute logarithm of the
density is at most \(12fr^2\rho\le1/10\). At the origin both
factors equal one. The holomorphic mean-value inequality on the
Euclidean ball gives the assertion since its volume is
\((\pi\rho^2)^b/b!\) and \(\exp(1/20+1/10)<2\).
\end{proof}

\begin{lemma}[Overlaps at a frame with no short lattice vector]\label{bergvol:returns}
Suppose the arithmetic group preserves a Hermitian \(O_E\)-lattice
\(\Lambda'\subset E^r\) self-dual for the \(O_E\)-valued form.
Let \(\|v\|_x^2\) be the sum of the positive majorants of that form
over the fixed CM type, obtained by changing the negative signs to
positive in \(x\). If its minimum on \(\Lambda'\setminus\{0\}\)
is at least \(\varepsilon\), at most \(A\) arithmetic elements
return \(B_x\) to itself, meaning \(\gamma B_x\cap B_x\ne\varnothing\).
\end{lemma}
\begin{proof}
Set \(m=2fr\), the real lattice dimension. The Euclidean dual of the
embedded lattice is \(2J\) times its dual for
\(\operatorname{Tr}_{E/\mathbb Q}h\), since that trace is twice
the real inner product after applying the sign involution \(J\).
The trace dual lies in \(\Delta^{-1}\Lambda'\): testing against
\(O_E\Lambda'\) puts each Hermitian pairing in the codifferent,
which is carried into \(O_E\) by \(\Delta\), as seen by inverting
the integral trace matrix. Hence the Euclidean dual minimum in length
is at least \(2\sqrt\varepsilon/\Delta\).
The Euclidean transference inequality
\(\lambda_m(L)\lambda_1(L^*)\le m\) applies to this full rank lattice.
It gives \(m\) independent vectors of length at most
\(m\Delta/(2\sqrt\varepsilon)\). They span \(E^r\) over \(E\):
otherwise their real span would have dimension less than \(m\).
Select an \(E\)-basis from them.
Banaszczyk,
\emph{New bounds in some transference theorems in the geometry of numbers},
\S2~\cite{banaszczyk}.

In the chart the positive subspace is the graph of a matrix \(T\)
with \(\|T\|\le\rho/(1-\rho)\), by inverting its upper block.
Its majorant is between one half and twice the central squared norm.
Indeed singular value decomposition reduces its change to independent
positive-negative planes with graph slope \(t\le\|T\|\); the
eigenvalues of the new majorant there are
\((1+t)/(1-t)\) and its inverse, by using the orthogonal lines
spanned by \((1,t)\) and \((t,1)\). Unpaired directions retain
eigenvalue one. A returning \(\gamma\) carries one chart majorant
isometrically onto another; hence \(\|\gamma v\|_x\le2\|v\|_x\).
All images of the selected basis vectors have length at most
\(m^2\Delta/\sqrt\varepsilon\), allowing the extra factor \(m\ge1\).
Lattice points in a ball of radius \(T_0\) number at most
\((1+2T_0/\sqrt\varepsilon)^m\), by packing disjoint open balls of
radius \(\sqrt\varepsilon/2\). The \(r\) basis images determine
\(\gamma\). Taking the \(r\)-th power of this count yields \(A\).
\end{proof}

\begin{proposition}[From cusp and mass data to a normalized kernel]\label{bergvol:global}
Assume the following on every analytic flag component at the chosen
embedding.
\begin{enumerate}
\item The automorphic sections under consideration are bounded in
Petersson norm. For each nonzero such cusp section there is an
orthonormal frame \(x\) where its norm is at least half the component
supremum and the minimum in Lemma~\ref{bergvol:returns} is at least
\(\varepsilon\), for the same \(\Gamma\)-stable self-dual lattice.
\item There are invariant flag measures in the chart normalization
of Lemma~\ref{bergvol:local} and measurable arithmetic fundamental
sets \(F_j\). With \(V=\sum_j\mu_j(F_j)\), assume \(0<V\le V_*\).
The pushforward of the adelic probability measure on arithmetic
and frame-invariant functions is \(\sum_j d\mu_j|_{F_j}/V\).
\end{enumerate}
Then for each cusp section under consideration the global squared supremum
is bounded by \(\beta_*\) times its squared Petersson norm for that probability. On any nonempty
algebraic component \(Z\) which corresponds to a union of these
analytic components of positive mass, the same bound holds with
the conditional probability on \(Z\). In particular the scalar
evaluation Bergman kernel on \(Z\) is at most \(\beta_*\).
\end{proposition}
\begin{proof}
At the chosen half-maximum frame, project the flag ball onto the
arithmetic quotient. Each fiber in the ball has at most \(A\)
points, since mapping one of them to the others requires elements
which return the ball. Equivalently, cut a measurable fundamental
domain into its arithmetic translates and integrate its invariant
integrand there; the number of copies of a point in the ball is
bounded by the return count, including stabilizers if needed.
Lemma~\ref{bergvol:local} and the half-maximum then give
\(\sup_{x\in F_j}|P(x)|_{\rm Pet}^2\le C_{\rm loc}
\int_{F_j}|P|_{\rm Pet}^2\,d\mu_j\).
Adding the nonnegative integrals and inserting \(V\) proves the
global result. For \(Z\) use the identical calculation summing only
over its analytic components and multiplying the conditional
integral by \(V_Z=\sum_{j\subset Z}\mu_j(F_j)\le V_*\).
Taking the supremum of the evaluation ratio over sections gives
the Bergman bound. Its integral against the conditional probability
bounds the rank as well: in any finite orthonormal system the sum
of the squared pointwise norms equals the squared norm of scalar
evaluation on its span, and each summand integrates to one.
\end{proof}

For clarity, the half-maximum hypothesis admits a quantitative test
which isolates the cusp input. Put \(H=2M\sqrt d\). If at a
half-maximum frame a lattice vector satisfies \(\|v\|_x^2<
\varepsilon\le1\), every conjugate of the algebraic integer
\(h(v,v)\) has absolute value less than one, since it is bounded
by the corresponding majorant. Thus \(h(v,v)=0\) by its norm.
For each such isotropic vector use the flow parameter scaled by
\(k_0=2|\tau(\xi_0)|\), independent of \(\tau\).
In this parameter the input needed is a transvection
period \(0<H'\le H\), acting with trivial scalar multiplier, and its
holomorphic one-variable cusp restriction \(G(\zeta)\). In the
normalization with scalar metric factor one at the center it
suffices that \(G(0)=P(x)\), that \(G\) be \(H'\)-periodic and
holomorphic on a neighborhood of \(\Im\zeta\ge-t\) for
\(t=1/(2\varepsilon)\), and that it tend uniformly to zero as
\(\Im\zeta\to+\infty\). The flag metric on the flow gives
\[
       |G(\zeta)|\le 2\cdot 2^{w/2}|G(0)|
                     \qquad(\Im\zeta=-t).
\]
To verify the metric assertion, write \(u\) for the coordinates of
\(\tau v\) at the orthonormal central frame, with positive columns
first. Since \(\tau\kappa_0=-i k_0\) on the chosen CM type, the
flow in the scaled parameter is \(U=I-i\zeta uu^*J\).
Isotropy gives \(u^*Ju=0\) and
\(U^*JU=J+2\Im\zeta\,Juu^*J\).
For \(S_i=\sum_{a=1}^i J_{aa}|u_a|^2\), the rank-one determinant
formula on each leading block now gives
\(D_{\tau i}=|1+2\Im\zeta\, S_i|\).
The \(S_i\) increase from zero to \(\|u\|^2/2\), then decrease to
zero, since the total positive and negative sums coincide.
Here \(\|u\|^2\) is the local majorant, less than
\(\varepsilon\) by hypothesis. All signed determinants retain their
central signs for \(\Im\zeta>-1/\varepsilon\). Thus the flow stays
inside the flag domain in this entire half-plane and the restriction
of a holomorphic section is holomorphic there.
At \(\Im\zeta=-1/(2\varepsilon)\), the \(D_{\tau i}\) decrease from
one to a number greater than \(1/2\), then increase to one.
The total variation of \(\log D_{\tau i}\) is at most \(2\log2\).
Centering the weights at their midrange in the squared metric formula
of Lemma~\ref{bergvol:local} bounds its absolute logarithm by
\(\sum_\tau\operatorname{spread}(\lambda_\tau)\log2\le w\log2\).
The boundary raw coefficient is therefore bounded by \(2^{w/2}\)
times its Petersson norm, which itself is at most \(2|G(0)|\).

For \(y=\Im\zeta\ge0\), the same variation argument gives a raw
coefficient of size at most \(2|G(0)|(1+\varepsilon y)^{w/2}\).
Thus for \(q=\exp(2\pi i\zeta/H')\) the growth is at most a polynomial
in \(\log(1/|q|)\). By Cauchy's estimate on circles shrinking to
zero, every negative Laurent coefficient vanishes. It suffices to
show that the constant Laurent coefficient vanishes, equivalently
that the average of this restriction along one real period is zero
at any fixed height.

For the analytic cusp space of this construction the constant is
indeed zero. Its definition imposes bounded Petersson norm and,
for every nonzero isotropic \(\xi\in\Lambda\) and every frame,
zero average of the coefficient on the following \(f\)-dimensional
real torus. Put \(F=E^+\); its period lattice in the raw flow
parameters is \(L=m e_\chi\,\iota(O_F)\), where \(\iota\) is the
real embedding vector, \(m\) is the integer transvection period of
the flag datum and \(e_\chi\) kills its finite order multiplier.
This is a full lattice by the nonzero discriminant of \(F\).
The real period in direction \(\iota(a)\) follows from the
rational transvection for \(a\in O_F\) and the multiplier is
trivial after its \(e_\chi\)-th power. Both boundedness and
zero average persist under complex linear combinations, so they
hold on the full span used for the analytic Bergman kernel.
\src{src/ReciprocityTheory/ReductiveGroups/ArithmeticGeometry/FundamentalDomains/ModerateCoefficientPositivity.lean}{58}{Definition of the analytic cusp space and its zero Fourier coefficient}.

Choose an integer \(c_1>0\) with \(\xi=c_1v\in\Lambda\), possible
since \(\Lambda\) is a full lattice in the rational vector space
underlying \(E^r\). Apply the flow in independent raw
parameters \(z_\tau\) to \(\xi\), and call the resulting holomorphic
coefficient \(A(z)\). For \(y_\tau=\Im z_\tau\ge0\), the same
principal minors are \(1+2 k_0 y_\tau S_i\), now with \(S_i\)
computed for \(\tau\xi\). Their variation and bounded Petersson norm
give
\[
 |A(t+iy)|\le C_P\prod_\tau
 (1+k_0 y_\tau\|\tau\xi\|_x^2)^{\operatorname{spread}(\lambda_\tau)/2}
 \qquad(t\in\mathbb R^f).
\]
Here the norm is the single-place majorant. In particular growth is
at most polynomial when one imaginary coordinate tends to infinity
and the others stay zero. For \(\eta\in L^*\) let \(a_\eta(y)\)
be the Fourier coefficient of \(t\mapsto A(t+iy)\) using
\(\exp(-2\pi i\eta\cdot t)\) and probability Haar measure on the
real torus. Holomorphy and integration by parts on this torus give
\(\partial_{y_j}a_\eta=-2\pi\eta_j a_\eta\), hence
\(a_\eta(y)=e^{-2\pi\eta\cdot y}a_\eta(0)\).
The polynomial bound implies \(a_\eta(0)=0\) if any \(\eta_j<0\).
The defining cusp average gives \(a_0(0)=0\).

The smooth periodic function at \(y=0\) has absolutely summable
Fourier coefficients: repeated integration by parts with
\(1-\sum_j\partial_{t_j}^2\) bounds them by
\(C_N(1+|\eta|^2)^{-N}\) for every \(N\).
At each \(y\ge0\) the resulting series represents \(A(t+iy)\)
by Fourier uniqueness. On the diagonal ray its nonzero modes
therefore tend uniformly to zero as the height tends to infinity,
by dominated convergence since \(\sum_j\eta_j>0\) on their support.
Finally \(G(\zeta)=A((\zeta/(k_0c_1^2))_\tau)\), since the
rank-one operator of \(c_1v\) is \(c_1^2\) times that of \(v\).
This proves the required uniform vanishing and the claimed
one-variable constant-term identity.

This exponential identifies the half-plane above
the boundary modulo
periods with a punctured disk up to radius \(e^{2\pi t/H'}\).
The induced function extends with value zero at \(q=0\). Apply
the maximum principle to its quotient by \(q\) and evaluate at
\(q=1\). This would give
\[
  1\le 2^{1+w/2}\exp(-2\pi t/H')
       \le 2^{1+w/2-(fw+2)}<1 .
\]
Thus the cusp test implies the required minimum. Its application
on the actual deep covers needs the same period bound on every
component for the lattice used in the return count.

The period has a direct lattice calculation once the integral model
is identified. In the host normalization put
\(c=\mathrm N(\mathfrak c_K)>0\); the purely imaginary \(\xi_0\)
satisfies \(|\tau(\xi_0)|=1/(c\sqrt d)\), and the transvection
coefficient \(\kappa_0=\pm2\xi_0\) satisfies
\((cd\kappa_0)^2=-4d\). Thus \(cd\kappa_0\) is integral. Suppose
the principal level is imposed on the lattice
\(\Lambda=\mathfrak a\Lambda'\), as in the actual comparison, with
\(\Lambda'\) Hermitian self-dual. For \(v\in\Lambda'\) isotropic
and \(y\in\Lambda\), \(h(y,v)v\in\Lambda\), using the convention
that \(h\) is linear in its first argument. The transvection
\[
        y\longmapsto y+Mcd\kappa_0\,h(y,v)v
\]
therefore preserves \(\Lambda\) and is the identity modulo \(M\).
It is unitary by \(\bar\kappa_0=-\kappa_0\) and \(h(v,v)=0\);
its inverse is obtained by changing the sign. Its raw period is
\(Mcd\), and scaling by \(k_0=2/(c\sqrt d)\) gives precisely
\(H=2M\sqrt d\). The same rational transformation realizes the
period simultaneously at every \(\tau\), since the imaginary
coefficient has the same negative imaginary part \(-k_0\) there.
For the even model weights the required arithmetic
character on these transvections must be trivial. The identification
of these lattices with the actual congruence groups on every deep
component, and the identification of the sections under test with
the stated analytic cusp space, remain necessary for this application.
\src{src/Development/AbelianVarietyTheory/Proofs/CubicalLineBundles/IsIntegralLevelMulKappa.lean}{35}{Integral transvection in the principal congruence subgroup}.

In the total mass comparison, the more precise intermediate upper
bound has the form
\[
   V\le 16^{fr^2} M^{2fr^2}\operatorname{vol}(U(r))^f
                            \operatorname{mass}_I .
\]
Here the compact group volume uses the real normalization in the
adelic tube comparison and \(\operatorname{mass}_I\) is the
definite standard lattice mass for the host. A bound
\(\operatorname{vol}(U(r))\le13^{r^2}\) and
\(\operatorname{mass}_I\le[3(2f)^2r^4R_E^2]^{6fr^2}\)
give \(V\le V_*\) since \(16\cdot13\le15^2\).
The principal-level index factor \(M^{2fr^2}\) itself follows by
reducing integral lattice automorphisms modulo \(M\). Their image
lies in the endomorphism ring of \(\Lambda/M\Lambda\), of order
\(M^{2fr^2}\) by local freeness of rank \(r\) over \(O_E\).
It remains to compare the adelic component tubes, with their disjointness
and volumes, and to identify the pushforward probability in the second
hypothesis. We calculate the compact factor and the definite mass below.
\src{src/ReciprocityTheory/ReductiveGroups/ArithmeticGeometry/FundamentalDomains/MaxBoundPf.lean}{12}{Isotropic flow and periodic cusp restriction};
\src{src/Development/WeilBound/RepresentationTheoryGL/TubeMeasures/DisjointTubes/MeasureSumBounds/SumMuFLe.lean}{110}{Principal level tube comparison for all components};
\src{src/Development/AutomorphicArithmetic/ThetaCorrespondence/FockModels/TamagawaPeterssonLambdaSqNonneg.lean}{371}{Restriction to normalized algebraic components}.

\paragraph{The compact real factor and the definite lattice mass.}
Put \(n=[E:\mathbb Q]=2f\), and write \(O=O_E\). The two numerical
factors in the preceding comparison have the following bounds.
\begin{proposition}\label{bergvol:compactmass}
For the real density specified below and the genus of the standard
positive definite Hermitian \(O\)-lattice of rank \(r\ge1\),
\[
 c_r=\operatorname{vol}(U(r))
     =\prod_{j=1}^r\frac{2\pi^j}{(j-1)!}\le13^{r^2},
 \qquad
 \operatorname{mass}_I
     \le [3n^2r^2R_E^2]^{2nr^2}
     \le[3n^2r^4R_E^2]^{3nr^2}.
\]
\end{proposition}
\begin{proof}
In the skew-Hermitian Lie algebra use the density \(dx\) given by
the imaginary diagonal entries and the real and imaginary parts of
the entries strictly below the diagonal. To compare it with the real
measure of the tube calculation, let \(\lambda\) be the additive
measure for which \((X,Y)\mapsto X+iY\) carries
\(\lambda\otimes\lambda\) to entrywise Lebesgue measure on
\(M_r(\mathbb C)\). Each diagonal block has absolute real determinant
one. Each off-diagonal pair gives the two complex entries
\(z+iw,-\bar z-i\bar w\), a block of absolute real determinant four.
Hence \(\lambda=2^{\binom r2}dx\). The specified real Haar density
at the identity, in exponential coordinates, is
\(2^{-\binom r2}\lambda=dx\).
\src{src/Formalism/LocalRingCardinality/DataCD.lean}{42}{Complexification measure on the real Lie algebra};
\src{src/Formalism/LocalToGlobal/BesselHermitianTheta/MinfDensity.lean}{13}{Real density relative to complexification}.

The map \(U(j)\to S^{2j-1}\), \(g\mapsto ge_1\), has fibre
\(U(j-1)\). At the identity its differential selects the first
column, whose independent coordinates give Euclidean tangent
measure on the sphere at \(e_1\). The remaining coordinates give
exactly the chosen density on \(U(j-1)\). By invariance the same
quotient measure is Euclidean surface measure on the whole sphere.
Starting with \(c_0=1\) gives the displayed product. For \(j\le r\)
each factor is at most \(2\cdot4^r\le8^r\).
In particular compact probability Haar has density
\(2^{-\binom r2}/c_r\) relative to \(\lambda\). For a real form
\(U(p,q)\), \(p+q=r\), the compact-dual convention prescribes this
same scalar relative to its own complexification measure. Thus
multiplication by \(c_r\) converts that convention to the specified
real density at each place.
\src{src/Formalism/LocalRingCardinality/DataCD.lean}{57}{Matching of the real density across signatures};
\src{src/ReciprocityTheory/ReductiveGroups/DiscreteSeries/FormalDegree/VolUFactProdMulLe.lean}{14}{Compact real volume factor}.

For the second assertion, every lattice \(L\) in the standard genus
is Hermitian self-dual. Indeed this holds locally for \(O^r\), and
localization commutes with the dual of a projective lattice. One may
use the conjugation-stable semilocal rings at pairs
\(\{\mathfrak p,\bar{\mathfrak p}\}\) here, including at split primes.
We bound the number of isometry classes of all totally positive
definite self-dual Hermitian \(O\)-lattices of rank \(r\).
Give the underlying real space the scalar product
\(\tfrac12\operatorname{Tr}_{E/\mathbb Q}h(x,y)\). Its squared norm
is \(\sum_\tau h_\tau(x_\tau,x_\tau)\), with one summand per conjugate
pair. For the absolute different \(\mathfrak D\) the Euclidean
dual lattice is \(L^*=2\mathfrak D^{-1}L\). Indeed testing the
pairing on all \(O\)-multiples of each lattice vector gives precisely
the codifferent condition, and \(L\) is Hermitian self-dual.
For \(0\ne y\in\mathfrak D^{-1}L\) we have
\(h(y,y)\in\mathfrak D^{-1}\bar{\mathfrak D}^{-1}\), whence
\(\lvert N_{E/\mathbb Q}h(y,y)\rvert\ge\Delta^{-2}\).
The arithmetic-geometric mean inequality gives
\(\|y\|^2\ge f\Delta^{-2/n}=fR_E^{-2}\).
Thus \(\lambda_1(L^*)\ge 2\sqrt f/R_E\).
By Banaszczyk's Euclidean transference theorem in dimension \(nr\),
\(\lambda_{nr}(L)\lambda_1(L^*)\le nr\)
\cite[Theorem~2.1]{banaszczyk}. Among real-linearly independent
vectors realizing this last minimum we may choose an \(E\)-basis
\(x_1,\ldots,x_r\) with \(\|x_i\|\le T=nrR_E\).

Write \(G=(h(x_i,x_j))\) and \(L_0=\sum O x_i\). The entries of
\(G\) are integral and have modulus at most \(T^2\) at every
embedding. There are at most \((1+2T^2)^n\) choices per entry.
Indeed the points of \(O\) in its Minkowski space have mutual sup distance
at least one by the norm inequality, so open polydiscs of radius
\(1/2\) around these points are disjoint. For a fixed \(G\) the marked
rational Hermitian space and \(L_0\) are fixed up to isometry.
The inclusions \(L_0\subseteq L\subseteq L_0^\#\), where \(\#\)
denotes the Hermitian dual, locate \(L/L_0\) in a finite module
\(A=L_0^\#/L_0\) of order \(\lvert N_{E/\mathbb Q}\det G\rvert
\le T^{2nr}\). The last inequality follows by Gram's determinant
inequality at every embedding.

Each \(O\)-submodule of \(A\) has an ordered generating tuple of
length \(r\). Locally at each prime in the support its inverse image
in \((L_0^\#)_{\mathfrak p}\) is free of rank \(r\) over a discrete
valuation ring, so its image has \(r\) generators. Combine these
tuples coordinate by coordinate under the primary decomposition of
the finite module. Consequently there are at most \(|A|^r\)
intermediate lattices for the fixed Gram matrix. All isometry
classes together number at most
\[
 (1+2T^2)^{nr^2}T^{2nr^2}\le (3T^2)^{2nr^2}.
\]
The definite standard mass is the sum of
\(1/|\operatorname{Aut}_O(L,h)|\) over the classes in its genus.
Each group is finite, since it is a discrete lattice stabilizer
inside a compact real group. The class count therefore proves
the required bound, including for nonfree \(L\).
\src{src/Formalism/SharpRestriction/GlobalSchurBounds/GeneralLinearActions/GapAtLeast.lean}{20}{Definition of the definite standard mass};
\src{src/Formalism/HermitianLattices/StatementMass.lean}{11}{Reciprocal stabilizer weights}.
\end{proof}

These estimates give \(V\le V_*\) once the displayed adelic tube
inequality holds. Applying that inequality to a different signature
still requires the comparison of the adelic covolumes with their
finite lattice stabilizers and the stated real densities.


\begin{corollary}[Size of the component bound]\label{bergvol:size}
Suppose \(L\ge1,\ \log(fr)\le L,\ \log M\le30L,\ \log R_E\le3L,\)
\(d^f\le\Delta\), and \(\log w\le25L\). Then
\[
  1\le\beta_*,\qquad \log\beta_*\le 400 f^2 r^2 L
          \le4000(frL)^3 .
\]
\end{corollary}
\begin{proof}
We have \(\log d\le6L\), \(\log\Delta\le6fL\),
\(\rho^{-1}\le e^{35L}\), and \(\varepsilon^{-1}\le
3fwM\sqrt d\le e^{61L}\). The radius estimate uses
\(\log(w+f)\le\log2+25L\); the bound on \(\varepsilon^{-1}\)
uses \(fw+2\le3fw\) and \(2\log2\le\pi\).
Since \(1\le b\le fr^2/2\) and \(\log b\le2L\), the logarithm
of \(8b!/(\pi\rho^2)^b\) is at most \(3L+36fr^2L\).
The logarithm of the base in \(A\) is at most
\(\log16+2\log(fr)+\log\Delta+\log(\varepsilon^{-1})
\le(66+6f)L\).
In \(V_*\), the first factor contributes at most \(66fr^2L\);
the logarithm of the base of its second factor is bounded by
\(\log3+2\log(2f)+4\log r+2\log R_E\le16L\).
Adding gives \(\log\beta_*\le400 f^2r^2L\). Each factor in
\(\beta_*\) is at least one, in particular \(\pi\rho^2<1\).
The final inequality uses \(f,r,L\ge1\).
\end{proof}
For the cover weight take \(w=2[E:\mathbb Q]\sum_{\tau,i}
|\lambda_{\tau i}|+1\). If the weight sum is zero then \(w=1\).
Otherwise it is an integer at least one; if its logarithm is at
most \(20L\) and \(\log[E:\mathbb Q]\le3L\), then
\(\log w\le\log3+23L\le25L\). The stage's field size
\(\log[E:\mathbb Q]+\log R_E\le3L\) implies both field estimates.
\src{src/Development/AutomorphicArithmetic/ThetaCorrespondence/FockModels/TamagawaPeterssonLambdaSqNonneg.lean}{458}{Cubic cost with the host discriminant bound};
\src{src/Development/AutomorphicArithmetic/ThetaCorrespondence/ThetaPairLifts/HodgeBundleWeights/IntegralCuspSlopePowerBound.lean}{137}{Use of the same bound for the component kernel}.

\paragraph{Short harmonic polynomials on every column lattice.}
Let \(O=O_E\), \(e=[E:\mathbb Q]=2f\), \(R_E=|\Delta_E|^{1/e}\),
and let \(F\) denote the coefficient field of the chosen deep
cover (in particular it contains the labelled images of \(E\)).
Write \(\Phi\) for the chosen CM type. In the following calculation
\(\Lambda\subset E^n\) is the \emph{column} lattice, self-dual for
\(v^*u\). It may vary through the entire allowed genus. The period
lattices of the abelian varieties have rank \(r=p+q\) over \(O\).

\begin{proposition}[A common short family for all embeddings]
\label{bfam:harm}
Assume \(n\ge p+q\). At each \(\tau\in\Phi\) prescribe two
partitions of lengths at most \(p,q\), respectively. Let \(d_P\)
be the sum of their sizes. Consider row highest weight polynomials
in the two matrices of \(p,q\) rows and \(n\) columns at each
\(\tau\), annihilated by all mixed contractions
\(\sum_b\partial_{X_{ib}}\partial_{Y_{jb}}\) at each type.
This space over \(F\) is spanned by a finite set of polynomials
in the integral algebra generated by the linear forms in the rows
with coefficients \(\tau(u)\) in the first block and
\(\tau(\bar v)\) in the second, for \(u,v\in\Lambda\).
Put \(D=(2p+1)d_P,\ h=en,\ T=\max(1,hD\sqrt e\,nR_E)\). Every
polynomial \(P\) in the set satisfies, at every embedding of \(F\),
\[
       \|P^\sigma\|_1\le (r\sqrt n\,T^{2p+1})^{d_P},
 \qquad \#\{P\}\le (D+1)^{hr}.
\]
The norm is the sum of absolute coefficients in the original
column coordinates. For the stage parameters \(L\ge1\),
\(\log(fr)\le L,\ \log R_E\le3L,\ \log S\le20L\), where
\(S=d_P+frn/2\), this gives
\(\log^+\|P^\sigma\|_1\ll rL d_P\), with an absolute constant.
\end{proposition}
\begin{proof}
Use the real Euclidean structure with norm squared
\(\operatorname{Tr}_{E/\mathbb Q}(u^*u)\). Its dual lattice is
\(\mathfrak D_E^{-1}\Lambda\) by self-duality and the trace
description of the codifferent. For \(0\ne v\) in that dual,
\(v^*v\in\mathfrak D_E^{-2}\) since conjugation preserves the
different. Its absolute field norm is at least
\(|\Delta_E|^{-2}\) by inclusion of principal fractional ideals.
All its conjugates are positive. Their arithmetic-geometric mean
inequality gives \(\|v\|_{\rm tr}\ge\sqrt e/R_E\). Euclidean
transference in dimension \(h\) now gives \(h\) independent
\(b_i\in\Lambda\) of trace norm at most \(\sqrt e\,nR_E\).
They form a real basis of the product of the \(f\) complex
column spaces. \cite[Theorem~2.1]{banaszczyk}

Independently choose \(p+q\) vectors \(u_i,w_j\) on the grid
\(\sum_{\ell=1}^h a_\ell b_\ell,\ 0\le a_\ell\le D\),
with integral \(a_\ell\). Write \(U\) for the matrix of the \(u_i\)
as columns, \(G=U^*U\), and set
\[
       v_j=\det(G)w_j-U\,\operatorname{adj}(G)U^*w_j .
\]
The coefficients used to combine the lattice vectors are integral
by their Hermitian pairings. Also \(U^*v_j=0\) by the adjugate
identity. At any embedding the input vectors have Euclidean norm
at most \(T\). For independent columns the displayed operator is
\(\det G\) times an orthogonal projection. By continuity its norm
is at most \(T^{2p}\) even at singular tuples, by Hadamard's
inequality. Hence \(\|\tau v_j\|\le T^{2p+1}\) at every embedding.

At \(\tau\), contract each of the first rows with the ordered list
\(\tau u_i\), and each of the second rows with \(\tau\bar v_j\),
using the ordinary bilinear dot product. Take the determinants of
the first \(i\) rows and columns in each resulting square matrix,
to powers the successive differences of the respective partitions
(use the reversed column order in the standard opposite block).
Multiply over \(i\) and \(\tau\). The determinants give the
specified row character. Every mixed derivative contracts a pair
of mutually orthogonal coefficient vectors and is zero. A linear
form has coefficient norm at most \(\sqrt n\,T^{2p+1}\).
A determinant of order \(i\) costs at most an additional \(i!\le
r^i\), giving the asserted bound. Its coefficients lie in the
stated integral algebra, also for a nonfree projective \(\Lambda\).

To see spanning, a complex linear functional annihilating the
grid polynomials gives a polynomial in the \(hr\) real grid
coordinates of degree at most \(D\): conjugation is real-linear
and each projected vector has degree \(2p+1\). Successive
one-variable interpolation makes it zero for all real choices.
The \(b_i\) allow independent choices at the \(f\) archimedean
places. In particular at each place use the first \(p\) and last
\(q\) columns of any unitary matrix, in the orders of the
prescribed minors. Here \(G=1\) and the projection fixes the
last columns. The obtained products form the orbit of the
product of the standard opposite flag minors under \(U(n)^f\).
Kashiwara and Vergne identify this product as a highest weight
vector and its column representation as the entire harmonic
row highest weight space for the given partitions, with multiplicity
one \cite[unitary case, \S6, Proposition~6.1 and Theorem~6.3]{kashiwara-vergne}.
Taking products gives the assertion at the \(f\) places. Descent
of matrix rank gives spanning over \(F\), since the equations
and grid vectors are defined there. Finally \(n\le2S,\ d_P\le S\)
give \(\log T\ll L\) if \(d_P>0\). For \(d_P=0\) take just \(1\).
\end{proof}
The two sets of rows here are labelled by
\(\tau,\bar\tau\), with the dual rows reversed when passing to
the ordered flag weight. Thus the integral algebra is exactly
the one generated by lattice directions in the harmonic jet of
Theorem~\ref{intjet:semi}. Relabelling the letters by an embedding
permutes the column coordinates and preserves the coefficient
estimate. For each fixed auxiliary integer, lattice and level
tuple the theta sum is linear in \(P\): the jet is the evaluation
of the constant-coefficient operator \(P(\partial)\) and the
translations and finite sums are independent of \(P\). Consequently
the short family has exactly the same image span as all such
harmonic polynomials for these data. This costs no determinant
or denominator from a change of basis.
\src{src/Development/AbelianVarietyTheory/Proofs/CuspFormEvaluations/ExistsShortHarmonicFamilyUnimodular.lean}{96}{The integral algebra and its simultaneous coefficient bounds};
\src{src/Development/AbelianVarietyTheory/Proofs/DivisorProofs/PolynomialCoefficientEmbeddingNormLeOfL1Bound.lean}{399}{Use at the actual deep coefficient field}.

\paragraph{From a half maximum to the polarized period metric.}
The transfer from the lattice in Proposition~\ref{bergvol:global}
to the periods in Proposition~\ref{intjet:arch} must account for
the host fractional ideal. Write the actual period lattice as
\(\Pi=\mathfrak a L'\), where \(L'\subset E^r\) is self-dual
for the row Hermitian form \(h_0\). Write \(\xi\) for its
imaginary polarization coefficient and \(x=|\tau\xi|\),
independent of \(\tau\). The row polarization is
\(\operatorname{Tr}_{E/\mathbb Q}(-\xi h_0)\), with
\((\xi)\mathfrak a\bar{\mathfrak a}\mathfrak D_E=O\).

\begin{lemma}[Comparison of the two minima]\label{bfam:minimum}
Let \({\rm maj}_\tau\) denote the positive majorant of \(h_0\)
at an orthonormal flag. Assume that
\[
        \sum_{\tau\in\Phi}{\rm maj}_\tau(u)\ge\varepsilon
                 \qquad(0\ne u\in L').
\]
Then the shortest nonzero polarized period
has squared length at least \(2\varepsilon/R_E^2\).
\end{lemma}
\begin{proof}
Taking ideal norms gives \(N(\mathfrak a)^{-1/f}=xR_E\).
In the product of \(f\) complex planes, \(\mathfrak a^{-1}\)
has covolume \(2^{-f}\sqrt{|\Delta_E|}/N(\mathfrak a)\).
The product of discs of squared radius \(xR_E^2\) has volume
strictly greater than \(2^{2f}\) times this covolume, as
\(\pi>2\). Minkowski gives \(0\ne b\in\mathfrak a^{-1}\)
with \(|\tau b|^2\le xR_E^2\) simultaneously. For \(v\in\Pi\)
nonzero, apply the hypothesis to \(bv\in L'\).
The polarized squared norm of \(v\) is
\(2x\sum_\tau{\rm maj}_\tau(v)\): in the orthogonal positive
and negative coordinates one conjugates the positive block in the
period convention of Lemma~\ref{dg:gaussian}, and the two terms
in the trace give \(2x\) on both
blocks. This also follows from the imaginary part of the Hermitian
form, which is the displayed trace polarization. The estimate
follows on summing
\({\rm maj}_\tau(bv)=|\tau b|^2{\rm maj}_\tau(v)\).
\end{proof}
\src{src/Development/AbelianVarietyTheory/Proofs/DivisorProofs/FirstMinAtHalfMax.lean}{82}{The inverse ideal and the distinction between the two lattices}.

\begin{corollary}[Simultaneous estimate in the normalized cover metric]
\label{bfam:metric}
Use the framed uniformization of the deep cover at each good embedding
and the self-dual row lattice model above. Assume the analytic cusp
reading and the half-maximum condition of the first hypothesis of
Proposition~\ref{bergvol:global} for the theta sections in question.
Use the parameters of Corollary~\ref{bergvol:size} and
Proposition~\ref{intjet:arch}, with \(S=d_P+frn/2\) and
\(\log S\le20L\). The sections read from the polynomials of
Proposition~\ref{bfam:harm}, multiplied by the reflex-ideal
scalar \(s\) of Theorem~\ref{intjet:trace}, have logarithmic
pointwise supremum at most \(C(frL)^2S\) in the normalized cover
metric, for an absolute \(C\). The same bound holds for the full
cover Petersson norm and every positive-mass component Petersson
norm. Multiplication by the common integer of
Theorem~\ref{intjet:semi} adds only its logarithm.
\end{corollary}
\begin{proof}
At an orthonormal flag all absolute leading Gram determinants are
one. The framed uniformization takes the Hodge directions to the
coordinate vectors and takes the theta generator to one of squared
mean one in the rigidified Gaussian metric on the period torus.
Explicitly its normalization is
\[
  \frac{1}{\operatorname{vol}\mathcal F}\int_{\mathcal F}
  |\theta(z)|^2 e^{-\pi H(z,z)}\,dz=1 .
\]
This is Haar probability on the abelian fibre and therefore
\(\|s_0\|_2=1\) in Proposition~\ref{intjet:arch}. It is part of
the framed uniformization, and changing the scalar of the
generator alone would change this equation. The polarized form
in these coordinates is \(H(z,z)=2x\sum_j|z_j|^2\), so that
\(U=\sqrt{2x}\). The cover metric multiplies the absolute raw
coefficient by \((2\pi x)^{-d_P/2}\) at this flag. In particular
the factor \((\sqrt2 U)^{d_P}\) there becomes
\((2/\pi)^{d_P/2}\), independently of the size of \(x\).
By Lemma~\ref{bfam:minimum} the minimum in the radius in
Proposition~\ref{intjet:arch} satisfies
\(\log^+\lambda^{-1}\le\log R_E+\tfrac12\log(1/\varepsilon)\ll L\).
That proposition, the short coefficient bound above, and
\(\log|\Delta_F|/[F:\mathbb Q]\ll(frL)^2\) give the asserted
estimate at the half maximum. Doubling bounds every frame.

The full cover norm integrates the square of this pointwise
quantity against the probability on the adelic unitary quotient
\(U(W_j)(\mathbb Q)\backslash U(W_j)(\mathbb A)\), writing
the group by restriction of scalars. The component norm uses the
same integrand through the frame map and the component measure
(conditional probability on positive mass). Any such measure of
mass at most one gives the same upper bound by integration.
Embeddings of the other CM orientation are read at their conjugate
good embeddings; the coefficient estimate was simultaneous for both.
This direction of the estimate uses no comparison of two total
adelic volumes.
\end{proof}
\src{src/Formalism/LocalToGlobal/ProjectiveLattices/PetersonNormSubspace.lean}{111}{The squared mean in the framed uniformization};
\src{src/Development/AbelianVarietyTheory/Proofs/DivisorProofs/EvenDomainWeightBound.lean}{68}{The full and component cover norms}.

\paragraph{Descent of the raw rule in characteristic zero.}
One can check that the finite-jet formula gives algebraic values
before identifying it with any classical modular form. Work on
the framed abelian locus over \(F\), with its actual level
structure, eigenframe, theta generator and polarized biextension.

\begin{lemma}\label{bfam:rule}
Assume the column lattice is standard locally at every prime
above the deep level \(M\). The odd-order finite theta sum of Lemma~\ref{newtheta:index}
with the prescribed half-primary division points and harmonic \(P\)
defines a regular rule on these framed points, functorial in
the characteristic-zero base. The complex reading of natural
weight rules on reduced finite-type \(F\)-tests is injective after
extension to \(\mathbb C\), whenever the framed uniformization
surjects onto complex framed points up to the weight action.
\end{lemma}
\begin{proof}
Work locally on the base in logarithm coordinates of the smooth
commutative formal group; they are defined over a \(\mathbb Q\)-algebra
by integrating its invariant differentials, with identity linear term.
On the polarized biextension choose a frame restricting to the
rigidified frames on both axes. Such a frame exists formally
locally: an arbitrary frame at the origin extends over the
completion and its axis restrictions can be divided out.
Use the biextension identification of its pullback by
\(([2],[2])\) with its fourth power. The logarithm of the
discrepancy starts in total degree three. Indeed each first-variable
addition discrepancy vanishes on \(X_1=0,\ X_2=0,\ Y=0\);
the second-variable discrepancy has the analogous property.
Compute dilation by doubling the two variables successively.
Changing the log frame by a degree-\(j\) homogeneous correction
changes the degree-\(j\) dilation discrepancy by \(2^j-4\) times
that correction. Thus all discrepancies of degree at least three
can be removed. Each addition discrepancy of the resulting frame
is now homogeneous of total degree two, by applying simultaneous
dilation to the addition diagram. Its three axis vanishings force
it to be zero.

The Rosati adjunction for an integer \(a\in O\) in this biadditive
frame can differ by an exponential of a bilinear form. In
types \((\tau,\nu)\), let its log coefficient be \(c_a\).
Functoriality for addition and composition of endomorphisms gives
\[
 c_{a+b}=c_a+c_b,\qquad
 c_{ab}=\tau(b)c_a+\nu(\bar a)c_b .
\]
Indeed for composition apply adjunction first to \(a\), with
first input \(bX\), then to \(b\), with second input \(\bar aY\).
If \(\tau=\bar\nu\), this is a derivation on the finite separable
field \(E/\mathbb Q\) with the given coefficient action, hence
zero. If \(\tau\ne\bar\nu\), interchanging \(a,b\) shows
\(c_a\) is a constant multiple of \(\tau(a)-\nu(\bar a)\).
Changing the bilinear log coefficient of the frame removes it.
The remaining ambiguity is Rosati-compatible bilinear. On the
product of \(n\) identical factors it gives the sums of mixed
quadratics killed by harmonic contraction in
Lemma~\ref{intjet:der}. Thus the jet is regular in the base and
independent of these choices. Its evaluation uses only a finite
truncation of these formal constructions and commutes with base
change. The weight action follows by changing the eigenframe in
the polynomial and raising the theta generator's scalar to the
\(n\)-th power.

After an étale base change one may enumerate \(Z_\Lambda\) and
divide each level point by \(N\) in its prescribed half-primary
part. In characteristic zero the torsion of an abelian scheme
is finite étale, and multiplication by \(N\) on the appropriate
primary torsion over geometric points is surjective onto the
prescribed level part. The symmetric normalized lifts on the
odd torsion are algebraic: in the theta group the fibre over
a point is a multiplicative torsor; the two normalizations of
Lemma~\ref{newtheta:index} select a unique lift étale locally.
Indeed first make its \(N^2\)-th power one by taking a root,
then adjust by an \(N^2\)-th root of unity to make inversion
agree with its inverse. The discrepancy is such a root and
squaring is an automorphism on these roots. Uniqueness gives
descent of the lift.

If two choices of division point differ by \(u\), then \(u\)
is \(N\)-torsion supported on the prescribed half above \(M\).
Here the column lattice is standard, so \(u\in K_0\) in the
notation of Lemma~\ref{newtheta:index}, by its period description
at those primes. The Weil pairing between two points in the
same half is trivial: its adjoint involution interchanges the
two relatively prime half ideals, hence their idempotents pair
only across the halves. The translations by \(t\) and \(u\) thus
commute and their normalized product is the lift at \(t+u\).
The vector to which the last translation is applied is fixed by
\(K_0\), by Lemma~\ref{newtheta:index}. This proves independence
of division and of enumeration. The denominator is invertible
in the base, and the level character takes values in \(F\).
The resulting regular scalars agree on the étale double overlap
and descend. On a disconnected base the assertion is checked
componentwise; the same identities hold scheme-theoretically by
the finite étale torsion and the uniqueness of its lifts.

Finally let \(\sum_i c_i f_i\) have vanishing complex reading
at a fixed embedding \(F\to\mathbb C\), with the \(f_i\)
independent over \(F\). On a framed point of any reduced
finite-type test algebra \(R\), it follows that
\(\sum_i c_i\otimes f_i(R)=0\) in \(\mathbb C\otimes_F R\).
Indeed its value is zero at each complex point by naturality,
and the uniformization up to the common weight factor.
This algebra is reduced since
\(F\) is perfect, and its complex points detect zero by the
Nullstellensatz. Expand all \(c_i\) in a basis of their finite
\(F\)-span in \(\mathbb C\). Each coefficient yields an \(F\)-linear
relation among the \(f_i\) on every test, hence is zero. This
proves injectivity.
\end{proof}
\src{src/Formalism/LocalToGlobal/ProjectiveLattices/ThShort.lean}{244}{The faithfully flat evaluation, half-primary division and weight};
\src{src/Formalism/LocalToGlobal/ProjectiveLattices/AlgebraMapCAt.lean}{9}{Reduced finite-type test rings};
\src{src/Development/AbelianVarietyTheory/Proofs/UnitaryExistence/ReadingLinearIndependentAndInjective.lean}{22}{Reading after extension of the coefficient field}.

\paragraph{The operator on row lattices.}
There is a separate numerical reason for the lower bound on the
theta rank. It applies before any denominator or choice of
representatives on the column genus.

\begin{lemma}[Invertibility at the stated theta rank]\label{bfam:gram}
Let \(\mathcal R\) be any set of full lattices of rational rank
\(b=2fr\), self-dual for the fixed rational symplectic polarization.
Set \(m(L,L')=[L:L\cap L']\). Suppose a complex matrix \(C\)
indexed by \(\mathcal R\) has constant diagonal \(A>0\) and
\(|C_{LL'}|\le A\,m(L,L')^{-n/2}\).
If \(n\ge4b+4\), then on bounded lattice-indexed coefficients
\[
     \|A^{-1}C-1\|_\infty\le\pi^2/6-1,\qquad
     \|C^{-1}\|_\infty\le A^{-1}(2-\pi^2/6)^{-1}.
\]
The result holds for coefficients in any complex Banach space
and on any closed subspace preserved by \(C\).
\end{lemma}
\begin{proof}
All self-dual lattices have the same symplectic covolume, hence
\([L':L\cap L']=m(L,L')\). In a rank \(b\) free abelian group
there are \((\ell^b-1)/(\ell-1)\le\ell^b\) subgroups of prime
index \(\ell\), by their kernels in the reduction modulo \(\ell\).
Any index-\(m\) subgroup occurs at the end of a chain with
successive prime indices in a prescribed ordering of the prime
factors of \(m\), by the primary decomposition of the quotient.
There are at most \(m^b\) such chains. The same count holds for
overlattices by ordinary lattice duality. Thus at most \(m^{2b}\)
choices of \(L'\) have this index, by first choosing \(L\cap L'\)
and then its overlattice. Index one gives only \(L\). The sum of
absolute off-diagonal entries in each normalized row is bounded
by \(\sum_{m\ge2}m^{2b-n/2}\le\sum_{m\ge2}m^{-2}\).
The norm-convergent Neumann series proves both conclusions,
including on the indicated subspaces.
\end{proof}

Here is the finite vacuum calculation giving the entry estimate.
For unshifted row lattices \(L,L'\), put
\(I=L\cap L'\). For any locally free rank-\(n\) column lattice
\(\Lambda\) self-dual for its Hermitian form, write
\(A_L=\Lambda\otimes_O L\) for the corresponding tensor lattice
with its trace symplectic form. Then
\[
 A_L\cap A_{L'}=\Lambda\otimes_O I,\qquad
 [A_L:A_L\cap A_{L'}]=m(L,L')^n .
\]
The first identity is flatness and the second follows on choosing
local bases at the finite primes. The two tensor lattices are
self-dual for the symplectic form since the column pairing
is perfect over \(O\). In the Heisenberg representation take
unit vacuum vectors \(v,v'\) with the prescribed characters on
these two lattices. Distinct translates of \(v'\) by representatives
of \(A_L/(A_L\cap A_{L'})\) are orthogonal: on \(A_{L'}\) their
eigencharacters differ by the symplectic commutator, whose
annihilator is precisely \(A_{L'}\). Their inner products with
\(v\) all have the same absolute value, because \(v\) is an
eigenvector on \(A_L\). Bessel's inequality gives
\(|\langle v,v'\rangle|\le m(L,L')^{-n/2}\).
This argument works with any translates of the two vacua, since
translation just changes their eigencharacters.

For the actual finite marks write \(\mathfrak m\) for the nonzero
proper odd ideal prime to \(\bar{\mathfrak m}\), and \(L_*\) for
the reference row lattice. The row set consists of self-dual \(L\)
equal to \(L_*\) at primes dividing \(\mathfrak m\bar{\mathfrak m}\).
The column lattice is standard there. Set
\(B_L=\Lambda\otimes_O(\mathfrak m\bar{\mathfrak m}^{-1}L)\).
The twist ideal times its conjugate is \(O\); hence these \(B_L\)
are again self-dual and their relative index remains \(m(L,L')^n\).
Choose a basis of \(L_*/\mathfrak m L_*\). The marked vector
\(\phi_L\) is the sum of translates of the unit vacuum on \(B_L\).
They are indexed by \(k\in J\subset{\rm GL}_r(O/\mathfrak m)\)
and weighted by the unitary level character. In fact the quotient
used to choose the translating tuples is the reduction modulo
\(\mathfrak m\) of \(\Lambda\otimes_O(2\bar{\mathfrak m}L)\).
Since \(2\bar{\mathfrak m}\) is locally a unit on \(\mathfrak m\),
this gives precisely the \(n\) marked rows modulo \(\mathfrak m L_*\).
The first \(r\) rows give the basis transformed by \(k\), and
the remaining rows vanish modulo \(\mathfrak m\). There is one
class for each such \(k\).

Use the common vacuum semicharacter on intersections. For
translates by \(y,y'\), nonzero inner product implies triviality
of the commutator of \(y-y'\) on \(B_L\cap B_{L'}\), hence
\(y-y'\in (B_L\cap B_{L'})^\#=B_L+B_{L'}\).
Here the duality identity follows by dualizing the sum of the
two self-dual lattices. At every \(\mathfrak p\mid\mathfrak m\)
this sum has rows in \(\mathfrak m L_{*,\mathfrak p}\).
The marked residues coincide, so only pairs with the same \(k\)
contribute. Their diagonal squared norm is \(|J|\); the mixed
inner product is at most \(|J|m(L,L')^{-n/2}\).
This proves precisely the hypotheses of Lemma~\ref{bfam:gram}
for the positive weighted class sum, with \(A=|J|\) times the
total weight. No lower bound on that mass is needed for invertibility.

The same calculation detects the arithmetic group of the marks.
For \(\gamma\in U(h_0)(E^+)\) replace \(B_{L'}\) by \(\gamma B_{L'}\);
the vacuum semicharacter is preserved because it is defined by the
polarization, the fixed \(O\)-scalar \(w_0\) with \(w_0+\bar w_0=1\),
and the unitary action. Nonzero overlap now gives
\(y-b=\gamma(y'-b')\) for some \(b\in B_L,\ b'\in B_{L'}\),
using the same duality identity. At \(\mathfrak p\mid\mathfrak m\)
both adjusted tuples have first \(r\) rows forming an integral
basis of \(L_{*,\mathfrak p}\), by reduction modulo the maximal
ideal. Thus \(\gamma\) preserves the localized reference lattice
and its residue maps the basis marked by \(k'\) to that marked by
\(k\). Its residue is \(k(k')^{-1}\in J\). This proves vanishing
on the graph orbits outside the subgroup specified by these
congruences.

One also obtains the finite vanishing along a nonzero isotropic
row line \(E\xi\). For completeness take a prime
\(\mathfrak p\mid\mathfrak m\), put \(a={\rm ord}_{\mathfrak p}
(\mathfrak m)>0\), and work at this prime and its distinct conjugate.
The trace symplectic pairing identifies \(L_{*,\bar{\mathfrak p}}\)
with the additive character dual lattice of \(L_{*,\mathfrak p}\)
in \(E_{\bar{\mathfrak p}}^r\). This follows by testing the
global self-duality locally on the two conjugate idempotents.
Scale \(\xi_{\bar{\mathfrak p}}\) to a primitive vector \(u\)
of \(L_{*,\bar{\mathfrak p}}\). The element
\(\varpi_{\bar{\mathfrak p}}^{-a}u\) belongs to the shifted row
lattice there but lies outside the additive dual of \(L_{*,\mathfrak p}\).
For any marked tuple \(y\) its first \(r\) rows form a basis at
\(\mathfrak p\). Some integral multiple (in the local integer
ring) of \(\varpi_{\bar{\mathfrak p}}^{-a}u\) therefore pairs
nontrivially under the additive character with one of those rows.
Put this vector in the corresponding column position of a tuple
\(\delta\in B_L\otimes\widehat{\mathbb Z}\) on \(E^n\otimes E\xi\),
and zero at the other components. Equivalently one can choose a
rational \(\delta\) giving the same character: the intersection of
a full lattice with a rational subspace spans a lattice in that
subspace, and its completion is the intersection of the completed
lattice with the completed subspace, by a basis adapted to the
saturated intersection. Approximation modulo a finite-index
sublattice then preserves the character.
On this subspace both the alternating form and the vacuum quadratic
phase vanish, since \(h_0(\xi,\xi)=0\). The Heisenberg translations
on it commute and their action on the marked translate has the
character of this pairing. For each of the finitely many labels
choose such a \(\delta\) with eigenvalue \(\eta\ne1\). The product
of the commuting operators \(T_\delta-\eta\) kills the sum of
marked translates and has nonzero augmentation
\(\prod(1-\eta)\). Hence its image in the coinvariants of the
isotropic row subgroup vanishes. This supplies the vanishing for
non-graph terms whose evaluation factors through those coinvariants.
\src{src/Development/AbelianVarietyTheory/Proofs/SiegelThetaTransformation/ThetaTestFunctionNormAndOverlapBound.lean}{133}{The localized congruence for a surviving pair of labels};
\src{src/Development/AbelianVarietyTheory/Proofs/SiegelThetaTransformation/HermitianThetaLevelDataFormsSystem.lean}{60}{The off-diagonal sum at rank \(8fr+4\)};
\src{src/Development/AbelianVarietyTheory/Proofs/SiegelThetaTransformation/ThetaTestFunctionNoCoinvariantsIsotropic.lean}{216}{Isotropic row coinvariants of the same marked tests}.

\paragraph{Which definite Siegel--Weil formula is needed.}
There is a way to use the ordinary adelic theorem without comparing
the volumes of different genera of self-dual Hermitian lattices.
Use just the genus of \(\Lambda_*=O^n\) in the positive definite
space \(V=E^n\), and require equality with the standard lattice
at a prescribed integer \(B_0\) divisible by the deep level and
the excluded rational primes. Every lattice used in this smaller
family is among the permitted columns. Its positive class mass
still gives the strict inequality in Lemma~\ref{bfam:gram}.

\begin{lemma}[Adelic realization of a lattice family]\label{bfam:genus}
Write \(H=U(V)\) over \(E^+\) and suppose \(n>4r\).
For each such column lattice let
\(\varphi_\Lambda\) be a function on \(V(E^+)^r\) (thus on
\(r\) columns in \(E^n\)). Assume the following conditions for
some positive integer \(R\).
\begin{enumerate}
\item The function is supported in \(R^{-1}\Lambda^r\) and
invariant under \(R\Lambda^r\).
\item If two lattices coincide at all primes over \(RB_0\),
their functions agree on the intersection of their indicated
supports \(R^{-1}\Lambda^r\cap R^{-1}\Lambda'^r\).
\item For a rational isometry \(\gamma\) integral and congruent
to the identity modulo \(B_0\) at its prime divisors,
\(\varphi_{\gamma\Lambda}(\gamma x)=\varphi_\Lambda(x)\).
\end{enumerate}
Then the Schwartz functions obtained by completion form the
translates of one finite adelic Schwartz function in this genus.
For a fixed \(H(E^+\otimes\mathbb R)\)-invariant archimedean
Schwartz function and the splitting in which \(H\) acts by
linear substitution in the doubled Schrödinger model, the
genus average of the theta functions with stabilizer weights,
divided by the sum of those weights, is the anisotropic
Siegel--Weil Eisenstein value for the standard lattice test.
The same is true of the genus average of the corresponding
Eisenstein values themselves.
\end{lemma}
\begin{proof}
Let \(\mathcal U\subset H(\mathbb A_f)\) be unrestricted away
from \(B_0\), and the principal congruence subgroup of the
standard lattice at its divisors. For \(g\in\mathcal U\)
consider \(\Lambda'=g\widehat\Lambda\cap E^n\). Write
\(\widehat\varphi\) for the function on the finite adeles
defined by the same finite quotient
\(R^{-1}\Lambda^r/R\Lambda^r\), extended by zero.
We claim
\(\widehat\varphi_{\Lambda'}(g x)=\widehat\varphi_\Lambda(x)\).
It suffices to check this on \(R^{-1}\widehat\Lambda^r\).
Choose a rational unitary \(\gamma\) sufficiently close to \(g\)
at all primes over \(RB_0\). Weak approximation here follows
from the Cayley transform between the skew-adjoint Lie algebra
and a dense open in \(H\); the latter is locally dense at every
finite place, since its complement is a proper algebraic closed
set in a smooth connected group. Linear weak approximation in
the Lie algebra proves the assertion. Thus \(\gamma\Lambda\)
and \(\Lambda'\) agree at those primes, \(\gamma\) meets the
congruences of the third hypothesis, and \(\gamma x,g x\)
have the same residues modulo the local period lattice there.
The lattice \(R^{-1}(\gamma\Lambda\cap\Lambda')^r\) maps onto
these residues by the Chinese remainder theorem. Choose a
rational \(z\) in it representing them. Outside \(R\) the
support lattice and period lattice coincide. Consequently the
second hypothesis compares \(\varphi_{\Lambda'}(z)\) and
\(\varphi_{\gamma\Lambda}(z)\), and the third identifies the latter
with \(\widehat\varphi_\Lambda(x)\). This proves the claim.

The same weak approximation at \(B_0\) gives
\(H(\mathbb A_f)=H(E^+)\mathcal U\).
Put \(K=\operatorname{Stab}_{\mathcal U}(\widehat\Lambda_*)\).
The compact right cosets on the finite double quotient
\(H(E^+)\backslash H(\mathbb A_f)/K\) correspond to the
lattice orbits with the indicated congruences. Representatives
can all be taken in \(\mathcal U\). Right invariance by \(K\)
follows from the claim. On each coset the compact archimedean
group integrates the invariant archimedean factor, and the
remaining mass is proportional to the reciprocal order of
the stabilizer. Indeed its chart is a quotient of
\(H(E^+\otimes\mathbb R)\times K\) by this finite group;
the numerator Haar volume is common to all cosets. Dividing by
the total mass gives the theta integral for Haar probability.
The doubled split group has Witt index \(r\); the source
dimension is \(n>4r\). In the notation of Gan--Qiu--Takeda the
anisotropic formula is at \(s=(n-r)/2>0\), with unitary
constant one for the probability theta integral
\cite[\S3.1 and Theorem~7.1]{ganQiuTakeda}. Finally each
finite transport has the same Siegel section: the \(H\)-action
commutes with the doubled group and linear substitution fixes
evaluation at zero. This proves the last assertion as well.
\end{proof}
For an archimedean polynomial insertion one takes the invariant
sum over a unitary orthonormal basis of the column representation
in Proposition~\ref{bfam:harm}, paired with its conjugate on
the second copy. This is the identity endomorphism of that
representation, so the resulting doubled Schwartz function is
invariant under the compact column group. Lemma~\ref{bfam:genus}
applies just as it does to the Gaussian. We verify the mixed
finite equivariance directly.
\src{src/Development/AbelianVarietyTheory/Definitions/SpanningFamilies/FirstCoordinateAdmissible.lean}{139}{Equivariance, periodicity and finite locality of the column family};
\src{src/Development/AbelianVarietyTheory/Proofs/SiegelThetaTransformation/AdmissibleDoubledHermitianThetaFamily.lean}{413}{Mixed vacuum comparison for the marked doubled test}.

\begin{lemma}[The finite phase for the projected marked vacua]\label{bfam:phase}
Use the finite adelic Schrödinger representation of the trace
symplectic space underlying \(V\otimes_E E^r\). For each row
lattice \(L_i\), \(i=1,2\), start with any unit vacuum \(v_i\) for
the standard column and the shifted row lattice
\(D_i=\mathfrak m\bar{\mathfrak m}^{-1}L_i\). For the column
\(\Lambda\) take the unit normalization \(u_i^\Lambda\) of the
orthogonal projection of \(v_i\) to the vacuum of
\(A_i^\Lambda=\widehat\Lambda\otimes_{\widehat O}\widehat D_i\).
The vacuum character at each lattice is the prescribed
semicharacter defined by \(w_0\). Form the marked sums
\(\phi_i^\Lambda\) above from these \(u_i^\Lambda\). For
\(h\in H(\mathbb A_f)\) preserving the marked residues (in
particular for \(h\in\mathcal U\) of the preceding proof),
\[
 W_{\phi_1^{h\Lambda},\phi_2^{h\Lambda}}(hx)
       =W_{\phi_1^\Lambda,\phi_2^\Lambda}(x),\qquad
 W_{u,v}(x)=\langle \rho(x)u,v\rangle .
\]
Here \(h\Lambda\) means the lattice with completion
\(h\widehat\Lambda\), and \(\rho\) denotes the Weyl operators.
In particular the finite doubled tests used in the ordinary
adelic theta identity are the translates of a single test for
each pair of row lattices, with their original complex coefficients.
\end{lemma}
\begin{proof}
Choose complementary rational Lagrangians in the underlying
symplectic space. Use dual Haar measures for the finite additive
character \(\psi_f\) and the pairing of the two Lagrangians. With
the convention \(\rho(a,b)u(t)=\psi_f(b(t+a/2))u(t+a)\) and the
inner product linear in its first argument, the map to the doubled
diagonal model is
\[
 W_{u,v}(a,b)=\int u(t+a/2)\overline{v(t-a/2)}\,\psi_f(bt)\,dt .
\]
Indeed one first changes variables from the two position coordinates
to their difference and midpoint, then Fourier transforms the
midpoint. The Jacobian of the first change is of absolute value
one, also at dyadic places. A change of basis for the dual
coordinate changes the corresponding dual Haar. Fourier inversion
and this change of variables give
\[
 \int W_{u,v}\overline{W_{s,t}}=
             \langle u,s\rangle\overline{\langle v,t\rangle}.
\]
The integrals here are integrals of Schwartz functions. In
particular \(W_{u,v}(0)=\langle u,v\rangle\).
In the paper's rational doubled coordinates the symbol is
\(a\,W_{u,v}(-x)\), where \(a\ne0\) is the common change of Haar
normalization. In fact translation of the two inputs by
\(y_1,y_2\) translates the position argument by \(-(y_1-y_2)\)
in those coordinates. Set \(y_1=-x,\ y_2=0\) and evaluate at
the origin. The remaining half-pairing has zero rational trace:
it is a purely imaginary scalar times the real Hermitian norm
of \(x\). Thus the value of the doubled symbol at \(x\) equals
its origin value on \(\rho(-x)u,v\). The origin pairing is the
displayed inner product by partial Fourier transformation. The
reflection commutes with every linear column isometry.

Write \(\epsilon(a)\rho(a)\) for the action fixing a vacuum on a
self-dual lattice \(A\). The formula for \(\epsilon\) depends on
the vector and \(w_0\), so its restrictions to intersections
agree. Explicitly it is the half quadratic character of the trace
pairing of \(a\) with \(w_0 a\), with the sign determined by
\(\psi_f\). On \(A\) it cancels the Weyl multiplication cocycle:
\(w_0+\bar w_0=1\) and the pairing of \(A\) with \(w_0A\) is
integral. The integral of \(\epsilon(a)\rho(a)\) against Haar
probability on \(A\) is the orthogonal projection \(P_A\) to its
fixed space. For a vacuum \(v\) on \(A_*\), \(W_{v,v}\) vanishes
outside \(A_*\): conjugation by \(b\in A_*\) multiplies the
coefficient by the commutator character and \(A_*^\#=A_*\).
Thus
\[
 \|P_Av\|^2=\langle P_Av,v\rangle
      =[A:A\cap A_*]^{-1}.
\]
The norm is positive and the unit projection has positive real
inner product with \(v\). This integral is also the normalized
finite average of the distinct translates used on the rational
lattice: the orbit on a Schwartz vector is finite and each fibre
of its orbit map has equal measure.

For vacua \(u_i\) on \(A_i\), the same commutator, tested on
\(A_1\cap A_2\), gives support in \(A_1+A_2\). At
\(x=a_1+a_2\), \(a_i\in A_i\), move \(\rho(a_1)\) to the
first vacuum and \(\rho(a_2)\) to the second. The result is
\(W_{u_1,u_2}(0)\) times a unit scalar determined by
\(\epsilon(a_i)\) and their Weyl cocycle. In particular its
absolute value is constant throughout the support. The preceding
Fourier identity on the diagonal gives
\[
 |W_{u_1,u_2}(0)|={\rm vol}(A_1+A_2)^{-1/2}.
\]
This volume is independent of the self-dual column: each
self-dual \(A_i\) has volume one and the relative index here
is \([D_2:D_1\cap D_2]^n\).

To compare the \emph{arguments} for columns \(\Lambda,\Lambda_*\),
observe that the intersection of the supports is
\[
 (\widehat\Lambda\cap\widehat\Lambda_*)\otimes_{\widehat O}
                        (\widehat D_1+\widehat D_2).
\]
This can be checked in a local basis of the row sum. Every
element of this intersection therefore has a decomposition
\(a_1+a_2\) with the two \(a_i\) in the corresponding vacuum
lattices for both columns. Its scalar relative to the value at
zero is identical for the two symbols. Their Fourier inner
product is consequently the positive volume of this intersection
times the product of the two values at zero, with the second
conjugated. On the other hand the Fourier identity identifies
this inner product with
\(\langle u_1^\Lambda,v_1\rangle
 \overline{\langle u_2^\Lambda,v_2\rangle}>0\).
The two values at zero have equal modulus, so they are equal.
This applies separately to every column and does not impose a
choice of phase on either initial \(v_i\).

An isometry \(h\) of the column carries a decomposition
\(a_1+a_2\) to \(ha_1+ha_2\). It preserves the half quadratic
character and the Weyl cocycle. Equality at zero now gives the
assertion for the unmarked vacua. For a marked summand translate
by \(y_i\in\widehat\Lambda\otimes 2\bar{\mathfrak m}\widehat L_i\).
Its doubled symbol is the unmarked symbol at a translate of \(x\)
multiplied by the Weyl phase of \(x,y_1,y_2\); the latter is
also invariant under \(h\). The lifts \(hy_i\) give the prescribed
residues for the new column. Any two lifts of the same residue
differ by an element of \(2A_i^{h\Lambda}\), as may be checked at
the primes of \(\mathfrak m,\bar{\mathfrak m}\) and their
complement. Its vacuum character is trivial, and its half
pairing with the chosen lift is integral since both lifts lie in
\(2(h\widehat\Lambda\otimes\widehat L_i)\).
Thus each marked translate is independent of this choice.
Rational lifts with these residues exist by the Chinese remainder
theorem. Summing with the given weights on \(J\times J\) proves
the formula.
\end{proof}
The proof uses the finite additive character trivial on
\(\widehat{\mathbb Z}\) obtained from the rational trace pairing.
It gives the finite diagram also at \(2\), with the specified
semicharacter. In particular the normalized class average of
these genuine doubled symbols is exactly the finite input of
Lemma~\ref{bfam:genus}. The choices of phases for the initial
row vacua are allowed to depend on \(L_i\).
\src{src/Development/AbelianVarietyTheory/Definitions/SpanningFamilies/FirstCoordinateAdmissible.lean}{81}{The sign in the doubled position coordinate};
\src{src/Development/AbelianVarietyTheory/Proofs/SiegelThetaTransformation/DoubledHermitianTestFunctionZeroAndNorm.lean}{275}{The common origin normalization};
\src{src/Development/AbelianVarietyTheory/Definitions/SpanningFamilies/IsSystem.lean}{148}{Vacuum semicharacter and the normalized projection};
\src{src/Development/AbelianVarietyTheory/Proofs/SiegelThetaTransformation/DoubledVacuumEqHermitianCharacter.lean}{554}{Constancy of the mixed scalar across columns}.

In the doubled model a rational maximal isotropic subspace of
\(W\oplus(-W)\) either is the graph of an isometry, or contains an
isotropic row line in \(W\oplus0\). Indeed if its first projection
has a kernel, its second projection is orthogonal to that kernel,
and hence has a nonzero kernel as well. If both projections are
invertible, isotropy is precisely the isometry equation. On a
graph the finite evaluation at zero is the inner product with the
second vacuum transported by the row isometry, since the base
graph has functional \(W_{u,v}(0)\). For a nongraph the evaluation
is invariant under Weyl translations in the indicated row line
tensored with \(V\). In the polarization of this maximal isotropic
subspace they are multiplication by additive characters of the
position variable and each such character is one at zero. Thus
the augmentation argument above kills this orbit term, including
after archimedean polynomial insertion. This accounts for all
rational isotropic subspaces.

\begin{corollary}[The Fourier mark on the identity cover]\label{bfam:mark}
At the deep level used for the raw spanning calculation the
auxiliary map of that calculation is the identity. Its auxiliary
ideal is \(O\), its integer is one and its auxiliary character is
trivial. Thus \(\mathfrak m\) is the same half of the level on
both sides. The subgroup \(J\) of residues is the identity:
its entries differ from the identity by reductions of elements
of \(\mathfrak m\) modulo \(\mathfrak m\). The algebraic
Fourier mark on the half-primary torsion selects exactly the
single translated shifted vacuum with the first \(r\) basis
rows prescribed and the remaining rows zero.
\end{corollary}
\begin{proof}
Put \(A=\widehat\Lambda\otimes\widehat L\) and
\(D=\widehat\Lambda\otimes\mathfrak m\bar{\mathfrak m}^{-1}
\widehat L\), and use a subscript \(*\) for the standard column.
The alternating commutator identifies \(A/(A\cap D)\) with
the character dual of \(D/(A\cap D)\): each induces a character
on the other group, and self-duality of both lattices gives
trivial kernel on both sides. The second group is supported on
\(\bar{\mathfrak m}\), where it is
\((\bar{\mathfrak m}^{-1} L/L)^n\); at those primes the column is
standard. It is the group \(Q\) of the \(n\) half-level torsion
points. To see the identification, use the level isomorphism
on \(M\)-torsion and the two Chinese remainder idempotents of
\(O/(M)\), since \((M)=\mathfrak m\bar{\mathfrak m}\).
The dual group similarly has residues in \((L/\mathfrak m L)^n\).
In these identifications the Fourier coefficient of the chosen
algebraic mark is precisely the pairing character for those
residues; this is the defining property of the dual marking.
One can lift its tuple by \(y\in\widehat\Lambda\otimes
2\bar{\mathfrak m}\widehat L\), since \(2\bar{\mathfrak m}\)
is a unit at \(\mathfrak m\). Let \(v\) be a unit unshifted
vacuum for the standard column and put \(z=P_Av\). The restriction
of \(\epsilon(d)\rho(d)\) to the intersection fixes \(z\).
Consequently averaging these operators on \(z\) amounts to averaging
over \(Q\). Conjugation by \(\rho(y)\) multiplies the summands
by their pairing character. Moreover \(\rho(y)z=z\), because
\(y\in2A\). With the character taken from this conjugation, the
normalized Fourier average is therefore \(\rho(y)P_D P_Av\).
Taking the opposite convention for the additive character conjugates
both the definition of the dual mark and the commutator.
The odd-order canonical Heisenberg lifts agree on \(Q\) with this
action on \(z\). Indeed the lifts on its isotropic subgroup are
additive, inversion acts by negation, and a change of such a lift
is a character of order dividing two on the odd group.

The normalization here can be compared on \emph{all} rows at
once. Pick the initial unit shifted vacuum to be the normalized
\(P_{D_*}v\). At the primes over \(M\), \(A=A_*,D=D_*\);
elsewhere \(D=A,D_*=A_*\). The local projectors at distinct
places commute, so on \(v\) we have
\(P_DP_Av=P_DP_{D_*}v\).
The product statement follows by factoring a vacuum for \(A_*\)
into local vacua. For completeness the vacuum line of a self-dual
lattice is one dimensional: the calculation of \(W_{u,v}\) above
gives \(|\langle u,v\rangle|=1\) for any two unit vectors fixed
with its character. The restricted tensor of local vacua supplies
one such vector, hence any choice differs only by a unit scalar.
Projector norms in Lemma~\ref{bfam:phase} now show that the
coefficient of the unit shifted marked vacuum is the positive
number
\[
 [D_*:D_*\cap A_*]^{-1/2}\,
 [D:D\cap D_*]^{-1/2}.
\]
The first index is \(|O/\mathfrak m|^{nr}\). The second is
\([\widehat\Lambda:\widehat\Lambda\cap\widehat\Lambda_*]^r\),
by using local bases of the shifted row lattice. Both are
independent of the row and its component. Changing the original
shifted unit vacuum for a fixed row multiplies all its projected
unit vacua by the same phase. Finally, the residue conditions on
the last \(n-r\) rows on both sides of the comparison are zero
modulo the identical ideal \(\mathfrak m\).
\end{proof}
\src{src/Formalism/LocalToGlobal/ProjectiveLattices/ThShort.lean}{300}{Identity auxiliary map on the deep model};
\src{src/Development/AbelianVarietyTheory/Definitions/SpanningFamilies/IsBaseClassLevelAdapted.lean}{78}{The two half-level ideals and the dual residue marking};
\src{src/Development/AbelianVarietyTheory/Proofs/SiegelThetaTransformation/LeSupOfCoprime.lean}{100}{Composition of the unshifted and shifted projections}.

\paragraph{Integral change by a flat character.}
Fix an algebraic component \(Z\) of the cover over \(F\).
For a rational trivialization \(t\) of a flat automorphic line on
\(Z\), suppose its pointwise metric is constant \(r_\sigma>0\)
on each complex fibre, with conjugate embeddings giving the same
constant. Suppose \(\rho_w\) is a lower bound for its vertical
orders at every test over the finite prime \(w\) of \(F\), in the
normalization which gives the order \({\rm ord}_w(a)\) to
\(a\in F^\times\). Assume \(\rho\) has finite support and is
integral off a finite set \(B\) of rational primes. Put \(h=[F:\mathbb Q]\)
and
\[
 \delta(t)=h^{-1}\left(\sum_w\rho_w\log Nw-
                         \sum_{\sigma:F\hookrightarrow\mathbb C}
                               \log r_\sigma\right).
\]
There is \(a\in F^\times\) such that \(at\) is integral at all
the tests and
\begin{equation}\label{bfam:flat}
 \max_\sigma\log\|at\|_\sigma
 \le-\delta(t)+\frac{\log|\Delta_F|}{2h}+\sum_{\ell\in B}\log\ell .
\end{equation}
Indeed take \(\mathfrak b=\prod_w w^{\lceil-\rho_w\rceil}\).
Every element of this ideal gives integrality. The rounding cost in
\(\log N\mathfrak b\) over \(\ell\in B\) is at most
\(\sum_{w\mid\ell}\log Nw\le h\log\ell\).
In the Minkowski embedding use the intervals and complex discs
with radii \(R/r_\sigma\) at real and complex places. Their volume
is \(2^{r_1}\pi^{r_2} R^h/\prod_\sigma r_\sigma\).
The covolume of \(\mathfrak b\) is
\(2^{-r_2}\sqrt{|\Delta_F|}\,N\mathfrak b\).
Taking \(R^h=\sqrt{|\Delta_F|}\,N\mathfrak b\prod_\sigma r_\sigma\)
achieves Minkowski's bound, including the equality case by compactness.
It gives \eqref{bfam:flat}.

If \(t\) is a nowhere vanishing section on the characteristic-zero
component, multiplication is a field-linear isomorphism to the
twisted cusp space there whenever the flat trivialization and
its inverse preserve cusp boundary conditions. Extension by zero
from the component and restriction back are inverse on its supported
sections. Thus any field-spanning family of global sections restricts
to a spanning family on \(Z\) under this component decomposition.
Multiplication by \(at\) preserves the vertical inequalities and
multiplies the component Petersson norm by \(|\sigma a|r_\sigma\):
this follows under the integral defining that norm, since the
absolute metric of the multiplier is constant on the component.
The inverse of \(at\) is needed only over \(F\). In particular no
division by a component mass or by a generator of a fractional
ideal is needed in the integral lattice.
\src{src/Development/AbelianVarietyTheory/Proofs/DivisorProofs/ExistsShortIntegralFlatLineTwistingSection.lean}{203}{Minkowski rescaling for fractional vertical orders};
\src{src/Development/AbelianVarietyTheory/Proofs/DivisorProofs/ExistsSpanningFamilyMulTwistingSection.lean}{97}{Component projection and multiplication on the field-spanning family}.

\paragraph{Passage from the marked family to the cover.}
Fix a good embedding of \(F\) and a component
\(k\). Let \(L_k\) be its row lattice after transport of its level
marking into the common rational Hermitian space. Write
\(\Theta_{\Lambda,P}(L_k,x)\) for the corresponding analytic
theta jet on a frame \(x\), using the marked vector of
Corollary~\ref{bfam:mark}, and write \(f_{\Lambda,N,P}\) for the
algebraic rule. Identify the polynomial variables by the
eigenframe and the framed periods. The required reading has the form
\begin{equation}\label{bfam:reading}
  f_{\Lambda,N,P}(k,x)=C_k\,u_{\Lambda,N}\,
                   \Theta_{\Lambda,P}(L_k,x).
\end{equation}
Here \(C_k\ne0\) is independent of \(\Lambda,N,P,x\); the
nonzero \(u_{\Lambda,N}\) is independent of \(k,P,x\).
A different normalization for the initial vacuum for \(L_k\)
can be absorbed in \(C_k\).

\begin{proposition}[The theta reading on a polarized period chart]
\label{bfam:unitread}
Use the principal PEL period chart of Corollary~\ref{bfam:metric}
on one connected component. Use its rigidified identification
\(\mathcal L=(1,w_0)^*\mathcal B\) with the pullback of its
polarized biextension, its given half-level marking and its dual,
and its labelled eigenframe as the tangent coordinates. Use the
identity auxiliary map of Corollary~\ref{bfam:mark}. For the raw
section of Lemma~\ref{newtheta:index}, \eqref{bfam:reading} holds
on unit flags up to a common factor of constant positive absolute
value in place of \(C_k\). With the shifted standard vacua chosen
as in Corollary~\ref{bfam:mark}, one may take
\[
        u_{\Lambda,N}=|Q|^{-1/2}
           [\Lambda:\Lambda\cap O^n]^{r/2}.
\]
For the convention that divides the raw rule by \(\beta_H\),
multiply this number by \(\sigma(\beta_H)^{-1}\).
On the good framed chart, \eqref{bfam:reading} holds on the
connected component with a single \(C_k\ne0\).
\end{proposition}
\begin{proof}
Write \(\Pi\) for the row period lattice and \(g=fr\).
In the Appell--Humbert frame over the Lie algebra the automorphy
factor of \(\mathcal L\) at a period \(v\) has unitary
semicharacter \(\chi(v)\) and Hermitian form \(H\).
The pullback by \((1,w_0)\) of the addition biextension has
Hermitian form \(H(z,w_0z')+H(w_0z,z')=H(z,z')\), by Rosati
adjunction. Its semicharacter is the quotient of the characters
at \((1+w_0)v,v,w_0v\), namely
\(\exp(\pi i E_L(v,w_0v))\) by the semicharacter law. Uniqueness
of the Appell--Humbert data of a rigidified line gives this same
semicharacter for \(\mathcal L\). Its values on periods are signs.

Choose a symplectic \(\mathbb Z\)-basis of \(\Pi\) and denote the
two period matrices in the unit tangent frame by \(P_0,Q_0\).
Set \(h=P_0^{-1}\), \(\Omega=h Q_0\), \(Y=\Im\Omega\). The
Riemann bilinear relations give \(Y>0\), \(\Omega^t=\Omega\)
and \(H(z,z)=(hz)^tY^{-1}\overline{hz}\). The signs in this
basis specify a half characteristic \((\alpha,\eta)\). Put
\[
 T(z)=e^{\pi(hz)^tY^{-1}(hz)/2}
          \sum_{\ell\in\mathbb Z^g+\alpha}
           e^{\pi i\ell^t\Omega\ell+2\pi i\ell^t(hz+\eta)}.
\]
Its multipliers follow by shifting the summation index on imaginary
periods and multiplying each term on real periods. The signs
produce the specified characteristic. The same period laws
force every holomorphic section to be a scalar times \(T\):
remove the first exponential, take the Fourier series in the real
periods, then use the other periods to recover every coefficient
from one. This also proves \(T\ne0\).

For a unit finite vacuum for \(A_*\), the \(n\)-fold Schrödinger
theta in these coordinates is
\[
       d_L(\det h)^{n/2}\prod_{j=1}^n T(z_j)
\]
with a fixed \(d_L\ne0\). In the basis adapted to \(\Pi\) the
integral translations and their dual characters force the finite
test to be supported on the translate
\(\alpha+\widehat{\mathbb Z}^{g}\) in each row; its values there
are fixed multiples of the character prescribed by \(\eta\).
Restriction of this function to the rational points gives exactly
the displayed characteristic series. The determinant is the
weight-\(n/2\) factor for the change \(z\mapsto hz\).
To compare to a fixed initial rational symplectic basis, it suffices
to change basis by a linear substitution, a symmetric shear, or
Fourier exchange. These generate the rational symplectic group by
symplectic elimination. Linear substitution reindexes the rational
sum. For a shear the archimedean and finite quadratic characters
at every rational point are inverse by the product formula. For
exchange, apply Poisson summation with dual Haar measures to the
finite test times the archimedean Gaussian with linear argument.
Completing the square gives in each row the factor
\(\det(-i\Omega)^{-1/2}
e^{-\pi i\zeta^t\Omega^{-1}\zeta}\) for the Fourier transform
with shifted argument \(\zeta\). The quadratic exponential is
absorbed by the change of the holomorphic Gaussian prefactor in
\(T\), and the half determinant by the change of \(\det h\).
Thus the new finite vacuum is multiplied by a fixed scalar
independent of the period point; choices of square root give
only a constant. This is the period-coordinate version of the
Schrödinger theta transformation.

The squared mean of \(T\) on the polarized torus is
\(\det(2Y)^{-1/2}\). In fact write \(hz=t+\Omega s\) in
the unit period cube. The integral over \(t\) kills all unequal
Fourier indices. The remaining integrand, including the metric
\(e^{-\pi H(z,z)}\), gives the translates of
\(e^{-2\pi s^t Y s}\); integrating the sum over the \(s\)-cube
gives its integral over \(\mathbb R^g\).
The squared mean of the \(n\)-fold Schrödinger product is
therefore
\[
    |d_L|^2|\det h|^n\det(2Y)^{-n/2}
      = |d_L|^2\bigl(\det H/2^g\bigr)^{n/2}.
\]
At unit flags \(\det H=(2x_0')^g\), where \(x_0'>0\) denotes
the common absolute polarization scalar called \(x\) in
Corollary~\ref{bfam:metric}; equivalently one may choose
orthonormal flags first and extend by the flag character.
The determinant is constant. The geometric generator has squared
mean one. Hence its \(n\)-fold product is the Schrödinger product
times a scalar of constant positive modulus, independent of all
column data and polynomial insertions.

It remains to check what the geometric translations and
differentiations do to this one product. In these coordinates
put
\[
  (U_t^{(m)}s)(z)=e^{-\pi mH(z+t/2,t)}s(z+t).
\]
It sends \([N]^*s=s(Nz)\) under \(U_t^{(N^2)}\) to
\((U_{Nt}^{(1)}s)(Nz)\). This pullback equality in the
Appell--Humbert frames follows since the factors of automorphy
agree and the isomorphism is rigidified at zero. When
\(N^2t\in2\Pi\), inversion sends this translate to its inverse
and its \(N^2\)-th power acts trivially: the residual period
sign at \(N^2t\) is one. Thus it is exactly the normalized
geometric lift of the rule. Such representatives can always
be chosen since \(N\) is odd. On an individual summand in
period coordinates \(u=a+\Omega b\) the finite Heisenberg
operator on a rational coefficient function is
\[
   v(s)\longmapsto
             e^{2\pi i a^t(s-b/2)}v(s-b).
\]
Reindex \(s\) by \(s+b\) in the characteristic sum. The linear
argument changes from \(\zeta\) to \(\zeta+u\), and the factor
is \(e^{\pi i(b^t\Omega b+2b^t\zeta+a^tb)}\).
Restoring the holomorphic Gaussian this is just \(U_{h^{-1}u}^{(1)}\).
Thus this comparison holds with phases before taking any norm.

For the first sum of the rule write the division points as
\(t=s/N\) with \(s\in\Lambda\otimes\Pi\) modulo \(N\Pi^n\),
as in Lemma~\ref{newtheta:index}. Representatives can be chosen
even in \(\Lambda\otimes\Pi\): its quotient by \(N\Pi^n\) has
odd order since the column is standard at two and \(N\) is odd.
Then \(Ns\in2\Pi^n\) and their semicharacters are trivial.
The finite sum on the initial vacuum is exactly
\(|Z_\Lambda|P_Av\), since \(N\Pi^n\) fixes it with the
prescribed semicharacter. The actions of the second, half-level
sum were computed in Corollary~\ref{bfam:mark}. Equivalently
its translating vectors modulo \(\Lambda\otimes\Pi\) exhaust
\(D/(D\cap A)\): at the primes of \(M\) this is the given
half-level isomorphism, and at its complement the primary torsion
class vanishes. Division representatives may be taken to be zero
modulo \(2N\Pi^n\) at the complement of \(\bar{\mathfrak m}\)
and even in the shifted lattice at \(\bar{\mathfrak m}\), by the
Chinese remainder theorem. Their symmetric lifts consequently
act by the normalized shifts used in that corollary.

In the universal-cover Appell--Humbert frames the constant
biextension frame on two Lie variables is biadditive and
Rosati-compatible. Indeed the factor of automorphy of the
addition biextension is obtained by dividing the factor at
\(z+z'\) by those at \(z,z'\). The resulting cross terms are
additive in the opposite variable, and
\(H(az,z')=H(z,\bar a z')\). Equivalently these statements follow
by taking the rigidified cube trivializations; in this frame the
third differences of the linear and quadratic exponents vanish.
The pullback frame along \((1,w_0)\) is the
Appell--Humbert frame of \(\mathcal L\) just used, by equality
of their Hermitian forms, semicharacters and rigidifications.
Consequently the algebraic jet may be evaluated by ordinary
differentiation on this analytic frame. Any change of compatible
frame is the common bilinear ambiguity killed by harmonicity
in Lemma~\ref{bfam:rule}. Each labelled Hodge derivative is
the corresponding coordinate derivative by the definition of the
framed period chart. Substitution \(z\mapsto Nz\) contributes
\(N^{d_P}\), cancelled by the same denominator in the rule.
All sums are finite except the normally convergent holomorphic
theta series, so these differentiations commute with summation.

The raw column factor is thus \(|Z_\Lambda|/|K_0|\) times the
two square-root indices in Corollary~\ref{bfam:mark}. Setting
\(I=[\Lambda:\Lambda\cap O^n]\), the three factors are
\(I^r,\ |Q|^{-1/2},\ I^{-r/2}\), respectively. This proves
the asserted formula, independent even of the auxiliary \(N\).

For the extension off unit flags it matters that the
Appell--Humbert Gaussian itself need not depend holomorphically
on the period. Work instead in the classical holomorphic
trivialization on the varying torus. The matrices \(P_0,Q_0\)
are holomorphic in the flag coordinate: their entries are either
entries of the inverse flag matrix applied to a fixed period,
or contractions of a flag column against a fixed vector by the
Hermitian form linear in the flag column. Thus \(h,\Omega\) are
holomorphic. The characteristic series without its first Gaussian
defines a holomorphic family of sections on the torus in these
coordinates. Its integral-period automorphy factors, read by
shifting the sum, are holomorphic in \(\Omega,\zeta\); the series
and its derivatives converge normally on compacta, by Gaussian
decay for \(Y\) uniformly positive. Rigidify its line at zero
in this holomorphic trivialization. On each fibre changing to the
Appell--Humbert frame multiplies that function by exactly the
Gaussian in \(T\): it has value one at zero and takes the
classical factors of automorphy to the Appell--Humbert factors.
On the tautological family of period tori this gives the same
relative rigidified polarized line as the pullback of its
biextension by \((1,w_0)\). Indeed their ratio is trivial on
every fibre by the computation of the two Appell--Humbert data;
its
pushforward is a line by cohomology and base change in degree
zero and evaluation is an isomorphism by fibrewise evaluation.
The rigidification trivializes that base line.

Apply the geometric torsion projections and harmonic jet to this
holomorphic characteristic section. These operations are
holomorphic in the period: the torsion points are locally
holomorphic sections of finite étale groups; the formal
biextension frame of Lemma~\ref{bfam:rule} can be constructed
over the local holomorphic base by the same finite-order
operations. The preceding fibre calculation identifies the
result with the analytic marked jet, up to the displayed
nonzero index factors, \(d_L\) and \((\det h)^{n/2}\), all
holomorphic in this chart. Lemma~\ref{dg:gaussian} also proves
this directly by removing the conjugate-type quadratic terms with
the harmonic contractions. This proves holomorphy of the
analytic jet without differentiating a nonholomorphic
\(Y^{-1}\) in the period. Its flag character follows in the
same coordinates. A triangular flag change gives the corresponding
triangular changes in the two labelled tangent blocks (first
columns in one block, last inverse rows in the other). It changes
both period matrices by the same tangent substitution, hence
fixes \(\Omega\); the determinant factor gives the scalar
weight \(n/2\), and the labelled derivatives give the character
of the row-highest polynomial. These are precisely the two parts
of the geometric weight. The geometric reading of a regular
weight rule is holomorphic with that character by the good framed
uniformization (its section property).
Finally all
cross products of the two rows of readings vanish on unit flags
and hence on every flag by the weight and Gram--Schmidt. Where
one analytic reading is nonzero its quotient with the geometric
reading is holomorphic; the quotient has the same absolute value
as on unit flags, hence is constant. Holomorphic continuation
gives the claimed equality everywhere on the component. When
all analytic entries vanish the same reasoning on unit flags
gives zero throughout for the geometric entries as well.
\end{proof}
\src{src/Formalism/LocalToGlobal/ProjectiveLattices/PetersonNormSubspace.lean}{75}{The section property and period identification in a good uniformization};
\src{src/Development/AbelianVarietyTheory/Proofs/SiegelThetaTransformation/PeriodCoordinateThetaOfSignedHeisenbergVacuum.lean}{451}{The half characteristic and determinant in an adapted basis};
\src{src/Development/AbelianVarietyTheory/Proofs/SiegelThetaTransformation/ExistsSymplOfSelfDual.lean}{226}{The probability normalization on the period torus};
\src{src/Development/AbelianVarietyTheory/Proofs/SiegelThetaTransformation/ShortThetaValueComplexPointEqClassicalTheta.lean}{194}{The product of generators and its covariant derivatives};
\src{src/Development/AbelianVarietyTheory/Proofs/SiegelThetaTransformation/ShortSeriesReadingsEqClassicalThetaSeries.lean}{88}{The separated nonzero constants in the actual chart comparison}.

\begin{theorem}[Integral field generation on the marked period charts]\label{bfam:oncharts}
Let \(p,q\ge1,\ m=fr+1,\ n=8m\). Suppose the actual level-\(M\)
PEL cover has at each good embedding a good framed uniformization
by the finitely many marked periods of its common rational
Hermitian space, as in Lemma~\ref{munf:marks} and
Proposition~\ref{bfam:unitread}. Use its identity auxiliary map
and the positive dominant harmonic weight of size \(S=d_P+frn/2\).
Assume the cusp space is read as bounded holomorphic forms of that
weight invariant under the marked principal stabilizers, the raw
rules have the prescribed cusp condition at all good embeddings,
and their readings satisfy the half-maximum hypothesis of
Corollary~\ref{bfam:metric}. Corollary~\ref{crev:good} verifies
these analytic conditions for the selected good covers and their
rule spaces. Use the semiabelian integral tests
of Theorem~\ref{intjet:semi} and the field, level and weight size
bounds of Corollary~\ref{bfam:metric}. Then each algebraic
component has a finite set of integral cusp sections spanning
its cusp space over \(F\). Their logarithmic Petersson norms at
every embedding are at most
\[
                 C(frL)^2 S+2\log D_M .
\]
The same estimate holds for simultaneous generators of the
untwisted full cover with its probability measure.

For the final weight \(\lambda\), put \(W=\sum_{\tau,i}
|\lambda_{\tau,i}|\), \(j_\tau=\lambda_{\tau,p}-1-4m\)
(algebraic weight rows numbered from one, with \(p\) the Lie
multiplicity at \(\tau\)). The first Hermitian block in its period
chart has \(p_0=q\) rows. Use the actual weight decomposition
with \(S\le(1+r)W+rf(1+4m)\), and the key section on the tests of
Theorem~\ref{ctwist:semi}. If \(j\ne0\), assume the horizontal
parallel determinant sections and their product have the common
vertical orders on each local component and the corrected degree
estimate of \eqref{ctwist:degree}, with the stated ramification-one
tests outside \(B\). Their regular trivializations and inverses
with constant pointwise norms preserve the rule cusp conditions
at all good embeddings by Corollary~\ref{crev:twist}. Under the stage bounds
for \(B\), the coefficient field and \(D_M\), the resulting
integral sections span on each component and satisfy
\[
           \log\|f\|_\sigma
            \le C'\bigl((frL)^3 W+(frL)^5\bigr).
\]
If every \(j_\tau=0\), the conclusion holds without the parallel
determinant or degree hypotheses. For \(j\ne0\) the same conclusion
also follows from the one-sided product criterion after
\eqref{ctwist:degree}, provided its lower degree bound gives the
same cost in \eqref{bfam:flat}. On the indefinite principal
family of Proposition~\ref{host:art} with its selected marked
PEL moduli and tests, Theorem~\ref{cmeig:splitmodel} supplies the
common orders and degree for every \(j\). Outside the excluded
primes the original level is hyperspecial, so integrality of
the common order already follows from
Proposition~\ref{cmeig:hyperspecial}.
\end{theorem}
\begin{proof}
First consider the row multiplier. A rational isometry maps the
unit vacuum of a row to a unit scalar times that of the transported
row, by preservation of the specified semicharacter and the
one-dimensional vacuum line. Conjugation also carries the
projectors to those for the transported row, so this scalar is
independent of the column. On the marked stabilizer it is one
at the rank used here. Indeed in a symplectic basis of the
\emph{shifted} row lattice the integral symplectic matrix
preserves its half characteristic, since the isometry commutes
with \(w_0\). For one characteristic theta the residual factor,
after the square root of the Siegel determinant and the quadratic
exponential, is an eighth root of unity. This follows on the
integral symplectic generators. A symmetric integral shear
changes the characteristic by the shear and its half diagonal;
the remaining exponential has argument \(\pi i\) times a
quarter-integer. An integral linear basis change reindexes the
sum, with at most a square root of \(-1\) from its determinant.
Fourier exchange gives \(e^{-\pi i g/4}\) times
\(e^{\pm 2\pi i\alpha^t\eta}\) and the displayed Gaussian
transform in Lemma~\ref{dg:gaussian}'s coordinates.
Integral changes of representatives of the characteristic give
signs, and choices in composing square roots give signs.
Thus the eighth-power assertion persists under composition.
The determinant factors cancel those of the period matrices in
the normalized theta, as in Proposition~\ref{dg:identity}. The
\(n\) identical rows contribute the \(n\)-th power of the
residual factor, equal to one since \(8\mid n\).

The marked translate introduces no extra central character here.
Indeed the residue label is the identity by
Corollary~\ref{bfam:mark}. Both \(y,\gamma y\) lie in
\(2\widehat\Lambda\otimes\bar{\mathfrak m}\widehat L\) for a
stabilizing \(\gamma\), and their residues agree at
\(\mathfrak m\). Their difference lies in \(2D\), by the lift
calculation in Lemma~\ref{bfam:phase}. The vacuum character on
this difference is one. The Weyl phase of the two translations is
also one since their half pairing is integral in the unshifted
self-dual lattice. Hence \(\rho(\gamma y)\) and \(\rho(y)\)
give the very same marked vector. This establishes compatibility
with the invariant geometric readings in the hypotheses. For
arbitrary even \(n\ge8fr+4\) the first conclusion also follows
whenever the geometric and vacuum stabilizer characters match;
divisibility by eight proves this particular matching at the
rank of the final weight composition.

Lemma~\ref{munf:marks} now supplies distinct marked row orbits.
Apply Proposition~\ref{munf:span} with the actual graph identity
of Proposition~\ref{dg:identity}, after dividing each component
by its constant \(C_k\). Equations~\eqref{bfam:reading} and the
nonvanishing of \(u_{\Lambda,N}\) give simultaneous generation.
The sum covers the entire genus of the standard column with its
prescribed level; a single representative need not suffice.
On each chosen column Proposition~\ref{bfam:harm} replaces the
polynomials by its finite spanning set. Lemma~\ref{bfam:rule}
and faithful scalar extension descend generation to \(F\).
Multiplication by the respective reflex scalars and the common
\(D_M^2\) preserves the field span. Theorem~\ref{intjet:semi}
gives integrality, and Corollary~\ref{bfam:metric} gives the
simultaneous norm bounds, including for the component probabilities.

On each component multiply further by its integral key section
of weight \(-4m\) and its integrally rescaled parallel determinant.
Theorem~\ref{ctwist:semi} and equations
\eqref{ctwist:degree}, \eqref{bfam:flat} give the combined
logarithmic cost at most
\[
  2frm\log(2\pi)+12rL\bigl(W+f(1+4m)\bigr)
  +\sum_{\ell\in B}\log\ell+
                    \frac{\log|\Delta_F|}{2[F:\mathbb Q]} .
\]
The last two quantities are \(O((frL)^2)\) under the stage
bounds. Also
\[
 2\log D_M=2(frn+d_P)\log M+2d_P\log\delta_E
                     \ll fL S
\]
by the level and discriminant bounds in
Theorem~\ref{intjet:semi}'s application.
Inserting the stated bound on \(S\) gives the result.
For \(j=0\) the parallel section is identically one and has
neither an order nor a metric cost. If only \(\sum_\tau j_\tau=0\),
formula~\eqref{ctwist:degree} cancels the common degree once
the common finite orders and common corrected degree have been
obtained. Component restriction and extension by zero preserve
the field decompositions and the integral evaluation tests.
\end{proof}
For every selected cover satisfying \(Y.\mathrm{Good}(e,o)\),
Corollary~\ref{crev:good} supplies the marked period, cusp,
half-maximum and measure clauses with the stated stage parameters.
For nonzero parallel weights
Theorem~\ref{cmeig:splitmodel} supplies the local orders and the
arithmetic degree on those marked components. Neither the finite
genus identity nor the positive real reproducing scalar depends
on that degree.
\src{src/Development/AbelianVarietyTheory/Proofs/SiegelThetaTransformation/KernelIdentity.lean}{181}{Doubled identity with the level tests and harmonic insertions};
\src{src/Development/AbelianVarietyTheory/Proofs/SiegelThetaTransformation/HermitianThetaSeriesSpanBoundedForms.lean}{168}{Inversion and reproduction on marked component families};
\src{src/Development/AutomorphicArithmetic/ThetaCorrespondence/ThetaPairLifts/HodgeBundleWeights/ShortSpanCompositionClosed.lean}{69}{The two actual component twists}.

\paragraph{The cusp test on the principal marked charts.}
Retain the period chart and identity Fourier mark of
Propositions~\ref{bfam:unitread} and \ref{bfam:oncharts}.
The following argument uses the mark at one prime over \(M\).
In particular it applies before summing over the definite column genus.

\begin{lemma}\label{crev:mode}
For every rational isotropic row line and every orthonormal frame on
the chart, the real diagonal transvection restriction of each marked
theta jet has mean zero along any period with trivial multiplier.
Its full real cusp-torus average is zero as well.
\end{lemma}
\begin{proof}
Write \(v\ne0\) for the isotropic row and \(I=(Ev)^n\) in the
rational trace symplectic space of the oscillator. Choose a rational
Lagrangian in one row containing \(Ev\), and take its \(n\)-fold sum
\(D\). Complete it to a symplectic polarization \(D\oplus D^*\),
using \(D^*\) for the position variable. Translations by \(a\in
D(\mathbb A_f)\) are multiplication by the finite additive character
of \(z(a)\) at a position \(z\). In particular they act as the identity
at all \(z\) annihilating \(I\).

Let \(\varphi\) be the marked finite vector in these coordinates.
The finite translation calculation preceding Lemma~\ref{bfam:genus}
gives an element of \(I(\mathbb A_f)\) acting on the single marked
vacuum with eigenvalue different from one. Indeed, at a prime of
\(\mathfrak m\) its first \(r\) residues are a basis and a primitive
vector on the conjugate line, divided by the local half-level,
detects one of them. The unshifted character of the isotropic line
is one. It follows pointwise that
\[
                 \varphi(z)=0\qquad (z|_{I(\mathbb A_f)}=0).
\]
This also holds for a sum of the indicated marks, term by term. The
calculation takes place after changing polarization, which intertwines
the Weyl translations by the Fourier and substitution laws in the
proof of Proposition~\ref{bfam:unitread}.

In this polarization the real diagonal flow is a shear. Its
infinitesimal action \(N\) has image \(I_{\mathbb R}\) and kernel
\(I_{\mathbb R}^{\perp}\), and it fixes \(D_{\mathbb R}\) pointwise.
The shear contributes the character with quadratic exponent,
up to the fixed Fourier sign, \(Q(z)=\tfrac12\langle \widetilde
z,N\widetilde z\rangle\), where \(\widetilde z\) is a lift of the
position and the pairing is symplectic. At each place the trace
polarization applied to the transvection
\(\kappa_0 h(y,v)v\) gives a nonzero real multiple of
\(\sum_{j=1}^n|h_\tau(y_j,v)|^2\). The multiples have the same sign:
they are obtained from the product of \(\tau\xi_0\) and
\(\tau\kappa_0=\pm2\tau\xi_0\) with one fixed choice of sign, and
the scaling \(k_0\) of the real parameter is positive. Consequently
\(Q(z)=0\) exactly on the annihilator of \(I_{\mathbb R}\).

More explicitly, on the real flow the first period matrix is fixed
because the shear fixes the chosen first Lagrangian. The second
matrix changes by the first times the real symmetric shear matrix.
Thus \(\Im\Omega\), the holomorphic argument \(hz_j\), the Gaussian
prefactor and \(\det h\) in Proposition~\ref{bfam:unitread} stay
fixed. The change in each summand is exactly its quadratic phase,
also after applying the polynomial in the \(z_j\)-derivatives.
The real restriction therefore has the expansion
\[
            \sum_{\ell\in D^*(\mathbb Q)}
                 \varphi(\ell)\,A_x(\ell)\,e^{2\pi i s Q(\ell)}
\]
with the Fourier sign absorbed in \(Q\). At the fixed frame the
archimedean coefficient is a polynomial times a decaying Gaussian,
and the compact finite support restricts rational positions to a
fixed denominator lattice. The series converges absolutely and
uniformly for real \(s\). If a rational position annihilates \(I\)
at infinity it annihilates \(I\) at all finite places. Hence every
zero-frequency summand vanishes. Averaging over \([0,T]\) and
taking \(T\to\infty\) gives zero by dominated convergence, even
when individual frequencies are not characters of the chosen
period circle. For a periodic sum this is its period mean.
Real translates of the starting frame give the same conclusion.
Integrating along the diagonal circles in the full cusp torus now
gives the torus average by translation invariance of Haar probability.
\end{proof}

\begin{proposition}\label{crev:half}
On each of these principal period charts the raw theta readings at
\(n=8(fr+1)\) are bounded in Petersson norm. For every nonzero
reading each orthonormal frame attaining at least half the component
supremum has row minimum squared at least \(\varepsilon\) of
Proposition~\ref{bergvol:global}. The same minimum conclusion applies
to every nonzero bounded cusp form with trivial transvection multiplier on
the chart. In particular the cusp and half-maximum tests on the raw
rules in Theorem~\ref{bfam:oncharts} follow from its stated marked
period uniformization.
\end{proposition}
\begin{proof}
Write \(a(x)=\min_{0\ne v\in L'}\sum_\tau{\rm maj}_{x,\tau}(v)>0\).
Lemma~\ref{bfam:minimum} and Proposition~\ref{intjet:arch}, at
orthonormalized frames, first give a bound of the form
\[
             \|P(x)\|_{\rm Pet}\le C(1+a(x)^{-1})^B .
\]
Indeed \(\lambda^2\ge 2a(x)/R_E^2\), the theta generator has
mean square one, and the directional norm and its weight correction
are fixed on orthonormal frames. Here \(C,B\) may depend on the
fixed section and stage. Multiplication by its nonzero fixed
reflex scalar does not affect this assertion.

Put \(H=2M\sqrt d\). For \(a=a(x)<\min(1,H^{-1})\) choose a
shortest \(v\). Its integral self-pairing vanishes by the norm
argument following Proposition~\ref{bergvol:global}. The
transvection there acts trivially modulo \(M\) on \(\mathfrak a L'\),
so lies in the actual marked stabilizer by Lemma~\ref{munf:marks}.
Its multiplier on the raw sections is trivial by the eighth-power
and identity-mark calculation in Theorem~\ref{bfam:oncharts},
which uses only the transformation formula. Take \(H\) as period
in the scaled coordinate, and write \(G(\zeta)\) for the raw
coefficient along this flow. It is holomorphic on
\(\Im\zeta>-1/a\) by the leading-minor calculation above.

We give the majorant comparison needed to apply the polynomial bound
before knowing a global supremum. At a single place set
\(l=\|\tau v\|_x^2\le a\). Rotate the positive and negative blocks
so the two components of \(\tau v\) lie on their first coordinate
axes, each of norm \(\sqrt{l/2}\). These rotations preserve the
central majorant. In this plane the new positive line has positive
and negative coordinates \(A=1-i\zeta l/2,\ B_1=-i\zeta l/2\),
with \(|A|^2-|B_1|^2=1+l\Im\zeta\). All orthogonal directions
retain their majorant. The two comparison eigenvalues in the plane
are \(K\) and \(K^{-1}\), where
\[
 K=\frac{(|A|+|B_1|)^2}{1+l\Im\zeta}
    \le\frac{(1+l|\zeta|)^2}{1+l\Im\zeta}.
\]
Indeed a unitary change of phase makes the normalized line coordinates
nonnegative and its orthogonal negative mate interchanges them.
At \(\zeta=s-i/(2a)\), \(0\le s\le H\), this gives \(K\le 13\):
the denominator is at least \(1/2\) and
\(l|\zeta|\le Ha+1/2\le3/2\).
Thus the new row minimum is at least \(a/13\).
The leading-minor variation estimate preceding
Lemma~\ref{bfam:harm} bounds the conversion from Petersson norm
to raw coefficient on this line by \(2^{w/2}\).
It follows that
\[
      \sup_{0\le s\le H}|G(s-i/(2a))|
                    \le 2^{w/2}C(1+13/a)^B .
\]
For \(y=\Im\zeta\ge0\), \(0\le s\le H\), the same comparison
gives \(K\le(1+a(H+y))^2\). The leading-minor formula bounds
the raw metric conversion by \((1+ay)^{w/2}\).
Together these give polynomial growth of \(G(s+iy)\) in \(y\)
for the fixed starting frame. In \(q=\exp(2\pi i\zeta/H)\),
Cauchy's coefficient estimates therefore remove all negative
Laurent powers. Lemma~\ref{crev:mode} removes the constant power.
The maximum principle for \(G/q\) on the disk of radius
\(\exp(\pi/(aH))\) gives
\[
      \|P(x)\|_{\rm Pet}=|G(0)|
       \le 2^{w/2}C(1+13/a)^B e^{-\pi/(aH)} .
\]
This is bounded as \(a\downarrow0\). On the complementary range
the original polynomial estimate is bounded. Hence every nonzero
section has a finite positive supremum and one can choose an
orthonormal frame above half of it. With boundedness established,
the half-maximum argument following Proposition~\ref{bergvol:global}
applies with its original \(\varepsilon\); the qualitative constants
\(B,C\) disappear. For a general bounded cusp form the real torus
and positive-frequency argument there supplies its diagonal zero
mode. Thus the same quantitative conclusion holds for that space
at every weight whose transvection character is trivial.
\end{proof}
The proof uses the short vector of the component under consideration.
It gives the absolute-constant bound of Corollary~\ref{bfam:metric}
for all selected parameters, all permitted column lattices and
auxiliary integers at every good embedding; it requires neither
a common fundamental domain as the rank changes nor a uniform
constant in the preliminary polynomial growth estimate.
\src{src/Development/AbelianVarietyTheory/Proofs/CuspFormEvaluations/ExistsExpansionOfEqSumThFr.lean}{26}{The source's different zero-mode test using positivity of the last row weight};
\src{src/Development/AbelianVarietyTheory/Proofs/DivisorProofs/FirstMinAtHalfMax.lean}{111}{Transfer to the minimum at the selected level}.

\begin{corollary}\label{crev:twist}
On the good marked charts the algebraic rule space defined by
testing all good complex readings requires no additional cusp
comparison for the raw spanning argument. A regular unit section on
a component with constant positive Petersson norm, in mutually dual
multiplicative metrics for it and its inverse, preserves this cusp
condition in both directions. This applies to the key and the
horizontal determinant whenever they exist on that component with
their stated constant metrics.
\end{corollary}
\begin{proof}
In the rule definition the cusp subspace is the intersection of
the regular weight rules with the inverse images, under every good
reading, of the complex spans of the analytic cusp forms. Each such
span is the analytic cusp space itself: holomorphy, transformation
law and vanishing of the torus integrals are linear, and each
finite sum of bounded forms is bounded. Conversely inclusion in the
span holds for each member. Lemma~\ref{bfam:rule},
Lemma~\ref{crev:mode} and Proposition~\ref{crev:half} place the
raw rules in this intersection. For generation one only needs to
express the reading of every member in the span of the raw readings,
which is the conclusion of Proposition~\ref{munf:span} on those
charts. Injectivity after scalar extension in Lemma~\ref{bfam:rule}
then descends the span to the coefficient field.

For the multiplication assertion use a common finite-index
transvection period lattice killing both characters. Passing to
it leaves each normalized constant term unchanged when applied to
a function periodic for the original lattice. Along a rational
isotropic row the positive-ray metric estimate following
Proposition~\ref{bergvol:global} and boundedness imply that
every nonzero Fourier coefficient of either factor has all
coordinates nonnegative, by growing one imaginary coordinate
at a time. At height zero the coefficients are absolutely summable
by smoothness on the real torus. Thus the product constant is the
product of the two constants, by convolution: two vectors with
nonnegative coordinates can sum to zero only if both are zero.
This proves cusp vanishing for the product, and boundedness follows
by multiplicativity of the metric. Apply the same argument to the
inverse section, which is holomorphic and has the reciprocal
constant norm. In particular its Fourier support is tested afresh
using holomorphy in the same upper tube, even though its \emph{weight}
has the opposite sign. The polynomial growth bound uses only the
weight spread. This treats the Fourier condition in the rule space;
a separate use of a subcanonical sheaf on a compactification must
identify its restriction with that space.
\end{proof}
\src{src/Formalism/LocalToGlobal/ProjectiveLattices/CuspFree.lean}{14}{All good embeddings in the definition of the cusp rules};
\src{src/Formalism/LocalToGlobal/ProjectiveLattices/JFDFree.lean}{42}{The span of the bounded holomorphic cusp readings}.

\paragraph{The flag measure of the selected unitary genus.}
Let \(F_0=E^+\). Use the host fields of Proposition~\ref{host:art},
the genus of the self-dual row, and its odd principal level \(M\ge3\).
Write \(K\) for the finite adelic stabilizer and \(K(M)\) for its
principal subgroup. Let \(c_r\) and \({\rm mass}_I\) be as in
Proposition~\ref{bergvol:compactmass}.

\begin{proposition}\label{crev:mass}
For the full adelic packet of marked rows, the sum of the component
flag volumes in the normalization of Lemma~\ref{bergvol:local} is
\[
        V=[K:K(M)]\,\frac{c_r^f\,{\rm mass}_I}{(2\pi)^{fr}}
           \le V_* .
\]
Adelic probability on flag-invariant functions pushes forward to
the sum of the fundamental-set measures divided by \(V\).
Conditioning on a nonempty union of components gives the same
formula with its own total flag volume.
\end{proposition}
\begin{proof}
In the marked period description of Lemma~\ref{munf:marks}, the
finite double cosets of \(G=U(h_0)\) represent the marked rows of
this genus. Indeed finite translates of the completed reference
lattice and its basis give a row by intersection with the rational
space. Conversely choose local isometries to a specified row in
the genus. At primes of \(M\) the extension is split and the mark
is a basis on one half with its pairing-dual on the other. An
invertible matrix over the local residue quotient lifts to an
integral invertible matrix by lifting entries. Thus the local
isometries may match the marks and the stabilizer is \(K(M)\).
In terms of the two algebraic Lie ranks \(p,q\), the real flag
space at a place is \(U(p_0,r-p_0)/U(1)^r\), with \(p_0=q\).

Use the rational form \(J=\operatorname{diag}(1_{p_0},-1_{r-p_0})\) given
by Lemma~\ref{cmeig:row}, and compare with \(G_c=U(1_r)\), the
group of the positive definite standard form. At every finite
place the two self-dual lattices are isometric, including residue
characteristic two, by the local part of that lemma. Keep track of
the Haar densities in this assertion. On the Lie algebra for a
Hermitian matrix \(J\), left translation identifies an invariant
top differential by identifying its tangent at one with the fixed
\(F_0\)-space of skew-Hermitian matrices via \(X\mapsto JX\).
Choose the same rational top differential on that vector space
for both \(J\) and \(1_r\). Under an isometry \(P^*JP=1_r\),
the comparison in these coordinates is \(A\mapsto P^*AP\).
Its determinant on the skew-Hermitian space is
\(N_{E/F_0}(\det P)^r\): extend scalars to the split algebra,
where the two factors multiply an unrestricted \(r\)-by-\(r\)
matrix on its left and right. Its absolute value is one at each
finite place because \(N(\det P)=1/\det J=\pm1\).
Thus the corresponding maximal stabilizers have the same local
volume. No choice of their local bases contributes a discriminant.

The global scalar and convergence factors in the Tamagawa measure
are the same: the base field and dimension agree and the geometric
character lattice in both cases is generated by determinant with
the quadratic Galois action of \(E/F_0\). At each real place the
common gauge in the \(A\)-coordinates is the same scalar times
Lebesgue measure on the imaginary diagonal and the complex
strictly lower entries. Multiplication by \(J\) only changes signs
in these coordinates. Divide out that same scalar for both groups.
Use volume one for the maximal stabilizers at finite places and
the resulting real coordinate densities. In these choices the
total adelic covolumes of the two groups are equal: they are
obtained from Tamagawa measure by a common scalar and both
Tamagawa numbers are two for a nonzero Hermitian space over a
quadratic extension. See \href{https://doi.org/10.1007/BF02391774}
{Weil, \emph{Sur la formule de Siegel dans la th\'eorie des groupes
classiques}}, and Ichino, \emph{A regularized Siegel--Weil formula
for unitary groups}, \S4. For the definite group the real volume is
\(c_r^f\). The finite double cosets at maximal level are its
standard lattice genus: cutting the quotient over such a coset
divides this common real volume by its finite automorphism group.
Summing gives \(c_r^f\,{\rm mass}_I\).

For the indefinite group, the quotient of the real density by
angle measure \(d\theta_1\cdots d\theta_r\) is the stated flag
measure. At the central flag the tangent projection takes the
strictly lower entries of the Lie algebra with Jacobian one.
At a lower unipotent matrix \(u=k b\), with \(k\) unitary and
\(b\) upper triangular, the projected tangent in this frame is
the strictly lower part of \(b(u^{-1}du)b^{-1}\).
The unipotent differential has determinant one in its entries.
Conjugation on matrices modulo upper triangular matrices has
complex determinant \(\prod_{i>j} b_{ii}/b_{jj}\). Its squared
absolute value is \(\prod_{i=1}^{r-1}D_i^{-2}\), using
\(|b_{ii}|^2=D_i/D_{i-1}\) and \(D_0=D_r=1\).
This recovers the chart density at every point.

Over each finite double coset \(gK(M)\), the quotient of
\(G_\infty\times gK(M)\) is by
\(\Gamma_g=G(F_0)\cap gK(M)g^{-1}\). Cutting the real factor
by a fundamental set for \(\Gamma_g\) gives the common finite
volume \(1/[K:K(M)]\) times real quotient volume. There is no
stabilizer on the flag: it would be discrete in a compact torus,
hence finite, and the principal subgroup on the underlying
integral period lattice is torsion-free. For the latter, at an
odd prime dividing \(M\) the powers of a nonidentity
\(1+\ell^a B\), \(a\ge1,\ B\not\equiv0\bmod\ell\), have depth
\(a+v_\ell(t)\) at exponent \(t>0\), by the binomial theorem.
Integrating along the angle tori therefore gives \((2\pi)^{fr}\)
times the flag integral. The same disintegration applies to any
nonnegative measurable function of the flag. Applying it first
to one proves the identity and then dividing by the total mass
proves the probability formula. The congruence index is at most
\(M^{2fr^2}\) by reduction on the row, so
Proposition~\ref{bergvol:compactmass} gives the required bound.
\end{proof}
In particular both analytic tests in
Proposition~\ref{bergvol:global} are available on the selected
principal packet for the bounded analytic cusp space of trivial
transvection character. Using this for an algebraic component
uses its identified framed period charts. For the rule space the
reading is as in Corollary~\ref{crev:twist}; agreement with a
separately defined toroidal sheaf is a geometric comparison.

\begin{corollary}[Analytic estimates on a selected good cover]\label{crev:good}
Take a stage in the selected range on the host of
Proposition~\ref{host:art}, a core with the prescribed host,
polarized datum and level, and the deep cover \(Y\) with
\(Y.\mathrm{Good}(e,o)\). Use its identity auxiliary rule
at \(n=8(fr+1)\), and the parameters and semiabelian tests of
Theorem~\ref{bfam:oncharts}. At every good embedding the raw
rules are cusp rules on this \(Y\). The short integral rules in
that theorem span the untwisted space on the cover and on each
component, with logarithmic Petersson norms at most
\(C(frL)^2S+2\log D_M\) for the stated stage size bounds.
For the component cusp spaces at the frame weights of
Corollary~\ref{bergvol:size} the normalized evaluation kernel
is at most \(\beta_*\).
\end{corollary}
\begin{proof}
Here \(Y.\mathrm{Good}(e,o)\) contains the following data. Its
\code{unf} gives the complex framed period uniformization at each
good embedding, including the section transformation law and the
Haar mean-square-one normalization of the theta generator on unit
frames. In these coordinates the first \(p_0=q\) Hermitian
directions are positive. Their conjugate, reversed tangent block
has Lie rank \(q\), and the other tangent block has Lie rank \(p\)
at the labelled \(\tau\). Thus the pivot \(\lambda_{\tau,p}\)
in Theorem~\ref{bfam:oncharts} is numbered in the algebraic
weights, while the positive block used for the majorant ends at
\(p_0\) in the Hermitian flag.

The \code{modelW} and \code{xiAbs} clauses give the signed host
Gram matrix and the absolute polarization coefficient at every
place. The common lattice embedding and alternating pairing in
\code{adPin} identify the original period lattice, its finite
adelic translates, the half-level mark and the eigenframe
coordinates on these same components. In particular stripping the
fixed host ideal from a translate of the reference period
\(\mathfrak a L'_*\) gives the translated Hermitian self-dual row:
finite unitary matrices commute with \(\mathfrak a\), and
localization detects self-duality. This is \(L'\) in the return
count and in Lemma~\ref{bfam:minimum}. The \code{lev} clause gives
auxiliary integer one on the deep variety; the spanning construction
there uses its identity auxiliary map of Corollary~\ref{bfam:mark}.
It does not replace the original map from the deep variety to the
base by an identity.

The \code{dictU} and \code{levelK} clauses identify the rational
and right-principal-level equivalences with the chart action and
exhaust the real flag orbits. Indeed the adelic frame map is
equivariant modulo arithmetic translates under the rational and
level actions and intertwines the real action; it hits each flag
modulo its upper triangular frame change. Equality of its frames
modulo the arithmetic action forces rational and finite-level
equivalence. The unit frames over each flag form a single orbit
under its diagonal angle group. Together with the marked-period
classification in Lemma~\ref{munf:marks}, this is the packet and
the probability disintegration of Proposition~\ref{crev:mass}.
The \code{unit} clause gives leading Gram determinants of absolute
value one at the adelic frames, as used in their metric reading.

Finally \code{chiTriv} gives trivial arithmetic multiplier at the
dominant natural weights on each good chart. It gives the same
for every integral flag weight: for fixed arithmetic element its
multiplier is multiplicative in addition of weights and never
zero. Take natural weights with \(n=0\) that are one on an
initial segment at one labelled embedding and zero elsewhere.
Consecutive differences give all individual weight coordinates,
with reversed order and a minus sign on the dual block. Their
integral span is the whole flag weight lattice. In particular
there is no character loss in the transvection period for the
frame weights. Apply Propositions~\ref{crev:half} and
\ref{crev:mass}, then Proposition~\ref{bergvol:global}.
The actual all-weight space \code{cuspW} is the intersection of
the regular rules with the good-reading tests of
Corollary~\ref{crev:twist}; at a natural weight it is the ordinary
rule space \code{cusp}, and \code{cuspOn} adds only the component
support condition. Thus Theorem~\ref{bfam:oncharts} applies with
its integral tests and no additional analytic cusp assumption.
\end{proof}
\src{src/Development/AbelianVarietyTheory/Definitions/Divisors/NormNegCmConstantEqAtAllEmbeddings.lean}{187}{The clauses of the selected deep cover};
\src{src/Formalism/LocalToGlobal/ProjectiveLattices/FamilyUniformisationBasePointOrthogonalBis.lean}{97}{The marked finite translates in the period identification};
\src{src/Development/AbelianVarietyTheory/Definitions/Foundations/IsProbabilityMeasureMuPart.lean}{425}{The actual all-weight rule space};
\src{src/Development/AbelianVarietyTheory/Counterexamples/LevelGroupExamples/AdelEqZeroOfOffDom.lean}{77}{Generation of the arithmetic characters by natural weights}.

Constructing such covers at all scales is a separate existence
statement. For example, the sized-cover construction starts with
the pinned core, a stage in range, the host splitting bound at least
five and \(15\) coprime to the base level. It constructs the half-level,
uses PEL moduli representability, forms the coefficient compositum
with a ray class field, and builds the marked uniformization; the
field-size bound uses the ray-class-field estimate. If the comparison
to a prescribed base frame requires \(Y.\mathrm{Match}\), that
condition concerns the same base embedding and base frames and is
not part of \(\mathrm{Good}\). The corollary asserts the estimates
on the selected \(Y\); it does not construct the core, the compatible
family or its toroidal compactifications.
\src{src/Development/AbelianVarietyTheory/Proofs/DivisorProofs/ExistsGoodCoverOverSizedField.lean}{111}{Inputs and construction of the sized good level};
\src{src/Development/AbelianVarietyTheory/Definitions/Divisors/NormNegCmConstantEqAtAllEmbeddings.lean}{207}{The common marked frames included in goodness}.


\subsection{Reproduction on the archimedean frame slab}
\label{arrep:section}

Write $\Psi$ for the CM type, $p=p_0$, $q=r-p$, and $G=G_\tau$
for the Hermitian matrix at $\tau$. The pairing is linear in its
first variable, so the Gram matrix of a column frame $x$ is
$x^tG\bar x$. Put $J=\operatorname{diag}(1_p,-1_q)$. The domain
$\mathcal D$ consists of invertible frames at all $\tau$ whose first $p$
columns span a positive space. This is equivalent to the prescribed
signs of all the leading Gram determinants: after the first $p$
columns, the additional orthogonal directions belong to a negative
definite space. Choose an orthoframe $x_0$ with Gram matrix $J$ at
every place. Let $B$ be the product of the invertible upper triangular
groups, and let $dy$ be ordinary Lebesgue measure in the original
complex matrix coordinates. Set
\[
 S_\delta=\{y:\ |(y_\tau^tG_\tau\bar y_\tau-J)_{ij}|<\delta
                         \text{ for every }\tau,i,j\}.
\]

For each place choose nonnegative integers $c_1,\ldots,c_q$,
$c'_1,\ldots,c'_p$ and write $a_i=\sum_{j\ge i}c_j$,
$a'_i=\sum_{j\ge i}c'_j$. In physical column order the character
$\chi_\lambda(b)=\prod_{\tau,i}b_{\tau,ii}^{\lambda_{\tau,i}}$ has weight
\[
 \lambda_\tau=(-n/2-a'_p,\ldots,-n/2-a'_1,\;
                          n/2+a_1,\ldots,n/2+a_q).
\]
The polynomial letters at each place have $n$ columns and two row
blocks. The $q$ rows are the negative rows in their physical order,
and the $p$ rows are the positive rows in reverse physical order.
A row-highest harmonic polynomial has row degrees $a_i,a'_i$, is
invariant under lower unitriangular row substitutions in each letter
block, and is killed by
$\sum_{k=1}^n\partial_{u_{ik}}\partial_{v_{jk}}$ for every pair
of opposite rows at the same place. These conventions apply to
differential symbols as well as ordinary polynomial variables.

\begin{proposition}\label{arrep:reproduce}
Let $n$ be even, $2r\le n$, and let $P_1,\ldots,P_m$ be any finite
family of such polynomials, with at least one nonzero member. Denote by
$K(x,y)$ the result of applying
$\sum_j P_j(\partial_z)\overline{P_j}(\partial_w)$ at $z=w=0$
to the doubled Gaussian below, where the bar conjugates coefficients.
If $0<\delta\le1/(4(r+1))$, then there is a positive real number
$\rho_\delta$ such that every holomorphic $F$ on $\mathcal D$ with
$F(xb)=\chi_\lambda(b)F(x)$ and bounded absolute value on the
orthoframes satisfies the absolutely convergent identity
\[
         \int_{S_\delta} K(x,y)F(y)\,dy=\rho_\delta F(x)
                               \qquad(x\in\mathcal D).
\]
On orthoframes the pointwise automorphic norm in the bounded-family
construction equals $|F|$. The scalar depends on the specified measure,
archimedean data, family and slab. It applies to every deep level
using those same data.
\end{proposition}

\begin{proof}
First compute what survives the insertion into the Gaussian. Suppress
$\tau$, let $P_+,P_-$ be the diagonal projections to the two physical
blocks, and set $S=-J$, $\kappa=\tau(\kappa_0)$,
$h=\kappa/4$, $t=\pi i\kappa=-\pi\operatorname{Im}\kappa>0$. Put
\[
\begin{aligned}
 L&=P_+\bar y^tG^t+P_-x^{-1},&
 H&=P_+x^tG+P_-\bar y^{-1},\\
 W&=(L^{-1}SL)^tG,& Y&=(hW-(hW)^*)/(2i).
\end{aligned}
\]
The kernel of the bottom rows of $L$ is the positive space spanned
by the first $p$ columns of $x$. The kernel of its top rows is
the negative orthogonal complement of the first $p$ columns of $y$.
They intersect trivially. Interchanging $x,y$ and conjugating gives
invertibility of $H$. The involution $L^{-1}SL$ acts by $-1$ on the
first space and $+1$ on the second. For $v=v_++v_-$ in this
decomposition, the real part of $v^tW\bar v$ is
$-\langle v_+,v_+\rangle+\langle v_-,v_-\rangle$ because the
cross terms have zero real part. Thus $W+W^*$ is negative definite.
As $h/(2i)<0$, the Hermitian matrix $Y=(h/(2i))(W+W^*)$ is positive.

Write $M=L^{-t}GH^{-1}$ and $N=L^{-t}GY^{-1}GH^{-1}$.
For each column $k$ put $\alpha=P_+z_k+P_-w_k$,
$\beta=P_+w_k+P_-z_k$. The actual undifferentiated Gaussian has
prefactor $\prod_\tau(\det L\,\det H)^{-n/2}$ and exponent
$-4\pi\sum_{\tau,k}h_\tau^2\beta^t N\alpha$.
The off-diagonal blocks of $M^{-1}=HG^{-1}L^t$ vanish because
the corresponding products before taking blocks are identity
matrices. Its diagonal blocks are the leading $p$ block of
$x^tG\bar y$ and the trailing $q$ block of $\bar y^{-1}G^{-1}x^{-t}$.
In particular $SM=MS$. Moreover
\[
                     \frac{h}{2i}(SN+NS)=M .
\]
Indeed $W G^{-1}W=G$ and $W^*G^{-1}W^*=G$. For $T=W+W^*$
these give $WG^{-1}T=TG^{-1}W^*$ and therefore
$GT^{-1}W=W^*T^{-1}G$. It follows that
$WT^{-1}G+GT^{-1}W=G$. Multiply this last identity on the left
by $L^{-t}$ and on the right by $H^{-1}$,
using $L^{-t}W=SL^{-t}G$ and $WH^{-1}=GH^{-1}S$.
The latter equality follows from $SM=MS$.

The off-diagonal terms in $\sum_k\beta^t N\alpha$ are linear
combinations of $\sum_kz_{ik}z_{jk}$ and $\sum_kw_{ik}w_{jk}$
for opposite rows $i,j$. Pairing with a polynomial differential
symbol at zero makes multiplication by each such quadratic adjoint
to its mixed Laplacian. Every term containing one of these quadratics
vanishes, also after both sets of derivatives are taken. Only finitely
many degrees of the exponential enter the pairing. Thus one may
differentiate the following explicit expression instead:
\[
 \prod_\tau (\det L\,\det H)^{-n/2}
 \exp\!\left(\sum_{\tau,k}t_\tau
      [\,w_{+,k}^t M_{++}z_{+,k}-z_{-,k}^t M_{--}w_{-,k}\,]\right).
\]
It is holomorphic in $x$ and antiholomorphic in $y$.
The two diagonal blocks of $M(y,x)$ are the adjoints of those
of $M(x,y)$, and the prefactors interchange by complex conjugation.
Consequently $K(y,x)=\overline{K(x,y)}$.

Let $\mathcal U=\prod_\tau\{u:u^tG_\tau\bar u=G_\tau\}$. Under
simultaneous left translation by $u$, $L,H$ acquire right factors
$u^{-1},\bar u^{-1}$, so $M$ and the prefactor are invariant.
Under $x\mapsto xb$, write the two diagonal blocks of $b$ as $A,D$.
Then $L$ acquires left factor $1\oplus D^{-1}$ and $H$ left factor
$A^t\oplus1$. The prefactor gains $(\det D/\det A)^{n/2}$ at this
place. In the exponent the substitutions in $z$ are $A^{-t}$ on
the positive rows and $D$ on the negative rows. In the differential
symbol the resulting row substitutions are $A^{-1}$ and $D^t$.
Both are lower triangular in letter order. The prescribed row
degrees give exactly the remaining factor of $\chi_\lambda(b)$.
This proves equivariance in $x$, and symmetry gives the conjugate
character in $y$.

At $x=y=x_0$ one has $L=x_0^{-1}$, $H=\bar x_0^{-1}$ and $M=J$.
The exponential becomes $\exp(\sum_{\tau,i,k}t_\tau z_{\tau ik}
w_{\tau ik})$. With $P_{j,\nu}$ the monomial coefficients, ordinary
differentiation gives
\[
 C_0=K(x_0,x_0)=
   \left(\prod_\tau|\det x_{0,\tau}|^n\right)
   \sum_{j,\nu}|P_{j,\nu}|^2\,\nu!\,
                       \prod_\tau t_\tau^{|\nu_\tau|}>0 .
\]
Here $|\nu_\tau|$ is the sum of all exponents at that place and
$\nu!$ is the product of their factorials.

We next identify the only holomorphic function that can survive
the compact torus average. In coordinates $g=x_0^{-1}y$, let
$\Delta_{\tau,i}$ be the leading $i$-row principal minor of $g_\tau$.
Define
\[
 H_\lambda(y)=\prod_\tau
    \Delta_{\tau,r}^{\,\lambda_{\tau,r}}
    \prod_{i=1}^{r-1}\Delta_{\tau,i}^{
                              \lambda_{\tau,i}-\lambda_{\tau,i+1}} .
\]
For $r=0$ take the empty product. The weight decreases within each
block. When $0<p<r$ the only internal negative exponent occurs at
$i=p$, where the minor cannot vanish: the projection of a positive
$p$-plane to the standard positive coordinate space has trivial
kernel. Hence $H_\lambda$ is holomorphic throughout $\mathcal D$ and
has value $1$ at $x_0$. Under left diagonal unitary translation in
the $g$ coordinates it has character $\chi_\lambda$, and under right
triangular multiplication it has the same character.

These two transformations determine a holomorphic function up to its
value at $x_0$. Indeed on the big cell near $g=1$ write $g=\ell b$
with $\ell$ lower unitriangular. The two character laws imply
invariance of the restriction to $\ell$ under diagonal unitary
conjugation. On every strictly lower entry $\ell_{ij}$ the simultaneous
circle $\operatorname{diag}(e^{is},e^{2is},\ldots,e^{ris})$ acts with
positive weight $i-j$. The holomorphic power series of the restriction
therefore has only its constant term. Right equivariance determines
the function on a neighborhood, and the identity theorem extends the
conclusion. Indeed indefinite Gram--Schmidt identifies the domain with
the product of the orthoframes and the upper triangular matrices with
positive real diagonal. Both are connected: the unitary group modulo
$U(p)\times U(q)$ is the matrix ball described below, with connected
compact fiber. It follows that
\[
               K(y,x_0)=C_0H_\lambda(y).
\]
For this application the left torus law follows from simultaneous
invariance and right equivariance of $K$ in its second variable.

We check absolute integrability on the slab. The preceding
Gram--Schmidt coordinates write $y=x_0 g b$ with $g$ standard
$J$-unitary and $b$ upper triangular with positive real diagonal.
For our $\delta$, each leading Gram matrix has operator distance
at most $r\delta<1/4$ from the corresponding signed identity. Thus
even the closure of the permitted Gram matrices is compact with all
pivots nonzero and of the required signs. The signed triangular
factorization $b^tJ\bar b$ has a continuous inverse there by the
Gram--Schmidt recursion. The permitted $b$ therefore have compact
closure inside the triangular group. In these smooth coordinates
Lebesgue measure is $dg\,j(b)\,db$, where $dg$ is Haar measure and
$j$ is smooth positive: the original measure is left invariant
under $\mathcal U$ since $|\det u_\tau|=1$. Both $|\chi_\lambda(b)|$
and $j(b)$ are bounded on this compact closure. In particular
$F$ is bounded throughout $S_\delta$.

At a place with $pq>0$ write the standard unitary matrix
$g=\left(\begin{smallmatrix}A&B\\ C&D\end{smallmatrix}\right)$.
All singular values of $A$ and $D$ are at least one by the unitary
relations. The complementary-minor identity and $g^{-1}=Jg^*J$
give $|\det A|=|\det D|$ and bound each leading minor of size
$i>p$ by $|\det D|$ using the complementary trailing minor of
$g^{-1}$. For $i\le p$ the exterior powers of $A$ bound it by
$|\det A|$. As $|\det g|=1$, the product of the minor estimates,
keeping the exact denominator at $i=p$, gives the one-place bound
\[
       |H_{\lambda,\tau}(x_{0,\tau}g)|
          \le |\det A|^{\lambda_{\tau,1}-\lambda_{\tau,r}}
          \le |\det A|^{-n}\le1 .
\]
If either block is empty the group is compact and the function bounded.
For the mixed signature the group modulo the compact block diagonal
subgroup is the ball $Z=C A^{-1}$, $Z^*Z<I_p$, the graph coordinates
of its positive planes. The group is transitive by orthogonal
completion, and $|\det A|=\det(I_p-Z^*Z)^{-1/2}$.
The invariant density on the ball is proportional to
$\det(I_p-Z^*Z)^{-r}\,dZ$. To see this, for an acting matrix with
blocks $a,b,c,d$, set $U=a+bZ$ and $Z'=(c+dZ)U^{-1}$.
Then $I_p-Z'^*Z'=U^{-*}(I_p-Z^*Z)U^{-1}$ and the differential is
$H\mapsto(d-Z'b)HU^{-1}$. Block elimination gives
$|\det U|\,|\det(d-Z'b)|=|\det g_{\rm acting}|=1$, so the real
Jacobian is $|\det U|^{-2r}$. Hence integration of $|\det A|^{-n}$
on the group reduces, up to a finite constant, to integration of
$\det(I_p-Z^*Z)^{n/2-r}$ on a bounded ball. Its exponent is
nonnegative. This proves that $H_\lambda$ is bounded and absolutely
integrable on $S_\delta$, including at $n=2r$. In particular both
$K(x_0,y)F(y)$ and $|K(x_0,y)|^2$ are integrable there.

Finally let $t$ range over the product diagonal unitary torus with
probability Haar measure. Its action on frames is
$t\cdot y=x_0 t x_0^{-1}y$. The holomorphic projection
$F_\lambda(y)=\int\chi_\lambda(t)^{-1}F(t\cdot y)\,dt$ has the two
character laws above and has value $F(x_0)$ at $x_0$. Thus
$F_\lambda=F(x_0)H_\lambda$. The slab and its measure are invariant
under this torus. By absolute convergence and the conjugate torus
law of $K(x_0,y)$, replacing $F$ with $F_\lambda$ in the integral
at $x_0$ leaves its value unchanged. The resulting constant is
\[
                \rho_\delta=\frac{1}{C_0}
                   \int_{S_\delta}|K(x_0,y)|^2\,dy>0 .
\]
Positivity follows from continuity, $K(x_0,x_0)=C_0>0$, and the
open neighborhood of $x_0$ in the slab (with mass one in dimension
zero). Left translation carries the identity and absolute convergence
to every orthoframe by applying it to the translated bounded
holomorphic function. Right equivariance then carries it to every
frame of $\mathcal D$.
\end{proof}

The required invariant finite family can be chosen in these same
variables. At each place form the product of the initial $i$-row
minors to powers $c_i$ in the negative letters and to powers $c'_i$
in the positive letters, using respectively the first $q$ and the
last $p$ columns, nested within each set. Since $n\ge p+q$, these
sets are disjoint. The product over places is a nonzero row-highest
polynomial and all mixed Laplacians kill it. The entire space of
the stated harmonics lies in total degree
$d=\sum_{\tau,i}i(c_{\tau,i}+c'_{\tau,i})$, extending the arrays by
zero. Its dimension is at most $\binom{nr|\Psi|+d}{d}$, the number
of monomials of degree at most $d$. The compact column groups $U(n)^\Psi$ act on opposite letters
by a unitary matrix and its contragredient, preserving the row
conditions and the contractions. An invariant Hermitian product
exists by averaging. An orthonormal basis $P_1,\ldots,P_m$ then has
invariant sum $\sum_j P_j(z)\overline{P_j(w)}$ and positive $C_0$.
Thus the insertion used for spanning bounded forms satisfies the
proposition with its actual weight and with finite $m$.

\src{src/Development/AbelianVarietyTheory/Definitions/SpanningFamilies/IsSystem.lean}{65}{The doubled Gaussian before the two polynomial insertions};
\src{src/Development/AbelianVarietyTheory/Proofs/SiegelThetaTransformation/SetIntegralCompOfStable.lean}{268}{The positive reproducing scalar on the full frame slab}; and
\src{src/Development/AbelianVarietyTheory/Proofs/SiegelThetaTransformation/HermitianThetaSeriesSpanBoundedForms.lean}{168}{Spanning bounded families in the range $4r<n$}.

\subsection{The graph terms of the doubled marked family}
\label{dg:section}

We compute the graph functional in the rational symplectic basis in
which the undoubled sections are defined. Write $g=r|\Psi|$,
$\Omega_x$ for the Siegel period matrix, $Y_x=\operatorname{Im}\Omega_x$,
and $h_x$ for the inverse of the first period matrix. Thus a real
period vector with symplectic coordinates $(u,v)$ has normalized
complex period $u+\Omega_x v$. Use Lebesgue measure on the $ng$ real
position coordinates and mass one for their integral finite lattice.
For one row the Gaussian, apart from the determinant, is
\[
 q_{x,z}(s)=\exp\!\left(\pi i\,s^t\Omega_xs+
          2\pi i\,s^t h_xz+\tfrac{\pi}{2}
                   z^t h_x^tY_x^{-1}h_xz\right).
\]
Put $Q_{x,z}=(\det h_x)^{n/2}\prod_{k=1}^nq_{x,z_k}(s_k)$
and $Q_{x,P}=P(\partial_z)Q_{x,z}|_{z=0}$. These give exactly
the rational theta sums defining $\theta_P(\phi;x)$ with finite
test $\phi$. We use inner products linear in their first variable.

\begin{lemma}\label{dg:gaussian}
For the family and weight of Proposition~\ref{arrep:reproduce},
\[
 \sum_j\langle Q_{x,P_j},Q_{y,P_j}\rangle_\infty
 =c_\infty K(x,y),\qquad
 c_\infty=2^{-ng/2}\prod_{\tau\in\Psi}
     \bigl(|\det G_\tau|\,[-\operatorname{Im}\tau(\kappa_0)]^r
                                                        \bigr)^{n/2}.
\]
This factor is independent of the weight, the polynomials, and both
frames.
\end{lemma}
\begin{proof}
First we show that $Q_{x,P}$ is holomorphic in $x$ with values in
Schwartz space. The physical period map at $\tau$ has components
$\langle x_i,s\rangle$ for $i\le p$ and $(x^{-1}s)_i$ for $i>p$.
Thus $h_x,\Omega_x$ are holomorphic. For the remaining quadratic put
$C=h_x^tY_x^{-1}h_x$. Choose a purely imaginary $\theta\in E$
which generates $E/\mathbf Q$. Such elements exist: writing the
purely imaginary subspace as $\eta E^+$ for any nonzero imaginary
$\eta$, distinct embeddings have distinct restrictions to this
subspace (evaluate first on $\eta$, then on its multiples).
Finitely many proper rational linear subspaces therefore suffice to
exclude all coincidences. Multiplication by $\theta$ in physical
period coordinates is diagonal $D$, with entries $\bar\tau(\theta)$
on the positive rows and $\tau(\theta)$ on the negative rows.
Write its rational blocks in the symplectic basis as
$\left(\begin{smallmatrix}a&b\\ c&d\end{smallmatrix}\right)$.
Applying it to first periods gives $h_xD=(a+\Omega_xc)h_x$.
Multiplication by $\theta$ is skew-adjoint for the alternating trace
pairing and complex-linear on the normalized periods. Hence it is
skew-adjoint for their Hermitian form with matrix $Y_x^{-1}$: the
imaginary part is the alternating pairing, and replacing one argument
by $i$ times that argument recovers the real part. With $A=a+\Omega_xc$,
this gives $A^tY_x^{-1}+Y_x^{-1}\bar A=0$. Consequently
\[
                  D^t C+C D=2i h_x^t c h_x .
\]
Whenever $D_i+D_j\ne0$, $C_{ij}$ is holomorphic by division.
The remaining pairs are exactly the opposite signs at the same
place. Their terms in the product exponential multiply the other
terms by exponentials of mixed row contractions. They disappear
under $P(\partial_z)$ by the mixed Laplacian, as in
Proposition~\ref{arrep:reproduce}. On each compact subset of
$\mathcal D$, $Y_x$ is bounded below by a positive matrix, so the
resulting polynomial Gaussians and their derivatives have locally
uniform Schwartz bounds. This proves holomorphy and hence
sesquiholomorphy of the displayed pairing.

On replacing $x$ by $xb$, with triangular diagonal blocks $A_0,D_0$,
the physical period map acquires left factor $A_0^t\oplus D_0^{-1}$.
Thus $\Omega_x$ is unchanged, and $h_x$ gains right factor
$A_0^{-t}\oplus D_0$. Ordinary differentiation and the row-highest
law give $Q_{xb,P}=\chi_\lambda(b)Q_{x,P}$. It suffices to compare
the two kernels on the diagonal of orthoframes and then on the full
diagonal by triangular equivariance.

For an orthoframe $x$, put $d_\tau=-\operatorname{Im}\tau(\kappa_0)$.
In its basis the physical periods of a vector $s$ are the conjugates
of its positive coefficients and its negative coefficients themselves.
The imaginary part of the Hermitian pairing with diagonal entries
$d_\tau$ in these coordinates is therefore the original alternating
form $\operatorname{Tr}_{E/\mathbf Q}(-\kappa_0\langle s,s'\rangle/2)$.
In normalized period coordinates the same form is
$\operatorname{Im}((u+\Omega_x v)^tY_x^{-1}
\overline{(u'+\Omega_x v')})=v^tu'-u^tv'$. Equality of imaginary
parts gives
$h_x^tY_x^{-1}\bar h_x=\operatorname{diag}(d_\tau 1_r)_\tau$.
Completing the real Gaussian square now gives
\[
 \langle Q_{x,z},Q_{x,\bar w}\rangle_\infty
 = |\det h_x|^n\det(2Y_x)^{-n/2}
          \exp\!\left(\pi\sum_{\tau,i,k} d_\tau z_{\tau ik}w_{\tau ik}\right).
\]
This is the orthoframe Gaussian of Proposition~\ref{arrep:reproduce}
times $c_\infty$, since $|\det x_\tau|^2=|\det G_\tau|^{-1}$.
Differentiate and use triangular equivariance. Sesquiholomorphic
kernels equal on the full diagonal are equal everywhere: the power
series of their difference at each diagonal point has coefficients
detected by independent Wirtinger derivatives on the diagonal,
and $\mathcal D$ is connected.
\end{proof}

\begin{proposition}\label{dg:identity}
Use the marked family in Lemmas~\ref{bfam:genus} and~\ref{bfam:phase},
even $n>4r$, and the column-invariant polynomial family of
Proposition~\ref{arrep:reproduce}. Let $\sum_\Lambda w_\Lambda$
denote the finite genus sum with reciprocal stabilizer weights.
Write $\phi_{L,\Lambda}$ for the marked vector on row $L$,
$\Delta$ for its prescribed rational row group, and
$w(\gamma)\phi_{L,\Lambda}=\mu(\gamma,L)\phi_{\gamma L,\Lambda}$.
For all permissible rows $L_1,L_2$ and all $x,y\in\mathcal D$,
\[
\begin{split}
 &\sum_\Lambda w_\Lambda\sum_j
    \theta_{P_j}(\phi_{L_1,\Lambda};x)
           \overline{\theta_{P_j}(\phi_{L_2,\Lambda};y)}\\
 &\quad=c_\infty\sum_{\gamma\in\Delta}\overline{\mu(\gamma,L_2)}
               K(x,\gamma y)
       \sum_\Lambda w_\Lambda
       \langle\phi_{L_1,\Lambda},\phi_{\gamma L_2,\Lambda}\rangle_f .
\end{split}
\]
The graph sum is absolutely convergent. This is the kernel identity
for this genus, including the phase of each marked row.
\end{proposition}
\begin{proof}
In the doubled oscillator use the rational diagonal polarization.
For a pure tensor $u\otimes\bar v$ the change to this polarization
takes difference and midpoint of the two position coordinates,
followed by Fourier transform in the midpoint. Evaluation at the
origin is therefore $\langle u,v\rangle$ in the self-dual position
measure, both at infinity and at the finite places. In the source
orientation the coordinate function is $W_{u,v}(-s)$ as in
Lemma~\ref{bfam:phase}. The product of the two rational theta sums
is the rational theta sum after this change: the rational pairs of
positions are in bijection with their rational differences and
midpoints, and Poisson summation in the latter variable replaces
it by its rational dual. This proof works for arbitrary Schwartz
inputs. In particular there is no primitivity restriction on the
rational differences. A common rescaling of the finite intertwiner
and the inverse rescaling of the archimedean intertwiner leave the
global evaluations unchanged.

Here it matters that the inserted archimedean test
$\sum_j Q_{x,P_j}\otimes\overline{Q_{y,P_j}}$ is invariant under the
compact column group in the doubled model. In physical period
coordinates a column unitary $u_\tau$ acts by $u_\tau$ in negative
rows and by $\bar u_\tau$ in positive rows. It preserves the
polarization, so its oscillator action fixes the Gaussian vacuum
line at each frame. One can see this by the first-order Gaussian
equations $(\partial_s-2\pi i\Omega_xs)q_{x,0}=0$: the operators
are the positive-polarization annihilators and their common solution
line is spanned by this Gaussian. Its unit multiplier is the same
at any two frames. Indeed their vacuum pairing is nonzero by the
Gaussian integral with matrix $-i(\Omega_x-\bar\Omega_y)$ of
positive real part. To check the action on the inserted vectors,
the Weyl translate of the vacuum with shift $v$ and modulation $u$
is $q_{x,z}\exp(-\pi\zeta^tY_x^{-1}\bar\zeta/2)$, where
$\zeta=u+\Omega_xv=h_xz$. This follows by replacing $s$ with $s+v$
and multiplying by $\exp(2\pi i u^t(s+v/2))$. The column action
preserves the Hermitian norm in this formula. Thus it transforms
the coherent parameters by the stated linear substitutions and the
common unit multiplier. The multipliers cancel between the two
copies, and column invariance of the sum of differential symbols
proves the assertion (the transposes in differentiation are again
paired unitary substitutions).

Apply Lemma~\ref{bfam:genus} to this invariant Schwartz test.
At the Siegel--Weil point the section for a Schwartz function $\Phi$
is $g\mapsto(\omega(g)\Phi)(0)$. Here the doubled group has split
rank $r$, the positive definite column space has dimension $n$ and
the parameter is $s_0=(n-r)/2$. The theta integral uses Haar
probability, giving coefficient one. Multiplying back by the
total mass gives the same weights on both sides.
The splitting character on the Siegel Levi can be taken trivial:
its restriction to $\mathbb A_{E^+}^\times$ is required to be the
$n$-th power of the quadratic extension character, and $n$ is even.
In this splitting the diagonal column group acts by substitution.
The same reference gives the doubled tensor intertwiner and its
value at zero; its general second-factor twist by the splitting
character is trivial here.
\cite[\S\S2.8, 3.1, 11.3 and Theorem~7.1]{ganQiuTakeda}

We specify the graph convention. Label the graph by the rational
row isometry $\gamma$ acting on the second member of the tensor to
give the diagonal functional. Thus for the seed graph one evaluates
at zero, and for this graph one first translates that member by
$\gamma$. In row coordinates a representative is the conjugate of
$\operatorname{diag}(1,\gamma^{-t})$ by the rational change of
polarization with blocks
$\left(\begin{smallmatrix}1&\kappa_0G/4\\-1&\kappa_0G/4\end{smallmatrix}\right)$.
Equivalently the action on the second undoubled theta test is the
rational oscillator action carrying its frame $y$ to $\gamma y$.
To check the scalar, let
$\left(\begin{smallmatrix}a&b\\c&d\end{smallmatrix}\right)$ be
the induced Siegel transformation in the symplectic period basis,
$U=c\Omega_y+d$, $\Omega'=(a\Omega_y+b)U^{-1}$.
The equality of the period maps on $\gamma s$ at $\gamma y$ and on
$s$ at $y$ gives $h_{\gamma y}=U^{-t}h_y$,
$\Omega_{\gamma y}=\Omega'$. The Gaussian action sends
$Q_{y,z}$ to $Q_{\gamma y,z}$ up to a constant phase independent
of $y,z$. Indeed under a real linear coordinate change this follows
by substitution, and under adding a real symmetric matrix to the
period it follows by multiplying by the corresponding quadratic
phase. Under Fourier inversion use the Gaussian integral
\[
 \int_{\mathbb R^g}e^{\pi i s^t\Omega s+2\pi i s^t t}\,ds
       =\det(-i\Omega)^{-1/2}e^{-\pi i t^t\Omega^{-1}t}.
\]
For the Fourier generator $\Omega'=-\Omega^{-1}$, the factor
$e^{-\pi i\zeta^t\Omega^{-1}\zeta}$ combines with the original
quadratic correction to give the one at
$\zeta'=\Omega^{-t}\zeta$, since
$\zeta'^tY'^{-1}\zeta'=\zeta^tY^{-1}\zeta-
2i\zeta^t\Omega^{-1}\zeta$. The determinant in $Q$ absorbs the
period-dependent determinant after taking $n$ rows. These generators
generate the rational symplectic group by symplectic elimination.
Choose the reciprocal scalar on the finite operator. Its action
$w(\gamma)$ is then exactly the one characterized by
$\theta(w(\gamma)\phi;\gamma y,z)=\theta(\phi;y,z)$: for the
Fourier generator this is rational Poisson summation, for rational
linear changes it is reindexing with the product formula, and for
the shear it is cancellation of the rational quadratic phases.
This characterization determines the finite operator, since at
fixed $y$ the distinct rational position frequencies of the theta
series can be separated by Fourier coefficients in $h_yz$ after
removing the common nonzero Gaussian prefactor.

These are also the normalizations of the graph section. Indeed the
diagonal-polarization change intertwines the rational row actions
with their conjugates in the doubled group. For the linear changes,
quadratic phases and Fourier generators this is just covariance
of Weyl translations under these operations. Any scalar ambiguity
between the two intertwiners is independent of the test by
irreducibility of the Schrödinger representation; its global value
is one since both rational actions preserve the nonzero rational
theta distribution by Poisson summation. The choice of another
rational Siegel-parabolic representative likewise preserves the
global evaluation, by the product formula on its determinant and
triviality of its unipotent character at zero. The term on the
graph consequently is precisely
\[
 \langle\phi_{L_1,\Lambda},w(\gamma)\phi_{L_2,\Lambda}\rangle_f\,
          \sum_j\langle Q_{x,P_j},Q_{\gamma y,P_j}\rangle_\infty .
\]
Graphs outside $\Delta$ vanish by the marked support condition.
For $\gamma\in\Delta$ the inner product in the second variable
contributes $\overline{\mu(\gamma,L_2)}$.

For a nongraph there is a nonzero rational isotropic row line
$E\xi$ in the intersection of its maximal isotropic with the first
copy. Indeed either projection being singular forces the other to
be singular by orthogonality and dimension. For every $v\in E^n$
the finite Weyl translation on the first input by $v\otimes\xi$ acts,
after passage to this polarization, by a pure modulation. It takes
value one at the origin: in difference and dual coordinates the
action is $\phi(t-d)\psi_f(\langle b,t-d/2\rangle)$ with $d=0$.
There is no residual scalar, since linear symplectic covariance
preserves the symmetric Weyl operators including their half
pairings. The resulting finite evaluation $\ell$ thus satisfies
$\ell(\rho(v\otimes\xi)\phi_1,\phi_2)=\ell(\phi_1,\phi_2)$.
The marked isotropic relation has finite coefficients $a_v$ with
$\sum a_v\ne0$ and $\sum a_v\rho(v\otimes\xi)\phi_1=0$.
It gives $(\sum a_v)\ell(\phi_1,\phi_2)=0$ before any
archimedean evaluation. This accounts for every nongraph.

Finally the polynomial differentiations are legitimate in this
expansion. On compact sets of complex $z,w$ the archimedean
Gaussians, with every fixed parameter derivative, are bounded in
Schwartz space, by completion of the square and $Y_x,Y_y>0$.
The same holds after the fixed partial Fourier transform.
For a Siegel representative with Iwasawa decomposition $pk$ and
Levi block $a_p$, evaluation is bounded by
$|\det a_p|_{\mathbb A_E}^{n/2}
\sup_k|(\omega(k)\Phi)(0)|$; the supremum is uniform on these
bounded test families by continuity and compactness. The finite
tests lie in finitely many compact orbits under the maximal compact
subgroup (each smooth test has open stabilizer there).
The scalar Siegel Eisenstein majorant converges absolutely for
$\operatorname{Re}s>r/2$, which contains $s_0$ since $n>4r$.
Thus either Cauchy's formula on a slightly larger parameter
polydisc or direct Schwartz differentiation permits the fixed
finite insertions under the sum. The genus formula above was only
needed for the invariant sum of inserted vectors. Gaussian theta
series on the left also converge normally since the finite test
restricts rational positions to a fixed fractional lattice.
Combining the surviving graph evaluations with Lemma~\ref{dg:gaussian}
proves the identity.
\end{proof}

\src{src/Development/AbelianVarietyTheory/Proofs/SiegelThetaTransformation/EisensteinSeriesTermUnitaryGraph.lean}{33}{The graph representative and both translated tests};
\src{src/Development/AbelianVarietyTheory/Proofs/SiegelThetaTransformation/KernelIdentity.lean}{181}{The class-weighted differentiated graph identity}.

\subsection{Marked period components and absolute unfolding}

Fix a rational Hermitian space \(W\) over \(E/F\), \(F=E^+\), its
polarizing trace form and CM type. Take the identity auxiliary cover
and its principal odd level \(M\), with
\((M)=\mathfrak m\bar{\mathfrak m}\) and
\(\mathfrak m+\bar{\mathfrak m}=O_E\).
In the period description a component has a self-dual row lattice
\(L\subset W\), the polarized half-level mark and the connected flag
domain of its prescribed signature. A marking on one half of \(L/ML\)
can equally be specified by its dual basis on the opposite half,
using the fixed primitive \(M\)-th root for the Weil pairing.

\begin{lemma}\label{munf:marks}
Fix a reference row \(L_*\) with a basis of \(L_*/\mathfrak mL_*\).
Every such marked period lattice in \(W\) can be carried by a rational
isometry to a row equal to \(L_*\) at the primes of \(M\), with the
specified basis on the \(\mathfrak m\)-half. Write \(\Delta\) for the
rational unitary elements integral there and acting identically on this
marked half of \(L_*\). Two transported rows lie in the same
\(\Delta\)-orbit exactly when their polarized marked period lattices
are isometric. The stabilizer of each row is its principal-level
arithmetic group.
\end{lemma}
\begin{proof}
Every place \(v\) of \(F\) above the level splits in \(E\), since a
prime fixed by conjugation could not divide the product of the
coprime conjugate halves. The polarization pairs the two summands of
\(L/ML\) perfectly: the CRT idempotents are interchanged by the
adjoint involution, so each summand is isotropic and is the dual of the
other. This proves the equivalence of the two ways of specifying
the marking, including at a prime power. Locally \(U(W)(F_v)\) is
\(\mathrm{GL}_r(F_v)\) on the first factor, acting by the inverse
adjoint on the second. Lift the indicated residue bases in both
first-factor lattices to bases over the DVR. The matrix carrying one
basis to the other carries the lattices onto one another. Its inverse
adjoint carries the second lattices and the dual marks, by perfection.

The resulting conditions on the tuple of local isometries form a
nonempty open set. Rational isometries meet it. Indeed the Cayley
map \(X\mapsto(1-X)(1+X)^{-1}\) identifies an open of the
\(F\)-vector space of skew-adjoint endomorphisms with the unitary
open where \(1+g\) is invertible. Every selected local open meets
this chart: multiply its matrix by a norm-one scalar arbitrarily
close to one, avoiding the finitely many prohibited eigenvalues.
At these split places such scalars are \((t,t^{-1})\).
Weak approximation on the vector space, in a fixed \(F\)-basis,
now gives a single rational isometry. The transported lattices agree
with \(L_*\) at the whole level. Globally, this is equivalent to the
existence of a common denominator prime to the level in the two
lattice inclusions: localize the finite torsion modules
\((L+L_*)/L\) and \((L+L_*)/L_*\), then choose an annihilator
congruent to one at those primes by CRT.

An element of \(\Delta\) carrying one transported lattice to another
preserves the mark on the first half, hence on the dual half.
Conversely a marked isometry between the lattices extends uniquely
by \(E\)-linearity to \(W\); its local action belongs to \(\Delta\).
When the two rows agree this gives precisely the integral isometries
acting identically modulo \(M\). To see the component statement in
period coordinates, a complex framed polarized torus has a period
lattice on its tangent space. An isomorphism lifts uniquely to its
complex-linear derivative on the universal covers, so it induces
the preceding integral isometry. Conversely transporting the
complex structure and flag by an isometry gives isomorphic marked
complex tori. For a fixed marked lattice the positive-plane domains
are matrix balls with connected compact flag fibres; their product
is connected. In a family of complex abelian varieties the integral
homology, its trace pairing and the finite marking are locally
constant in a topological trivialization. This gives exactly the
isometry classes as component indices in the marked period model.
\end{proof}
\src{src/Development/AbelianVarietyTheory/Proofs/SiegelThetaTransformation/ExistsMoveMarked.lean}{133}{Rational transport of the marked basis},
\src{src/Development/AbelianVarietyTheory/Proofs/SiegelThetaTransformation/ComponentEqOfMatchedLatticesMarkings.lean}{74}{identification of marked indices} and
\src{src/Development/AbelianVarietyTheory/Proofs/SiegelThetaTransformation/ExistsComponentsLatticesMovedToLevelFamily.lean}{162}{distinct orbits after simultaneous transport}.

\begin{proposition}[Unfolding the finite span]\label{munf:span}
On the countable set of admissible rows let \(C\) be the equivariant
matrix of Lemma~\ref{bfam:gram} and \(K\) the kernel of
Proposition~\ref{arrep:reproduce}. Suppose the marked periodized
kernel is the finite sum of paired theta series for the indicated
column genus, with a single nonzero normalization. Then every
bounded family of holomorphic forms of the prescribed weight and
unitary multiplier on these row orbits lies in the span of those
theta series. The conclusion also holds simultaneously on any finite
collection of distinct row orbits.
\end{proposition}
\begin{proof}
Write the input as \(F\) and put \(V=C^{-1}F\). Use the norm given
by the supremum over all rows and orthoframes. The Neumann series
converges in this norm and preserves the weight and multiplier,
since \(C\) commutes with the action. It converges locally uniformly
on the full domain by triangular equivariance and Gram--Schmidt,
so \(V\) is holomorphic there. Its absolute value is bounded on all
rows of the slab by one constant \(B\), by the compactness of the
triangular factors in the proof of Proposition~\ref{arrep:reproduce}.
For a fixed output \(a=(L,x)\) put
\[
          f_a(L',y)=C_{LL'}K(x,y)V_{L'}(y).
\]
On rows times the slab we have
\[
 \sum_{L'}\int_{S_\delta}|f_a(L',y)|\,dy
       \le B\sum_{L'}|C_{LL'}|\int_{S_\delta}|K(x,y)|\,dy<\infty .
\]
Thus we may unfold under the countable diagonal action on this
product with its counting and Lebesgue measures. The action
preserves both measures.

Here is an explicit reason a measurable fundamental set exists.
Choose one \(L'\) in each row orbit. Its stabilizer is discrete in
the real unitary group: on an integral basis of the underlying
rank-\(2fr\) lattice it has integral matrix entries. It acts freely
on the invertible frames and properly, since the isometries carrying
one compact set of frames into another form a subset with compact
closure (write \(g=y'y^{-1}\)). Cover the slab by countably many
open patches disjoint from their nonidentity stabilizer translates.
Keep in each patch the complement of the translates of all earlier
patches. Their union meets every stabilizer orbit once and is Borel.
The union on the representative rows is a diagonal fundamental set
\(\mathcal D_0\). Absolute integrability and Tonelli now give
\[
 \sum_{L'}\int_{S_\delta}f_a(L',y)\,dy
       =\int_{\mathcal D_0}\sum_{\gamma\in\Delta}
                           f_a(\gamma L',\gamma y)\,d(L',y).
\]
The left side is \(\rho_\delta(CV)_L(x)=\rho_\delta F_L(x)\).
On the right the multiplier law for \(V\) and the given periodized
identity express the inner sum as \(\sum_{j=1}^s q_j(a)h_j(L',y)\),
where the \(q_j\) form a basis selected from the finitely many theta
series. The functions \(h_j\) are integrable even if a single
uncombined term in the original genus expression has no known
integral. Indeed choose \(s\) evaluations \(a_i\) for which
\((q_j(a_i))\) is invertible: all evaluation vectors span
\(\mathbf C^s\) by independence of the functions. The sums at these
\(s\) points are integrable by the preceding absolute estimate, and
matrix inversion expresses each \(h_j\) as their linear combination.
The zero-dimensional span gives zero directly. Divide by the
positive \(\rho_\delta\) to obtain the asserted finite expansion.
For a collection of orbits extend the families by the equivariant
law from each representative and by zero on the complement. The
stabilizer law makes the extension well-defined and the unitary
multiplier preserves its bound.
\end{proof}
\src{src/Development/AbelianVarietyTheory/Proofs/SiegelThetaTransformation/ExistsFundamentalDomain.lean}{138}{A fundamental set on lattice--frame pairs} and
\src{src/Development/AbelianVarietyTheory/Proofs/SiegelThetaTransformation/HermitianThetaSeriesSpanBoundedForms.lean}{22}{finite rank after absolute unfolding}.

For the rows of Lemma~\ref{munf:marks} in their common rational
Hermitian space, use the graph sum of Proposition~\ref{dg:identity}.
The passage to algebraic cusp sections then uses \eqref{bfam:reading}
on that space with its marked multiplier.


\paragraph{The fourth power of the theta key.}
Write \(g=f(p+q)\). On the abelian locus let \(L\) be the rigidified
symmetric principal theta line, \(P=\pi_*L\), and
\(\omega=\det(\operatorname{Lie} A)^\vee\), where the dual applies to
the determinant. The line \(P\) commutes with base change. Indeed its
fibres have one section and vanishing higher cohomology. A local
generator \(s\) cuts out a relative effective theta divisor
\(\Theta\). The canonical evaluation identifies
\(\mathcal O(\Theta)=L\otimes\pi^*P^{-1}\), hence
\(e^*\mathcal O(\Theta)=P^{-1}\). In particular this identification
also works when the origin lies on \(\Theta\).

\begin{proposition}\label{ctwist:key}
There is a functorial unit section on the abelian locus, up to a
sign on each parity, of \(P^{\otimes8}\otimes\omega^{\otimes4}\).
Its scalar \(k(s,v)\) in a theta frame \(s\) and a tangent basis
\(v\) transforms by
\[
 k(a s,Mv)=a^{-8}(\det M)^4 k(s,v).
\]
In the metric on the cover its absolute norm is \((2\pi)^{2g}\).
In a primitive theta frame and a full integral tangent basis it is
a unit at good reduction, including residue characteristic two.
\end{proposition}
\begin{proof}
Apply Moret--Bailly, Theorems 0.2 and 0.5 and Corollary 1.5
\cite{moret-bailly-key}, to the effective principally polarized
symmetric divisor. His line is
\(\omega\otimes(e^*\mathcal O(\Theta))^{-2}\).
The scalar is taken relative to \(s^8\otimes\eta_v^4\), where
\(\eta_v(v_1,\ldots,v_g)=1\). This gives both transformation
exponents and the assertion over an abelian scheme over a DVR.
The multiplier-line description of Candelori, Theorem 5.1 and
Remarks 5.2--5.3 \cite{candelori-key}, is over
\(\mathbb Z[1/2,i]\).

Here is the metric calculation. On the even characteristic chart
take periods \((1,\Omega)\), \(Y=\operatorname{Im}\Omega>0\),
the usual Riemann theta, and
\(v_i=(2\pi i)^{-1}\partial/\partial z_i\). The analytic
trivialization of Theorem 0.5 gives \(|k|=1\) in these frames.
Write \(N_s\) for the Haar mean of the squared rigidified theta
norm. With \(z=x+\Omega y\), integration in \(x\) annihilates the
cross terms, and summation of the diagonal terms gives
\(\int_{\mathbb R^g}\exp(-2\pi y^tYy)\,dy=(\det(2Y))^{-1/2}\).
Put \(N_v=\operatorname{covol}(\Lambda)/
(\pi^g|\det(v)|^2)\) in the same complex coordinates. Then
\[
 |k|N_s^4 N_v^2=(2\pi)^{2g}.
\]
The two exponents just computed give invariance under changes of
theta and tangent frames. A symplectic basis puts every even
characteristic in the chosen chart. For an odd characteristic take
the product of the pair with itself. This has even parity, the
two metrics multiply, and Corollary 1.5 identifies the product
trivializations up to sign. Positive square roots give the claimed
identity in odd parity as well. On an orthonormal cover frame the
Hermitian matrix of the polarization is \(2\xi I\) and the theta
has \(N_s=1\). Unimodularity of the alternating form gives
\((2\xi)^g\operatorname{covol}(\Lambda)=1\), so
\(N_v=(2\pi\xi)^{-g}\). The cover's Lie factor for the four
determinants is \((2\pi\xi)^{-2g}\), giving the stated norm.
\end{proof}

There is a further integral change of frame even on the abelian
locus. Let \(R\) be a test DVR whose fraction field splits \(E\),
\(T\) its free Lie lattice, and \(T_\tau=T\cap(T_K)_\tau\) of
rank \(m_\tau\). If \(\mathfrak D_E\) is the different, then
\[
 0\le {\rm length}_R(T/\textstyle\bigoplus_\tau T_\tau)
 \le\sum_\tau m_\tau {\rm ord}_R(\tau\mathfrak D_E).
\]
Indeed for the projector \(e_\tau\), the trace-dual basis formula
\(e_\tau=\sum_i b_i\otimes\tau(b_i^*)\) gives
\(\tau(\mathfrak D_E)e_\tau T\subset T_\tau\):
the \(b_i\) are an integral basis and \(b_i^*\in\mathfrak D_E^{-1}\).
Projection embeds the displayed quotient in
\(\bigoplus_\tau e_\tau T/T_\tau\). Each summand is generated by
at most \(m_\tau\) elements and is killed by
\(\tau(\mathfrak D_E)R\). In the combined eigentangent basis
the valuation of the key increases by four times this length.
All normalized valuations follow by dividing by the ramification
index of the test DVR. Consequently \(k^m\) is integral at these
tests for every \(m\ge0\), with logarithmic complex cost
\(2gm\log(2\pi)\).
\src{src/Development/AbelianVarietyTheory/Proofs/DivisorProofs/ThetaKeySectionNormComplexPointsEqTwoPiPow.lean}{124}{The normalized cover metric};
\src{src/Development/AbelianVarietyTheory/Proofs/DivisorProofs/ExistsKeyUnitPowTwistingSectionOnComponent.lean}{135}{The power of the key and the separate root section}.

\begin{theorem}[Integral key at the semiabelian tests]\label{ctwist:semi}
For the rigidified graph line of Lemma~\ref{intjet:maximum},
\(k(s,v)\) has valuation zero at every semiabelian DVR test in
a primitive theta frame and a full integral tangent basis. The
different bound above therefore applies to the combined eigentangent
basis at all finite places. On each component \(k^m\) supplies the
weight-\(-4m\) multiplier with cost \(2gm\log(2\pi)\).
\end{theorem}
\begin{proof}
It suffices to compute an order after completion and finite
extension. Use the connected model, and split the semistable
reduction after extension. Write \(A/K\) for the generic fibre.
In an additional finite extension make \(A[8]\) rational. Let
\(\mathcal A'\) be the open subgroup of the Néron model generated
by its identity component and the sections of \(A[2]\).
The identity component is the given connected semistable model
(SGA~7, IX, Proposition 3.2, also used in
\cite[IV, \S3.4]{moret-bailly-pinceaux}); this identification
also applies to its semistable base change. The component group
of \(\mathcal A'\) is killed by two, and
\(\mathcal A'[2]\) is finite flat. For the latter assertion the
kernel of multiplication by two is quasi-finite flat on a
semistable group: on each geometric fibre multiplication factors
on the connected semiabelian variety through isogenies of its
torus and abelian quotient, and flatness follows from the
fibrewise flatness criterion. Its generic fibre is thus
schematically dense. Its support is contained in the union of
the finitely many sections extending the generic two-torsion,
which is closed and finite over the base. The kernel is proper
by the valuative criterion, hence is finite. This argument allows
nonreduced special two-torsion.

Moret--Bailly, II, Corollary 1.2.8 with \(n=2\), gives a symmetric
cubical extension \(L'\) on \(\mathcal A'\). Its restriction to
the identity component is the rigidified cubical line of the
test by uniqueness. It is relatively ample by VI, Proposition
2.1, and \(K((L')^2)=\mathcal A'[2]\), since both are the flat
closures of the generic principal two-torsion. The special
fibre of the connected model is geometrically integral. A
primitive theta section has order zero there: if it vanished
along the fibre, division by a uniformizer would give another
integral section. Its nonzero reduction persists under residue
field extension. Lemma~\ref{intjet:maximum} gives this section
canonical norm at most one on the entire analytification.
Consequently its order along every additional component \(C\)
of \(\mathcal A'\) is nonnegative, by the cubical model formula
\({\rm ord}_C(s)=-\log\|s(x_C)\|\) at its Shilov point
(\cite[Proposition 10.3]{deJongShokrieh}). Smoothness gives
normality, so the vertical inequalities and generic regularity
extend the section over \(\mathcal A'\). Restriction back to the
identity component proves that the two theta lattices agree.
This reasoning also compares a test with additional components
by restriction to the identity component. The Lie lattices agree
under open immersions.

Now apply Moret--Bailly, VIII, 1.0--1.3(i)
\cite{moret-bailly-pinceaux}. The base is a normal trait, the
abelian locus is schematically dense, \(K((L')^2)\) is finite,
and the rank \(d\) of the theta module is one. The asserted
canonical trivialization of \(\Delta^{a(g)}\),
\(\Delta=(\pi_*L')^2\otimes\omega\), therefore applies over
\(\mathbb Z[1/d]=\mathbb Z\). Denote it by \(u\). On the integral
stack of abelian schemes with their effective symmetric principal
theta divisors, \(k^{a(g)}/u^4\) is a global unit. Its square is
one by Moret--Bailly, Proposition 1.4(iii) and Corollary 1.5
\cite{moret-bailly-key}. Here one passes between rigidified lines
and effective divisors by the evaluation construction preceding
Proposition~\ref{ctwist:key}. Compatibility of \(u\) with
restriction to the generic point shows
\(a(g){\rm ord}(k)=4{\rm ord}(u)=0\) in the two agreed lattices.
Neither extracting an integral root nor bounding \(a(g)\) is
required. The eigentangent change and Proposition~\ref{ctwist:key}
now give both estimates. Taking the power on the component uses
the inverse only on the characteristic-zero space; extension by
its component idempotent gives the supported multiplier.
\end{proof}

The distinction between nonnegative order and order zero can be
checked on the Tate elliptic curve with parameter \(Q\) and odd
divisor \([0]\). In its connected model
\(e^*\mathcal O([0])=\omega^{-1}\), hence
\(\Delta^4=\omega^{12}\). The modular discriminant is a
nowhere-zero section of this line on the odd abelian stack.
It agrees up to sign with the key by the unit calculation above.
In the invariant Tate frame its scalar is
\(Q\prod_{n\ge1}(1-Q^n)^{24}\), of order \(v(Q)\).
More generally, on a connected semistable model without the
graph condition one always has \(P'\subseteq P\), with \(P'\)
the enlarged model's theta lattice. If \(P'=aP\), the same
canonical-power argument on \(\mathcal A'\) gives
\({\rm ord}(k(s,v))=8{\rm ord}(a)\ge0\). The odd Tate example
has \({\rm ord}(a)=v(Q)/8\) after the extension. Its toric
characteristic is \(1/2\): in the half-characteristic theta
series the toric exponents are shifted by \(1/2\).
In the graph case Lemma~\ref{intjet:maximum} removes this
index and gives the upper different bound as well.

On a normal toroidal chart with these primitive theta and
invariant differential lattices the same argument applies to
the trait at each generic boundary divisor in characteristic
zero. The key and its inverse are units there in full Lie
frames. Thus its multiplication preserves vanishing at those
boundary divisors. The scalar of this fourth-power section and
the existence of the auxiliary one-key section are distinct.
The latter asks for a fourth root up to a scalar in \(F^\times\)
on each component. Over a characteristic-zero component with
trivial geometric monodromy on the symplectic four-torsion this
follows also from Candelori, Definition 3.2 and Theorem 5.1
\cite{candelori-key}. The multiplier torsor is obtained from
the isometry torsor of the symplectic four-group with its theta
characteristic by pushing forward along a character into \(\mu_4\).
In a geometric marking both this torsor and the torsor of
fourth roots of the key are constant: their two fourth-power
trivializations differ by a constant on the universal complex
stack, by the unit calculation in \cite{moret-bailly-key} in
genus \(g\ge2\). Descent of a geometrically constant \(\mu_4\)-torsor
on a geometrically connected component amounts to a torsor over
\(F\), i.e. to a class in \(F^\times/(F^\times)^4\).
The deep level in the statement has odd \(M\ge3\) and a separate
auxiliary marking \(A\). The odd level alone gives no full
four-torsion marking. For the original even weight spaces the
unit \(k^m\) above is intrinsically defined.
\src{src/Development/AbelianVarietyTheory/Proofs/DivisorProofs/ExistsKeySectionFourthRootPerComponent.lean}{93}{The separate descent up to a component scalar}.

\paragraph{Determinants on a characteristic-zero component.}
Let \(\Lambda\) be the projective \(\mathcal O_E\)-lattice of a
component, with level group contained in its unitary isometries
acting as the identity on \(\Lambda/M\Lambda\), \(M\ge3\).
Every such isometry has determinant one over \(E\). In fact its
determinant \(u\) preserves the invertible ideal
\(\bigwedge_{\mathcal O_E}^{p+q}\Lambda\), is \(1\bmod M\)
(localize to check this on free bases), and satisfies
\(|\tau u|=1\) at every embedding. If \(u\ne1\),
\(M^{[E:\mathbb Q]}\le |N_{E/\mathbb Q}(u-1)|
\le2^{[E:\mathbb Q]}\), an impossibility.
For the split half-marking, fixing \(\Lambda/\mathfrak m\Lambda\)
first gives the congruence modulo \(\mathfrak m\). Taking the
adjoint with respect to the perfect alternating integral pairing
gives the conjugate-ideal congruence for the inverse, and hence
for the isometry itself. Indeed the dual of
\(\mathfrak m\Lambda\) is
\(\overline{\mathfrak m}^{-1}\Lambda\). The two ideals are
coprime with product \(M\mathcal O_E\), giving the asserted
principal congruence.

On such a classical complex component the determinant of the
\(\tau\)-part of the standard Betti local system thus has a global
horizontal generator, including for noncompact signatures. On the
period domain the polarization gives the Hodge exact sequence
\[
 0\longrightarrow\omega_\tau\longrightarrow
 H^1_{\rm dR,\tau}\longrightarrow\omega_{\bar\tau}^{\vee}
 \longrightarrow0.
\]
Dualizing its determinant gives
\(\det\operatorname{Lie}_\tau/\det\operatorname{Lie}_{\bar\tau}\),
the line of the elementary section \(s_\tau\).
The norm of a horizontal volume is constant: in a unitary
orthonormal basis the absolute determinant of the Hermitian form
on the full space is one, independently of its positive and
negative Hodge subspaces. Changing to the cover normalization
multiplies this by a constant fixed on the component.
For a smooth geometrically connected algebraic component over
\(F\) carrying the indicated abelian scheme and splitting of \(E\),
the Gauss--Manin determinant is regular singular. Full faithfulness
for regular singular connections gives an algebraic horizontal
generator over \(\mathbb C\) (Deligne, II, Theorem 5.9
\cite{deligne-reg-sing}). It descends up to a scalar to \(F\):
horizontal sections are the kernel of the connection on global
sections, and both global-section spaces and the kernel commute
with extension of scalars on this quasi-compact separated
component. The descended one-dimensional space generates the
line everywhere by faithful flatness. This proves rational
triviality with its archimedean constant for these components.

There are two further consequences for a fixed classical component.
Take a smooth characteristic-zero toroidal compactification with
unipotent local monodromy and with the Hodge bundles given by the
canonical extension of the Gauss--Manin bundle. The horizontal
determinant and its inverse extend as units in the ratio of
the two extended Hodge determinants. Indeed, near a normal
crossings boundary with parameters \(z_i\), the canonical
extension in a multivalued horizontal basis is obtained by
multiplying by
\(\exp(-\sum_i(\log z_i)\log T_i/(2\pi i))\), where \(T_i\)
are the local monodromies. Each \(\log T_i\) is nilpotent on
the \(\tau\)-part, so this matrix has determinant one. Taking
determinants of the extended Hodge sequence identifies its
volume line with the ratio in question. Thus multiplication
in either direction preserves the subcanonical cusp ideal
\(\mathcal O(-D_\infty)\) on these charts. The extended Hodge
sequence and the boundary ideal in the classical PEL model are
given in \href{https://www-users.math.umn.edu/~kwlan/articles/cpt-ram-ord.pdf}
{Lan, Proposition 1.4.2.1 and Definition 1.4.2.9}.

On any fixed finite-type model of this compactification and line
over an open of \(\operatorname{Spec}\mathcal O_F\), the same
section is an integral generator after deleting finitely many more
places. To see this, choose a finite affine trivializing cover.
The coordinates of \(s\) and \(s^{-1}\) on the generic fibre and
the finitely many gluing identities all descend after inverting a
single nonzero base integer. On this open the order is zero at
every test whose lattice is the pullback of that model line.
Away from \(\Delta_E\), the embedding idempotents of the split
étale algebra \(\mathcal O_E\otimes\mathcal O_{F,w}\) are integral,
so formation of its Lie eigensummands commutes with pullback.
In particular this gives the order for the abelian tests extending
to the fixed PEL chart with its universal abelian scheme. The
exceptional primes produced by this argument depend on the
component, horizontal generator and model. Controlling all the
semiabelian tests and the total arithmetic degree in a varying
family requires the comparisons below.

\paragraph{One vertical order on one component.}
Fix a component \(Z\), a host embedding \(\tau\), and a finite place
\(w\) of its coefficient field \(F\). Orders on test traits are
normalized to restrict to \({\rm ord}_w\) on \(F\). The following
divisor calculation explains what an integral model of the parallel
line has to establish. The special fibre may have several irreducible
components with unequal multiplicities.

\begin{lemma}[Vertical order on a proper model]\label{ctwist:vertical}
Let \(\mathcal X\) be integral, normal, proper and flat over a DVR
\(R\) with uniformizer \(\pi\) and nonempty connected special fibre.
Let \(\mathcal L\) be an invertible sheaf, trivialized by \(s\) on the
entire generic fibre. Suppose some positive power of \(\mathcal L\)
is trivial on the reduced special fibre. Write
\({\rm div}(\pi)=\sum_i m_i C_i\) and
\({\rm div}(s)=\sum_i a_i C_i\). Then all \(a_i/m_i\) are the same
rational number \(\alpha\). For every local test \(R\to V\) with
\(\operatorname{Spec}V\to\mathcal X\), where \(V\) is a valuation
ring extending \(R\), the coordinate of \(s\) in an integral generator
has order \(\alpha\,v(\pi)\). If \(R=\mathcal O_{F,w}\) and one
such test is a DVR of ramification index one, then \(\alpha\in\mathbb Z\).
\end{lemma}
\begin{proof}
Set \(\alpha=\min_i a_i/m_i\). Choose \(M>0\) divisible by its
denominator and by the order of the restricted line in the Picard
group. The rational section \(u=\pi^{-M\alpha}s^M\) has nonnegative
order along every prime divisor and order zero at a minimizing
component. It extends over \(\mathcal X\) by normality, checked in
local frames by the height-one intersection criterion.
In a trivialization of \(\mathcal L^M\) on the reduced special
fibre, its restriction is a global function nonzero at that
component's generic point. Global functions on a proper reduced
connected scheme over a field form a field: their algebra is
finite-dimensional and reduced, and connectedness excludes
nontrivial idempotents. Hence \(u\) is a unit on the reduced fibre,
and therefore generates at every point of the special fibre by
Nakayama. It also generates on the generic fibre. Pulling
\(s^M=\pi^{M\alpha}u\) to \(V\) proves the assertion. At an
unramified DVR, \(v(\pi)=1\) and the coordinate has integral order.
\end{proof}
The two scheme facts used here are
\href{https://stacks.math.columbia.edu/tag/0AVB}{Stacks Project,
Tag 0AVB} and
\href{https://stacks.math.columbia.edu/tag/0366}{Section 33.9}.
If the special fibre has just one irreducible component, the same
proof dispenses with its Picard group: after clearing the single
ratio both \(u\) and \(u^{-1}\) extend by normality. In the case
of a reduced irreducible fibre \(M=1\) suffices. For a model over
a ramified base \(R/\mathcal O_{F,w}\), divide its orders by the
ramification index before comparing tests. The rational choice of
frame alone cannot determine \(\alpha\): for the lattice
\(\ell^a\mathbb Z_\ell\) in \(\mathbb Q_\ell\), the frame \(1\)
has integral coordinate \(\ell^{-a}\).

Here is the specific model comparison needed to apply this reasoning.
On the good deep cover write \(U\) for the framed uniformization.
For each fixed \((Z,\tau,w)\), applying the vertical divisor lemma
requires comparing the test frame with local charts on a normal
projective flat model after a finite level and base change. The
comparison must lift
\emph{every} admissible test, match its frame coordinate up to a
unit, and use the same normalized order on the base uniformizer.
Triviality on the reduced fibre, or an equivalent constancy assertion
in these compatible charts, is essential on reducible fibres.
On the dependency path of the exported theorem the generator on
the component is obtained using resolution and GAGA. The source
then uses the associated period frame to compute the degree.
The finite correction between the two frame conventions at a fixed
\((\tau,w)\) is \(q\) times the logarithmic difference between the
different and \(d_0\) in the prescribed local absolute value. Adding
this fixed correction preserves equality of the readings on that
component. The final change of the global degree follows by summing
the finite corrections. Lemma~\ref{ctwist:vertical} proves the
vertical divisor step from its stated model hypotheses. The splitting
model below computes the ordinary integral lattices directly.
\src{src/Development/AutomorphicSpectrum/Proofs/GeometryProofs/FlatFirstFromResolutionAndGagaStatement.lean}{8}{Generator from resolution and GAGA};
\src{src/Development/AutomorphicSpectrum/Proofs/GeometryProofs/KroneckerCharacterFrameBound.lean}{23}{Passage to the component period frame};
\src{src/Development/AutomorphicSpectrum/Proofs/GeometryProofs/KroneckerCharacterPeriodFormula.lean}{122}{Finite correction in the ordinary metric};
\src{src/Development/AutomorphicSpectrum/Proofs/DivisibilityProofs/FlatDegreeStOfResKGaga.lean}{23}{Passage from the constant degree to the corrected section}.

\paragraph{Computing the common degree at a CM point.}
Write \(A_\chi=\operatorname{Re}\log\Omega(E,\chi)\). Assume the
local equality just discussed, and the constant archimedean norm
in the prescribed metric. One point with everywhere good abelian
reduction now suffices to read the uncorrected degree on \(Z\).
In the original construction one chooses an orthogonal \(E\)-frame
in the \emph{specified} complex uniformizing component, constructs
its integral CM object over some \(K/F\), and identifies its framed
generic fibre with a point of \(Z\). This choice works for every
\(\tau\) of the chosen CM orientation on that \(Z\). Write \(T_+\)
and \(T_-\) for the integral Lie lattice intersected with the
\(\tau\)- and \(\bar\tau\)-eigenspaces there, with the induced
polarization metrics. Their ranks are \(p,q\). The frame law
identifies the uncorrected line degree of \(s_\tau\) with
\(\widehat{\deg} T_+-\widehat{\deg} T_-\), using normalized
determinant degrees. Indeed integral local bases compute its finite
coordinate by a ratio of exterior determinants. The same ratio in
orthonormal frames computes its norm at infinity. Changing rational
bases adds principal divisors of degree zero. Passing to \(K\)
preserves the answer because
\(\sum_{v\mid w}e(v/w)f(v/w)=[K:F]\), and each complex embedding
has \([K:F]\) extensions. This calculation uses the integral
eigensublattices even at primes dividing the different.
\src{src/Development/AbelianVarietyTheory/Proofs/DivisorProofs/ExistsCmGoodReductionPointOnComponent.lean}{57}{The CM point on the prescribed component};
\src{src/Development/AbelianVarietyTheory/Proofs/DivisorProofs/FlatLineDegreeEqEigenLatticeDegreeSub.lean}{20}{Finite and infinite determinant readings}.

The cancellation of the component's indices can be seen explicitly.
Let \(S_\pm\subseteq T_\pm\) be the full-rank images spanned by the
differentials of the \(p\) and \(q\) CM factors of the two types,
\(I_\pm\) their determinant index ideals in \(\mathcal O_K\), and
\(g_\pm\) their respective polarization Gram matrices over
\(\mathcal O_E\). Put \(G_\pm=|N_{E/\mathbb Q}\det g_\pm|\).
The CM period comparison for the images gives
\[
 \widehat{\deg}S_+=pA_\chi/2-\frac{\log G_+}{2[E:\mathbb Q]},
 \qquad
 \widehat{\deg}S_-=qA_\chi/2-\frac{\log G_-}{2[E:\mathbb Q]} .
\]
For this mixed-isogeny calculation suppose that the finite comparison
supplies an ideal \(\mathfrak n\) of \(E\) satisfying, with ideals
extended via \(\tau\) when necessary,
\[
 \mathfrak n\bar{\mathfrak n}=(\det g_+\det g_-),\qquad
 I_-(\mathfrak n\mathcal O_K)=I_+(\tau\det g_-).
\]
The definite specialization of the index comparison occurs on the
exported dependency path. The displayed mixed ideal identities are
the hypotheses of the following cancellation calculation.
Proposition~\ref{cmeig:degree} below gives an evaluating point on
the indefinite family for which both Lie lattices are integral
direct sums and the mixed ideal hypotheses can be dispensed with.
\src{src/Development/AbelianVarietyTheory/Proofs/DivisorProofs/DegreeImagesEigenLinesCmUnitaryShimuraObject.lean}{122}{Period and Gram determinant of each image};
\src{src/Development/AbelianVarietyTheory/Proofs/DivisorProofs/CmDefiniteLieIndexIdealEqNormIdeal.lean}{20}{Index comparison for the definite specialization};
\src{src/Development/AbelianVarietyTheory/Proofs/DivisorProofs/CmObjectDefiniteLieEigenLatticeDegreeDifference.lean}{21}{Its degree consequence}.

Taking norms gives \(N\mathfrak n=(G_+G_-)^{1/2}\) and
\(NI_+/NI_-=(N\mathfrak n/G_-)^{[K:\tau E]}\). The finite index
adds \(\log NI_\pm/[K:\mathbb Q]\) to the normalized degree in
passing from \(S_\pm\) to \(T_\pm\). Consequently
\[
 \widehat{\deg}T_+-\widehat{\deg}T_-=(p-q)A_\chi/2 .
\]
All dependence on the two Gram matrices has cancelled. This
explains why the degree is independent of \(Z,\tau\) at fixed
host and signature, although the individual local orders may vary.
It also explains why using the CM point \emph{before} proving local
constancy would leave the orders at the other tests undetermined.

\paragraph{The finite degree in the parallel twist.}
Use \(j_\tau=\lambda_{\tau,p}-1-4m\), \(m=fr+1\). Here the
one-based index \(p\) is the multiplicity at \(\tau\) in the
algebraic weight. The positive row block in the period chart has
size \(p_0=q\). In the
source the elementary flat section \(s_\tau\) has entry \(-1\)
on the entire \(\tau\)-block, so the required product is
\(\prod_\tau s_\tau^{-j_\tau}\). For each \(w\) the comparison
needed with its integral model is the equality of the orders at
all tests over \(w\), with common value
\[
 \rho_{\tau w}+D_{\tau w},\qquad
 D_{\tau w}=q\bigl({\rm ord}_w(\tau\mathfrak D_E)
                            -{\rm ord}_w(d_0)\bigr),
\]
where \(d_0\) is the host prime and \(\rho_\tau\) has finite
support. The source defines the corrected degree using
\(\rho_\tau\); denote it by \(-c_1\) if it is independent of
\(\tau,Z\). The norm of the different and the product formula
give
\[
 [F:\mathbb Q]^{-1}\sum_w D_{\tau w}\log Nw
      =q(\log {\rm rd}(E)-\log d_0),\qquad
 {\rm rd}(E)=|\Delta_E|^{1/[E:\mathbb Q]} .
\]
Indeed the norm of the extended different is
\(|\Delta_E|^{[F:E]}\), while that of \(d_0\mathcal O_F\) is
\(d_0^{[F:\mathbb Q]}\). The uncorrected degree of the product
in \eqref{bfam:flat} is therefore exactly
\begin{equation}\label{ctwist:degree}
 \delta\left(\prod_\tau s_\tau^{-j_\tau}\right)
   =\bigl(c_1-q(\log{\rm rd}(E)-\log d_0)\bigr)
       \sum_\tau j_\tau .
\end{equation}
If \(j_\tau=0\) for every \(\tau\), the multiplier is \(1\) and its
integral cost is zero before any comparison. If only
\(\sum_\tau j_\tau=0\), cancellation in \eqref{ctwist:degree}
uses equality of the corrected degrees for the embeddings with
nonzero exponents. For general \(j_\tau\), suppose as in the composition
that \(|c_1|\le9rL\), \(\log{\rm rd}(E)\le3L\), and
\(|\log{\rm rd}(E)-\log d_0|\le\log{\rm rd}(E)\).
The last inequality follows when \(E\) contains the host quadratic
field of absolute discriminant \(d_0\): the discriminant tower
gives \(2\log{\rm rd}(E)\ge\log d_0\).
At each \(w\) above a prime outside the bad set \(B\), a test of
ramification index one makes the common order integral. Take the
product first, then apply \eqref{bfam:flat} over \(F\) once.
Writing \(W=\sum_{\tau,i}|\lambda_{\tau i}|\), the logarithmic
cost for both twists together is at most
\[
 2gm\log(2\pi)+12rL\,\bigl(W+f(1+4m)\bigr)
 +\sum_{\ell\in B}\log\ell
 +\frac{\log|\Delta_F|}{2[F:\mathbb Q]} .
\]
Here \(q\le r\) and \(\sum_\tau|j_\tau|\le W+f(1+4m)\).
The estimate holds for every component and embedding at the same
threshold as the field, level and degree bounds. On the selected
classical family Theorem~\ref{cmeig:splitmodel} gives the
uncorrected degree \((p-q)A_\chi/2\) of \(s_\tau\). Subtracting
the degree of \(D_\tau\) gives the corrected degree
\[
        -c_1=\tfrac12(p-q)A_\chi-q(\log{\rm rd}(E)-\log d_0).
\]
Thus the product coefficient in \eqref{ctwist:degree} is
\(-\tfrac12(p-q)A_\chi\), including zero in balanced signature.
On the chosen family
the bound on \(c_1\) follows numerically from
\(|A_\chi|\le3L\), \(0\le\log d_0\le L\),
\(0\le\log{\rm rd}(E)\le2L\), and \(p+q=r\).
In fact these bounds give \( |c_1|\le(3/2+2)rL\).
\src{src/Development/AbelianVarietyTheory/Proofs/DivisorProofs/FlatDegreeLogFormAbsLe.lean}{55}{Period and discriminant bounds on the selected stage};
\src{src/Development/AbelianVarietyTheory/Definitions/Foundations/ExistsFlatSecRankZero.lean}{751}{The corrected degree in Lie multiplicities};
\src{src/Development/AbelianVarietyTheory/Proofs/DivisorProofs/ExistsTwSecOfConst.lean}{87}{Conversion to the degree of the actual product};
\src{src/Development/AutomorphicSpectrum/Proofs/GeometryProofs/KroneckerCharacterFrameBound.lean}{20}{The period in the degree of the uncorrected frame};
\src{src/Development/AutomorphicArithmetic/ThetaCorrespondence/ThetaPairLifts/HodgeBundleWeights/ShortSpanCompositionClosed.lean}{54}{The algebraic weight index in the parallel exponent}.

There is a weaker sufficient test when only a lower bound for the
product degree is sought. On a common normal proper flat model at
\(w\) carrying the actual lines and tests of Lemma~\ref{ctwist:vertical},
let \(a_{\tau i}\) be the order of \(s_\tau\) along \(C_i\), and use
the base \(\mathcal O_{F,w}\). Set
\[
 b_w(j)=\min_i\frac{-\sum_\tau j_\tau a_{\tau i}}{m_i}.
\]
Then \(b_w(j)\) is a lower bound at \emph{every} test for this
product without any special-fibre triviality assumption. Indeed,
raise the product to a power \(M\) clearing the denominator of
\(b_w(j)\), and divide by \(\pi^{Mb_w(j)}\). Its divisor is
effective, so normality extends it over the model. Pullback to
each test yields the inequality. A generic trait of a minimizing
vertical divisor realizes equality among all trait maps. On a
reduced special fibre \(m_i=1\) and \(b_w(j)\) is integral.
Thus, if these lower bounds have finite support, are integral
outside \(B\), and give a sufficiently large normalized degree
in \eqref{bfam:flat} with the product's actual metric, the same
integral spanning estimate follows. This criterion applies on each
chosen model separately; a bound uniform in the window still has
to be supplied for the indicated weighted minima.

\paragraph{A product CM point on each indefinite marked row.}
For the fields of Proposition~\ref{host:art} there is another
choice of evaluating point. Write \(k=\mathbb Q(\sqrt{-d_0})\),
\(F_0=E^+\), \(O=O_E\), and \(\Phi\) for the embeddings inducing
the chosen embedding of \(k\). For \(\tau\in\Phi\) write \(p,q\)
for the Lie multiplicities at \(\tau,\bar\tau\), respectively.
Write \(p_0\) for the number of positive Hermitian row directions.
The selected period map conjugates the first \(p_0=q\) coordinates
and is \(\tau\)-linear on the last \(r-p_0=p\) coordinates.
Keep the period convention \(\Pi=\mathfrak a L\) of
Lemma~\ref{bfam:minimum}, so \(L\) is Hermitian self-dual.
\src{src/Formalism/LocalToGlobal/ProjectiveLattices/PetersonNormSubspace.lean}{67}{The conjugated positive block in the period map};
\src{src/Formalism/LocalToGlobal/ProjectiveLattices/PetersonNormSubspace.lean}{75}{The positive row index on a good chart};
\src{src/ReciprocityTheory/ReductiveGroups/ArithmeticGeometry/BirationalGeometry/ToroidalCompactification/FlagBasis.lean}{76}{The Lie multiplicities of the two embeddings}.

\begin{lemma}\label{cmeig:row}
If \(p,q\ge1\), every marked row in this family has, up to rational
isometry, underlying lattice
\[
        J e_1\ \perp\ O e_2\ \perp\cdots\perp\ O e_r,\qquad
        J\bar J=O,\qquad
        h(e_i,e_i)=
        \begin{cases}1&i\le p_0,\\-1&i>p_0.\end{cases}
\]
The transported mark need not split into marks of the summands.
\end{lemma}
\begin{proof}
The two imaginary quadratic subfields of discriminants
\(-d_0,-q_0\) in Proposition~\ref{host:art} have coprime
discriminants (here \(q_0\) is the auxiliary prime). At each
rational prime at least one of them is unramified. Its inertia
restriction is trivial, whereas complex conjugation acts on it
nontrivially. Thus \(E/F_0\) is split or unramified at every finite
place, since relative inertia is the intersection with
\(\operatorname{Gal}(E/F_0)\).

At a split place a basis of the first factor and its pairing-dual
in the second identify a self-dual lattice with the standard one.
At an unramified quadratic place its reduction is a nondegenerate
Hermitian space over the quadratic residue extension. There is a
vector of nonzero norm: otherwise polarization of \(h(x+ay,x+ay)\)
and the perfect residue trace pairing would give \(h=0\).
Lifting gives a unit norm, which can be scaled to either sign,
since the norm maps unramified units onto units (use the residue
norm and lift successively with the residue trace). Its integral
orthogonal complement is again self-dual. This gives the local
orthogonal basis by induction, including in residue characteristic
two. At infinity the signatures agree with the displayed form.
For clarity, the rational isometry follows already from
Hasse--Minkowski: the underlying quadratic form \(h(x,x)\)
represents the prescribed first sign at every place, hence globally.
For the latter implication apply isotropy to \(h(x,x)-a t^2\);
if \(t=0\) gives a nonzero isotropic vector, the original quadratic
space contains a hyperbolic plane and represents \(a\).
Split off the Hermitian line of norm \(a\); by Witt cancellation its
complement is locally isometric to the remaining diagonal form.
Induction now identifies the two rational Hermitian spaces.

Put \(L_0=\bigoplus O e_i\), \(G=U(h)\). Choose
\(g\in G(\mathbb A_{F_0,f})\) carrying \(\widehat L_0\) onto
\(\widehat L\) using the local isometries, choosing the identity
where the two lattices coincide. Write \(t=\det g\) and
\(d(t)=\operatorname{diag}(t,1,\dots,1)\). The determinant-one
groups have strong approximation away from the infinite places:
\(\mathrm{SU}(h)\) is simply connected, absolutely almost simple
(\(r\ge2\)), and noncompact at a real place by \(p,q\ge1\).
This is \href{https://www.mathnet.ru/links/02c4d3528af06a248862fb1aa3ba53f9/im2220_eng.pdf}
{Platonov, Theorem A}. Thus the open set
\(g K_0 d(t)^{-1}\cap\mathrm{SU}(h)(\mathbb A_{F_0,f})\)
contains a rational element \(\gamma\): it contains
\(g d(t)^{-1}\), and \(K_0\) is the open stabilizer of
\(\widehat L_0\). Then \(\gamma^{-1}L\) has completion
\(d(t)\widehat L_0\). Its first ideal is
\(J=E\cap t\widehat O\); \(t\bar t=1\) gives \(J\bar J=O\).
Transport the entire given finite marking by \(\gamma^{-1}\).
The resulting polarized marked row is isomorphic to the one
started with, as in Lemma~\ref{munf:marks}. No approximation to
the identity of the marking, or change of its level, is needed.
\end{proof}

\begin{proposition}[The degree on the product point]\label{cmeig:degree}
On each classical marked period component of
Lemma~\ref{cmeig:row} there is a point over a finite extension of
the coefficient field with everywhere good abelian reduction. For
every \(\tau\in\Phi\) its two integral Lie determinants in the
polarization metric of the cover satisfy
\[
 \widehat{\deg}T_+-\widehat{\deg}T_-=
   \frac{p-q}{2}\left(\gamma_{\rm Eul}+
     \log\frac{2\pi}{{\rm rd}(E)}
       -\frac{L'(1,\chi)}{L(1,\chi)}\right).
\]
All finite places, including those ramified in \(E\), occur with
their ordinary integral eigensublattices.
\end{proposition}
\begin{proof}
In the representative of Lemma~\ref{cmeig:row} take the flag of
the positive and negative coordinate axes. The period lattice
is the orthogonal direct sum of the fractional ideals
\(I_1=\mathfrak a J,\ I_i=\mathfrak a\ (i>1)\).
The last \(p\) factors have CM type \(\Phi\), and the first \(q\)
have its conjugate, by the period map above.
Let \(C\) be the elliptic curve with complex lattice \(O_k\),
with the indicated action (conjugate the action for the opposite
type). It is defined with its endomorphisms over a number field
and acquires good reduction everywhere after finite extension:
its CM \(j\)-value is an algebraic integer, hence gives potential
good reduction at the finitely many places of bad reduction.
The Serre construction \(I_i\otimes_{O_k} C\) exists integrally
since \(I_i\) is projective over \(O_k\): an idempotent on a
finite direct power of \(C\) realizes it as a direct factor.
Its periods at the chosen embedding are exactly the images of
\(I_i\) at the extensions of that \(k\)-embedding. The alternating
trace polarization on each factor is induced from a positive
Hermitian form \(\operatorname{Tr}_{E/k}(\beta u\bar v)\).
Indeed divide its purely imaginary trace coefficient (with the
coordinate sign) by the coefficient of the standard principal
elliptic polarization, a signed generator of
\(\mathfrak D_{k/\mathbb Q}^{-1}\); the ratio \(\beta\in E^+\)
is totally positive by the Riemann positivity condition.
Perfection is exactly
\[
       (\beta)\,I_i\bar I_i\,\mathfrak D_{E/k}=O .
\]
This also follows by taking the trace-dual ideal. Hence the
polarized product has exactly the required period lattice and
polarization. For the polarized tensor construction see
\href{https://arxiv.org/abs/1507.07607}{Amir-Khosravi, \S2}.
The transported finite torsion marking is defined after another
finite extension, irrespective of how it mixes factors. It is
needed on the characteristic-zero point; the underlying product
extends over the full ring of integers. In a classical PEL moduli
chart this constructs the point of the prescribed component,
including its canonical rigidified graph line. An additional
finite étale marking cover is handled by lifting the point to
the given component of the cover, whose map to the underlying
connected component is finite étale and surjective.

Here is the integral and metric calculation for a single factor
with ideal \(I\) and tensor coefficient \(\beta\). Enlarge the
number field \(K\) of definition to split \(E\), and put
\(\mathfrak D=\mathfrak D_{E/k}\). In
\(I\otimes_{O_k}O_K\), the intersection with the coordinate line
at an active embedding \(\nu\) is \(\nu(I\mathfrak D)O_K\).
In fact test membership in the lattice by its dual. The trace
dual of \(I\) over \(O_k\) is \(I^{-1}\mathfrak D^{-1}\);
this remains the dual after flat base change. On the split
generic algebra the trace is the sum of coordinates. The vector
supported at \(\nu\) with entry \(z\) thus pairs integrally
exactly when \(z\nu(I^{-1}\mathfrak D^{-1})\subset O_K\), as
asserted. Tensoring with the invertible \(O_K\)-module
\(\operatorname{Lie} C\) gives the actual Lie eigenline of the
abelian scheme, since taking Lie commutes with the idempotent
defining the tensor construction.

On that coordinate its squared metric at \(\sigma:K\to\mathbb C\)
is the squared elliptic Lie metric multiplied by
\(\sigma\nu(\beta)\): diagonalizing the trace Hermitian form in
the embedding coordinates gives coefficients \(\sigma\nu(\beta)\);
the polarization convention multiplies both Hermitian forms by
the same \(\pi\). For \(e=[E:\mathbb Q]\), normalized ideal and
metric degrees therefore give
\[
\begin{split}
 2\widehat{\deg}T_\nu(I)
 &=2\widehat{\deg}\operatorname{Lie}C
   -\frac{2\log N(I\mathfrak D)+\log|N_{E/\mathbb Q}\beta|}{e}\\
 &=2\widehat{\deg}\operatorname{Lie}C-\frac{\log N\mathfrak D}{e}.
\end{split}
\]
The second equality uses perfection, so even the fractional ideal
on the first coordinate disappears. The norm is the absolute
ideal norm and \(N\mathfrak D=|\Delta_E|/d_0^{e/2}\).

For \(C(\mathbb C)=\mathbb C/\Pi_C\) with Euclidean period area
\(a_C\), principal polarization gives the Hermitian coefficient
\(1/a_C\). The tangent metric is \(\pi/a_C\), so its dual on
\(dz\) is \(a_C/\pi=(i/2\pi)\int dz\wedge d\bar z\).
Thus \(-\widehat{\deg}\operatorname{Lie}C\) is exactly the
stable elliptic Faltings height with the normalization of
\href{https://people.math.wisc.edu/~tonghaiyang/KRY2.pdf}
{Kudla--Rapoport--Yang, (10.14)--(10.15), Proposition 10.10}.
Their quadratic CM formula gives
\[
 2\widehat{\deg}\operatorname{Lie}C
 =\gamma_{\rm Eul}+\log(2\pi/\sqrt{d_0})
                  - L'(1,\chi)/L(1,\chi).
\]
This applies to either elliptic orientation. Substitution gives
the same \(2\widehat{\deg}T_\nu(I)=A_\chi\) on every active
coordinate; orthogonality and the integral direct product give
the two determinants by addition over \(p,q\) coordinates.
\end{proof}

In particular no comparison between a mixed isogeny image and
its saturation is needed to evaluate this point. If the vertical
orders on a selected component agree at all tests and its
archimedean norm is constant, the degree is that of this point
by the normalized base-change calculation preceding
\eqref{ctwist:degree}.

\paragraph{Vertical orders at hyperspecial places.}
There is a direct good-reduction comparison for the ordinary Lie
frames of the classical PEL problem. It also distinguishes a prime
in the level from a prime ramified in the endomorphism algebra.

\begin{proposition}\label{cmeig:hyperspecial}
Suppose the characteristic-zero cover on the chosen component
is a classical finite marking level for its principally polarized
\(O_E\)-abelian scheme, and its finite tests use the integral Lie
eigensublattices of connected semiabelian extensions of that same
scheme. Let \(w\mid\ell\) in the splitting coefficient field \(F\),
with \(\ell\nmid 2\Delta_E\). Each rational horizontal determinant
generator has the same normalized order at all semiabelian DVR
tests on the component above \(w\). If its marked level itself is
hyperspecial at \(\ell\), this common order belongs to \(\mathbb Z\).
\end{proposition}
\begin{proof}
First take a neat prime-to-\(\ell\) marking with full hyperspecial
level at \(\ell\). The algebra is unramified at \(\ell\), its
involution is of unitary type and the principal alternation is
perfect. Thus this is a good reduction PEL problem. With smooth
projective cone data it has a proper smooth toroidal model with
universal semiabelian scheme by
\href{https://www.kwlan.org/articles/cpt-PEL-type-thesis-revision-dense.pdf}
{Lan, Theorems 6.4.1.1 and 7.3.3.4}.
Base change to the completed coefficient DVR. The geometric
components of a proper smooth family form a finite étale scheme;
the selected geometrically connected generic component defined
over the coefficient field extends to an open and closed component
with geometrically connected special fibre. That fibre is smooth,
so is reduced and irreducible. The split embedding idempotents of
\(O_E\otimes_{\mathbb Z}O_{F,w}\) are integral. The indicated Lie
determinant ratio is therefore a line bundle on the model.
Its characteristic-zero horizontal generator extends as a unit at
the toroidal boundary by the unipotent determinant calculation
above Lemma~\ref{ctwist:vertical}. On the whole model its divisor
is an integral multiple of the single reduced special fibre.

Every test extends to the toroidal model by properness. The
pullback of its connected semiabelian group is the given connected
semiabelian extension by uniqueness over a trait (SGA~7 I, IX,
Proposition 3.2). This is an isomorphism respecting the endomorphisms
by generic-fibre uniqueness. Hence its Lie eigensummands, including
those of a toric degeneration, are exactly the test lattices.
Pulling back the divisor in an integral local generator gives the
same normalized integer at every such test.

For general finite marking choose another full neat level of order
prime to \(\ell\) and a geometric connected component of the common
refinement lying over the original component. Over a finite field
extension it and both projection maps are defined. Fix one place
of this field above \(w\). The common component is finite étale
and surjective over each of its two characteristic-zero image
components, as follows by forgetting finite torsion bases. Its two
pulled-back determinant generators differ by a scalar in the
field of definition: both are horizontal for the same connection,
so their quotient has zero differential on the smooth geometrically
connected component. On the auxiliary prime-to-\(\ell\) component
the generator exists by determinant-one monodromy at its full neat
level, as proved above Lemma~\ref{ctwist:vertical}. Any test on the
original component first extends to the fixed base place and
then lifts generically to the common component after a finite
extension of its fraction field, by surjectivity. Connected
semiabelian pullbacks agree by the same uniqueness. The order on
the lifted test is the order on the auxiliary component plus the
valuation of the fixed scalar, both independent of the test after
normalization to \(F_w\). This proves the claim without requiring
an integral level-\(\ell\) marking of a semiabelian group.
\end{proof}

\begin{theorem}[Integral determinant comparison at the ramified primes]
\label{cmeig:splitmodel}
For the fields of Proposition~\ref{host:art}, under the same
classical moduli and test identifications as in
Proposition~\ref{cmeig:hyperspecial}, the common normalized order
exists at every finite place. For each indefinite component the
uncorrected arithmetic degree of the elementary horizontal
determinant is therefore \((p-q)A_\chi/2\).
\end{theorem}
\begin{proof}
Fix \(w,\tau\), and use a neat marking away from its residue
characteristic \(\ell\), with maximal level at \(\ell\) stabilizing
the single self-dual lattice (and its periodic translates).
As in Proposition~\ref{cmeig:hyperspecial} we can pass to this
level and to a finite local splitting field. At every place of
\(E^+\) above \(\ell\) the quadratic algebra is split or
unramified by Lemma~\ref{cmeig:row}. These are exactly the
unitary cases (AL) and (AU) of the Pappas--Rapoport splitting
model. Its order \(O_E\) is maximal, and its alternating lattice
has index one in its dual. These are the hypotheses at the
maximal level of \href{https://stephane-bijakowski.perso.math.cnrs.fr/PRgeometry.pdf}
{Bijakowski--Hernandez, \S2.4 and Theorem 2.30}, which gives a
smooth splitting model in these cases. At two the imaginary
quadratic subfield \(k\) already splits by Lemma~\ref{host:different},
so only the linear case (AL) occurs there.

Use the ordering of the \(E\)-embeddings in
\href{https://www.kwlan.org/articles/cpt-spl.pdf}
{Lan, Choices 2.3.1, Definition 2.3.3}: put \(\tau,\bar\tau\)
first as a consecutive pair in their \(E^+\)-place, and order the
other conjugate pairs consecutively. In the split case the pair
is in distinct integral field factors. In the inert unramified
quadratic case its restrictions to the maximal unramified subfield
are distinct, so the two are again separated by integral
idempotents over the splitting base. On these two different
summands the first quotient of the decreasing Hodge filtration
is labelled respectively by \(\tau,\bar\tau\).
In Lan's interleaved ordering the intervening graded piece on
the other summand is zero. In the increasing notation of
Bijakowski--Hernandez these are the two \emph{last} quotients.
Reversing the ordering within each list, with conjugate lists
paired term by term, is permitted in their \S2.4.

Here are the model comparisons used with the smoothness theorem.
For the single periodic self-dual lattice the open splitting model
parametrizes the polarized Hodge flags of the given ranks. Separating
the unramified idempotents turns Lan's interleaved flags into
the flags for the separate unramified embeddings; conversely they
are interleaved by direct sum. The polarization conditions can
also be compared on the common generic fibre, without identifying
their scheme structures on nonflat thickenings. Indeed both flag
problems are closed in the same relative product of flag varieties
over the naive moduli problem, by the endomorphism and pairing
conditions. They give the same unique flags on the generic fibre.
The Bijakowski--Hernandez model is smooth, hence equals the flat
closure of this generic fibre in that flag space, and is normal.
Lan's normalized open model is the normalization of that closure
on the selected generic components: the forgetful map of the
normalized underlying moduli to the naive moduli is finite and
generically open and closed at full lattice-stabilizer level
(Proposition 2.2.11). Finiteness persists after pullback to the
flag space and after normalization over the excellent base
(Definition 2.4.5 and Proposition 2.3.10). Consequently the normalized
open model on those components is the smooth one, by finite
birationality over a normal scheme.

Lan's Theorem 3.4.1 and Proposition 3.4.14 give a proper flat
toroidal compactification with geometrically regular fibres for
smooth cone data: geometric regularity of each fibre is
equivalent to that of the open splitting model. It is thus smooth.
By Proposition 3.1.2 and Definition 3.1.3 the extended filtration
is a filtration by subbundles on the Hodge bundle of the
semiabelian extension (identified with that of its dual by the
principal polarization, with conjugate labels if using the
dual-action convention). In particular the two locally free
first quotients above exist along the whole boundary.

For any trait write \(H\) for its integral invariant
differentials in the relevant unramified summand. If \(H'\) is
the kernel of its first quotient, then \(H'\) is saturated and
its generic fibre is the sum of all the other eigenspaces in
that summand. Hence
\[
   (H/H')^\vee=\{u\in H^\vee:u(H')=0\}
           =H^\vee\cap (H^\vee_{\rm frac})_\tau .
\]
The identical argument on the conjugate summand gives the
second test lattice. This uses the end quotient of the
filtration: in general an intermediate graded piece has a
different lattice. The determinant ratio of these dual quotients
is therefore an invertible sheaf which gives exactly the
ordinary test metric at \(\ell\). Its generic horizontal section
is a unit across the characteristic-zero boundary, as before.
Over the fixed geometric component the smooth proper model
has a single reduced special fibre. Lemma~\ref{ctwist:vertical}
and the common-refinement argument of
Proposition~\ref{cmeig:hyperspecial} give equality of all orders,
including on semiabelian tests. A different ordering or model
can be used for a different pair \(\tau,\bar\tau\), since orders
are calculated pairwise before multiplying sections.

Finally compare with the CM point of Proposition~\ref{cmeig:degree}
at each place and with its archimedean norm. Its finite extension
preserves normalized arithmetic degrees. This gives the claimed
global degree and hence the uncorrected product degree in
\eqref{ctwist:degree}.
\end{proof}

\paragraph{The selected good cover.}
The level condition on \(Y\) fixes its auxiliary marking to have
order one. Its finite marking is the half-marking at the odd integer
\(M\ge3\).
The characteristic-zero comparison can be made from the moduli
property. The universal object is proper and semiabelian, hence
abelian. Its rigidified line is relatively ample and its
polarization kernel is trivial, hence its polarization \(\lambda\)
is principal. The adjunction property for its biextension gives
the specified Rosati involution on \(O_E\). Conversely from these
PEL data pull the normalized Poincar\'e bundle back along
\((1,\lambda w_0)\). Its polarization is
\(\lambda(w_0+\bar w_0)=\lambda\); the rigidification, symmetry
and cube law are the canonical ones. Its biextension is the
one for \(\lambda\). This reconstruction also recovers the
original line from its graph condition. The rigidified isomorphism
is unique: an automorphism is multiplication by a unit on the
abelian scheme, which comes from the base by properness and
geometric connectedness and equals one at the identity.

Over the splitting field the determinant condition is precisely
the \(p,q\) multiplicity condition. On a reduced
characteristic-zero base the Hodge eigensummands are vector bundles
and their ranks can be tested at geometric points. The mark kills
\(\mathfrak m\) and on the complementary summand is the reduction
of a polarized Tate marking. By CRT these restrictions determine
it on the whole lattice. Its compatibility at all torsion orders
in the definition of an object is a condition on the local
lattice genus, not a choice of additional sections. More explicitly,
in a connected complex PEL family the Betti lattice, its pairing,
action and finite mark are locally constant; a compatible marked
Tate isometry at one point transports to every point. Equivalently
at a geometric point the reduction of a compatible tower is a
point of the prescribed compact-open marking orbit. Conversely
any representative of that orbit gives the required compatible
reductions at all \(n\). To check the roots of unity, fix the
primitive \(M\)-th root in the coefficient field and, at a geometric
point, extend it to a compatible system \((\zeta'_n)\). Relative
to the tower \((\zeta_n)\) of a Tate marking there is a unique
\(u\in\widehat{\mathbb Z}^{\times}\) with \(\zeta'_n=\zeta_n^u\)
for all \(n\ge1\). The reference lattice has an
integral similitude of multiplier \(u\) fixing the half-mark.
At \(\ell\mid M\), use \(1\) on the marked half and \(u_\ell\) on
its pairing-dual half. At \(\ell\nmid M\), use a central scalar of
norm \(u_\ell\) at each place of \(F_0\) over \(\ell\), using the
surjectivity of split and unramified unit norms in Lemma~\ref{cmeig:row}.
Precomposing a Tate marking by this similitude gives the prescribed
tower with the same marked sections. At order \(M\) the Weil pairing
with the marked basis and the fixed root then determines the basis
on the opposite half uniquely. This construction works over the
base at order \(M\) by perfect duality of finite étale groups. The
higher towers occur only in the geometric-point existence test.
This identifies the half-marking with the corresponding fine PEL
level on the selected components. After finite extension the
source moduli property gives a classifying map for the classical
universal object on its reduced smooth base, and the classical
moduli property gives the reverse map on the reduced smooth
source base. The reconstructed universal objects are isomorphic,
so both composites are the identity by uniqueness. The maps
preserve the underlying group and polarization, also after
restricting to a specified marked row.
\src{src/ReciprocityTheory/ReductiveGroups/ArithmeticGeometry/BirationalGeometry/ToroidalCompactification/TwoWeightEqOneAddRoot.lean}{131}{The compatible torsion isometries and primitive roots};
\src{src/ReciprocityTheory/ReductiveGroups/ArithmeticGeometry/BirationalGeometry/ToroidalCompactification/TwoWeightEqOneAddRoot.lean}{144}{The two marks, signature and geometric-point existence condition};
\src{src/ReciprocityTheory/ReductiveGroups/ArithmeticGeometry/BirationalGeometry/ToroidalCompactification/PushInvertible.lean}{29}{The universal object and its reduced-base moduli property};
\src{src/Development/AbelianVarietyTheory/Definitions/Divisors/NormNegCmConstantEqAtAllEmbeddings.lean}{187}{The selected level, period charts and polarization scale}.

On a test DVR the definition supplies a semiabelian group whose
generic fibre is cartesian over this universal group. Its fibres
are geometrically connected. The test vectors form a basis for
all integral eigenvectors at each labelled embedding, so they give
exactly the intersections used in Theorem~\ref{cmeig:splitmodel}.
Indeed the images of algebraic integers in the coefficient field
belong to the valuation ring; the equality expressing the
eigenvector property therefore holds integrally if and only if
it holds over the fraction field, by torsion-freeness of the
smooth Lie lattice. The theta generator in the test contributes
no factor to the parallel character. At infinity the comparison
uses the polarization metric \(\pi H\) on Lie vectors; this is
the metric used in Proposition~\ref{cmeig:degree}. Thus the
ordinary parallel determinant comparison follows on the selected
principal family with its given period readings.
\src{src/Formalism/LocalToGlobal/ProjectiveLattices/VAtVNeBot.lean}{53}{Cartesian semiabelian tests and the universal integral eigenbasis condition};
\src{src/ReciprocityTheory/ReductiveGroups/ArithmeticGeometry/SemiabelianSchemes/AbelianSchemeQuotients/IsAbelianScheme.lean}{41}{Connected geometric fibres of the semiabelian group};
\src{src/Development/AbelianVarietyTheory/Definitions/CMLattices/IsDiffOf.lean}{101}{Polarization norm in the CM comparison}.

\begin{corollary}[Short sections in the final weight on a good cover]
\label{cmeig:final}
Fix one good deep cover \(Y\) as in Corollary~\ref{crev:good}
on a selected host, with \(p,q\ge1\). Take the admissible final
weight \(\lambda\) and the level, denominator, field and bad-prime
size bounds of Theorem~\ref{bfam:oncharts}.
On each marked component \(Z\), its regular cusp rule space of
weight \(\lambda\) is spanned over the coefficient field by
integral sections \(s_i\) with
\[
     \log\|s_i\|_{\sigma}\le
          C\bigl((frL)^3 W+(frL)^5\bigr),\qquad
     W=\sum_{\tau,i}|\lambda_{\tau i}|,
\]
at every embedding in the theorem. The same constant works over
these choices under the common stage bounds.
\end{corollary}
\begin{proof}
Start with the integral generators of Corollary~\ref{crev:good}.
Multiply by the key of Theorem~\ref{ctwist:semi} and by the
product with exponents \(-j_\tau\) in \eqref{ctwist:degree}.
Theorem~\ref{cmeig:splitmodel} computes the latter degree at the
actual integral eigenlattices. Outside the excluded primes the
level is hyperspecial, so the common order is integral by
Proposition~\ref{cmeig:hyperspecial}. At the remaining primes
round the product order once as in \eqref{bfam:flat}.
Proposition~\ref{host:period} supplies the uniform numerical
degree bound. Both unit sections and their inverses preserve the
cusp condition by Corollary~\ref{crev:twist}. The final-weight
estimate of Theorem~\ref{bfam:oncharts} now gives the claim.
This argument starts with \(Y\) and the stated numerical bounds.
Selecting such covers compatibly at all stages and identifying the
restriction maps on the modified flags enter separately in the
Hilbert--Samuel estimate.
\end{proof}



Once \(\beta\) is known, the required archimedean multiplication
estimate has no further dimension factor. For the shared Petersson
measure on the component, let \(B(y)\) be the supremum of
\(|u(y)|^2/\|u\|_2^2\) over nonzero sections of the first frame space,
using the fiber norm of its weight. Suppose \(B(y)\le\beta\) on the
support of the measure, and the fiber metric on a product is the
tensor product of the fiber metrics. For any frame section \(u\)
and higher weight section \(v\), pointwise multiplication and
integration give
\[
 \|uv\|_2^2=\int |u(y)|^2|v(y)|^2\,d\mu(y)
       \le \beta\|u\|_2^2\|v\|_2^2.
\]
Using the original frame norm for \(\|u\|_2\) gives the operator
bound in the determinant argument. This requires the quotient
measure and fiber normalizations for the three weights to agree
under tensor product. A bound for the kernel relative to a
separately rescaled measure does not imply the same \(\beta\).
\src{src/Development/AbelianVarietyTheory/Proofs/DivisorProofs/PeterssonNormMulLeSqrtBergmanMulComponent.lean}{79}{Component product estimate}.

\begin{proposition}[Absorption of the component errors]\label{newtheta:absorption}
Assume \eqref{newtheta:budgets} and \eqref{newtheta:finiteHS},
\(1\le D,\ 3D+5\le T,\ \delta\ge1\), and
\[
 b\ge1,\quad \log b\le P_0,\quad TP_0\le\sqrt b,\quad
 z=7(D^2/T+D\rho)\le\tfrac12,\quad \rho\ge0,\quad
 zP_0\kappa\le b .
\]
If \((1-z)d_0\le B\) and \(g_M\ge g-d_0c\), \(c\ge0\), then
\[
 g_M\ge B\big[Th-3(D+1)|h|
       -26(26D^2+2zT)\log(2+|h|)-4000b-2c\big].
\]
\end{proposition}
\begin{proof}
Write \(P=P_0,\ v=\kappa/T\) and \(L_*=27Pv+92P\). Then
\(TP\le b,\ (TP)^2\le b,\ \log D\le\log T\le P/2\).
The expression \(2+|h|+\log\Lambda\) is at most
\((2+|h|)(1+\log\Lambda)\). Since
\(\log(1+\log\Lambda)\le\log\Lambda\), \eqref{newtheta:budgets}
gives \(L_h\le L_*+26\log(2+|h|)\). Specifically, the terms free
of \(h\) are at most \(27P(v+1)+26(P+P/2+P)=L_*\).
The rank comparisons give \(0\le u\le z\le1/2\),
\(d_0(1-u)=B\), and \(d_0u/B=u/(1-u)\le2z\).
The bracket multiplied by \(u\) is nonnegative. Substitution in
\eqref{newtheta:finiteHS} leaves the claimed height coefficient
and the following upper bound on the error independent of \(h\):
\begin{align*}
 J+26D^2(L_*+\log T)&+(T+D+1)\log\beta+(T+D^2)(G+2)\\
 &+2z\big[T(\log2+L_*+\tfrac12\log\beta)
                   +\log\Lambda_D+\log\Lambda_T\big].
\end{align*}
For explicit bounds use
\[
 D^2P\le b/14,\qquad D^2Pv\le b/7,\qquad zTPv\le b .
\]
Indeed \(D^2\le zT/7\le T/14\); the last inequality displayed
is \(zP\kappa\le b\). Since \(D\le D^2\) and \(DP\le TP\le b\),
the first cost \(J\) is at most \(6b/7+8b<9b\).
The second is at most \(26(27/7+92.5/14)b<300b\).
The third is at most \(2TP\le2b\). For the fourth,
\(G+2\le5DP\) by \(D,P\ge1\), and \(T+D^2\le2T\);
hence it is at most \(10TDP\le10(TP)^2\le10b\).
Finally \(zT\log2\le b/2,\ zT L_*\le27b+46b,\)
\(zT\log\beta/2\le b/4\), and
\(z(\log\Lambda_D+\log\Lambda_T)\le2b\).
Thus the last cost is less than \(200b\). Their sum is bounded by
\(4000b\). Since \(d_0\le2B\), the transfer costs at most \(2Bc\).
The argument nowhere requires a sign for \(h\).
\end{proof}

\begin{proposition}[Summing the component degrees]\label{newtheta:aggregate}
Let \(I\) be a finite nonempty set of components. For \(i\in I\) let
\(d_i>0\) be its rank, \(h_i,\ell_i,g_i\) its height, comparison
slope, and transferred degree, and \(B_i>0\) its degree share.
Let \(0\le \rho<1,\ z\ge0,\ D\ge0,\ T\ge3D+5,\ A>0,\ R\ge0\).
Assume for some \(b_0>0\), \(b_0\le B_i\le b_0/(1-\rho)\) and
\(B_i\le d_i\). Suppose
\[
 g_i\ge B_i[Th_i-3(D+1)|h_i|-A\log(2+|h_i|)-R],
 \qquad 2\sum_i g_i/\sum_i d_i\le\mu .
\]
If \(s,c,X\ge0\) satisfy
\[
 h_i+s\ge0,\qquad
 (1-\rho)(\ell_i+s)_+\le h_i+s+c,\qquad
 |I|^{-1}\sum_i\ell_i\ge-X ,
\]
and for \(E\ge1,\ C_1,C_2\ge0\)
\[
 \max\{TX,Tc,s(6D+10+T(z+2\rho))\}\le C_1E,\qquad
 A\log\max(A,1)+R\le C_2E ,
\]
then \(\mu\ge-2(3C_1+C_2+4)E\).
\end{proposition}
\begin{proof}
Set \(K=3D+4\), \(U=T-K>0\), \(B_\Sigma=\sum_i B_i\).
For \(v\ge0\), \(\log t\le t-1\) at \(t=(2+v)/A\) gives
\(A\log(2+v)\le v+2+A\log\max(A,1)\).
Since \(|h_i|\le h_i+2s\), the bracket in the component bound is
at least \(U(h_i+s)-(T+K)s-2-C_2E\).
Using nonnegativity of the positive parts and their unweighted mean,
\[
 \frac{\sum_i B_i(\ell_i+s)_+}{B_\Sigma}
 \ge(1-\rho)|I|^{-1}\sum_i(\ell_i+s)_+
 \ge(1-\rho)(s-X)_+ .
\]
Consequently the weighted average of \(h_i+s\) is at least
\((1-\rho)^2(s-X)_+-c\ge(1-2\rho)s-X-c\).
The last inequality follows from \((s-X)_+\ge s-X\),
\((1-\rho)^2\ge1-2\rho\), and \((1-\rho)^2\le1\).
It follows that
\[
 \frac{\sum_i g_i}{B_\Sigma}
 \ge-[2K+2\rho U]s-U(X+c)-2-C_2E
 \ge-(3C_1+C_2+2)E .
\]
Indeed \(2K+2\rho U\le6D+10+T(z+2\rho)\) and \(U\le T\).
As \(B_\Sigma\le\sum_i d_i\) and the last bound is nonpositive,
the same lower bound holds with \(\sum_i d_i\) in the denominator.
Multiplication by two proves the assertion, including the stated
slightly larger constant.
\end{proof}
\src{src/Development/AbelianVarietyTheory/Proofs/DivisorProofs/UniformLowerBoundOfSums.lean}{16}{Component shares and covering slope}.

\paragraph{The geometric hypotheses on the family used in Section~\ref{main:assembly}.}
Fix one selected scale, a stage in the prescribed range and its
deep cover \(Y\) with coefficient field \(F\), level \(M\), and
component \(Z\). Its auxiliary marking has order \(A=1\) by the
selected good-level condition. In the indefinite short-section
step one has \(p,q\ge1\), \([E:\mathbb Q]/2>2\), \(p+q=r\).
In the notation of the source, the stage hypothesis is
\(Y.\mathrm{Good}(e,o)\), with
\(\log M\le30L\), host field size at most \(3L\), \(\log(fr)\le L\)
and \(o\le391(frL)^2\). The dominant final weight has its block
gap at least \(8fr+10\) and logarithmic size at most \(10L\).
The following applications on this same \(Y,Z\) distinguish the
points where the paper uses geometric theorems of the original proof.

\begin{enumerate}
\item On this \(Y\), Corollary~\ref{crev:good} identifies the
principal framed charts with their marked periods at each good
embedding. The identity Fourier mark kills the zero transvection
mode, and Proposition~\ref{crev:half} places the positive half
maximum on the thick part for the same self-dual row. The theta
line has Haar mean square one. Proposition~\ref{crev:mass}
computes the total flag volume and the density \(1/V\) of adelic
probability, giving the conditional Petersson measure and the
normalized Bergman bound on each component. These are deductions
on the selected good cover. The construction of such covers
with their field-size bounds is the separate existence step
specified after Corollary~\ref{crev:good}.
\src{src/Development/AutomorphicArithmetic/ThetaCorrespondence/ThetaPairLifts/HodgeBundleWeights/ExistsGoodCoverAtMinimalLevel.lean}{29}{The field-spanning family and its two metric estimates};
\src{src/ReciprocityTheory/ReductiveGroups/ArithmeticGeometry/FundamentalDomains/MaxBoundPf.lean}{12}{Cusp bound};
\src{src/Development/AbelianVarietyTheory/Proofs/DivisorProofs/FirstMinAtHalfMax.lean}{111}{Minimum at a positive half maximum on the same component};
\src{src/Development/WeilBound/RepresentationTheoryGL/TubeMeasures/DisjointTubes/MeasureSumBounds/SumMuFLe.lean}{110}{Mass in the chosen normalization}.

\item Use the rigidified cubical graph line on the semiabelian tests
of Theorem~\ref{ctwist:semi}. For the other line the relevant lattice
is the determinant ratio for the two integral Lie eigenspaces.
Theorem~\ref{cmeig:splitmodel} gives the local comparison for every
test on \(Z\), separately for each \(w,\tau\), and evaluates the
degree on its marked product CM point. For each \(w\)
outside the primes dividing \(2d_0|\Delta_E|\operatorname{lcm}(M,A)\),
the source supplies on this \(Z\) an unramified test, hence a test
taking integral values on nonzero evaluations. At those places
the integral order also follows from the single reduced fibre
in Proposition~\ref{cmeig:hyperspecial}. At the bad primes the
single rounding of the product costs at most \(\sum_{\ell\in B}\log
\ell\). The source then applies the component CM degree formula,
the key and the product to the same spanning set. For the
characteristic-zero rule spaces Corollary~\ref{crev:twist} proves
cusp compatibility in both directions for these constant-norm
units.
\src{src/Development/AbelianVarietyTheory/Definitions/Divisors/MulSmulIdemMem.lean}{264}{Exact excluded primes and integral-valued test};
\src{src/Development/AutomorphicArithmetic/ThetaCorrespondence/ThetaPairLifts/HodgeBundleWeights/UnramifiedPointsAtDeepLevel.lean}{11}{Unramified CM tests on this cover};
\src{src/Development/AutomorphicArithmetic/ThetaCorrespondence/ThetaPairLifts/HodgeBundleWeights/ShortSpanCompositionClosed.lean}{74}{Twists on the chosen component}.

\item For the Hilbert--Samuel step take the modified toroidal flag
component for \(Z\) and the two birational weights of its weight
path. Restriction in degree zero must be surjective for the
transverse flags, and the successive symbols and deepest-space
generators must be those of the same restriction. The short integral
generators and column norms of the first frame, the Bergman bound
\(\beta\), the short sections in the subsequent weights, and the
multiplication bound \(\sqrt\beta\|v\|_{\rm frame}\) use that
component's metric. The paper derives the finite-degree inequality
from these hypotheses in Corollary~\ref{newtheta:effective}.
Their simultaneous realization for the family is invoked at the
following source steps.
\src{src/Development/AutomorphicArithmetic/ThetaCorrespondence/ThetaPairLifts/HodgeBundleWeights/ExistsWeightHodgeBundleGenerated.lean}{22}{Toroidal flag and modified bundles on \(Z\)};
\src{src/Development/AbelianVarietyTheory/Proofs/DivisorProofs/HilbertSamuelFiltrationAtGoodSections.lean}{51}{Restriction and filtration};
\src{src/Development/AutomorphicArithmetic/ThetaCorrespondence/ThetaPairLifts/ThirdPlaceSeries/ShortVectorsSpanCuspLattice.lean}{33}{Short vectors at the family weights};
\src{src/Development/AutomorphicArithmetic/ThetaCorrespondence/ThetaPairLifts/ThirdPlaceSeries/HilbertSamuelBergmanFrameBound.lean}{26}{Frame and kernel bounds at the family stage};
\src{src/Development/AutomorphicArithmetic/ThetaCorrespondence/ThetaPairLifts/HodgeBundleWeights/ExistsDegreeParametersBound.lean}{33}{Combination on a common cover and component}.
\end{enumerate}

\subsection{The coefficient field and the two flag systems on a family stage}
\label{zetageo:section}

\begin{proposition}[A coefficient field for the deep level]\label{zetageo:field}
Write \(s(F)=\log[F:\mathbf Q]+\log\operatorname{rd}(F)\), where
\(\operatorname{rd}(F)=|\Delta_F|^{1/[F:\mathbf Q]}\).
At a stage as above let \(F_0\) be the initial coefficient field,
\(K=\mathbf Q(\sqrt{-d})\) the imaginary quadratic field of the
marking, and \(g_0,g_1\) the two host level factors. Suppose that
\(\Delta_K=-d\), that \(3,5\) split in \(K\), and that
\((15,g_0g_1)=1\). Put \(m=g_0g_1\,15^u\) for \(u\ge0\), \(m\ge3\).
There is a choice of a half of the split \(15^u\)-level, and a
field \(F\supset F_0K(\zeta_m)\) containing the CM \(j\)-invariants
and normalized Weber values of the \(m\)-torsion for the full order
\(\mathcal O_K\). It satisfies
\[
 s(F)\le s(F_0)+\log 2+\tfrac32\log d+3\log m.
\]
On the family with \(u=\lfloor3L/\log3\rfloor\), \(L\ge1\),
\(\log g_i\le9L+1,\ \log d\le L,\ s(F_0)\le217Z_0^2\) and
\(Z_0=frL\ge8L\), this gives \(s(F)\le391Z_0^2\).
\end{proposition}
\begin{proof}
Choose one prime over three and one over five. Equivalently, choose
a root of \(-d\) at each prime and combine the roots to a unit
\(a\bmod 15\). The ideal \(I=(15,a+\sqrt{-d})^u\) is coprime
to its conjugate and \(I\bar I=(15^u)\), by reduction at the four
distinct primes. This is the required half level.
Let \(R\) be the ray class field of \(K\) of modulus \(m\mathcal O_K\).
The ray class exact sequence gives
\([R:K]\le h_K|(\mathcal O_K/m)^\times|\le d m^2\).
Indeed each ideal class has an integral representative of norm
at most \(2\sqrt d/\pi\) by Minkowski. There are at most
\(\sum_{n\le\sqrt d}\sum_{a\mid n}1\le\sqrt d(1+\log\sqrt d)\le d\)
such ideals, by the quadratic Euler factors. The conductor of each
ray class character divides \(m\mathcal O_K\). The
conductor--discriminant formula in \(R/K\) thus gives
\(\operatorname{rd}(R)\le\sqrt d\,m\). Conjugation preserves the
modulus, so \(R/\mathbf Q\) is normal. The Artin action on
\(\zeta_m\) factors through this ray class group by the ideal norm.
The main theorem of complex multiplication in its Weber form gives
the stated CM values in \(R\)~\cite{coxCM}.

Take \(F=F_0R\). Degrees are submultiplicative under compositum, and
\(\operatorname{rd}(F)\le\operatorname{rd}(F_0)\operatorname{rd}(R)\).
For the latter inequality, base change of the different gives
\(\mathfrak D_{R/\mathbf Q}\mathcal O_F\subseteq
\mathfrak D_{F/F_0}\). Apply the tower discriminant formula and
take the \([F:\mathbf Q]\)-th root. The two bounds for \(R\) give the
claimed estimate. Since \(15\le3^3\), one has \(u\log15\le9L\) and
\(\log m\le27L+2\). The increase in \(s\) is at most \(90L\)
for \(L\ge1\). In fact \(90L\le(90/64)Z_0^2\), giving
\(s(F)<219Z_0^2\) inside the stated budget. This estimate uses
the same \(L,f,r\) for every component of the stage.
\end{proof}

Propositions~\ref{addh:constants} and~\ref{addh:level} realize the
component constants and the split half-level map on the principal
marked packet over this \(F\). The compatible charts and the integral
comparisons on modified flags enter the construction of the good cover \(Y\).
Their combination and the bounds on the selected stage are recorded
in the \src{src/Development/AbelianVarietyTheory/Proofs/DivisorProofs/ExistsGoodCoverOverSizedField.lean}{111}{construction of the marked cover}
and the \src{src/Development/AbelianVarietyTheory/Proofs/DivisorProofs/DeepExistenceStatement.lean}{29}{uniform field-size calculation}.

\begin{lemma}[One modification and the cost of changing the lattice]
\label{zetageo:lattice}
On a smooth compactified flag component over its coefficient field
suppose the two systems of Proposition~\ref{newtheta:modified} have
no base points on the marked open \(U\). There is one smooth proper
modification, an isomorphism over \(U\), on which both moving systems
are generated by their given sections. If the rational map of either
system is birational onto its image, its moving system on the
modification has that same property.

Let \(M,M'\) be full \(\mathcal O_F\)-lattices in the same space of
sections of rank \(a>0\), with Hermitian metrics at infinity. Put
\(H=M\cap M'\). At each embedding \(\sigma\), let \(d_\sigma\)
be the determinant norm of the identity from the metric on \(M\)
to that on \(M'\). For normalized Arakelov slopes,
\[
 \widehat\mu(M')-\widehat\mu(M)
 =\frac{\log|M'/H|-\log|M/H|-\sum_\sigma\log d_\sigma}
        {a[F:\mathbf Q]} .
\]
Here both embeddings over a complex place are counted.
\end{lemma}
\begin{proof}
Principalize the product of the two base ideals in characteristic
zero, with centres outside \(U\)~\cite{wlodarczyk}. On the resulting
integral variety each factor is invertible: for \(IJ\) invertible
the fractional ideal \(J(IJ)^{-1}\) is an inverse to \(I\).
Divide the pulled-back sections by a local generator of each
factor. Their coefficients now generate the unit ideal. The two
maps agree with the original rational maps on the dense open.
Proposition~\ref{newtheta:modified} then gives the adjoint
inclusion whenever its extension and cusp identifications hold.
For the slope formula, compute the degree of the determinant
of \(H\) first with each metric. Passing from \(H\) to either
superlattice adds the logarithm of its finite index in the
unnormalized degree. Changing the metric subtracts the logarithm
of the determinant norm at each embedding. Divide by the common
rank and field degree.
\end{proof}

In particular equality of the two rational section spaces determines
neither finite index in this formula. For a comparison with a
separately defined integral flag lattice, these indices and the
determinant norms must be evaluated on the same \(Y,Z\) that
produced the sections. This is one place where a pointwise proof
of \eqref{scope:missing} needs more than an archimedean norm estimate.

\begin{proposition}[A finite test for the component mean height]\label{zetageo:mean}
Fix a good cover, a weight \(u\) with the required gap, and the
finite valuation conditions used at that gap. On each component
\(Z\) let \(D\) be the flag dimension, \(\delta_Z>0\) its prescribed
top degree, \(d_F=[F:\mathbf Q]\), and \(L_{Z,m}\) its mixed integral
cusp lattice in weight \(mu\). Give it the supremum of the marked
norms over all good embeddings and frames.
Suppose the following sections belong to these actual lattices and
their products belong to the lattices of the summed weights, with
submultiplicative norms.
\begin{itemize}
\item There are \(D+1\) nonzero sections \(s_0,\ldots,s_D\) in
degree \(m_0\in8\mathbf Z_{>0}\). The ratios \(s_i/s_0\)
\((1\le i\le D)\) are algebraically independent over \(F\) and
the function field of \(Z\) has degree \(r=m_0^D\delta_Z\) over
the field of these ratios.
\item There are \(r\) nonzero sections \(b_j\) in one degree
\(m_0 k_0\) such that \(b_j/s_0^{k_0}\) form a basis for this field
extension. All sections used here have finite positive norm.
\end{itemize}
Then the mixed component volume slope satisfies
\[
 \operatorname{compVolSlope}_Z(u)\ \ge\
 a_Z:=-\frac{\sum_{i=0}^D\log\|s_i\|}{(D+1)m_0}.
\]
If these finite systems are available on every component and
\(|I|^{-1}\sum_{Z\in I}a_Z\ge-X\), they give the required component
mean with error \(X\). The degrees \(m_0,k_0\) and systems may
depend on the weight and the component.
\end{proposition}
\begin{proof}
In degree \(m=m_0(k_0+k)\) take all products \(b_j s^\alpha\) with
\(\alpha_i\ge0,\ \sum_i\alpha_i=k\). Dividing by \(s_0^{k+k_0}\)
and using the two independence assumptions proves that these products
are \(F\)-linearly independent. Multiply them by a fixed integral
basis of \(\mathcal O_F\) over \(\mathbf Z\). We obtain
\(R_k=d_F r\binom{k+D}{D}\) independent integral sections \(v_j\).
In particular \(D!R_k/(d_Fm^D)\to\delta_Z\).
On averaging over all the monomials each exponent has mean
\(k/(D+1)\), by symmetry. Thus, writing
\(\lambda_j=-m^{-1}\log\|v_j\|\), one has
\(\liminf_k R_k^{-1}\sum_j\lambda_j\ge a_Z\).
The fixed sections \(b_j\) and the coefficient basis cost \(O(1/m)\)
in this average.

For any real \(t\), form the integer combinations with coefficients
between zero and \(\lfloor e^{m(t+\lambda_j)}/R_k\rfloor\) inclusive.
They are distinct and have norm at most \(e^{mt}\). Their log count
is at least both zero and
\(mR_k(t+R_k^{-1}\sum_j\lambda_j)-R_k\log R_k\).
The mixed volume \(V_Z(t)\) takes the limsup of the log count on
degrees \(m=8k'\), multiplied by \(D!/(d_Fm^{D+1})\).
Our degrees are a subsequence and \(\log R_k=o(m)\). Hence
\[
 V_Z(t)\ge\delta_Z\max(0,t+a_Z),\qquad
 \inf_{t\in\mathbf R}\bigl(V_Z(t)/\delta_Z-t\bigr)\ge a_Z .
\]
The infimum is the component volume slope. Taking the \(a_Z\) as
witnesses gives the mean assertion. No interchange of a component
average with a limsup is needed.
\end{proof}

For the family's \src{src/Development/AbelianVarietyTheory/Definitions/Divisors/SemiposAtOfPositive.lean}{101}{mean height condition},
the lattices in this test use the height valuations at the primes
selected by the weight and gap and the recorded valuations elsewhere.
The numerical requirement is the displayed mean with
\(X=C\,\mathrm{errX}(s,u,g)\), for a single \(C\) over the admitted
stages and weights. If the \(s_i\) give a finite map to
\(\mathbf P^D\) from a model of degree \(\delta_Z\) in weight \(u\),
the projection formula gives the required degree \(r\).
A finite map also supplies the generic basis from sufficiently high
twists of its coherent direct image on projective space. Pulling that
basis back into the specified cusp spaces and placing it in the mixed
lattices are part of the stated hypotheses. Only the norms of the
\(D+1\) parameter sections enter the final mean.

There is a quantitative way to choose the parameters. Suppose \(s\)
integral sections in degree \(m_0\), of norm at most \(H\), define
a base-point-free map onto a projective variety of dimension \(D\)
and degree \(\delta'\). Its Chow form in \(D+1\) hyperplanes is
nonzero and has degree \(\delta'\) in each block. Polynomial
interpolation in each coefficient supplies hyperplanes with
coefficients in \(\{0,\ldots,\delta'\}\) where this form is nonzero.
Their common center misses the image. Projection from that center
is finite on the image: the pullback of a hyperplane is ample, so
no curve is contracted. Its degree is \(\delta'\) by the projection
formula. The resulting parameter sections have norm at most
\(s\delta'H\). If the original map has generic degree \(g\) with
\(g\delta'=m_0^D\delta_Z\), and the generic basis belongs to the
mixed lattices as above, this gives the explicit lower bound
\(-m_0^{-1}\log(s\delta'H)\) for the component volume slope.

To pass from one stage to the family, first fix a common threshold
at which the host, core-field and size bounds of
Proposition~\ref{zetageo:field} hold for every admitted point and
stage. On the resulting good covers, suppose the marked flag
identifications, integral restrictions, short frames and their
metric bounds specified above hold with uniform constants. Suppose
also the height and rank comparisons following Lemma~\ref{geo:component}
hold there, in particular the lower bound on the component average
of \(\ell_Z\), which can now be supplied by
Proposition~\ref{zetageo:mean}. Proposition~\ref{newtheta:aggregate} then gives
\(\mu_{\rm cover}\ge-2(3C_1+C_2+4)E_P\) at every such stage.
After the stated descent and Gram comparisons,
Proposition~\ref{geo:covolume} absorbs the exponential error with
one threshold depending only on the fixed constants. Only the
finitely many thresholds for uniform estimates are combined here.
The marked integral comparisons and the mean height of the actual
components remain inputs to this application. The formal proof
combines the height inputs at the
\src{src/Development/AutomorphicArithmetic/ThetaCorrespondence/ThetaPairLifts/ThirdPlaceSeries/HeightLowerBoundGeneral.lean}{21}{component height theorem}
and passes from the covers to the grid at the
\src{src/Development/AutomorphicArithmetic/ThetaCorrespondence/ThetaPairLifts/ThirdPlaceSeries/TopCovolG.lean}{49}{covolume descent}.
